%=====================================================================
%  微分方程 · 期中之后（L13--L23）完整讲义
%  由六份分册讲义合并：L13 / L14-15 / L16-17 / L17-18 / L18-L20 / L21-L23
%  合并只做结构与编号，不改各章推导；各章 § 编号在章内自成一套（§0--§12）
%=====================================================================
\documentclass[11pt]{ctexart}
%=====================================================================
% theme-violet.tex —— 讲义精讲 · 紫花主题
% 用法：\documentclass[11pt]{ctexart}
%       \input{assets/theme-violet}
%       \begin{document} ... \end{document}
% 只用 ctexart 自带能力 + xcolor/enumitem/fancyhdr/pifont/amsmath，
% 不依赖 tikz/tcolorbox，保证 MiKTeX 可直接编译。
%=====================================================================

\usepackage[a4paper,left=2.0cm,right=2.0cm,top=2.1cm,bottom=2.3cm,%
            headsep=12pt,footskip=22pt]{geometry}
\usepackage{amsmath,amssymb}
\usepackage{needspace}
\usepackage{xcolor}
\usepackage{array}
\usepackage{enumitem}
\usepackage{pifont}
\usepackage{fancyhdr}
\usepackage[hidelinks]{hyperref}

%---------------------------------------------------------------- 配色
\definecolor{violet}{RGB}{107,78,143}      % 主色：紫
\definecolor{violetdeep}{RGB}{63,44,92}    % 深紫：标题/强调
\definecolor{petal}{RGB}{183,159,208}      % 花瓣：装饰线、花饰
\definecolor{bandviolet}{RGB}{247,243,251} % 浅紫底：档案/提示框
\definecolor{bandcream}{RGB}{252,249,242}  % 米白底：易错/注意框
\definecolor{rulepurple}{RGB}{217,207,230} % 分隔线
\definecolor{inkgray}{RGB}{92,96,110}      % 次要文字

%---------------------------------------------------------------- 字体
% 有 Noto CJK 用 Noto，没有就退回 ctex 默认（Windows 下为宋体/黑体）
\IfFontExistsTF{Noto Serif CJK SC}{\setCJKmainfont{Noto Serif CJK SC}}{}
\IfFontExistsTF{Noto Sans CJK SC}{\setCJKsansfont{Noto Sans CJK SC}}{}

%---------------------------------------------------------------- 版式
\setlist{nosep,leftmargin=1.9em}
\linespread{1.28}
\setlength{\parskip}{3.5pt}
\setlength{\parindent}{0pt}
\setlength{\tabcolsep}{5pt}
\renewcommand{\arraystretch}{1.28}

%---------------------------------------------------------------- 花饰
% \ding{95} = 五瓣小花，\ding{96} = 星形花；都来自 pifont，必定存在
\newcommand{\orn}{\textcolor{petal}{\ding{95}}}          % 行内小标
\newcommand{\ornstar}{\textcolor{petal}{\ding{96}}}      % 星形花
\newcommand{\flowerrule}{%                                % 居中花饰分隔线
  \par\vspace{7pt}%
  \noindent\makebox[\linewidth]{\color{petal}%
    \rule{0.20\linewidth}{0.6pt}\hspace{8pt}{\small\ding{95}}\hspace{4pt}%
    \raisebox{1.1pt}{\tiny\ding{95}}\hspace{4pt}{\small\ding{95}}\hspace{8pt}%
    \rule{0.20\linewidth}{0.6pt}}%
  \par\vspace{7pt}}

%---------------------------------------------------------------- 标题与节头
\newcommand{\doctitle}[3]{%   #1 主标题  #2 副标题  #3 元信息（可用 \\ 分行）
  \begin{center}
    {\color{petal}\small \orn\enspace #2\enspace \ornstar}\\[7pt]
    {\LARGE\bfseries\color{violetdeep} #1}\\[10pt]
    {\color{petal}\rule{0.55\linewidth}{1pt}}\\[4pt]
    {\small\color{inkgray} #3}
  \end{center}\vspace{2pt}}

\newcommand{\hdr}[2]{%   #1 章节号（如 §0）  #2 标题
  \par\vspace{12pt}%
  \noindent\colorbox{violetdeep}{\color{white}\bfseries\small\;#1\;}\hspace{7pt}%
  {\color{violetdeep}\large\bfseries #2}\par\nobreak
  \vspace{3pt}{\color{petal}\hrule height 1.2pt}\vspace{6pt}}

\newcommand{\subhdr}[1]{%
  \par\vspace{6pt}%
  \noindent{\color{violet}\bfseries #1}\par\vspace{1pt}}

%---------------------------------------------------------------- 提示框
\newcommand{\callout}[1]{% 浅紫底：讲义档案、读法说明
  \par\medskip\noindent
  \fcolorbox{rulepurple}{bandviolet}{\parbox{\dimexpr\linewidth-2\fboxsep-2\fboxrule}{\vspace{2pt}#1\vspace{2pt}}}%
  \par\medskip}

\newcommand{\warnbox}[1]{% 米白底：易错、边界条件
  \par\medskip\noindent
  \fcolorbox{rulepurple}{bandcream}{\parbox{\dimexpr\linewidth-2\fboxsep-2\fboxrule}{\vspace{2pt}#1\vspace{2pt}}}%
  \par\medskip}

\newcommand{\gapitem}[2]{% #1 缺口类型（紫） #2 内容
  \par\vspace{2pt}\noindent
  \hangindent=1.4em\hangafter=1
  {\color{petal}\small\ding{96}}\enspace{\bfseries\color{violetdeep}#1}\enspace #2\par}

% 题干块（\qstem）：每道例题／作业的「详解」开头必须先复述题干与出处。
% 起因：2026-10-25 第五章讲义把七道作业的题面集中放在一页、详解另起一页，
% 读者翻到「第 4 题」时看不到条件，等于一道没有题干的题目。
% 规则：详解块内自带题干，读者不翻页就能读懂题在问什么。
\newcommand{\qstem}[1]{%
  \par\medskip\noindent
  \fcolorbox{violet}{bandviolet}{%
    \parbox{\dimexpr\linewidth-2\fboxsep-2\fboxrule}{%
      \vspace{2pt}%
      {\small\bfseries\color{violet}题干}\hspace{0.8em}#1%
      \vspace{2pt}}}%
  \par\medskip}

% v2: 保护题干与开头，长题可跨页，续页补上下文。
\newcommand{\qbreak}{\needspace{6\baselineskip}}
% 独立练习：ID、来源、题干、考点、答案标签（通常 ans:ID）。
\newcommand{\practice}[5]{%
  \qbreak\subhdr{练习 #1}\qstem{（#2）#3}%
  \textbf{考点}：#4\par
  详解见附录第~\pageref{#5}~页。\par\medskip}
% 附录开头：ID、来源、条件摘要；自动产生 ans:ID 标签。
\newcommand{\answerstart}[3]{%
  \qbreak\subhdr{#1\quad 详解}\label{ans:#1}\qstem{（#2）#3}}
\newcommand{\answercontinue}[2]{%
  \needspace{5\baselineskip}\subhdr{#1\quad 续解}%
  {\small\color{inkgray}本页接续：#2}\par\medskip}
% 延迟提供默认值，让旧讲义可在导言区保留自己的 \newcommand{\step}。
\AtBeginDocument{\providecommand{\step}[1]{\needspace{3\baselineskip}\par\medskip\textbf{#1}\par}}

% 前提旁注：只放在「需要理解的那一步」后面，给陈述 + 本题代入，不重证。
\newcommand{\aside}[1]{%
  \par\vspace{1pt}\noindent
  \hangindent=1.6em\hangafter=1
  {\small\color{inkgray}%
    \hspace*{0.6em}{\color{petal}$\triangleright$}\hspace{2pt}#1}%
  \par\vspace{2pt}}

%---------------------------------------------------------------- 表格
% 本机 MiKTeX 的 colortbl 与 array 版本不匹配（\rowcolor 会报错），
% 故表格统一用「上粗紫线 + 表头 + 中细线」的三线式，不依赖 colortbl。
\newcommand{\tabtop}{\par\vspace{3pt}{\color{violet}\hrule height 1.2pt}\vspace{1pt}}
\newcommand{\tabmid}{\par\vspace{1pt}{\color{petal}\hrule height 0.5pt}\vspace{1pt}}
\newcommand{\tabbot}{\par\vspace{3pt}{\color{rulepurple}\hrule height 0.5pt}\vspace{8pt}}
\newcommand{\thead}[1]{{\bfseries\color{violetdeep}#1}}

%---------------------------------------------------------------- 页眉页脚
\pagestyle{fancy}
\fancyhf{}
\renewcommand{\headrulewidth}{0.4pt}
\renewcommand{\headrule}{\hbox to\headwidth{\color{rulepurple}\leaders\hrule height \headrulewidth\hfill}}
\fancyhead[L]{\small\color{inkgray}\songti \doctag}
\fancyhead[R]{\small\color{inkgray}\songti \coursetag}
\fancyfoot[C]{\small\color{inkgray}\orn\enspace{\color{violet}\thepage}\enspace\orn}

% 每篇讲义在导言区设这两个值
\newcommand{\doctag}{}
\newcommand{\coursetag}{}
\newcommand{\settags}[2]{\renewcommand{\doctag}{#1}\renewcommand{\coursetag}{#2}}

%---------------------------------------------------------------- 数学记号
\newcommand{\vect}[1]{\boldsymbol{#1}}
\newcommand{\dd}{\,\mathrm{d}}
\newcommand{\ddx}{\frac{\mathrm{d}}{\mathrm{d}x}}
\addtolength{\jot}{2pt}



\usepackage{tikz}
\usetikzlibrary{decorations.pathmorphing,arrows.meta}
\usepackage{graphicx}

\settags{期中之后 · 完整讲义}{微分方程}

% ---- 各分册共用的小宏（合并时去重）----
\renewcommand{\qbreak}{\needspace{0.26\textheight}}
\newcommand{\qpagebreak}{\needspace{0.92\textheight}}
\newcommand{\mat}[1]{\begin{pmatrix}#1\end{pmatrix}}
\newcommand{\tr}{\operatorname{tr}}
\newcommand{\rank}{\operatorname{rank}}
\newcommand{\diag}{\operatorname{diag}}
\newcommand{\ii}{\mathrm{i}}
\newcommand{\ee}{\mathrm{e}}
\newcommand{\R}{\mathbb{R}}
\newcommand{\C}{\mathbb{C}}
\newcommand{\Z}{\mathbb{Z}}
\newcommand{\N}{\mathbb{N}}
\newcommand{\X}{\mathbf{X}}\newcommand{\Y}{\mathbf{Y}}
\newcommand{\U}{\mathbf{U}}
\newcommand{\F}{\mathbf{F}}
\newcommand{\Jac}{J}
\newcommand{\eps}{\varepsilon}
\newcommand{\abs}[1]{\left|#1\right|}
\newcommand{\norm}[1]{\left\lVert#1\right\rVert}
\newcommand{\cnum}[1]{\ding{\the\numexpr171+#1\relax}}

% ---- 名词英文标注：淡墨斜体，紧跟在中文名词后的括号里 ----
% 用 \textcolor 而非 {\color ...}：后者在 p{} 单元格开头会插一个竖直模式节点，
% 使该格文字整体下沉一行（xcolor + xdvipdfmx 实测）。
\newcommand{\EN}[1]{\textcolor{inkgray}{\itshape #1}}

% ---- 图块（第 6 章沿用）----
\newcounter{hwfig}
\newcommand{\figblock}[4][\relax]{%
  \par\medskip\begin{center}%
  \refstepcounter{hwfig}%
  \includegraphics[width=#3\linewidth]{figs-hw22/#2}\\[3pt]
  {\small\color{inkgray}\textbf{图\,\thehwfig}\quad#4}%
  \ifx#1\relax\else\label{c6:#1}\fi
  \end{center}\par\medskip}

% ---- 书的部件：分部页与章题 ----
\newcommand{\bigpart}[3]{%
  \clearpage
  \thispagestyle{empty}
  \par\vspace*{0.20\textheight}
  \noindent\makebox[\linewidth]{\color{petal}\rule{0.62\linewidth}{1.2pt}}
  \par\vspace{12pt}
  \noindent\makebox[\linewidth]{\color{violetdeep}\Huge\bfseries 第 #1 部}
  \par\vspace{8pt}
  \noindent\makebox[\linewidth]{\color{violetdeep}\LARGE #2}
  \par\vspace{14pt}
  \noindent\makebox[\linewidth]{\parbox{0.62\linewidth}{\centering\color{inkgray}\small #3}}
  \par\vspace{12pt}
  \noindent\makebox[\linewidth]{\color{petal}\rule{0.62\linewidth}{1.2pt}}
  \clearpage}

\newcommand{\chapterbanner}[4]{%
  \clearpage
  \settags{第 #1 章\quad #2}{微分方程 · 期中之后}
  \par\vspace*{0.9cm}
  \noindent{\color{petal}\small 第 #1 章}
  \par\vspace{2pt}
  \noindent{\color{violetdeep}\LARGE\bfseries #2}
  \par\vspace{3pt}
  \noindent{\color{inkgray}\large #3}
  \par\vspace{6pt}
  \noindent{\color{petal}\rule{\linewidth}{1.2pt}}
  \par\vspace{4pt}
  \noindent{\small\color{inkgray}#4}\par\vspace{8pt}}

\begin{document}

\doctitle{微分方程 · 期中之后}%%
{L13--L23 完整讲义}%%
{依据：本课程 8 份课件，共 334 页（Lecture 13--23，2026-09-24 至 2026-11-03）\\
范围：期中考试（2026-10-08）之后的全部内容；此前的 L1--L12 不在本书内\\
全书由六份分册讲义合并而成，正文推导与各章附录详解均保留原稿}


\callout{\textbf{先读这一段。}期中之后老师换了三个题目，但其实是同一件事的三层：先问“多个未知函数怎么一起解”（L13--L18），再问“系数不是常数怎么办”（L18--L20），最后问“根本解不出来时看什么”（L21--L23）。这本书按这三层分成三部；每一章都可以独立读，但第 1--4 章是一条链，后面几章直接踩在前面几章上。}

\hdr{§0}{全书导览}

\subhdr{0.1 主线：十一讲在回答同一个问题}

期中之前解的是\textbf{一个}未知函数的高阶方程。从 L13 起，问题变成\textbf{一组}未知函数的一阶方程：

\[
\vect{Y}'=A\vect{Y}+\vect{F}(x),\qquad \vect{Y}\in\mathbb{R}^{n}.
\]

整段内容的走法，就是把这个问题一层层剥开。下面这张图是全书骨架，后面每一章都在解释其中一格。

\begin{center}
\renewcommand{\arraystretch}{1.02}
\begin{tabular}{@{}c@{}}
\fcolorbox{rulepurple}{bandviolet}{\parbox{0.86\linewidth}{\centering
  \textbf{\ding{172}\;一个高阶方程}\quad $y^{(n)}+a_1y^{(n-1)}+\cdots+a_ny=f(x)$}}\\[1pt]
{\color{petal}$\Downarrow$}\enspace{\small\color{inkgray}第 1 章：用一个变换，把它拆成一阶组}\\[8pt]
\fcolorbox{rulepurple}{bandviolet}{\parbox{0.86\linewidth}{\centering
  \textbf{\ding{173}\;一阶线性方程组}\quad $\vect{Y}'=A\vect{Y}+\vect{F}(x)$}}\\[1pt]
{\color{petal}$\Downarrow$}\enspace{\small\color{inkgray}第 2 章：$A$ 是常矩阵、$\vect{F}=0$ 时，用特征值拼出解}\\[8pt]
\fcolorbox{rulepurple}{bandviolet}{\parbox{0.86\linewidth}{\centering
  \textbf{\ding{174}\;齐次常系数}\quad 三种特征值情形，各给一套基本积木}}\\[1pt]
{\color{petal}$\Downarrow$}\enspace{\small\color{inkgray}第 3 章：三种情形合成为一个对象 $\mathrm{e}^{Ax}$}\\[8pt]
\fcolorbox{rulepurple}{bandviolet}{\parbox{0.86\linewidth}{\centering
  \textbf{\ding{175}\;统一写法}\quad $\vect{Y}(x)=\mathrm{e}^{Ax}\vect{Y}_0$，三种情形都收进一个矩阵}}\\[1pt]
{\color{petal}$\Downarrow$}\enspace{\small\color{inkgray}第 4 章：$\vect{F}\neq0$ 时，再补一个特解}\\[8pt]
\fcolorbox{rulepurple}{bandviolet}{\parbox{0.86\linewidth}{\centering
  \textbf{\ding{176}\;有外力}\quad $\vect{Y}=\vect{Y}_c+\vect{Y}_p$，两种求 $\vect{Y}_p$ 的办法}}\\[1pt]
{\color{petal}$\Downarrow$}\enspace{\small\color{inkgray}第 5 章：$A$ 不再是常矩阵，怎么办}\\[8pt]
\fcolorbox{rulepurple}{bandviolet}{\parbox{0.86\linewidth}{\centering
  \textbf{\ding{177}\;变系数}\quad 常点用幂级数，正则奇点用 Frobenius 级数}}\\[1pt]
{\color{petal}$\Downarrow$}\enspace{\small\color{inkgray}第 6 章：方程非线性、解不出来，怎么办}\\[8pt]
\fcolorbox{rulepurple}{bandviolet}{\parbox{0.86\linewidth}{\centering
  \textbf{\ding{178}\;非线性}\quad 不解方程，画相图定性；在临界点附近借用线性}}\\
\end{tabular}
\end{center}

读这张图要抓住一处：\textbf{\ding{173}\ding{174}\ding{175}\ding{176} 是一条竖线}（同一个问题的四种情形），\textbf{\ding{177}\ding{178} 是两次换问题}——一次换“系数”，一次换“能不能解出来”。所以复习时不要把六章平铺，先走完竖线，再看两次换问题。

\subhdr{0.2 三部六章}

\begin{minipage}{\linewidth}
\tabtop
\begin{tabular}{@{}p{0.06\linewidth}p{0.07\linewidth}p{0.24\linewidth}>{\raggedleft\arraybackslash}p{0.11\linewidth}p{0.30\linewidth}@{}}
\thead{部} & \thead{章} & \thead{内容} & \thead{课件页} & \thead{这一部在换什么}\\
\end{tabular}
\tabmid
\begin{tabular}{@{}p{0.06\linewidth}p{0.07\linewidth}p{0.24\linewidth}>{\raggedleft\arraybackslash}p{0.11\linewidth}p{0.30\linewidth}@{}}
II & 第 1 章 & 方程组入门（L13） & 37 & 把高阶方程换成一阶组：换的是\textbf{语言}\\
II & 第 2 章 & 常系数齐次组（L14--15） & 64 & 特征值三情形：换的是\textbf{解的形态}\\
II & 第 3 章 & 矩阵指数（L16--17） & 44 & $\mathrm{e}^{Ax}$：换的是\textbf{写法}\\
II & 第 4 章 & 非齐次线性组（L17--18） & 49 & 补一个特解：换的是\textbf{有外力}\\
III & 第 5 章 & 幂级数解法（L18--20） & 79 & 系数不是常数：换的是\textbf{假设}\\
IV & 第 6 章 & 临界点与稳定性（L21--23） & 61 & 解不出来：换的是\textbf{目标}\\
\hline
 & & 合计 & 334 & \\
\end{tabular}
\tabbot
\end{minipage}

第 II 部（194 页）占全段的六成，是主战场；第 III、IV 部各是一次“换问题”。

\subhdr{0.3 依赖链：哪一章塌了，会连累哪一章}

\begin{minipage}{\linewidth}
\tabtop
\begin{tabular}{@{}p{0.09\linewidth}p{0.28\linewidth}p{0.52\linewidth}@{}}
\thead{这一章} & \thead{供给什么} & \thead{不熟会卡在哪}\\
\end{tabular}
\tabmid
\begin{tabular}{@{}p{0.09\linewidth}p{0.28\linewidth}p{0.52\linewidth}@{}}
第 1 章 & 解的结构、叠加原理 & 第 4 章的 $\vect{Y}=\vect{Y}_c+\vect{Y}_p$ 变成口号；说不出为什么齐次解要两个\\
第 2 章 & 特征值三情形 & 第 3 章“Case 1 对角化”看不懂条件从哪来；第 6 章的类型判据只能背表\\
第 3 章 & $\mathrm{e}^{Ax}$ 与 $\Phi$ & 第 4 章参数变易法没有可用的积分核；非齐次题只能做“好猜”的那类\\
第 5 章 & 幂级数递推 & Frobenius 连第一步都下不了手（指标方程是从递推的常数项读出来的）\\
第 6 章 & 相图与稳定性定义 & 线性化定理没有几何直觉，只能记结论\\
\end{tabular}
\tabbot
\end{minipage}

\subhdr{0.4 怎么读这本书}

\begin{enumerate}
  \item \textbf{每章自成一套节号}：全书的节号写作 §0、§1……，它们\textbf{只在章内有效}（第 3 章的 §4 与第 4 章的 §4 不是同一节）。跨章引用一律写明“第 $k$ 章 §$n$”。
  \item \textbf{正文不给答案}：各章正文只给题面与考点，例题与作业的详解统一在\textbf{附录 A}（连续编号 A.1--A.41，按章排列）。正文里的“详解：附录 A.$n$，第 $m$ 页”可以直接跳。
  \item \textbf{三条路线}：想通读，按第 1$\to$2$\to$3$\to$4$\to$5$\to$6 章；想速通，只抓第 2、4、6 章（约 5 小时）；想查工具，直接翻各章的“公式/判据”小节与附录 C。
  \item \textbf{每章末尾都有四件东西}：易错点、缺口清单（课件哪里跳了步、本章补了什么）、三句话总结、与课件的页码对照。复习时优先看后两件。
\end{enumerate}

\warnbox{\textbf{一条编号上的坑。}课件文件名里 \texttt{17\_18} 是“非齐次线性方程组”，而 \texttt{18\_power-series} 是“幂级数”，两份都标了 10 月 20 日——L18 这天老师同时收尾方程组、开启幂级数。\textbf{按标题认，不要按文件名里的数字认}：本书把 \texttt{17\_18} 归入第 4 章，把 \texttt{18} 归入第 5 章。}

\subhdr{0.5 中英术语对照}

课件是英文的，讲义是中文的，两边名词先对上，读课件才不会卡在字面上。

\begin{minipage}{\linewidth}
\tabtop
\begin{tabular}{@{}p{0.25\linewidth}p{0.17\linewidth}p{0.26\linewidth}p{0.17\linewidth}@{}}
\thead{课件写法} & \thead{中文} & \thead{课件写法} & \thead{中文}\\
\end{tabular}
\tabmid
\begin{tabular}{@{}p{0.25\linewidth}p{0.17\linewidth}p{0.26\linewidth}p{0.17\linewidth}@{}}
system of ODEs & 微分方程组 & fundamental matrix & 基本解矩阵\\
linear system & 线性系统 & matrix exponential & 矩阵指数\\
homogeneous & 齐次的 & Cayley--Hamilton & 凯莱--哈密顿\\
nonhomogeneous & 非齐次的 & variation of parameters & 参数变易\\
fundamental set & 基本解组 & undetermined coefficients & 待定系数\\
Wronskian & 朗斯基行列式 & ordinary point & 常点\\
regular singular point & 正则奇点 & indicial equation & 指标方程\\
Frobenius method & 弗罗贝尼乌斯法 & Bessel equation & 贝塞尔方程\\
phase portrait & 相图 & critical point & 临界点\\
saddle & 鞍点 & node & 结点\\
focus & 焦点 & center & 中心\\
hyperbolic & 双曲的 & first approximation & 第一近似\\
asymptotically stable & 渐近稳定 & linearization & 线性化\\
\end{tabular}
\tabbot
\end{minipage}

这张表是名词台账，不承载推导，逐格可回课件核对。少数术语各教材译法不同（focus 也译“螺旋点”、node 也译“节点”），本书统一用上表的译名。

\bigpart{II}{线性方程组 · L13--L18}{四章一条竖线：同一个 $\vect{Y}^{\prime}=A\vect{Y}+\vect{F}$，从“怎么建”讲到“怎么解”。课件 194 页，占全段六成，是主战场。}
\chapterbanner{1}{方程组入门}{从高阶方程到一阶组}{课件 \texttt{13\_DE\_systems.pdf}（37 页，2026-09-24）}
\callout{\textbf{先读这一段。}期中之前你解的是“一个未知函数的高阶方程”，从这一讲起变成“一组未知函数的一阶方程”。这一讲不教怎么解方程，只做三件事：\textbf{说清为什么需要组}（§1 三个物理模型）、\textbf{把“解”这件事重新定义一遍}（§2--§5）、\textbf{给出两条互相逆行的通道}（§6：高阶方程怎么化成组；§7：常系数组怎么解——只开个头，留给 L14--15）。这一讲是整章的“语法课”，后面四讲都在用这里的词。}

%=====================================================================
\hdr{§0}{这一讲在整门课里的位置}

期中考试（2026-10-08）之前，L1--L12 处理的是\textbf{一个}未知函数 $y(x)$ 的 $n$ 阶方程。L13 起，未知量变成一个向量

\[
\vect{Y}(x)=\bigl(y_1(x),\dots,y_n(x)\bigr)^{\mathsf T},
\]

方程变成 $\vect{Y}'=A(x)\vect{Y}+\vect{F}(x)$。这一转变不是换记号，而是换了问题的结构：高阶方程要配 $n$ 个初值，一阶组也要配 $n$ 个，两边恰好对应——§6 就把这件事写成一步变换。

本讲课件 37 页，按内容分五块：

\begin{minipage}{\linewidth}
\tabtop
\begin{tabular}{@{}p{0.10\linewidth}p{0.18\linewidth}p{0.56\linewidth}@{}}
\thead{页} & \thead{板块} & \thead{这一块在干什么}\\
\end{tabular}
\tabmid
\begin{tabular}{@{}p{0.10\linewidth}p{0.18\linewidth}p{0.56\linewidth}@{}}
p3--p6 & 物理来源 & 三个模型：两质量三弹簧、并联 LRC、弹簧-质量。说明为什么会有“组”\\
p7--p10 & 定义与例子 & 系统、正规形、解、通解/特解/奇解；例 1、例 2\\
p11--p18 & 线性系统 & 矩阵写法；存在唯一性；线性算子；叠加原理；例 3、例 4\\
p19--p23 & 解的结构 & 线性相关、Wronskian、基本解组；定理 2.3--2.5；例 5\\
p24--p31 & 变换 & $n$ 阶方程化成一阶组的五步；例 6、例 7\\
p32--p36 & 常系数入口 & 试探解 $\vect{V}\mathrm{e}^{\lambda x}$、特征方程、代数重数\\
p37 & 作业 & 一道题（化 $y'''=y''-x^2(y')^2$ 为一阶组）\\
\end{tabular}
\tabbot
\end{minipage}

这张表是页码台账，不承载推导；逐格都能在课件里对上。全讲的主线只有一句：\textbf{先承认“组”是自然的，再把“解方程组”翻译成“向量空间里挑 $n$ 个基”}。

\subhdr{0.1 记号}

本讲义与课件记号一致，只有一处需要提前声明：

\begin{itemize}
  \item $\vect{Y}$ 是\textbf{列向量} $(y_1,\dots,y_n)^{\mathsf T}$；$\vect{Y}_k$ 表示第 $k$ 个解向量，不是第 $k$ 个分量。
  \item $A(x)$ 是 $n\times n$ 系数矩阵，$a_{ij}(x)$ 是第 $i$ 个方程里 $y_j$ 的系数（\textbf{先行后列}）。
  \item $\mathcal{L}[\vect{Y}]=\vect{Y}'-A\vect{Y}$ 是线性算子，$\mathcal{L}[\vect{Y}]=\vect{F}$ 即原方程。
  \item $W(\vect{Y}_1,\dots,\vect{Y}_n)$ 是由这 $n$ 个\textbf{列}向量拼成的行列式，见 §5.2。
\end{itemize}

%=====================================================================
\hdr{§1}{三个模型：为什么需要“组”}

课件用三页各给一个物理系统（p4、p5、p6），目的不是教你建模，而是让你接受一件事：\textbf{真实系统从来不止一个未知函数}。两质量系统有两个位移，LRC 电路有电压和电流两个状态量，弹簧-质量只有一个位移但它是一阶导也出现的二阶方程——三种情形最后都写成同一个形状。

\subhdr{1.1 两质量三弹簧（p4）}

两个质量 $m_1,m_2$ 在光滑水平面上，三根弹簧常数 $k_1,k_2,k_3$：左墙—$k_1$—$m_1$—$k_2$—$m_2$—$k_3$—右墙。设 $x_1,x_2$ 是两质量相对各自平衡位置的位移，\textbf{向右为正}；外力 $F_1(t),F_2(t)$ 也向右为正。

\begin{center}
\begin{tikzpicture}[font=\small,>=Stealth]
  \draw[gray!55] (-0.2,-1.42) -- (7.9,-1.42);
  \foreach \x in {0.1,0.7,...,7.7} {\draw[gray!55] (\x,-1.42) -- (\x-0.16,-1.64);}
  \node[gray!70,anchor=west] at (7.95,-1.55) {\footnotesize 光滑水平面};
  \draw[thick] (0,0.55) -- (0,-1.42);
  \foreach \y in {0.4,0.15,...,-1.30} {\draw[gray!45] (0,\y) -- (-0.22,\y-0.18);}
  \draw[decorate,decoration={zigzag,segment length=5pt,amplitude=3pt}] (0,-0.42) -- (1.7,-0.42);
  \node at (0.85,0.58) {$k_1$};
  \draw[thick,fill=bandviolet] (1.7,-0.95) rectangle (2.72,-0.05);
  \node at (2.21,-0.52) {$m_1$};
  \draw[decorate,decoration={zigzag,segment length=5pt,amplitude=3pt}] (2.72,-0.42) -- (4.62,-0.42);
  \node at (3.67,0.58) {$k_2$};
  \draw[thick,fill=bandviolet] (4.62,-0.95) rectangle (5.64,-0.05);
  \node at (5.13,-0.52) {$m_2$};
  \draw[decorate,decoration={zigzag,segment length=5pt,amplitude=3pt}] (5.64,-0.42) -- (7.5,-0.42);
  \node at (6.57,0.58) {$k_3$};
  \draw[thick] (7.5,0.55) -- (7.5,-1.42);
  \foreach \y in {0.4,0.15,...,-1.30} {\draw[gray!45] (7.5,\y) -- (7.72,\y-0.18);}
  \draw[->] (2.21,-1.16) -- (2.86,-1.16) node[right] {\footnotesize $x_1$};
  \draw[->] (5.13,-1.16) -- (5.78,-1.16) node[right] {\footnotesize $x_2$};
  \draw[->,violet] (2.73,-0.52) -- (3.21,-0.52) node[right,black] {\footnotesize $F_1(t)$};
  \draw[->,violet] (5.65,-0.52) -- (6.13,-0.52) node[right,black] {\footnotesize $F_2(t)$};
\end{tikzpicture}\\[2pt]
{\small\color{inkgray}图 1\quad 两质量三弹簧系统（按课件 p4 重绘）：三根弹簧的伸长量分别是 $x_1$、$x_2-x_1$、$x_2$。}
\end{center}

\textbf{受力怎么逐项读出来}（课件 p4 直接给了结果，没写这一步）。三根弹簧的“伸长量”分别是：

\begin{itemize}
  \item 左弹簧：$m_1$ 的位移就是它的伸长量，$x_1$；它把 $m_1$ 往左拉，力为 $-k_1x_1$。
  \item 中弹簧：两端分别是 $m_1$ 与 $m_2$，伸长量是 $x_2-x_1$；它拉 $m_1$ 向右、拉 $m_2$ 向左，对 $m_1$ 的力是 $+k_2(x_2-x_1)$，对 $m_2$ 是 $-k_2(x_2-x_1)$。
  \item 右弹簧：$m_2$ 的位移是它的伸长量，$x_2$；把 $m_2$ 往左拉，力为 $-k_3x_2$。
\end{itemize}

于是

\[
m_1\ddot x_1=-(k_1+k_2)x_1+k_2x_2+F_1(t),\qquad
m_2\ddot x_2=k_2x_1-(k_2+k_3)x_2+F_2(t),
\]

与课件 p4 的两个式子逐项相同。写成矩阵就是

\[
\mat{\ddot x_1\\[2pt]\ddot x_2}
=\mat{-(k_1+k_2)/m_1 & k_2/m_1\\[3pt] k_2/m_2 & -(k_2+k_3)/m_2}
\mat{x_1\\[2pt]x_2}+\mat{F_1/m_1\\[2pt]F_2/m_2}.
\]

\textbf{注意}：这还是二阶的组（两个二阶方程），不是课件最后要的一阶组。要变成一阶组，还要再降一次阶——这正是 §1.3 与 §6 要做的事。

\subhdr{1.2 并联 LRC 电路（p5）}

$R,L,C$ 三者并联。设 $V$ 是电容两端电压，$I$ 是流过电感的电流。课件直接给出

\[
\frac{\mathrm{d}I}{\mathrm{d}t}=\frac{V}{L},\qquad
\frac{\mathrm{d}V}{\mathrm{d}t}=-\frac{I}{C}-\frac{V}{RC},
\]

但没写这两步怎么来。补上：并联意味着三者两端\textbf{电压相同}，都是 $V$；元件的本构关系是

\[
V=L\frac{\mathrm{d}I}{\mathrm{d}t}\ \ (\text{电感}),\qquad
I_C=C\frac{\mathrm{d}V}{\mathrm{d}t}\ \ (\text{电容}),\qquad
I_R=\frac{V}{R}\ \ (\text{电阻}).
\]

在三个支路汇合的节点上用基尔霍夫电流定律：无外接电流时，流出节点之和为零，

\[
C\frac{\mathrm{d}V}{\mathrm{d}t}+\frac{V}{R}+I=0
\ \Longrightarrow\
\frac{\mathrm{d}V}{\mathrm{d}t}=-\frac{I}{C}-\frac{V}{RC},
\]

第一式由 $V=L\,\mathrm{d}I/\mathrm{d}t$ 直接除 $L$ 得到。两个式子合起来是\textbf{两个一阶方程}，未知函数是 $I$ 和 $V$——天然就是正规形。

\begin{center}
\begin{tikzpicture}[font=\small,>=Stealth]
  \draw (0,0) -- (6,0);
  \draw (0,-2.7) -- (6,-2.7);
  \draw (1.2,0) -- (1.2,-0.75);
  \draw (0.85,-0.75) rectangle (1.55,-1.35);
  \node at (1.2,-1.05) {$R$};
  \draw (1.2,-1.35) -- (1.2,-2.7);
  \draw (3,0) -- (3,-0.75);
  \draw[decorate,decoration={coil,aspect=0.55,segment length=5pt,amplitude=3.6pt}] (3,-0.75) -- (3,-1.55);
  \node at (3.42,-1.15) {$L$};
  \draw (3,-1.55) -- (3,-2.7);
  \draw[->] (3,-1.95) -- (3,-2.28);
  \node at (3.3,-2.12) {\footnotesize $I$};
  \draw (4.8,0) -- (4.8,-1.05);
  \draw (4.45,-1.05) -- (5.15,-1.05);
  \draw (4.45,-1.32) -- (5.15,-1.32);
  \node at (5.38,-1.18) {$C$};
  \draw (4.8,-1.32) -- (4.8,-2.7);
  \draw[<->] (0.42,0) -- (0.42,-2.7) node[midway,left] {$V$};
  \node at (0.64,0.24) {\footnotesize $+$};
  \node at (0.64,-2.94) {\footnotesize $-$};
  \fill (0,0) circle (1.6pt); \fill (0,-2.7) circle (1.6pt);
  \fill (6,0) circle (1.6pt); \fill (6,-2.7) circle (1.6pt);
\end{tikzpicture}\\[2pt]
{\small\color{inkgray}图 2\quad 并联 LRC（按课件 p5 重绘）：三条支路电压同为 $V$，$I$ 由 $L$ 支路自上而下。}
\end{center}

\textbf{符号约定先定死}：$I$ 向下流过电感为正，$V$ 上端相对下端为正。换一种约定，两个方程的符号会同时变，但“一阶组”这个结构不变。

\subhdr{1.3 弹簧-质量（p6）}

单个质量，阻尼 $\gamma$，弹簧常数 $k$，外力 $F(t)$，$u$ 为位移（平衡位置为原点，向下为正）：

\[
m u''+\gamma u'+k u=F(t).
\]

课件写“can be transformed into a system of first-order differential equations”，然后直接给出 $x_1=u,\ x_2=u'$ 与结果。补上中间这一步：由 $x_1=u$ 得 $x_1'=u'=x_2$；由 $x_2=u'$ 得 $x_2'=u''$，而原方程解出 $u''=\bigl(F-\gamma u'-ku\bigr)/m$，代入 $u=x_1,u'=x_2$：

\[
x_1'=x_2,\qquad
x_2'=-\frac{k}{m}x_1-\frac{\gamma}{m}x_2+\frac{F(t)}{m},
\]

与课件 p6 一致。\textbf{这一步就是 §6 的一般变换在 $n=2$ 时的样子}，只是这里手工做了一遍。

\subhdr{1.4 三个模型的共同点}

三个系统的物理内容完全不同，落到纸上却是同一个形状：\textbf{若干未知函数的一阶导数，被写成这些未知函数的线性组合加一个已知项}。这就是本讲的主对象：

\[
\vect{Y}'=A\vect{Y}+\vect{F}(t).
\]

后面 §2--§5 建立“解”的语言，§6 说明任何 $n$ 阶方程都能化到这个形状，§7 开始处理 $A$ 为常矩阵的情形。

%=====================================================================
\hdr{§2}{系统、正规形、解}

\subhdr{2.1 一般系统与正规形（p7）}

课件给两个定义。第一个是\textbf{一般系统}：$n$ 个一阶方程，每个都写成隐式

\[
F_k(x,y_1,\dots,y_n,y_1',\dots,y_n')=0,\qquad k=1,\dots,n .
\]

第二个是\textbf{正规形}（normal system）：每个方程都已经解出对应的导数

\[
y_k'=f_k(x,y_1,\dots,y_n),\qquad k=1,\dots,n ,
\]

再配初值 $y_k(x_0)=y_k^0$。

\textbf{这里有一个课件没说的前提}：从一般形写到正规形，要能\textbf{解出} $y_k'$。这件事不是永远做得到的——例如 $(y_1')^2+y_1-1=0$ 就只能写成两个分支 $y_1'=\pm\sqrt{1-y_1}$。本课程后面处理的都是正规形，所以这个前提实际上被默认了；但你写答案时如果遇到“化不成正规形”的情形，要说明一句。

\subhdr{2.2 解与初值问题（p8）}

一组函数 $y_k=\varphi_k(x)$ 在区间 $I$ 上称为解，如果把它代进系统后\textbf{在 $I$ 上每一点}都成立。几何上，解是 $n$ 维空间里的一条曲线，初值决定这条曲线的起点；系统加初值就是\textbf{初值问题}（IVP）。

\textbf{“在区间上每点成立”不是废话}。例 2 的通解 $y_1=(x+c_1)^2$ 只在 $x+c_1\ge0$ 时满足 $y_1'=2\sqrt{y_1}$：因为 $\sqrt{(x+c_1)^2}=|x+c_1|$，当 $x+c_1<0$ 时右边是 $-2(x+c_1)$，符号反了。所以这道题的解只能定义在半区间上——这是非线性方程特有的麻烦，线性系统不会出现（§4.1 的 Remark 2.1）。

\subhdr{2.3 通解、特解、奇解（p9）}

课件的定义是：含 $n$ 个任意常数 $c_1,\dots,c_n$ 的一组函数，如果对区域 $D$ 内每一点都能取到，就称通解；取定常数值得到特解；不能由通解取常数得到的解叫奇解。

课件原文在“for every point $(x,y_1,\dots,y_n)\in D$”这一句被排版截断了，读起来不知道条件是什么。这里给一个\textbf{可操作的判据}：把初值 $(x_0,\vect{y}^0)$ 代进通解，得到关于 $c_1,\dots,c_n$ 的 $n$ 个方程；若这个方程组对 $D$ 内每一点都有唯一解，则这组函数是通解。换句话说，\textbf{常数个数够、而且能反解出初值}，两条都满足才算通解。

\textbf{“$n$ 个常数”不能多也不能少}：少了覆盖不了所有初值，多了则常数之间不独立、本质上还是 $n$ 个。例 1 的 $c_1\cos(x-c_2)$ 与 $A\cos x+B\sin x$ 都是两个常数，等价，见附录 A.1。

\subhdr{2.4 例 1 与例 2（p10）}

两道题都在这页，题干与考点如下。

\qbreak
\qstem{（课件 p10 例 1）$\dfrac{\mathrm{d}y_1}{\mathrm{d}x}=y_2,\quad \dfrac{\mathrm{d}y_2}{\mathrm{d}x}=-y_1$。课件给通解 $y_1=c_1\cos(x-c_2)$，$y_2=-c_1\sin(x-c_2)$。\textbf{验证它是解，并说明常数怎么写都行。}}
\par\textbf{考点}：解的定义（代回恒成立）；通解里常数的等价写法。
\par\textbf{关键一步}：把两个式子分别求导，代回右端；再令 $A=c_1\cos c_2,\ B=c_1\sin c_2$ 看两种写法是否互推。
\par\textbf{详解}：附录 A.1，第 \pageref{c1:app:a1} 页。

\qbreak
\qstem{（课件 p10 例 2）$\dfrac{\mathrm{d}y_1}{\mathrm{d}x}=2\sqrt{y_1},\quad \dfrac{\mathrm{d}y_2}{\mathrm{d}x}=y_1$。课件给通解 $y_1=(x+c_1)^2$，$y_2=\tfrac13(x+c_1)^3+c_2$，以及奇解 $y_1=0,\ y_2=c$。\textbf{验证通解与奇解。}}
\par\textbf{考点}：非线性方程的定义域限制；“奇解”为什么真的是新的解。
\par\textbf{关键一步}：先解只含 $y_1$ 的第一式；再检查 $x+c_1<0$ 时 $\sqrt{y_1}$ 的符号；最后看 $y_1\equiv0$ 能否由通解取常数得到。
\par\textbf{详解}：附录 A.2，第 \pageref{c1:app:a2} 页。

%=====================================================================
\hdr{§3}{线性系统与矩阵写法}

\subhdr{3.1 一般线性系统（p12）}

线性系统的正规形是

\[
y_k'=\sum_{j=1}^{n}a_{kj}(x)y_j+f_k(x),\qquad k=1,\dots,n,
\]

即每个方程的右端对未知函数是\textbf{一次齐次}的（可以带 $x$ 的函数系数），外加一个只依赖 $x$ 的项。矩阵形式是

\[
\vect{Y}'(x)=A(x)\vect{Y}+\vect{F}(x),\qquad
A=(a_{kj}),\quad \vect{Y}=(y_1,\dots,y_n)^{\mathsf T},\quad \vect{F}=(f_1,\dots,f_n)^{\mathsf T}.
\]

$\vect{F}\equiv\vect{0}$ 时称齐次，否则非齐次。

\subhdr{3.2 课件的一处排版笔误（p13）}

p13 的矩阵 $A$ 第二行最后一个元素印成了 $a_{1n}$，按定义应为 $a_{2n}$。这不影响理解，但抄公式时容易照抄错。正确的写法是

\[
A=\mat{
a_{11} & a_{12} & \cdots & a_{1n}\\
a_{21} & a_{22} & \cdots & a_{2n}\\
\vdots & \vdots & & \vdots\\
a_{n1} & a_{n2} & \cdots & a_{nn}}.
\]

\textbf{记法只有一条}：$a_{kj}$ 是第 $k$ 个方程里 $y_j$ 的系数，先行后列。写矩阵时逐行抄系数，比“先写列”不容易错。

\subhdr{3.3 例 3（p14）}

\qbreak
\qstem{（课件 p14 例 3）设 $x(t),y(t),z(t)$ 未知：$x'=y-z,\ y'=x+y,\ z'=x+z$。\textbf{写成矩阵形式，并判断是否齐次。}}
\par\textbf{考点}：逐行读系数；齐次的判断。
\par\textbf{关键一步}：第一行里 $x$ 的系数是 $0$（没出现），$y$ 是 $1$，$z$ 是 $-1$；按 $x,y,z$ 顺序排成一行。
\par\textbf{详解}：附录 A.3，第 \pageref{c1:app:a3} 页。

%=====================================================================
\hdr{§4}{存在唯一性、算子、叠加原理}

\subhdr{4.1 定理 2.1（p15）}

设 $a_{kj}(x),f_k(x)$ 在 $[a,b]$ 上连续，则对任意 $x_0\in[a,b]$ 与任意初值 $\vect{y}^0$，初值问题

\[
y_k'=\sum_{j=1}^{n}a_{kj}(x)y_j+f_k(x),\qquad y_k(x_0)=y_k^0
\]

在 $x_0$ 的某个邻域内\textbf{存在且唯一}解。

课件紧接着给了一条 Remark：\textbf{线性系统没有奇解}。这句话的理由课件没写，补在这里：只要系数连续，解可以沿区间一路延拓，不会在中途“断掉”或“分叉”；既然每个初值都唯一对应一条解，就不可能存在“通解之外另有一条解”的情形。这与例 2 的非线性方程形成对照：那里 $y_1\equiv0$ 是一条真实的、通解取不到的轨道。

\subhdr{4.2 线性算子与性质 2.1（p16）}

定义算子

\[
\mathcal{L}[\vect{Y}]=\frac{\mathrm{d}\vect{Y}}{\mathrm{d}x}-A\vect{Y}.
\]

齐次方程就是 $\mathcal{L}[\vect{Y}]=\vect{0}$，非齐次是 $\mathcal{L}[\vect{Y}]=\vect{F}$。性质 2.1 说 $\mathcal{L}$ 是线性的：

\[
\mathcal{L}[c\vect{Y}]=c\,\mathcal{L}[\vect{Y}],\qquad
\mathcal{L}[\vect{Y}_1+\vect{Y}_2]=\mathcal{L}[\vect{Y}_1]+\mathcal{L}[\vect{Y}_2].
\]

由此立刻得到：齐次解的任意线性组合仍是齐次解。

最后一条要单独说：\textbf{若 $A$ 是实矩阵，而复值解 $\vect{Y}=\vect{U}+\mathrm{i}\vect{V}$ 满足 $\mathcal{L}[\vect{Y}]=\vect{0}$，则实部 $\vect{U}$ 与虚部 $\vect{V}$ 也都是解}。理由是 $\mathcal{L}$ 只做求导与乘实矩阵，实部虚部互不干扰：

\[
\vect{0}=\mathcal{L}[\vect{U}+\mathrm{i}\vect{V}]=\mathcal{L}[\vect{U}]+\mathrm{i}\,\mathcal{L}[\vect{V}]
\ \Longrightarrow\
\mathcal{L}[\vect{U}]=\vect{0},\quad \mathcal{L}[\vect{V}]=\vect{0}.
\]

这一条在 L14--15 处理共轭复特征值时会直接用到：先得复值解 $\vect{V}\mathrm{e}^{(\alpha+\mathrm{i}\beta)x}$，再取实部虚部得到两个实解。

\subhdr{4.3 叠加原理与定理 2.2（p17）}

叠加原理：若 $\vect{Y}_1$ 是 $\mathcal{L}[\vect{Y}]=\vect{F}$ 的解、$\vect{Y}_0$ 是齐次解，则 $\vect{Y}_0+\vect{Y}_1$ 仍是 $\mathcal{L}[\vect{Y}]=\vect{F}$ 的解。一行验证：

\[
\mathcal{L}[\vect{Y}_0+\vect{Y}_1]=\mathcal{L}[\vect{Y}_0]+\mathcal{L}[\vect{Y}_1]=\vect{0}+\vect{F}=\vect{F}.
\]

定理 2.2（课件只给结论，这里补证明）：非齐次方程的通解是

\[
\vect{Y}(x)=\vect{Z}(x)+\tilde{\vect{Y}}(x),
\]

其中 $\vect{Z}$ 是齐次通解、$\tilde{\vect{Y}}$ 是非齐次的一个特解。

\textbf{证明}：设 $\vect{Y}$ 是任一解。令 $\vect{W}=\vect{Y}-\tilde{\vect{Y}}$，由线性性

\[
\mathcal{L}[\vect{W}]=\mathcal{L}[\vect{Y}]-\mathcal{L}[\tilde{\vect{Y}}]=\vect{F}-\vect{F}=\vect{0},
\]

所以 $\vect{W}$ 是齐次解，即 $\vect{W}=\vect{Z}$，于是 $\vect{Y}=\vect{Z}+\tilde{\vect{Y}}$。反方向由叠加原理直接得到。$\square$

这个证明只有一个动作：\textbf{两个解相减消掉右端}。后面 L17--18 求非齐次特解时，反复用的就是这一条。

性质 2.2（右端可拆）：若 $\vect{Y}_i$ 是 $\mathcal{L}[\vect{Y}]=\vect{F}_i$ 的解，则 $\sum_i\vect{Y}_i$ 是 $\mathcal{L}[\vect{Y}]=\sum_i\vect{F}_i$ 的解。它说明非齐次项可以一项一项地攻。

\subhdr{4.4 例 4（p18）}

\qbreak
\qstem{（课件 p18 例 4）$\vect{Y}'=\mat{1&1\\4&1}\vect{Y}$。验证 $\vect{Y}_1=\mat{1\\2}\mathrm{e}^{3x}$、$\vect{Y}_2=\mat{1\\-2}\mathrm{e}^{-x}$ 都是解，从而 $c_1\vect{Y}_1+c_2\vect{Y}_2$ 是解。}
\par\textbf{考点}：解的定义；叠加原理。
\par\textbf{关键一步}：算 $A\vect{Y}_1$，看它是否等于 $3\vect{Y}_1$（因为 $\vect{Y}_1'=3\vect{Y}_1$）。
\par\textbf{详解}：附录 A.4，第 \pageref{c1:app:a4} 页。

%=====================================================================
\hdr{§5}{线性相关、Wronskian、基本解组}

这一节回答一个问题：\textbf{齐次方程要几个解才够？}答案是 $n$ 个，而且这 $n$ 个必须线性无关。判据是 Wronskian。

\subhdr{5.1 线性相关与线性无关（p19）}

$n$ 个向量函数 $\vect{Y}_1,\dots,\vect{Y}_n$ 在 $[a,b]$ 上线性相关，指存在不全为零的常数 $c_1,\dots,c_n$ 使

\[
c_1\vect{Y}_1+\cdots+c_n\vect{Y}_n=\vect{0},\qquad x\in[a,b].
\]

注意是\textbf{在整个区间上恒为零}，不是某一点为零。否则任何两个在某点碰头的解都会被判成相关。

\subhdr{5.2 Wronskian（p20）}

把 $n$ 个列向量排成一个矩阵，取行列式：

\[
W(\vect{Y}_1,\dots,\vect{Y}_n)(x)=\det\begin{pmatrix}
y_{11}(x) & y_{21}(x) & \cdots & y_{n1}(x)\\
y_{12}(x) & y_{22}(x) & \cdots & y_{n2}(x)\\
\vdots & \vdots & & \vdots\\
y_{1n}(x) & y_{2n}(x) & \cdots & y_{nn}(x)
\end{pmatrix},
\]

其中 $\vect{Y}_k=(y_{k1},\dots,y_{kn})^{\mathsf T}$，\textbf{按列排}。课件 p20 把同一个行列式先按列写、又按行写，两种写法互为转置，数值相同（行列式转置不变），但下标容易看错。本讲义统一按上面的写法：\textbf{第 $k$ 列就是 $\vect{Y}_k$}。

课件 p20 还有一处小笔误：线性组合写成 $\sum_{j=0}^{n}y_{kj}c_j$，下标应从 $j=1$ 起。

\warnbox{\textbf{一个必须点明的坑}：$W\equiv0$ 并\textbf{不}等于线性相关。这条等价关系只对“齐次线性方程组的解”成立（定理 2.3），对任意向量函数不成立。反例：}

\[
\vect{Y}_1=\mat{1\\ x},\qquad \vect{Y}_2=\mat{x\\ x^2}.
\]

它们的 Wronskian 是 $1\cdot x^2-x\cdot x\equiv0$，但两者线性无关：若 $c_1\vect{Y}_1+c_2\vect{Y}_2=\vect{0}$ 对一切 $x$ 成立，看第一个分量得 $c_1+c_2x\equiv0$，只能 $c_1=c_2=0$。

\warnbox{考试里若用 Wronskian 判无关，\textbf{前提必须先说清}：“这两个是齐次线性方程的解，所以 $W\not\equiv0$ 与线性无关等价”。少了这句前提，反例就能推翻你的论证。}

\subhdr{5.3 定理 2.3（p21）}

齐次系统 $\vect{Y}'=A\vect{Y}$ 的 $n$ 个解 $\vect{Y}_1,\dots,\vect{Y}_n$ 线性无关，\textbf{当且仅当}存在一点 $x_0$ 使 $W(x_0)\neq0$；此时对区间内任意 $x$ 都有 $W(x)\neq0$。

后半句“一旦非零就处处非零”不是显然的，它来自 Abel 公式（一阶线性方程的 Wronskian 满足 $W'=(\tr A)W$，故 $W(x)=W(x_0)\exp\int_{x_0}^{x}\tr A$，指数因子永不为零）。L13 只引用结论不证，这里给一个指针：学到 L16--17 时可以把 $\tr A$ 与矩阵指数的行列式对上。

\subhdr{5.4 基本解组与定理 2.5（p22）}

定理 2.4：齐次解的任意线性组合仍是解（就是性质 2.1 的复述）。

定义：齐次系统的 $n$ 个\textbf{线性无关}的解，称为\textbf{基本解组}（fundamental set of solutions, FSS）。

定理 2.5：若 $\{\vect{Y}_1,\dots,\vect{Y}_n\}$ 是基本解组，则齐次方程的通解可写成

\[
\vect{Y}=\sum_{k=1}^{n}c_k\vect{Y}_k .
\]

\textbf{为什么恰好是 $n$ 个？}因为解空间是 $n$ 维的：给定 $x_0$，每个初值向量 $\vect{Y}(x_0)\in\mathbb{R}^n$ 由定理 2.1 唯一决定一条解，反过来每条解也由它的初值唯一决定。于是“解”与“$\mathbb{R}^n$ 中的向量”一一对应，$\mathbb{R}^n$ 的一组基（$n$ 个线性无关向量）就对应一组基本解。这就是“$n$ 个”的来历，也是定理 2.5 能说“任意解都是线性组合”的理由。

\subhdr{5.5 例 5（p23）}

\qbreak
\qstem{（课件 p23 例 5）仍取 $\mathrm{d}y_1/\mathrm{d}x=y_2,\ \mathrm{d}y_2/\mathrm{d}x=-y_1$。给定 $\vect{Y}_1=(\cos x,-\sin x)^{\mathsf T}$、$\vect{Y}_2=(\sin x,\cos x)^{\mathsf T}$，算 Wronskian 并写通解。}
\par\textbf{考点}：Wronskian 的计算；由基本解组写通解。
\par\textbf{关键一步}：$W=\cos x\cdot\cos x-(-\sin x)\cdot\sin x$，用 $\cos^2+\sin^2=1$ 收尾。
\par\textbf{详解}：附录 A.5，第 \pageref{c1:app:a5} 页。

%=====================================================================
\hdr{§6}{变换：$n$ 阶方程化成一阶组}

这是本讲最“工程”的一节：给定一个 $n$ 阶方程，怎么机械地写出一阶组。课件的例 6 用五步走完一般情形，例 7 给一个具体方程。

\subhdr{6.1 一般构造（p25）}

设方程已经解出最高阶导数：

\[
y^{(n)}=f\bigl(x,y,y',\dots,y^{(n-1)}\bigr).
\]

令

\[
y_1=y,\quad y_2=y',\quad y_3=y'',\quad \dots,\quad y_n=y^{(n-1)} .
\]

则前 $n-1$ 个方程是“下一行就是上一行的导数”，最后一个方程把 $y^{(n)}$ 用 $f$ 写成 $y_1,\dots,y_n$ 的函数：

\[
y_1'=y_2,\quad y_2'=y_3,\quad\dots,\quad y_{n-1}'=y_n,\quad y_n'=f(x,y_1,\dots,y_n).
\]

初值也跟着走：$y(x_0)=c_0,\dots,y^{(n-1)}(x_0)=c_{n-1}$ 对应 $y_k(x_0)=c_{k-1}$。

\textbf{这个变换是可逆的}：从组回到方程，只需把前 $n-1$ 个式子逐次求导代进最后一个；从方程到组，就是上面的定义。所以“解 $n$ 阶方程”与“解这个特殊形状的一阶组”是同一件事。

\subhdr{6.2 例 6：五步与那个“倒序”（p26--p29）}

课件把一般线性方程

\[
b_0(t)y^{(n)}+b_1(t)y^{(n-1)}+\cdots+b_{n-1}(t)y'+b_n(t)y=g(t),\qquad b_0\neq0
\]

化成 $\vect{Y}'=A(t)\vect{Y}+\vect{F}(t)$，走了五步。

\textbf{第 1 步，归一化}：两边除 $b_0$，把最高阶导数解出来：

\[
y^{(n)}=a_1(t)y^{(n-1)}+\cdots+a_{n-1}(t)y'+a_n(t)y+f(t),
\qquad a_j=-\frac{b_j}{b_0},\quad f=\frac{g}{b_0}.
\]

\textbf{第 2 步，设变量}：$y_1=y,\ y_2=y',\ \dots,\ y_n=y^{(n-1)}$。

\textbf{第 3 步，写前 $n-1$ 个方程}：$y_k'=y_{k+1}$（$k=1,\dots,n-1$）。

\textbf{第 4 步，写最后一个方程}：$y_n'=y^{(n)}$，把归一化后的右端用 $y_1,\dots,y_n$ 代掉：

\[
y_n'=a_n y_1+a_{n-1}y_2+\cdots+a_1y_n+f(t).
\]

\textbf{第 5 步，拼矩阵}：

\[
A=\begin{pmatrix}
0&1&0&\cdots&0\\
0&0&1&\cdots&0\\
\vdots& & &\ddots&\vdots\\
0&0&0&\cdots&1\\
a_n&a_{n-1}&a_{n-2}&\cdots&a_1
\end{pmatrix},
\qquad
\vect{F}=\mat{0\\ \vdots\\ 0\\ f(t)},\qquad
\vect{Y}(t_0)=\mat{c_0\\ \vdots\\ c_{n-1}}.
\]

\textbf{为什么最后一行是倒序 $(a_n,a_{n-1},\dots,a_1)$？}这是最容易写错的地方，课件 p28 只给结果没解释。看第 4 步的式子：$y_1$ 乘的是 $a_n$，$y_2$ 乘的是 $a_{n-1}$，……，$y_n$ 乘的是 $a_1$。而矩阵乘法里第 $n$ 行第 $j$ 列的元素正好乘 $y_j$，所以按 $y_1,\dots,y_n$ 的顺序排出来就是 $a_n,a_{n-1},\dots,a_1$。\textbf{一句话：系数下标与变量下标反着走。}

\subhdr{6.3 例 7（p30--p31）}

\qbreak
\qstem{（课件 p30 例 7）$t^4\dfrac{\mathrm{d}^2x}{\mathrm{d}t^2}+\sin t\dfrac{\mathrm{d}x}{\mathrm{d}t}-4x=\ln t,\quad t>0$。\textbf{化为一阶线性组。}}
\par\textbf{考点}：归一化后再设变量；系数含 $t$ 时矩阵 $A(t)$ 怎么写。
\par\textbf{关键一步}：先除 $t^4$ 解出 $x''$；再令 $y=x'$。注意这里的 $x,y$ 是原方程的未知函数与其导数，不是 §6.2 里的 $y_1,y_2$。
\par\textbf{详解}：附录 A.7，第 \pageref{c1:app:a7} 页。

\subhdr{6.4 一个 $n=3$ 的模板}

把 §6.2 的 $n=3$ 写出来，方便对照作业：

\[
y'''=a_1y''+a_2y'+a_3y+f
\ \Longrightarrow\
A=\mat{0&1&0\\0&0&1\\ a_3&a_2&a_1},\qquad \vect{F}=\mat{0\\0\\f}.
\]

最后一行按 $y_1,y_2,y_3$ 排，是 $(a_3,a_2,a_1)$——再一次是倒序。作业里那道题不是线性的，所以拼不成这个 $A$，但前两行完全一样。

%=====================================================================
\hdr{§7}{常系数齐次组的入口（p32--p36）}

这一节只开个头：$A$ 是常矩阵时，解长什么样。完整的三种情形在 L14--15。

\subhdr{7.1 从 $n=1$ 找类比（p33）}

$n=1$ 时方程退化成 $y'=ay$，解是 $y=c\mathrm{e}^{ax}$。这里 $a$ 是常数、$c$ 由初值定。这一行是全部猜想的来源。

\subhdr{7.2 试探解与特征方程（p34--p35）}

对 $\vect{Y}'=A\vect{Y}$，猜

\[
\vect{Y}(x)=\vect{V}\mathrm{e}^{\lambda x},\qquad \vect{V}\neq\vect{0}\ \text{常向量}.
\]

代进方程：左边 $\vect{Y}'=\lambda\vect{V}\mathrm{e}^{\lambda x}$，右边 $A\vect{V}\mathrm{e}^{\lambda x}$；约掉 $\mathrm{e}^{\lambda x}$（它恒不为零）得

\[
A\vect{V}=\lambda\vect{V}.
\]

于是\textbf{微分方程问题被换成了代数问题}：找 $A$ 的特征值 $\lambda$ 与特征向量 $\vect{V}$。把 $A\vect{V}=\lambda\vect{V}$ 移项得 $(A-\lambda E)\vect{V}=\vect{0}$；这是关于 $\vect{V}$ 分量的齐次线性方程组，要它有非零解，必须

\[
\det(A-\lambda E)=0 .
\]

这叫\textbf{特征方程}，它的左端叫特征多项式，次数是 $n$。

\subhdr{7.3 代数重数与几何重数（p36）}

特征多项式可分解为

\[
\det(A-\lambda E)=(\lambda-\lambda_1)^{m_1}\cdots(\lambda-\lambda_k)^{m_k},
\qquad m_1+\cdots+m_k=n .
\]

$m_j$ 叫特征值 $\lambda_j$ 的\textbf{代数重数}；$m=1$ 时叫单特征值，此时不同的特征值对应的特征向量自动线性无关。

课件 p36 最后一句是：“若 $\lambda$ 的重数 $k>1$，则 $\lambda$ 有 $1$ 到 $k$ 个线性无关的特征向量（或广义特征向量）。”这句话是全章最重要的伏笔，但它没给名字、没给算法。补上：

\[
g=\dim\ker(A-\lambda E)=n-\rank(A-\lambda E),\qquad 1\le g\le m .
\]

$g$ 叫\textbf{几何重数}，即“特征向量够不够多”。$g=m$ 时特征向量够用，解全部是 $\mathrm{e}^{\lambda x}$ 乘向量；$g<m$ 时不够，要补广义特征向量，解里会出现 $x\mathrm{e}^{\lambda x}$ 这样的项。这两种情形就是 L14--15 的 Case 3a 与 Case 3b。

%=====================================================================
\hdr{§8}{作业（p37）}

\qbreak
\qstem{（课件 p37 作业 I）把方程 $y'''=y''-x^2(y')^2$ 化为等价的一阶方程组。（课件原注：It's not a linear system.）}
\par\textbf{考点}：§6.1 的一般变换在非线性方程上的用法；为什么这次拼不出矩阵 $A$。
\par\textbf{关键一步}：照 §6.4 的模板设 $y_1=y,\ y_2=y',\ y_3=y''$，前两行照抄，第三行把 $y'''$ 用原方程换掉。
\par\textbf{详解}：附录 A.8，第 \pageref{c1:app:a8} 页。

%=====================================================================
\hdr{§9}{易错点}

\begin{enumerate}
  \item \textbf{通解的常数必须“够且能反解”}。例 1 的 $c_1\cos(x-c_2)$ 与 $A\cos x+B\sin x$ 都是两个常数；但写在答卷上时，若用前者，要能说清 $c_1,c_2$ 与初值的关系（见 A.1）。
  \item \textbf{$W\equiv0$ 不等于线性相关}。只对“齐次线性方程的解”成立；反例见 §5.2。答卷上必须写出“这是齐次解”这个前提。
  \item \textbf{变换矩阵的最后一行是倒序}。$A$ 的最后一行按 $y_1,\dots,y_n$ 排是 $(a_n,\dots,a_1)$，不是 $(a_1,\dots,a_n)$（§6.2）。
  \item \textbf{非线性方程化成的组不能写成 $A\vect{Y}$}。作业那题的第三式含 $x^2y_2^2$，是二次项，只能写成 $\vect{Y}'=\vect{F}(x,\vect{Y})$。
  \item \textbf{奇解只出现在非线性方程}。例 2 有奇解，线性系统没有（Remark 2.1 的理由见 §4.1）。
  \item \textbf{$a_{kj}$ 先行后列}。写矩阵时逐行抄系数，别按“先写各变量”去拼；课件 p13 第二行末元素的笔误说明这条线也会被印错。
\end{enumerate}

%=====================================================================
\hdr{§10}{缺口清单}

下表的“缺口”指课件跳过的步骤或表述不清处，逐条给页码与原文短语。

\gapitem{机制缺口}{p4 “The two masses move on a frictionless surface … = −(k1+k2)x1 + k2x2 + F1” 只给结果，未说明三根弹簧的伸长量分别是 $x_1$、$x_2-x_1$、$x_2$。本讲义 §1.1 逐项补受力。}
\gapitem{机制缺口}{p5 “Let V be the voltage drop … dV/dt = −I/C − V/(RC)” 直接给方程，未写基尔霍夫与元件本构。本讲义 §1.2 补推导。}
\gapitem{代数跳步}{p6 “The above equation can be transformed … Let x1 = u and x2 = u'.” 只声明可化，未演示。本讲义 §1.3 补出中间一步。}
\gapitem{表述不清}{p9 “is called the general solution on the domain D, if for every point (x, y1, …, yn) $\in$ D” 条件被截断。本讲义 §2.3 给可操作判据：常数个数 + 初值可反解。}
\gapitem{平凡解／奇解}{p10 “Singular Solution: y1 = 0, y2 = c” 未说明为何不是通解取常数。本讲义 A.2 补理由：需 $(x+c_1)^2\equiv0$。}
\gapitem{排版笔误}{p13 矩阵第二行末元素印成 $a_{1n}$，应为 $a_{2n}$。本讲义 §3.2 更正。}
\gapitem{排版笔误}{p20 线性组合写成 $\sum_{j=0}^{n}y_{kj}c_j$（下标应从 1 起），且 Wronskian 行列式先按列后按行两种写法并列。本讲义 §5.2 统一为按列排列。}
\gapitem{适用区间}{p21 定理 2.3 未说明“$W\neq0\Leftrightarrow$ 线性无关”只对解成立。本讲义 §5.2 给反例并标明前提。}
\gapitem{代数跳步}{p28 矩阵最后一行写成 $(a_n,a_{n-1},\dots,a_1)$ 未解释倒序。本讲义 §6.2 说明“系数下标与变量下标反着走”。}
\gapitem{术语未定义}{p36 “anywhere from 1 to k linearly independent eigenvectors” 只描述现象，未定义几何重数。本讲义 §7.3 补 $g=n-\rank(A-\lambda E)$。}

%=====================================================================
\hdr{§11}{三句话}

\begin{enumerate}
  \item 方程组不是新东西：真实系统本来就是多个未知函数，L13 只是把它们写成向量。
  \item 解的结构是线性的：齐次解的任意组合仍是解，非齐次解 $=$ 齐次通解 $+$ 一个特解；齐次只需 $n$ 个线性无关的解，判据是 Wronskian（\textbf{只对解}成立）。
  \item 任何 $n$ 阶方程都能化成一阶组：设 $y_k=y^{(k-1)}$，前 $n-1$ 行是 $0,1$ 的移位，最后一行是系数倒序；$A$ 为常矩阵时，再把问题换成特征值。
\end{enumerate}

%=====================================================================
\hdr{§12}{与课件的页码对照}

\begin{minipage}{\linewidth}
\tabtop
\begin{tabular}{@{}p{0.12\linewidth}p{0.40\linewidth}p{0.38\linewidth}@{}}
\thead{本讲义} & \thead{内容} & \thead{课件页}\\
\end{tabular}
\tabmid
\begin{tabular}{@{}p{0.12\linewidth}p{0.40\linewidth}p{0.38\linewidth}@{}}
§1.1 & 两质量三弹簧与受力 & p4\\
§1.2 & 并联 LRC 与基尔霍夫 & p5\\
§1.3 & 弹簧-质量化一阶组 & p6\\
§2.1 & 一般系统与正规形 & p7\\
§2.2 & 解与初值问题 & p8\\
§2.3 & 通解、特解、奇解 & p9\\
§2.4 & 例 1、例 2 & p10\\
§3.1 & 线性系统一般形 & p12--p13\\
§3.2 & 矩阵写法与笔误 & p13\\
§3.3 & 例 3 & p14\\
§4.1 & 定理 2.1 与 Remark 2.1 & p15\\
§4.2 & 线性算子与性质 2.1 & p16\\
§4.3 & 叠加原理、定理 2.2 & p17\\
§4.4 & 例 4 & p18\\
§5.1 & 线性相关与无关 & p19\\
§5.2 & Wronskian 与反例 & p20\\
§5.3 & 定理 2.3 & p21\\
§5.4 & 定理 2.4、FSS、定理 2.5 & p22\\
§5.5 & 例 5 & p23\\
§6.1 & 一般变换 & p25\\
§6.2 & 例 6 五步与倒序 & p26--p29\\
§6.3 & 例 7 & p30--p31\\
§7 & 常系数入口：试探解与特征方程 & p32--p36\\
§8 & 作业 & p37\\
\end{tabular}
\tabbot
\end{minipage}

%=====================================================================
\newpage
\chapterbanner{2}{常系数齐次组}{三种特征值情形与广义特征向量}{课件 \texttt{14\_15\_Linear-DE-systems-\ldots.pdf}（64 页，2026-09-29／10-08；全名见该册原封面）}
\callout{\textbf{先读这一段。}这一讲只解决一个方程：$\vect{Y}'=A\vect{Y}$，$A$ 是常矩阵。全部内容由特征值分成三支：\textbf{相异实根}（§3）、\textbf{共轭复根}（§4）、\textbf{重根}（§6–§7，还要再分 $g=m$ 与 $g<m$）。三种情形的解长得不一样，但来源只有一个：猜 $\vect{Y}=\vect{V}\mathrm{e}^{\lambda x}$（§2）。你要练的不是"背三种公式"，而是\textbf{拿到一个矩阵，先算出 $\lambda$ 与 $m,g$，再决定走哪一支}。§8 把这件事做成一张六步检查表。}

%=====================================================================
\hdr{§0}{三种情形一览，与三个要一直盯住的量}

先把结论摆出来，后面的每一节都是在解释这张表怎么来的。

\begin{minipage}{\linewidth}
\tabtop
\begin{tabular}{@{}p{0.16\linewidth}p{0.20\linewidth}p{0.42\linewidth}p{0.14\linewidth}@{}}
\thead{情形} & \thead{特征值} & \thead{解的基本形状} & \thead{课件页}\\
\end{tabular}
\tabmid
\begin{tabular}{@{}p{0.16\linewidth}p{0.20\linewidth}p{0.42\linewidth}p{0.14\linewidth}@{}}
Case 1 & 相异实根 $\lambda_k$ & $\vect{V}_k\mathrm{e}^{\lambda_k x}$，各管各的 & p9--p15\\
Case 2 & 共轭复根 $a\pm\mathrm{i}b$ & $\mathrm{e}^{ax}\bigl(\vect{P}\cos bx-\vect{Q}\sin bx\bigr)$ 与其虚部 & p16--p25\\
Case 3a & 重根，$g=m$ & 仍是 $\mathrm{e}^{\lambda x}$ 乘 $m$ 个独立特征向量 & p30--p38\\
Case 3b & 重根，$1\le g<m$ & 出现 $x\mathrm{e}^{\lambda x}$、$\tfrac{x^2}{2}\mathrm{e}^{\lambda x}$ 的链 & p39--p45\\
\end{tabular}
\tabbot
\end{minipage}

这张表是路线台账，不承载推导；它回答的是"算完特征值之后走哪条路"，每一行的理由在后面各节。

\textbf{三个要一直盯住的量}：

\begin{itemize}
  \item $\lambda$：特征值，解里指数的底数。
  \item $m$：\textbf{代数重数}，$\lambda$ 在特征多项式里出现的次数。
  \item $g$：\textbf{几何重数}，$g=n-\rank(A-\lambda E)$，即"这个 $\lambda$ 能拿出几个线性无关的特征向量"。
\end{itemize}

永远有 $1\le g\le m$。三条路的分岔点不是"$m$ 等于几"，而是"\textbf{$g$ 够不够}"：$g=m$ 走 Case 3a，$g<m$ 走 Case 3b。这一点课件在 p7 埋了一句，到 p31 才给出定义，中间夹着一整个 Case 1 与 Case 2，很容易读丢。

%=====================================================================
\hdr{§1}{特征值问题：从 $A\vect{V}=\lambda\vect{V}$ 到特征方程}

\subhdr{1.1 定义与特征方程（p3、p6）}

设 $A$ 是 $n\times n$ 实常矩阵。若存在数 $\lambda$ 与非零向量 $\vect{V}$ 使

\[
A\vect{V}=\lambda\vect{V},
\]

则 $\lambda$ 叫特征值、$\vect{V}$ 叫对应的特征向量。把它移项：

\[
(A-\lambda E)\vect{V}=\vect{0}.
\]

这是一个关于 $\vect{V}$ 各分量的\textbf{齐次线性代数方程组}。齐次方程组只有两种命运：要么只有零解，要么有无穷多非零解。我们要的是后者，而后者成立的充要条件是系数行列式为零：

\[
\det(A-\lambda E)=0 .
\]

这就是\textbf{特征方程}；左端是关于 $\lambda$ 的 $n$ 次多项式，叫特征多项式。

\subhdr{1.2 代数重数与几何重数（p7）}

把特征多项式分解：

\[
\det(A-\lambda E)=(\lambda-\lambda_1)^{m_1}\cdots(\lambda-\lambda_k)^{m_k},
\qquad m_1+\cdots+m_k=n .
\]

$m_j$ 叫 $\lambda_j$ 的\textbf{代数重数}。对每个 $\lambda_j$，解 $(A-\lambda_jE)\vect{V}=\vect{0}$ 得到的\textbf{线性无关特征向量的个数}，叫\textbf{几何重数}

\[
g_j=\dim\ker(A-\lambda_jE)=n-\rank(A-\lambda_jE),
\]

并且总有 $1\le g_j\le m_j$。

\textbf{为什么 $g\le m$？}直觉上，$g$ 是"解空间维数"，$m$ 是"根在这点重复了几次"；特征值每重复一次，最多多给一个独立方向，给不出更多。这个不等式是后面一切分类的基础，但课件 p7 只写了结论，没解释，也没告诉你怎么算——算的办法就是"数秩"，$g=n-\rank$。

\subhdr{1.3 三个量与三条路的对应}

拿到 $A$ 之后，先算特征多项式，得到两样东西：\textbf{有哪些 $\lambda$}、\textbf{每个 $\lambda$ 的 $m$}。然后对每个 $\lambda$ 算 $g=n-\rank(A-\lambda E)$。

\begin{itemize}
  \item 所有 $\lambda$ 都是单根（$m=1$，于是 $g=1$）：\textbf{Case 1}。
  \item 出现一对共轭复根：\textbf{Case 2}（复根一定成对出现，见 §4.4）。
  \item 某个 $\lambda$ 的 $m>1$：看它的 $g$。$g=m$ 走 \textbf{Case 3a}；$g<m$ 走 \textbf{Case 3b}。
\end{itemize}

一个矩阵里可以同时出现多种情形：例如 $\lambda=2$（$m=1$）与 $\lambda=3$（$m=2$，$g=2$）并存，就一半走 Case 1、一半走 Case 3a。各支互不干扰，最后把解加起来（叠加原理）。

%=====================================================================
\hdr{§2}{为什么猜 $\vect{Y}=\vect{V}\mathrm{e}^{\lambda x}$}

\subhdr{2.1 $n=1$ 的类比（p4）}

$n=1$ 时方程退化成 $y'=ay$（$a$ 常数），解是 $y=c\mathrm{e}^{ax}$。这一条是全部猜想的来源：方程右端是"常数乘自己"，解就该是"指数乘常数"。

\subhdr{2.2 代入与约化（p5）}

设 $\vect{Y}(x)=\vect{V}\mathrm{e}^{\lambda x}$，其中 $\vect{V}$ 是非零常向量。求导：

\[
\vect{Y}'=\lambda\vect{V}\mathrm{e}^{\lambda x},\qquad
A\vect{Y}=A\vect{V}\mathrm{e}^{\lambda x}.
\]

两边相等，而 $\mathrm{e}^{\lambda x}$ 恒不为零，可以约掉，剩下

\[
\lambda\vect{V}=A\vect{V}.
\]

\textbf{这一步是整章的关键}：微分方程问题被换成了代数问题——找矩阵的特征值。约掉指数之所以合法，是因为 $\mathrm{e}^{\lambda x}$ 处处非零；这也解释了为什么解里必然带指数。

\subhdr{2.3 为什么偏偏是指数}

从另一头看：$\vect{Y}'=A\vect{Y}$ 的解，每个分量都是若干个 $\mathrm{e}^{\lambda x}$（或 $x^k\mathrm{e}^{\lambda x}$）的线性组合。理由在于常系数线性方程的解空间是有限维的，而指数函数是求导算子的"特征函数"——求导把它送回自己乘一个常数。这一条是直觉，严格证明要用矩阵指数（L16--17 的 $\mathrm{e}^{Ax}$），这里先用着。

%=====================================================================
\hdr{§3}{Case 1：相异实根（p9--p15）}

\subhdr{3.1 定理与通解（p10）}

设 $A$ 的 $n$ 个特征值 $\lambda_1,\dots,\lambda_n$ 互不相同，对应的特征向量为 $\vect{V}_1,\dots,\vect{V}_n$。则

\[
\vect{Y}_k(x)=\vect{V}_k\mathrm{e}^{\lambda_k x},\qquad k=1,\dots,n
\]

是 $n$ 个线性无关的解，通解为

\[
\vect{Y}(x)=\sum_{k=1}^{n}c_k\vect{V}_k\mathrm{e}^{\lambda_k x}.
\]

\textbf{为什么线性无关}（课件 p10 的定理 2.1 只给结论，这里补）：设

\[
c_1\vect{V}_1\mathrm{e}^{\lambda_1x}+\cdots+c_n\vect{V}_n\mathrm{e}^{\lambda_nx}=\vect{0}
\]

对一切 $x$ 成立。两边乘 $\mathrm{e}^{-\lambda_1x}$，再令 $x\to+\infty$：若某个 $c_k\neq0$ 且 $\lambda_k>\lambda_1$，则对应项会发散，等式不可能成立；于是只能 $c_1=0$。逐个消去，得所有 $c_k=0$。所以无关。

\textbf{更直接的判据}：这 $n$ 个解拼成的 Wronskian 在 $x=0$ 处等于 $\det(\vect{V}_1,\dots,\vect{V}_n)$。而分属不同特征值的特征向量必线性无关（否则 $A\vect{V}_i=\lambda_i\vect{V}_i$ 与 $A\vect{V}_j=\lambda_j\vect{V}_j$ 会推出 $\lambda_i=\lambda_j$），故行列式非零。

\subhdr{3.2 操作三条}

\begin{enumerate}
  \item 写出 $A$，解 $\det(A-\lambda E)=0$ 得全部 $\lambda_k$（要检查是否互异）。
  \item 对每个 $\lambda_k$ 解 $(A-\lambda_kE)\vect{V}_k=\vect{0}$，取一个非零解（通常把某个自由变量取 $1$）。
  \item 拼 $\vect{Y}=\sum c_k\vect{V}_k\mathrm{e}^{\lambda_k x}$。
\end{enumerate}

\textbf{注意}：特征向量可以乘任意非零常数，不同人取的 $\vect{V}_k$ 可能差一个倍数，但通解集合相同。答卷上不必统一"最美"的写法，只要非零即可。

%=====================================================================
\hdr{§4}{Case 2：共轭复根（p16--p25）}

\subhdr{4.1 复值解怎么来的（p17--p18）}

设 $\lambda_1=a+\mathrm{i}b$（$b\neq0$）是特征值，$\vect{V}_1$ 是对应特征向量。形式上

\[
\vect{Y}_1(x)=\vect{V}_1\mathrm{e}^{(a+\mathrm{i}b)x}
=\mathrm{e}^{ax}\bigl(\cos bx+\mathrm{i}\sin bx\bigr)\vect{V}_1
\]

确实是解——但它是\textbf{复值}的，而原方程是实系数的，我们要实值解。

\subhdr{4.2 取实部与虚部：为什么合法（p19、p21）}

把 $\vect{V}_1$ 拆成实部与虚部：$\vect{V}_1=\vect{P}+\mathrm{i}\vect{Q}$（$\vect{P},\vect{Q}$ 实向量）。则

\[
\vect{Y}_1(x)=\underbrace{\mathrm{e}^{ax}\bigl(\vect{P}\cos bx-\vect{Q}\sin bx\bigr)}_{\vect{U}(x)}
+\mathrm{i}\underbrace{\mathrm{e}^{ax}\bigl(\vect{P}\sin bx+\vect{Q}\cos bx\bigr)}_{\vect{W}(x)}.
\]

\textbf{$\vect{U}$ 与 $\vect{W}$ 都是解}。理由是线性算子 $\mathcal{L}[\vect{Y}]=\vect{Y}'-A\vect{Y}$ 只做求导与乘实矩阵，对复值的实部虚部互不干扰：

\[
\vect{0}=\mathcal{L}[\vect{Y}_1]=\mathcal{L}[\vect{U}]+\mathrm{i}\mathcal{L}[\vect{W}]
\ \Longrightarrow\ \mathcal{L}[\vect{U}]=\vect{0},\ \mathcal{L}[\vect{W}]=\vect{0}.
\]

课件 p19 用的是另一种说法：$\vect{U}=(\vect{Y}_1+\bar{\vect{Y}}_1)/2$、$\vect{W}=(\vect{Y}_1-\bar{\vect{Y}}_1)/(2\mathrm{i})$，再由叠加原理得它们是解。两种说法等价，本质都是"线性算子把实部虚部分开"。

\textbf{它们线性无关}：课件 p19 直接算了 Wronskian（例 2 里 $W=\mathrm{e}^{-x}\neq0$）。一般地，$\vect{U},\vect{W}$ 的线性组合若为零，等价于 $\vect{V}_1\mathrm{e}^{\lambda_1x}$ 与其共轭的某个实组合为零，由 $\vect{P},\vect{Q}$ 与 $\cos,\sin$ 的独立性可推出系数全零。

\subhdr{4.3 一般公式（p20）}

设 $\lambda_1=a+\mathrm{i}b$、$\lambda_2=a-\mathrm{i}b$（$b\neq0$）是一对共轭复根，其余 $\lambda_3,\dots,\lambda_n$ 为实数且互异。则通解为

\[
\vect{Y}(x)=c_1\vect{U}(x)+c_2\vect{W}(x)+\sum_{k=3}^{n}c_k\vect{V}_k\mathrm{e}^{\lambda_k x},
\]

其中

\[
\vect{U}(x)=\mathrm{e}^{ax}\Bigl(\operatorname{Re}(\vect{V}_1)\cos bx-\operatorname{Im}(\vect{V}_1)\sin bx\Bigr),
\qquad
\vect{W}(x)=\mathrm{e}^{ax}\Bigl(\operatorname{Im}(\vect{V}_1)\cos bx+\operatorname{Re}(\vect{V}_1)\sin bx\Bigr).
\]

\textbf{实操只有一句}：\textbf{只算 $\lambda=a+\mathrm{i}b$ 这一个特征值}，取 $\vect{V}=\vect{P}+\mathrm{i}\vect{Q}$，代进上面两式。共轭的那一半不用算，它给出的正是 $\vect{U},\vect{W}$ 的同一对组合。

\subhdr{4.4 为什么复根一定成对（p21）}

$A$ 是实矩阵，特征多项式的系数全是实数。实系数多项式的复根必成共轭对出现，所以 $\lambda_2=\bar\lambda_1$ 也是特征值；并且 $\bar{\vect{V}}_1$ 就是它的特征向量。这一条保证"只算一半"永远够用。

\subhdr{4.5 一个必须提防的坑}

复根情形下，\textbf{不要}把 $\vect{V}\mathrm{e}^{\lambda x}$ 直接当答案交上去——那不是实值解。也不要只写"$\cos$ 那一半"：$\vect{U}$ 与 $\vect{W}$ 都要，它们是两个独立解，缺一个就少一维。

%=====================================================================
\hdr{§5}{四维的例子：多自由度弹簧-质量（p26--p29）}

\subhdr{5.1 建模（p26）}

两个质量、三根弹簧的系统，化成四个一阶方程：令 $y_1=x_1$、$y_2=x_2$、$y_3=x_1'$、$y_4=x_2'$，则

\[
\vect{Y}'=\begin{pmatrix}
0&0&1&0\\
0&0&0&1\\
-2&3/2&0&0\\
4/3&-3&0&0
\end{pmatrix}\vect{Y}.
\]

矩阵里 $2\times2$ 的零块与单位块，是"两个二阶方程降阶"留下的痕迹；下面 $2\times2$ 的块是刚度矩阵除以质量。

\warnbox{\textbf{课件 p26 有一处符号笔误}：它写 $m_2x_2''=-k_2x_1-(k_2+k_3)x_2$，但按 L13 p4 的受力约定应为 $+k_2x_1$；p27 的矩阵第四行也确实是 $+4/3$。以矩阵为准，p26 的 $-\tfrac43x_1$ 应为 $+\tfrac43x_1$。}

\subhdr{5.2 特征方程里只有 $\lambda$ 的偶次幂（p27）}

\[
\det(A-\lambda E)=\lambda^4+5\lambda^2+4=(\lambda^2+1)(\lambda^2+4),
\]

所以 $\lambda=\pm\mathrm{i},\ \pm2\mathrm{i}$：\textbf{四个纯虚根}。

\textbf{纯虚根意味着什么？}指数是 $\mathrm{e}^{\pm\mathrm{i}bt}$，没有实部，振幅不增不减——这是\textbf{无阻尼的简谐振动}。这正对应模型里没有阻尼项：方程组里没有 $x'$ 的一次项，矩阵的上右块是单位块、下右块是零块，于是特征方程只含 $\lambda$ 的偶次幂。这个结构不是巧合，是"无阻尼"在代数上的样子。

\subhdr{5.3 复根在四维里的用法}

$\lambda=\mathrm{i}$ 的特征向量是 $\vect{V}_1=(3,2,3\mathrm{i},2\mathrm{i})^{\mathsf T}$（课件 p27）。取实部虚部得到两条实解；$\lambda=2\mathrm{i}$ 同样处理，再得两条。四条拼起来就是通解（课件 p29）。注意：$b=1$ 与 $b=2$ 对应两个不同的振动频率，$\cos t$ 与 $\cos2t$ 各是一阶模态。

%=====================================================================
\hdr{§6}{Case 3a：重根但 $g=m$（p30--p38）}

\subhdr{6.1 定义与定理（p31）}

设 $\lambda_1$ 的代数重数为 $m_1$，几何重数

\[
g_1=n-\rank(A-\lambda_1E).
\]

\textbf{定理 2.3}：若 $g_1=m_1$，则恰有 $m_1$ 个线性无关的特征向量 $\vect{V}_1,\dots,\vect{V}_{m_1}$ 对应 $\lambda_1$。

\subhdr{6.2 通解（p32）}

\[
\vect{Y}(x)=\bigl(c_1\vect{V}_1+\cdots+c_{m_1}\vect{V}_{m_1}\bigr)\mathrm{e}^{\lambda_1x}
+\sum_{k=m_1+1}^{n}c_k\vect{V}_k\mathrm{e}^{\lambda_kx}.
\]

注意它和 Case 1 的差别只在"同一个 $\lambda$ 配了好几个特征向量"：\textbf{指数还是 $\mathrm{e}^{\lambda x}$，没有 $x$ 因子}。$g=m$ 时特征向量够用，不需要广义特征向量。

\subhdr{6.3 判别 $g=m$ 的算法}

\begin{enumerate}
  \item 算 $A-\lambda_1E$，用初等行变换化成阶梯形，数非零行数 = $\rank$。
  \item $g=n-\rank$。若 $g=m$，进入本情形。
  \item 解 $(A-\lambda_1E)\vect{V}=\vect{0}$，把自由变量逐个取 $1$（其余取 $0$），得到 $g$ 个独立特征向量。
\end{enumerate}

\subhdr{6.4 初值问题的六步（p37--p38）}

课件把整套流程写成六步，适用于所有情形（复根时在第 III 步取实虚部）：

\begin{enumerate}
  \item 求全部特征值 $\lambda_1,\dots,\lambda_k$；
  \item 求对应的特征向量（或广义特征向量）$\vect{V}_1,\dots,\vect{V}_n$；
  \item 造解 $\vect{Y}_k=\vect{V}_k\mathrm{e}^{\lambda_kx}$；复根时改取实部、虚部；
  \item 验 $W=\det(\vect{Y}_1,\dots,\vect{Y}_n)\neq0$（课件注明"通常自动满足"）；
  \item 写通解 $\vect{Y}=\sum c_k\vect{Y}_k$；
  \item 若给初值 $\vect{Y}(0)$，代入定 $c_1,\dots,c_n$。
\end{enumerate}

第 IV 步不能省：它保证你找到的 $n$ 个解真的构成基本解组。$g=m$ 时这一条自动成立（$m$ 个独立特征向量）；$g<m$ 时就要靠链的构造来保证。

%=====================================================================
\hdr{§7}{Case 3b：重根且 $1\le g<m$（p39--p45）}

这一节是全讲最绕的地方。课件用"试错"的方式引出广义特征向量，中间的逻辑跳了一步，这里补全。

\subhdr{7.1 第一次尝试为什么失败（p40）}

$\lambda$ 是二重根、却只有一个特征向量 $\vect{V}_1$。自然的想法是再乘一个 $x$：设

\[
\vect{Y}_2=x\mathrm{e}^{\lambda x}\vect{V}_1 .
\]

求导得 $\vect{Y}_2'=(1+\lambda x)\mathrm{e}^{\lambda x}\vect{V}_1$，而 $A\vect{Y}_2=\lambda x\mathrm{e}^{\lambda x}\vect{V}_1$。要相等，必须

\[
1+\lambda x=\lambda x,
\]

对任何 $x$ 都不成立。\textbf{失败的原因}：$\vect{V}_1$ 已经在 $A-\lambda E$ 的核里，乘一个 $x$ 只改变系数，不产生新的方向。

\subhdr{7.2 修正：加一个待定向量（p41）}

改设

\[
\vect{Y}_2=x\mathrm{e}^{\lambda x}\vect{V}_1+\mathrm{e}^{\lambda x}\vect{V}^{(1)},
\]

其中 $\vect{V}^{(1)}$ 待定。求导：

\[
\vect{Y}_2'=\mathrm{e}^{\lambda x}\bigl[\lambda x\vect{V}_1+(\vect{V}_1+\lambda\vect{V}^{(1)})\bigr],
\qquad
A\vect{Y}_2=\mathrm{e}^{\lambda x}\bigl[\lambda x\vect{V}_1+A\vect{V}^{(1)}\bigr].
\]

比较两项（$x\mathrm{e}^{\lambda x}$ 的系数自动相等，剩下的常数项必须相等）：

\[
\vect{V}_1+\lambda\vect{V}^{(1)}=A\vect{V}^{(1)}
\ \Longleftrightarrow\
(A-\lambda E)\vect{V}^{(1)}=\vect{V}_1 .
\]

\textbf{这个方程就是把"缺的那一维"补出来的地方}。它一般有解（只要 $\vect{V}_1$ 落在 $A-\lambda E$ 的像里），$\vect{V}^{(1)}$ 就叫 $\lambda$ 对应的\textbf{广义特征向量}。

\subhdr{7.3 链：$g$ 与 $m$ 差多少，就接多长（p43）}

一般情形：若 $\lambda_1$ 有 $g_1$ 个独立特征向量、代数重数 $m_1$，则需要 $k=m_1-g_1$ 个广义特征向量，串成一条链：

\[
(A-\lambda_1E)\vect{V}=\vect{0},\quad
(A-\lambda_1E)\vect{V}^{(1)}=\vect{V},\quad
(A-\lambda_1E)\vect{V}^{(2)}=\vect{V}^{(1)},\ \dots,\
(A-\lambda_1E)\vect{V}^{(k)}=\vect{V}^{(k-1)}.
\]

\textbf{一个容易忽略的条件}：链的起点 $\vect{V}$ 必须是"能接上"的特征向量，即 $\vect{V}\in\operatorname{Range}(A-\lambda_1E)$。判别法是比较

\[
\rank(A-\lambda_1E)\quad\text{与}\quad\rank\bigl(A-\lambda_1E\mid\vect{V}\bigr).
\]

两者相等，才有 $\vect{V}^{(1)}$；不相等就换一个特征向量（或换成几个特征向量的组合）。课件 p61 在三维 $g=2$ 的例子正是这么做的：取 $\vect{V}=(0,0,1)$ 时两个秩不等，接不上；改成 $\vect{V}=(2,1,1)=(2,1,0)+(0,0,1)$ 就能接上。

\subhdr{7.4 通解的一般形式（p44、p45）}

把这些解加起来：

\[
\vect{Y}=\bigl(c_1\vect{V}_1+\cdots+c_{g}\vect{V}_{g}\bigr)\mathrm{e}^{\lambda x}
+\sum_{j=1}^{k}c_{g+j}\Bigl(\vect{V}\frac{x^{j}}{j!}+\vect{V}^{(1)}\frac{x^{j-1}}{(j-1)!}+\cdots+\vect{V}^{(j)}\Bigr)\mathrm{e}^{\lambda x}
+\sum_{\text{其余 }\lambda}c_k\vect{V}_k\mathrm{e}^{\lambda_kx}.
\]

\textbf{读法}：括号里是"$\vect{V}$ 配 $x^j/j!$、$\vect{V}^{(1)}$ 配 $x^{j-1}/(j-1)!$、……、$\vect{V}^{(j)}$ 配常数"，系数正好是泰勒展开的系数。$k=1$ 时退化成 $(\vect{V}x+\vect{V}^{(1)})\mathrm{e}^{\lambda x}$，就是 §7.2 的结果。

\subhdr{7.5 一张对照表（p45）}

\begin{minipage}{\linewidth}
\tabtop
\begin{tabular}{@{}p{0.20\linewidth}p{0.30\linewidth}p{0.34\linewidth}@{}}
\thead{} & \thead{$g=m$} & \thead{$1\le g<m$（$k=m-g$）}\\
\end{tabular}
\tabmid
\begin{tabular}{@{}p{0.20\linewidth}p{0.30\linewidth}p{0.34\linewidth}@{}}
特征向量 & $\vect{V}_1,\dots,\vect{V}_m$ & $\vect{V}_1,\dots,\vect{V}_g$\\
广义特征向量 & 不需要 & $\vect{V}^{(1)},\dots,\vect{V}^{(k)}$\\
解 & 全乘 $\mathrm{e}^{\lambda x}$ & $\vect{V}\mathrm{e}^{\lambda x}$、$(\vect{V}x+\vect{V}^{(1)})\mathrm{e}^{\lambda x}$、$\bigl(\vect{V}\tfrac{x^2}{2}+\vect{V}^{(1)}x+\vect{V}^{(2)}\bigr)\mathrm{e}^{\lambda x}$、……\\
解够不够 & 够（$m$ 条） & 单独几条不够，要接成链才够\\
\end{tabular}
\tabbot
\end{minipage}

这张表是形态对照，不承载推导：它能告诉你"该长什么样"，但链怎么接、$\vect{V}^{(1)}$ 怎么解，仍要回到 §7.2、§7.3。

%=====================================================================
\hdr{§8}{拿到一个矩阵时的六步检查表}

\begin{enumerate}
  \item \textbf{写 $A$}，检查是否常系数、是否齐次。
  \item \textbf{算特征多项式} $\det(A-\lambda E)$，因式分解，列出所有 $\lambda$ 与各自的 $m$。
  \item \textbf{算 $g$}：对每个重根算 $n-\rank(A-\lambda E)$。这一步决定后面走 3a 还是 3b；单根不用算（$g=m=1$）。
  \item \textbf{求向量}：单根与 $g=m$ 的情形解 $(A-\lambda E)\vect{V}=0$；复根只算一个，取实虚部；$g<m$ 时先求特征向量，再按链解 $(A-\lambda E)\vect{V}^{(j)}=\vect{V}^{(j-1)}$（并检查秩条件）。
  \item \textbf{拼解}，数一数总共是不是 $n$ 条。
  \item \textbf{验}：把解代回原方程（至少验一条），或算 Wronskian 看是否非零；有初值再代初值。
\end{enumerate}

\textbf{第 5 步的数数是最便宜的检查}：解的总数必须恰好等于 $n$。少了，说明某条链断了或某个复根只取了一半；多了，说明有两列线性相关。

%=====================================================================
\hdr{§9}{解长什么样：三种情形的几何与稳定性}

前面三节解决"怎么算出解"，这一节回答"解画出来是什么"。这一节不是可有可无的补充：\textbf{下一章（L21--23 相图与稳定性）整章都建立在这里}，那章用的 $\tau$--$\Delta$ 判据，就是把本节三种情形翻译成迹与行列式的语言。

\subhdr{8.1 一个实根：轨道是一条射线}

若解里只有一项 $c\vect{V}\mathrm{e}^{\lambda x}$（$c\neq0$），把它写成参数式：$\vect{Y}(x)=c\mathrm{e}^{\lambda x}\vect{V}$。无论 $x$ 取什么值，$\vect{Y}$ 永远是 $\vect{V}$ 的常数倍，所以\textbf{轨道落在过原点、方向为 $\vect{V}$ 的直线上}：

\begin{itemize}
  \item $\lambda<0$：$c\mathrm{e}^{\lambda x}\to0$，轨道沿直线缩向原点（$c>0$ 从 $\vect{V}$ 那一侧来，$c<0$ 从对侧来）。
  \item $\lambda>0$：轨道沿直线离开原点，趋向无穷。
  \item $\lambda=0$：$c\mathrm{e}^{0}\vect{V}$ 不动——这是平衡点，不出现在一般情形。
\end{itemize}

两个相异实根时，通解是两条直线方向的叠加，轨道一般是一条穿过原点的曲线，被两个特征方向"夹"住：$x\to+\infty$ 时由 $|\lambda|$ 较大的那项主导，轨道趋于那个特征方向；$x\to-\infty$ 时由另一项主导。若一正一负（$\Delta<0$），轨道从一个方向进来、从另一个方向出去——这就是\textbf{鞍点}。

\subhdr{8.2 一对复根：轨道是螺旋}

复根情形下 $\vect{U}(x)=\mathrm{e}^{ax}\bigl(\vect{P}\cos bx-\vect{Q}\sin bx\bigr)$。在 $\vect{P},\vect{Q}$ 张成的平面里看，它是一条\textbf{对数螺旋}：

\begin{itemize}
  \item 半径按 $\mathrm{e}^{ax}$ 变化：$a<0$ 向内卷，$a>0$ 向外卷；
  \item 角度按 $bx$ 匀速增加：每走 $\Delta x=2\pi/b$ 转一圈；
  \item $a=0$（纯虚根）：半径不变，轨道是闭合椭圆——这就是\textbf{中心}，对应无阻尼振动（§5 的弹簧系统正是这种）。
\end{itemize}

\textbf{注意旋转方向}：它由 $b$ 的符号与 $\vect{P},\vect{Q}$ 的取向共同决定，不能只看 $b$。要判断，取 $x=0$ 处的点 $\vect{P}$，看导数 $\vect{U}'(0)=a\vect{P}-b\vect{Q}$ 指向哪边。

\subhdr{8.3 重根：$g=m$ 还是刚才的直线，$g<m$ 会弯}

\begin{itemize}
  \item $g=m$：解全是 $\mathrm{e}^{\lambda x}$ 乘特征向量，轨道与 8.1 一样是射线，只是同一个 $\lambda$ 下有好几条。
  \item $g<m$：解里多出 $x\mathrm{e}^{\lambda x}\vect{V}$。$x\to\pm\infty$ 时 $x$ 项压过常数项，所以轨道\textbf{最终平行于 $\vect{V}$}，但中途会偏离；$\lambda<0$ 时整体仍缩向原点，只是"绕着直线弯进来"；$\lambda>0$ 时向外甩。
\end{itemize}

\textbf{一句话区分}：$g<m$ 的重根不会给你新的方向，只给你"沿同一方向多走一会儿"的自由——这就是为什么解里出现 $x$，而轨道不出现新的直线。

\subhdr{8.4 与 L21--23 的 $\tau$--$\Delta$ 判据对接}

$n=2$ 时，两个特征值满足

\[
\tau=\lambda_1+\lambda_2=\operatorname{tr}A,\qquad
\Delta=\lambda_1\lambda_2=\det A,\qquad
D=\tau^2-4\Delta=(\lambda_1-\lambda_2)^2\ge0 .
\]

于是本节的三种情形可以逐条翻译成代数判据：

\begin{itemize}
  \item $D>0$：两个不同实根 → 两条特征方向（结点，或 $\Delta<0$ 时的鞍点）；
  \item $D<0$：一对共轭复根 → 螺旋（焦点；$\tau=0$ 时是中心）；
  \item $D=0$：重根 → 星形结点（$g=2$）或退化结点（$g=1$）。
\end{itemize}

\textbf{这不是新知识，是同一件事换了语言}：L22--23 作业里要求你算 $\tau,\Delta,D$ 再查表，查的就是这里 8.1–8.3 的三种情形；螺旋的朝向由 $\tau$ 的符号定（$\tau<0$ 向内、$\tau>0$ 向外），与 $a=\tau/2$ 的符号一致。

\subhdr{8.5 稳定性的一句话}

把上面收成一句：

\[
\text{渐近稳定}\iff\text{所有 }\operatorname{Re}\lambda<0;
\qquad
\text{不稳定}\iff\text{存在 }\operatorname{Re}\lambda>0 .
\]

若出现 $\operatorname{Re}\lambda=0$（纯虚根或零根），线性化\textbf{给不出结论}——轨道可能是中心，也可能是缓慢旋出/旋进的细焦点，要看高阶项。这一条是 L22--23 的入口，本讲义只点到，不展开。

%=====================================================================
\hdr{§10}{作业（p64）}

课件作业四题，都是"解 $\vect{Y}'=A\vect{Y}$"。题干与考点如下，详解在附录。

\qbreak
\qstem{（课件 p64 第 1 题）$A=\mat{1&1&-3\\3&2&-3\\1&0&1}$，求 $\vect{Y}'=A\vect{Y}$ 的通解。}
\par\textbf{考点}：重根 + $g<m$ 的链；三维链长 2。
\par\textbf{关键一步}：特征值 $1$（$m=2$）与 $2$（$m=1$）；对 $\lambda=1$ 算 $g=3-\rank(A-E)$，若 $g=1<m$ 就要接一个广义特征向量。
\par\textbf{详解}：附录 A.18，第 \pageref{c2:app:hw1} 页。

\qbreak
\qstem{（课件 p64 第 2 题）$A=\mat{-2&1\\-5&2}$，求通解。}
\par\textbf{考点}：纯虚特征值与实值化的写法。
\par\textbf{关键一步}：先看迹（$=0$）与行列式（$=1$），猜出 $\lambda=\pm\mathrm{i}$；只算 $\lambda=\mathrm{i}$ 一个，取实虚部。
\par\textbf{详解}：附录 A.19，第 \pageref{c2:app:hw2} 页。

\qbreak
\qstem{（课件 p64 第 3 题）$A=\mat{-2&-1&1\\-3&-3&2\\-4&-3&2}$，求通解。}
\par\textbf{考点}：三重根、$g=1$，需要长度 3 的链。
\par\textbf{关键一步}：$\lambda=-1$，$m=3$；算 $g=3-\rank(A+E)=1$，于是要解两次 $(A+E)\vect{W}=\vect{V}$、$(A+E)\vect{U}=\vect{W}$。
\par\textbf{详解}：附录 A.20，第 \pageref{c2:app:hw3} 页。

\qbreak
\qstem{（课件 p64 第 4 题）$A=\mat{1&1\\9&1}$，求通解。}
\par\textbf{考点}：相异实根，Case 1 的基本操作。
\par\textbf{关键一步}：特征方程 $\lambda^2-2\lambda-8=0$，两根 $4$ 与 $-2$；各配一个特征向量。
\par\textbf{详解}：附录 A.21，第 \pageref{c2:app:hw4} 页。

%=====================================================================
\hdr{§11}{易错点}

\begin{enumerate}
  \item \textbf{重根不一定出现 $x$}。$g=m$ 时全是纯指数；只有 $g<m$ 才出现 $x\mathrm{e}^{\lambda x}$。判据是 $g$，不是 $m$（§6、§7）。
  \item \textbf{复根要把实虚部都留下}。只写 $\vect{U}$ 少一维；只写复值解不合题意（§4.5）。
  \item \textbf{链只能从"能接上"的特征向量起头}。要比较 $\rank(A-\lambda E)$ 与 $\rank(A-\lambda E\mid\vect{V})$（§7.3）。
  \item \textbf{广义特征向量有无穷多个，选一个就行}。方程 $(A-\lambda E)\vect{V}^{(1)}=\vect{V}$ 的解不唯一，差一个特征向量；任取一个非零解都可以。
  \item \textbf{解的总数必须等于 $n$}。这是最便宜的自检（§8 第 5 步）。
  \item \textbf{特征向量可以乘任意非零常数}，但同一个 $\lambda$ 的多个特征向量必须线性无关，不能取成倍数。
\end{enumerate}

%=====================================================================
\hdr{§12}{缺口清单}

\gapitem{机制缺口}{p4--p5 “Seek a solution in the form $Y(x)=Ve^{\lambda x}$” 直接猜形状，未说明为什么可以这样猜。本讲义 §2.1--§2.2 补 $n=1$ 类比与代入约化。}
\gapitem{术语未定义}{p7 “If $\lambda_j$ has $g_j$ linearly independent eigenvectors, then $\lambda$ has geometric multiplicity $g_j$” 只给名字，未给算法，且“anywhere from 1 to k”没有解释。本讲义 §1.2 补 $g=n-\rank(A-\lambda E)$ 与 $1\le g\le m$ 的理由。}
\gapitem{证明缺口}{p10 定理 2.1 只陈述“distinct $\Rightarrow$ independent”，未证。本讲义 §3.1 补两条证明（指数线性无关论证与 Wronskian 论证）。}
\gapitem{机制缺口}{p19 “According to the principle of superposition” 一句带过实虚部是解，未说明为什么。本讲义 §4.2 补线性算子对实部虚部分离。}
\gapitem{形式改写}{p20 一般公式只用 $\operatorname{Re}$、$\operatorname{Im}$ 记号给出。本讲义 §4.3 展开成显式 $\cos b x$、$\sin b x$ 形式，便于代初值。}
\gapitem{排版笔误}{p26 写 $m_2x_2''=-k_2x_1-(k_2+k_3)x_2$，与 L13 p4 的 $+k_2x_1$ 及 p27 矩阵第四行 $+4/3$ 矛盾。本讲义 §5.1 指出应为 $+k_2x_1$。}
\gapitem{排版笔误}{p36 通解首行把 $c_1\vect{V}_1$ 配 $\mathrm{e}^{3x}$、$c_2\vect{V}_2$ 配 $\mathrm{e}^{2x}$，与 $\lambda_1=2$、$\lambda_2=3$ 矛盾（下一行分量式才对）。本讲义 A.5 按 $\lambda_1=2$ 重写。}
\gapitem{代数跳步}{p41 “Comparing terms” 跳过“$x\mathrm{e}^{\lambda x}$ 项自动相等、只剩常数项”这一步。本讲义 §7.2 展开比较过程。}
\gapitem{机制缺口}{p40 只说 $1+\lambda x=\lambda x$ “doesn't work”，未解释失败原因。本讲义 §7.1 补：$\vect{V}_1$ 已在核里，乘 $x$ 不产生新方向。}
\gapitem{适用条件缺失}{p43 未说明链的起点必须在 $\operatorname{Range}(A-\lambda E)$ 中；p61 的秩比较没有文字解释。本讲义 §7.3 补秩判据。}

%=====================================================================
\hdr{§13}{三句话}

\begin{enumerate}
  \item 全部内容来自一个猜测：$\vect{Y}=\vect{V}\mathrm{e}^{\lambda x}$，代进方程就变成 $A\vect{V}=\lambda\vect{V}$。
  \item 分三支看特征值：相异实根各乘各的指数；共轭复根取实部虚部；重根看 $g$——$g=m$ 仍是纯指数，$g<m$ 要接广义特征向量的链。
  \item 判断走哪支的钥匙是 $\rank(A-\lambda E)$ 给出的 $g$，不是 $m$；链能否接上，也要看秩。
\end{enumerate}

%=====================================================================
\hdr{§14}{与课件的页码对照}

\begin{minipage}{\linewidth}
\tabtop
\begin{tabular}{@{}p{0.12\linewidth}p{0.44\linewidth}p{0.34\linewidth}@{}}
\thead{本讲义} & \thead{内容} & \thead{课件页}\\
\end{tabular}
\tabmid
\begin{tabular}{@{}p{0.12\linewidth}p{0.44\linewidth}p{0.34\linewidth}@{}}
§1 & 特征值、特征方程、$m$ 与 $g$ & p3--p7\\
§2 & 为什么猜 $\vect{V}\mathrm{e}^{\lambda x}$ & p4--p5\\
§3 & Case 1：相异实根 & p9--p15\\
§4 & Case 2：共轭复根 & p16--p25\\
§5 & 四维弹簧-质量例 & p26--p29\\
§6 & Case 3a：$g=m$ & p30--p38\\
§7 & Case 3b：$1\le g<m$ & p39--p45\\
§9 & 六步检查表 & p37--p38\\
§10 & 作业四题 & p64\\
\end{tabular}
\tabbot
\end{minipage}

%=====================================================================
\newpage
\chapterbanner{3}{矩阵指数}{三种情形压成一个 $\mathrm{e}^{Ax}$}{课件 \texttt{16\_17\_matrix\_exponential.pdf}（44 页，2026-10-13）}
\callout{\textbf{先读这一段。}L14--15 把 $\vect{Y}'=A\vect{Y}$ 分三种情形解完了。这一讲问一个更狠的问题：\textbf{能不能只写一个式子，把三种情形都盖住？}能，就是 $e^{Ax}$。它的定义不依赖特征值（§1），所以对相异根、复根、重根一视同仁；但它的\textbf{计算}要分路线：可对角化时用 §3，不可对角化或维度小用 §4 的 Cayley--Hamilton，已经有一组解时用 §5 的 $\Phi$。§6 给一张选路表。}

%=====================================================================
\hdr{§0}{三条路线，一张图}

\begin{minipage}{\linewidth}
\tabtop
\begin{tabular}{@{}p{0.12\linewidth}p{0.24\linewidth}p{0.36\linewidth}p{0.10\linewidth}@{}}
\thead{路线} & \thead{前提} & \thead{做什么} & \thead{课件页}\\
\end{tabular}
\tabmid
\begin{tabular}{@{}p{0.12\linewidth}p{0.24\linewidth}p{0.36\linewidth}p{0.10\linewidth}@{}}
定义 & 无 & $e^{Ax}=\sum A^kx^k/k!$，证明它是解 & p3--p6\\
对角化 & 所有特征值 $g=m$ & $A=V\Lambda V^{-1}$，$e^{Ax}=Ve^{\Lambda x}V^{-1}$ & p8--p15\\
Cayley--Hamilton & 无（但计算量随 $n$ 涨） & 化成 $n-1$ 次矩阵多项式，用特征值定系数 & p16--p32\\
基本解矩阵 & 已有一组解 & $e^{Ax}=\Phi(x)\Phi^{-1}(0)$ & p39--p43\\
\end{tabular}
\tabbot
\end{minipage}

这张表是路线台账，不承载推导；每一行的条件与算法在对应小节展开。

\textbf{一句话读法}：定义只负责"存在且正确"，不负责"算得出来"（p6 自己说了：级数对实际计算没用）。真正的计算是后三条，它们给出的是同一个矩阵的不同表达式。

%=====================================================================
\hdr{§1}{定义：从标量指数到矩阵指数}

\subhdr{1.1 标量的类比（p4）}

标量初值问题 $x'=ax$、$x(0)=x_0$ 的解是 $x=x_0\mathrm{e}^{at}$。矩阵情形 $\vect{Y}'=A\vect{Y}$、$\vect{Y}(0)=\vect{Y}_0$ 能不能照抄成 $\vect{Y}=\mathrm{e}^{Ax}\vect{Y}_0$？

能——但前提是把 $\mathrm{e}^{Ax}$ 定义成一个\textbf{矩阵}。标量指数有泰勒级数 $\mathrm{e}^{at}=\sum a^nt^n/n!$，把这个式子里的 $a$ 换成矩阵 $A$ 就行。

\subhdr{1.2 级数定义（p5，定理 1.1）}

对 $n\times n$ 常矩阵 $A$，定义

\[
\mathrm{e}^{Ax}=E+Ax+\frac{A^2x^2}{2!}+\cdots+\frac{A^kx^k}{k!}+\cdots=\sum_{k=0}^{\infty}\frac{A^kx^k}{k!},
\qquad A^0=E .
\]

这个级数对每个 x 都收敛。课件只说了 “convergent”，理由补在这里：设 A 的矩阵范数记为 ‖A‖，则
\[
\left\|\frac{A^kx^k}{k!}\right\|\le\frac{\|A\|^k|x|^k}{k!},
\]
右端是标量指数级数的通项，处处收敛，故矩阵级数绝对收敛。

\textbf{注意 $\mathrm{e}^{Ax}$ 是 $n\times n$ 矩阵}，不是数，也不是向量；$\vect{Y}=\mathrm{e}^{Ax}\vect{Y}_0$ 是"矩阵乘向量"。

\subhdr{1.3 它为什么是解（p5）}

把 $\vect{Y}=\mathrm{e}^{Ax}\vect{C}$ 代进方程，只需验

\[
\frac{\mathrm{d}}{\mathrm{d}x}\Bigl(\sum_{k=0}^{\infty}\frac{A^kx^k}{k!}\Bigr)\vect{C}
=\sum_{k=1}^{\infty}\frac{kA^kx^{k-1}}{k!}\vect{C}
=A\sum_{k=1}^{\infty}\frac{A^{k-1}x^{k-1}}{(k-1)!}\vect{C}
=A\,\mathrm{e}^{Ax}\vect{C}.
\]

第一处等号是逐项求导（级数在任意有限区间上一致收敛，可以逐项求导），第二处把 $k/k!=1/(k-1)!$ 并用 $A^k=AA^{k-1}$。于是

\[
\bigl(\mathrm{e}^{Ax}\vect{C}\bigr)'-A\mathrm{e}^{Ax}\vect{C}=\vect{0},
\]

对任意常向量 $\vect{C}$ 成立。所以\textbf{通解就是} $\vect{Y}=\mathrm{e}^{Ax}\vect{C}$，其中 $\vect{C}$ 是 $n$ 维任意常向量。

\subhdr{1.4 初值问题（p6）}

若给的是 $\vect{Y}(x_0)=\vect{Y}_0$，把 $x_0$ 代进通解得 $\vect{C}=\mathrm{e}^{-Ax_0}\vect{Y}_0$，于是

\[
\vect{Y}(x)=\mathrm{e}^{Ax}\mathrm{e}^{-Ax_0}\vect{Y}_0=\mathrm{e}^{A(x-x_0)}\vect{Y}_0 .
\]

这里用了性质 $\mathrm{e}^{At}\mathrm{e}^{As}=\mathrm{e}^{A(t+s)}$（同底数的矩阵指数可交换，见 §2）。

\subhdr{1.5 级数不能直接算（p6 末句）}

$\mathrm{e}^{Ax}$ 的元素是无穷级数，逐项算不现实。所以后面三条路线都是"把无穷级数换成有限表达式"：对角化把它换成 $n$ 个指数，C--H 把它换成 $n$ 个矩阵的线性组合，$\Phi$ 直接跳过级数。

%=====================================================================
\hdr{§2}{性质：能用的五条与一条容易忽略的}

\subhdr{2.1 五条基本性质（p7，性质 1.1）}

\begin{enumerate}
  \item $\mathrm{e}^{A\cdot0}=E$。代 $x=0$，级数只剩第一项。
  \item $\mathrm{e}^{At}\mathrm{e}^{As}=\mathrm{e}^{A(t+s)}$。两个同底数的级数相乘，二项式系数正好凑成 $(t+s)^k/k!$。
  \item 若 $AB=BA$，则 $\mathrm{e}^{A+B}=\mathrm{e}^A\mathrm{e}^B$。\textbf{条件不能去}：见下面的反例。
  \item $\mathrm{e}^{-At}=(\mathrm{e}^{At})^{-1}$。由第 2 条令 $s=-t$ 得 $\mathrm{e}^{At}\mathrm{e}^{-At}=E$。
  \item $\mathrm{e}^{rtE}=\mathrm{e}^{rt}E$。因为 $(rtE)^k=(rt)^kE$，级数逐项都带 $E$。
\end{enumerate}

\textbf{第 3 条的反例}（课件没给）：取

\[
A=\mat{0&1\\0&0},\qquad B=\mat{0&0\\1&0},\qquad AB=\mat{1&0\\0&0},\quad BA=\mat{0&0\\0&1}.
\]

计算得 $\mathrm{e}^{A}\mathrm{e}^{B}$ 与 $\mathrm{e}^{A+B}$ 的 $(1,1)$ 元素分别为 $2$ 与 $\cosh1\approx1.543$，不相等。所以 \textbf{$A,B$ 不可交换时，指数不能拆}。

\subhdr{2.2 特征向量也是 $\mathrm{e}^{A}$ 的特征向量（p7，性质 1.2）}

若 $A\vect{V}=\lambda\vect{V}$，则

\[
\mathrm{e}^{A}\vect{V}=\sum_{k=0}^{\infty}\frac{A^k}{k!}\vect{V}
=\sum_{k=0}^{\infty}\frac{\lambda^k}{k!}\vect{V}
=\mathrm{e}^{\lambda}\vect{V}.
\]

中间一步用了 $A^k\vect{V}=\lambda^k\vect{V}$（每次作用一次 $A$ 就提出一个 $\lambda$）。所以 $\vect{V}$ 是 $\mathrm{e}^{A}$ 的特征向量，特征值是 $\mathrm{e}^{\lambda}$。

\textbf{这条性质很重要}：它说明"矩阵指数不改变特征方向，只改变特征值（$\lambda\mapsto\mathrm{e}^{\lambda}$）"。L14--15 的三种情形，在这里变成"$\mathrm{e}^{Ax}$ 的三个特征值 $1$、$\mathrm{e}^{\lambda_1x}$、……"。

%=====================================================================
\hdr{§3}{路线一：对角化（要求 $g=m$）}

\subhdr{3.1 前提（p9）}

若 $A$ 的每个特征值都满足 $g_j=m_j$（即 $A$ 有 $n$ 个线性无关的特征向量），把特征向量按列排成

\[
V=(\vect{V}_1,\dots,\vect{V}_n),
\]

则

\[
V^{-1}AV=\Lambda=\diag(\underbrace{\lambda_1,\dots,\lambda_1}_{m_1},\lambda_{m_1+1},\dots,\lambda_n).
\]

\textbf{这就是 L14--15 的 Case 1 与 Case 3a}：特征向量够用时才能这样写。$g<m$ 时 $V$ 不满秩，$V^{-1}$ 不存在，这条路走不通——那是 §4 的活。

\subhdr{3.2 从级数到 $Ve^{\Lambda x}V^{-1}$（p10）}

把 $A=V\Lambda V^{-1}$ 代进级数。关键是幂次：

\[
A^k=(V\Lambda V^{-1})(V\Lambda V^{-1})\cdots(V\Lambda V^{-1})=V\Lambda^kV^{-1},
\]

中间相邻的 $V^{-1}V$ 逐个抵消，只剩两头的 $V$ 与 $V^{-1}$。于是

\[
\mathrm{e}^{Ax}=\sum_{k=0}^{\infty}\frac{A^kx^k}{k!}
=V\Bigl(\sum_{k=0}^{\infty}\frac{\Lambda^kx^k}{k!}\Bigr)V^{-1}
=V\mathrm{e}^{\Lambda x}V^{-1}.
\]

课件 p10 把这一步写成一长串，核心就是"$V$ 与 $V^{-1}$ 成对抵消"。

\subhdr{3.3 对角矩阵的指数（p11）}

$\Lambda$ 是对角矩阵，$\Lambda^k=\diag(\lambda_1^k,\dots,\lambda_n^k)$，逐项求和：

\[
\mathrm{e}^{\Lambda x}=\diag(\mathrm{e}^{\lambda_1x},\dots,\mathrm{e}^{\lambda_nx}).
\]

所以

\[
\boxed{\ \mathrm{e}^{Ax}=V\,\diag(\mathrm{e}^{\lambda_1x},\dots,\mathrm{e}^{\lambda_nx})\,V^{-1}\ }.
\]

\subhdr{3.4 通解（p12）}

\[
\vect{Y}=\mathrm{e}^{Ax}\vect{C}=V\mathrm{e}^{\Lambda x}V^{-1}\vect{C}
=V\mathrm{e}^{\Lambda x}\tilde{\vect{C}}
=\sum_{k=1}^{n}\tilde c_k\vect{V}_k\mathrm{e}^{\lambda_kx}.
\]

这正是 L14--15 的答案。所以对角化路线没有新东西，它只是把"拼解"写成了矩阵乘法。

\textbf{代价}：要算 $V^{-1}$。三维以上手算很烦，而且只有可对角化时才能用。

%=====================================================================
\hdr{§4}{路线二：Cayley--Hamilton（不需要可对角化）}

\subhdr{4.1 定理 3.1：方阵满足自己的特征多项式（p17）}

设 $A$ 的特征多项式为

\[
\Delta(\lambda)=\det(A-\lambda E)=b_n\lambda^n+b_{n-1}\lambda^{n-1}+\cdots+b_1\lambda+b_0,
\]

则

\[
b_nA^n+b_{n-1}A^{n-1}+\cdots+b_1A+b_0E=[0].
\]

\textbf{它的用处}：$A^n$ 可以用 $E,A,\dots,A^{n-1}$ 表示。于是\textbf{任何} $A$ 的多项式都能压到 $n-1$ 次以内（p19 的降阶）：

\[
p(A)=q(A)\Delta(A)+r(A)=r(A),\qquad \deg r\le n-1 .
\]

\subhdr{4.2 定理 3.2：$\mathrm{e}^{Ax}$ 是 $n-1$ 次矩阵多项式（p18）}

把上面的结论用到 $Ax$ 上：

\[
\mathrm{e}^{Ax}=\alpha_{n-1}A^{n-1}x^{n-1}+\alpha_{n-2}A^{n-2}x^{n-2}+\cdots+\alpha_1Ax+\alpha_0E,
\]

其中 $\alpha_0,\dots,\alpha_{n-1}$ 是 $x$ 的函数，要由特征值定。$n=2$ 时是 $\alpha_1Ax+\alpha_0E$，$n=3$ 时是 $\alpha_2A^2x^2+\alpha_1Ax+\alpha_0E$。

\subhdr{4.3 定系数：把特征值代进去（p21--p23）}

设 $r(\lambda)=\alpha_{n-1}\lambda^{n-1}+\cdots+\alpha_1\lambda+\alpha_0$。由 $f(A)=r(A)$（$f$ 是解析函数），对每个特征值有

\[
f(\lambda_i)=r(\lambda_i).
\]

取 $f(\lambda)=\mathrm{e}^{\lambda}$，得到

\[
\mathrm{e}^{\lambda_i}=r(\lambda_i)=\sum_{k=0}^{n-1}\alpha_k\lambda_i^k .
\]

$n$ 个特征值给 $n$ 个方程，解出 $n$ 个 $\alpha_k$。

\warnbox{\textbf{这里有一个必须点明的记号坑}：上式里的 $\lambda_i$ 是\textbf{矩阵 $Ax$ 的特征值}，不是 $A$ 的特征值。$A$ 的特征值是 $\mu_i$，则 $Ax$ 的特征值是 $\lambda_i=\mu_ix$。课件 p22 写 $\mathrm{e}^{\lambda_i}$ 时没明说这一点，读者容易把 $x$ 丢掉。\textbf{记住：}代进方程的指数是 $\mathrm{e}^{\mu_ix}$，而方程里的多项式变量是 $\mu_ix$。}

\subhdr{4.4 重根怎么办（p28--p29，定理 3.3）}

若 $\lambda_i$ 是 $m$ 重根，$n$ 个特征值给不出 $n$ 个独立方程（重复的方程相同）。补的办法是\textbf{求导}：

\[
\mathrm{e}^{\lambda_i}=r(\lambda_i),\quad
\mathrm{e}^{\lambda_i}=r'(\lambda_i),\quad
\mathrm{e}^{\lambda_i}=r''(\lambda_i),\ \dots,\
\mathrm{e}^{\lambda_i}=r^{(m-1)}(\lambda_i).
\]

\textbf{为什么可以求导}：$\Delta(\lambda)$ 在 $m$ 重根处满足 $\Delta(\lambda_i)=\Delta'(\lambda_i)=\cdots=\Delta^{(m-1)}(\lambda_i)=0$，而 $f(\lambda)=\Delta(\lambda)q(\lambda)+r(\lambda)$ 两边求导后，含 $\Delta$ 的项在 $\lambda_i$ 处全为零，只剩 $f^{(j)}(\lambda_i)=r^{(j)}(\lambda_i)$。

\subhdr{4.5 四步总结（p32）}

\begin{enumerate}
  \item 求 $Ax$ 的特征值（即 $A$ 的特征值乘 $x$），记下各自的重数。
  \item 设 $\mathrm{e}^{Ax}=\sum_{k=0}^{n-1}\alpha_kA^kx^k$，并写出 $r(\lambda)=\sum\alpha_k\lambda^k$。
  \item 对每个特征值按重数写方程：单根写 $1$ 个，$m$ 重根写 $m$ 个（含导数）。
  \item 解出 $\alpha_k$，代回第 2 步。
\end{enumerate}

课件 p30--p31 的例 6 就是"重数 3 与 1"的方程组建法：$4$ 个方程解 $4$ 个 $\alpha$。它只到"列方程"为止，没解下去——因为没给具体的 $A$。

%=====================================================================
\hdr{§5}{路线三：基本解矩阵 $\Phi$}

\subhdr{5.1 定义（p40）}

若 $\{\vect{Y}_1,\dots,\vect{Y}_n\}$ 是 $\vect{Y}'=A\vect{Y}$ 的基本解组，把它们按列排成

\[
\Phi(x)=(\vect{Y}_1(x),\dots,\vect{Y}_n(x)),
\]

则称 $\Phi$ 为\textbf{基本解矩阵}。通解是 $\vect{Y}=\Phi(x)\vect{C}$。

\textbf{它与 L14--15 的关系}：L14--15 找到的 $n$ 个解，按列排起来就是 $\Phi$。所以 $\Phi$ 不是新对象，只是把"解组"写成矩阵。

\subhdr{5.2 初值问题（p42）}

$\vect{Y}(x_0)=\Phi(x_0)\vect{C}=\vect{Y}_0$。因为 $\det\Phi(x_0)\neq0$（基本解组的 Wronskian 非零），可以解出 $\vect{C}=\Phi^{-1}(x_0)\vect{Y}_0$，于是

\[
\vect{Y}(x)=\Phi(x)\Phi^{-1}(x_0)\vect{Y}_0 .
\]

\subhdr{5.3 与 $\mathrm{e}^{Ax}$ 的关系（p42--p43）}

同一个初值问题，§1.4 给出的解是 $\vect{Y}=\mathrm{e}^{A(x-x_0)}\vect{Y}_0$。\textbf{初值问题的解唯一}（L13 定理 2.1），所以两个表达式必须相等：

\[
\Phi(x)\Phi^{-1}(x_0)=\mathrm{e}^{A(x-x_0)},\qquad\text{特别地}\quad \Phi(x)\Phi^{-1}(0)=\mathrm{e}^{Ax}.
\]

若再取 $\Phi(0)=E$，就有 $\mathrm{e}^{Ax}=\Phi(x)$。

\textbf{定理 4.1}（p43）：$\mathrm{e}^{Ax}$ 的列构成基本解组，所以 $\mathrm{e}^{Ax}$ 本身就是基本解矩阵。这条定理是上面等式的直接推论：$\mathrm{e}^{Ax}$ 在 $x=0$ 处等于 $E$，取 $\Phi=\mathrm{e}^{Ax}$ 即满足 $\Phi(0)=E$。

\subhdr{5.4 三种写法的关系}

\begin{minipage}{\linewidth}
\tabtop
\begin{tabular}{@{}p{0.22\linewidth}p{0.36\linewidth}p{0.32\linewidth}@{}}
\thead{写法} & \thead{形式} & \thead{什么时候方便}\\
\end{tabular}
\tabmid
\begin{tabular}{@{}p{0.22\linewidth}p{0.36\linewidth}p{0.32\linewidth}@{}}
级数 & $\sum A^kx^k/k!$ & 只用于证明性质，不用于计算\\
对角化 & $V\mathrm{e}^{\Lambda x}V^{-1}$ & 可对角化且维数小\\
C--H & $\sum_{k=0}^{n-1}\alpha_kA^kx^k$ & 任何矩阵，$n\le3$ 最舒服\\
基本解矩阵 & $\Phi(x)\Phi^{-1}(0)$ & 已经有一组解，或要接非齐次\\
\end{tabular}
\tabbot
\end{minipage}

这张表是对照，不承载推导；四条给的是同一个矩阵，选哪条只看手头有什么。

\subhdr{5.5 三种情形下 $\mathrm{e}^{Ax}$ 长什么样}

三条路线算的是同一个矩阵，但它在三种特征值情形下的"外貌"不同。这一小节把外貌说清，也顺便解释为什么 L14--15 的答案长那样。

\textbf{(a) 相异实根}：对角化成立，

\[
\mathrm{e}^{Ax}=V\diag(\mathrm{e}^{\lambda_1x},\dots,\mathrm{e}^{\lambda_nx})V^{-1}
\ \Longrightarrow\
(\mathrm{e}^{Ax})_{ij}=\sum_{k=1}^{n}c_{ijk}\mathrm{e}^{\lambda_kx},
\]

每个元素都是 $n$ 个纯指数的线性组合。系数 $c_{ijk}$ 由 $V$ 与 $V^{-1}$ 的元素相乘得到——这就是附录里那些矩阵的元素总是"几个 $\mathrm{e}^{\lambda x}$ 相加"的原因。

\textbf{(b) 共轭复根}：$V$ 与 $V^{-1}$ 都含复数，但乘出来的 $\mathrm{e}^{Ax}$ 是\textbf{实矩阵}，元素形如

\[
\mathrm{e}^{ax}\bigl(A\cos bx+B\sin bx\bigr).
\]

共轭的两列配共轭的两行，虚部成对抵消。例 5（A.4）的旋转矩阵、例 8（A.6）的元素都是这个形状。

\textbf{(c) 重根，$g=m$}：对角化仍成立，只是对角线上有重复的 $\mathrm{e}^{\lambda x}$。外貌与 (a) 相同。

\textbf{(d) 重根，$g<m$}：$V$ 不可逆，不能对角化，这时

\[
\mathrm{e}^{Ax}=\mathrm{e}^{\lambda x}P(x),
\]

其中 $P(x)$ 是\textbf{元素为多项式}的矩阵，次数不超过 $m-1$。例 7（A.5）的 $\mathrm{e}^{-2x}\mat{2x+1&-4x&0\\x&1-2x&0\\x&-2x&1}$ 就是：$P$ 是一次多项式矩阵，$\lambda=-2$ 的重数是 $3$，次数上限 $3-1=2$（这里实际只到 $1$）。

\textbf{一句话总结}：$\mathrm{e}^{Ax}$ 的每个元素都是"\textbf{多项式乘指数} $\mathrm{e}^{\lambda x}$"，次数\textbf{小于}该 $\lambda$ 的重数。相异根时多项式是常数，重根时才会出现 $x$、$x^2$——这正是 L14--15 里 $x\mathrm{e}^{\lambda x}$ 项的来源，只是那里按列写，这里按矩阵写。

%=====================================================================
\hdr{§6}{拿到 $A$ 时怎么选}

\begin{enumerate}
  \item \textbf{先看 $n$}。$n=2$ 且特征值好算：直接 C--H，两个方程两个未知数，最快。
  \item \textbf{再看是否可对角化}。$g=m$ 时对角化省事（尤其 $n\ge3$ 且有重根时，C--H 要解含导数的方程组）。
  \item \textbf{$g<m$ 时只能走 C--H}（或直接用 L14--15 的链式解，再按列排成 $\Phi$）。
  \item \textbf{已经有解组}（例如刚用 L14--15 的方法解过）：直接排成 $\Phi$，再用 $\Phi\Phi^{-1}(0)$ 得 $\mathrm{e}^{Ax}$。求逆比重解一遍便宜。
\end{enumerate}

%=====================================================================
\hdr{§7}{作业（p44）}

课件作业三题，都是"求 $\mathrm{e}^{Ax}$ 并解 $\vect{Y}'=A\vect{Y}$"。

\qbreak
\qstem{（课件 p44 第 1 题）$A=\mat{-6&-4\\8&6}$。求 $\mathrm{e}^{Ax}$ 与通解。}
\par\textbf{考点}：C--H 在 $2\times2$、相异实根上的标准流程。
\par\textbf{关键一步}：$\operatorname{tr}A=0$、$\det A=-4$，特征方程 $\lambda^2-4=0$，$\lambda=\pm2$（乘 $x$ 后为 $\pm2x$）；设 $\mathrm{e}^{Ax}=\alpha_1Ax+\alpha_0E$。
\par\textbf{详解}：附录 A.28，第 \pageref{c3:app:hw1} 页。

\qbreak
\qstem{（课件 p44 第 2 题）$A=\mat{-4&4\\-5&4}$。求 $\mathrm{e}^{Ax}$ 与通解。}
\par\textbf{考点}：纯虚特征值；指数方程组的复数解法。
\par\textbf{关键一步}：$\operatorname{tr}A=0$、$\det A=4$，$\lambda=\pm2\mathrm{i}$；解 $\mathrm{e}^{\pm2\mathrm{i}x}=\pm2\mathrm{i}x\alpha_1+\alpha_0$，用欧拉公式取实虚部。
\par\textbf{详解}：附录 A.29，第 \pageref{c3:app:hw2} 页。

\qbreak
\qstem{（课件 p44 第 3 题）$A=\mat{-24&0&0\\4&-2&0\\1&0&2}$。求 $\mathrm{e}^{Ax}$ 与通解。}
\par\textbf{考点}：$3\times3$、三个相异实根；下三角矩阵的指数仍是下三角。
\par\textbf{关键一步}：下三角阵的特征值就是对角元 $-24,-2,2$，互异；设 $\mathrm{e}^{Ax}=\alpha_2A^2x^2+\alpha_1Ax+\alpha_0E$，三个方程解三个 $\alpha$。
\par\textbf{详解}：附录 A.30，第 \pageref{c3:app:hw3} 页。

%=====================================================================
\hdr{§8}{易错点}

\begin{enumerate}
  \item \textbf{$\mathrm{e}^{Ax}$ 是矩阵，不是数}。写答案时不能把 $\mathrm{e}^{Ax}$ 当成标量去"乘开"。
  \item \textbf{$\mathrm{e}^{A+B}=\mathrm{e}^A\mathrm{e}^B$ 只在 $AB=BA$ 时成立}（§2.1 的反例）。
  \item \textbf{代进方程的指数是 $\mathrm{e}^{\mu_ix}$，多项式变量是 $\mu_ix$}。C--H 里最容易把 $x$ 丢掉（§4.3 的警告）。
  \item \textbf{重根要补导数方程}，否则方程数不够，解不唯一（§4.4）。
  \item \textbf{对角化只对 $g=m$ 有效}。$g<m$ 时 $V$ 不可逆，硬写 $V^{-1}$ 会出错（§3.1）。
  \item \textbf{$\Phi(x)\Phi^{-1}(x_0)=\mathrm{e}^{A(x-x_0)}$ 依赖"初值问题解唯一"}，不是恒等式（§5.3）。
\end{enumerate}

%=====================================================================
\hdr{§9}{缺口清单}

\gapitem{证明缺口}{p5 “The sum of (2) is convergent” 只给结论。本讲义 §1.2 补范数比较的收敛理由。}
\gapitem{适用条件缺失}{p7 “If $AB=BA$, then $\mathrm{e}^{A+B}=\mathrm{e}^A\mathrm{e}^B$” 给了条件但没给反例，读者容易以为总成立。本讲义 §2.1 补 $2\times2$ 反例。}
\gapitem{证明缺口}{p7 性质 1.2 未证。本讲义 §2.2 补 $A^k\vect{V}=\lambda^k\vect{V}$ 的推导。}
\gapitem{代数跳步}{p10 级数推导中 $(V\Lambda V^{-1})^k=V\Lambda^kV^{-1}$ 被写成一带而过。本讲义 §3.2 说明相邻 $V^{-1}V$ 抵消。}
\gapitem{证明缺口}{p17 定理 3.1（Cayley--Hamilton）只陈述。本讲义 §4.1 给出它的用处与降阶结论，未重证定理本身（属线性代数经典结论）。}
\gapitem{记号混乱}{p22 “$\mathrm{e}^{\lambda_i}=\sum\alpha_k\lambda_i^k$” 中的 $\lambda_i$ 是 $Ax$ 的特征值，未说明它含 $x$。本讲义 §4.3 明确写 $\lambda_i=\mu_ix$。}
\gapitem{机制缺口}{p28 “the first $(m-1)$ derivatives of $\Delta(s)$ all vanish at the eigenvalues, therefore $f^{(j)}(\lambda_i)=r^{(j)}(\lambda_i)$” 跳过推导。本讲义 §4.4 补“含 $\Delta$ 的项在根处为零”。}
\gapitem{解答不完整}{p30--p31 例 6 只列四个方程，没有解出 $\alpha$。本讲义 §4.5 说明这是“列方程训练”，并给出重根补导数的一般规则。}
\gapitem{排版笔误}{p14 的 $V^{-1}$ 第 1 行第 2 列印成 $0$，应为 $-1$（用 $V\cdot V^{-1}=E$ 可验）。本讲义 A.1 给出正确逆矩阵。}
\gapitem{排版笔误}{p27 例 5 通解末行写成 $c_1\binom{\cos x}{-\sin x}+c_2\binom{-\sin x}{\cos x}$，第二项应为 $\binom{\sin x}{\cos x}$——用 $\mathrm{e}^{Ax}\vect{C}$ 直接乘可得。本讲义 A.4 更正。}

%=====================================================================
\hdr{§10}{三句话}

\begin{enumerate}
  \item $\mathrm{e}^{Ax}$ 由级数定义，对任何 $A$ 都有意义，且 $\vect{Y}=\mathrm{e}^{Ax}\vect{C}$ 就是通解——三种情形被它一个式子盖住。
  \item 算它有两条主路：能对角化就 $V\mathrm{e}^{\Lambda x}V^{-1}$；不能就 Cayley--Hamilton 化成 $n-1$ 次多项式，用特征值（重根加导数）定系数。
  \item 已经有解组 $\Phi$ 时，$\mathrm{e}^{Ax}=\Phi(x)\Phi^{-1}(0)$；三者是同一个矩阵的不同写法。
\end{enumerate}

%=====================================================================
\hdr{§11}{与课件的页码对照}

\begin{minipage}{\linewidth}
\tabtop
\begin{tabular}{@{}p{0.12\linewidth}p{0.46\linewidth}p{0.32\linewidth}@{}}
\thead{本讲义} & \thead{内容} & \thead{课件页}\\
\end{tabular}
\tabmid
\begin{tabular}{@{}p{0.12\linewidth}p{0.46\linewidth}p{0.32\linewidth}@{}}
§1 & 定义、收敛、通解、初值问题 & p3--p6\\
§2 & 五条性质与 $\mathrm{e}^A$ 的特征值 & p7\\
§3 & 对角化：$Ve^{\Lambda x}V^{-1}$ & p8--p15\\
§4 & Cayley--Hamilton：降阶与定系数 & p16--p32\\
§5 & 基本解矩阵 $\Phi$ 与定理 4.1 & p39--p43\\
§7 & 作业三题 & p44\\
\end{tabular}
\tabbot
\end{minipage}

%=====================================================================
\newpage
\chapterbanner{4}{非齐次线性组}{参数变易与待定系数}{课件 \texttt{17\_18\_Nonhomogeneous-linear-systems.pdf}（49 页，2026-10-15／10-20）}
\callout{\textbf{先读这一段。}前三讲都在解 $\vect{Y}'=A\vect{Y}$（齐次）。这一讲加一个外力项：$\vect{Y}'=A(x)\vect{Y}+\vect{F}(x)$。结构上只有一句话——\textbf{通解 $=$ 齐次通解 $+$ 一个特解}（§1）。力气全花在"怎么找一个特解"上，两条路：\textbf{参数变易法}（把常数换成函数，万能但要积分，§2--§3）与\textbf{待定系数法}（猜形状、比较系数，快但只对常系数与特定 $\vect{F}$，§4）。§5 给出选路的判据，§6 用一道题把三条路走一遍。}

%=====================================================================
\hdr{§0}{问题、结构与两条路}

\begin{minipage}{\linewidth}
\tabtop
\begin{tabular}{@{}p{0.16\linewidth}p{0.26\linewidth}p{0.34\linewidth}p{0.10\linewidth}@{}}
\thead{方法} & \thead{前提} & \thead{代价} & \thead{课件页}\\
\end{tabular}
\tabmid
\begin{tabular}{@{}p{0.16\linewidth}p{0.26\linewidth}p{0.34\linewidth}p{0.10\linewidth}@{}}
参数变易（逐分量） & 已知基本解组 & 解 $n$ 元一次方程组 + 积分 & p6--p10\\
参数变易（矩阵） & 已知 $\Phi$ 或 $\mathrm{e}^{Ax}$ & 求 $\Phi^{-1}$ + 积分 & p11--p16\\
待定系数 & $A$ 常矩阵，$\vect{F}$ 是"多项式×指数×三角" & 猜形 + 解矩阵方程 & p18--p28\\
\end{tabular}
\tabbot
\end{minipage}

这张表是路线台账；两条参数变易本质是同一条（§3.1 说明），待定系数则是另一条思路。

%=====================================================================
\hdr{§1}{解的结构：通解 = 齐次通解 + 一个特解}

\subhdr{1.1 定理（p4）}

设

\[
\vect{Y}'=A(x)\vect{Y}+\vect{F}(x)\tag{N}
\]

的系数矩阵与 $\vect{F}$ 连续。则 (N) 的通解是

\[
\vect{Y}(x)=\vect{Y}_c(x)+\vect{Y}_p(x),
\]

其中 $\vect{Y}_c$ 是齐次方程 $\vect{Y}'=A\vect{Y}$ 的通解（$\sum c_k\vect{Y}_k$ 或 $\Phi(x)\vect{C}$），$\vect{Y}_p$ 是 (N) 的\textbf{任意一个}特解。

\subhdr{1.2 为什么（L13 定理 2.2 的复习）}

设 $\vect{Y}$ 是 (N) 的任一解，$\vect{Y}_p$ 是已知的一个特解。令 $\vect{W}=\vect{Y}-\vect{Y}_p$，则

\[
\vect{W}'-A\vect{W}=(\vect{Y}'-A\vect{Y})-(\vect{Y}_p'-A\vect{Y}_p)=\vect{F}-\vect{F}=\vect{0},
\]

所以 $\vect{W}$ 是齐次解，即 $\vect{W}=\vect{Y}_c$。反方向由叠加原理直接得到。

\textbf{这个证明只有一个动作：两个解相减，右端消掉。}后面所有"求特解"的工作，都是为了拿到那一个 $\vect{Y}_p$；齐次部分不用重算。

\subhdr{1.3 一个常见的误解}

$\vect{Y}_p$ \textbf{不唯一}：任何特解加一个齐次解仍是特解。所以两条方法给出的"特解"可能长得不一样，但通解一定相同。答卷上不必追求与答案逐字相同，只要代回原方程成立即可。

%=====================================================================
\hdr{§2}{参数变易法 I：逐分量版（p6）}

\subhdr{2.1 设形}

设 $\{\vect{Y}_1,\dots,\vect{Y}_n\}$ 是齐次方程的基本解组。把通解 $\sum c_k\vect{Y}_k$ 里的\textbf{常数换成函数}：

\[
\vect{Y}_p(x)=c_1(x)\vect{Y}_1(x)+\cdots+c_n(x)\vect{Y}_n(x).
\]

这正是"参数变易"（variation of parameters）这个名字的来历——L8 的标量常数变易法是它在 $n=1$ 时的样子。

\subhdr{2.2 关键一步：$c_k'$ 满足的方程}

对 $\vect{Y}_p$ 求导：

\[
\vect{Y}_p'=\sum_k c_k'\vect{Y}_k+\sum_k c_k\vect{Y}_k'
=\sum_k c_k'\vect{Y}_k+A\sum_k c_k\vect{Y}_k
=\sum_k c_k'\vect{Y}_k+A\vect{Y}_p .
\]

最后一步用了 $\vect{Y}_k'=A\vect{Y}_k$（它们是齐次解）。要 $\vect{Y}_p$ 满足 (N)，必须

\[
\vect{Y}_p'=A\vect{Y}_p+\vect{F}
\ \Longleftrightarrow\
\sum_{k=1}^{n}c_k'(x)\vect{Y}_k(x)=\vect{F}(x).\tag{$\star$}
\]

\textbf{课件 p6 从"$\vect{Y}_p'=A\vect{Y}_p+\vect{F}$"直接跳到 ($\star$)，中间省略的就是"$c_k$ 那一半正好等于 $A\vect{Y}_p$、与右边抵消"}。

\subhdr{2.3 为什么 ($\star$) 一定有唯一解}

把 ($\star$) 看成关于 $c_1',\dots,c_n'$ 的线性方程组，系数矩阵正是 $\Phi=(\vect{Y}_1,\dots,\vect{Y}_n)$：

\[
\Phi(x)\vect{C}'(x)=\vect{F}(x),\qquad \vect{C}'=(c_1',\dots,c_n')^{\mathsf T}.
\]

由于 $\{\vect{Y}_k\}$ 是基本解组，$\det\Phi(x)=W(x)\neq0$，所以 $\vect{C}'=\Phi^{-1}\vect{F}$ 唯一确定。这就是逐分量版与矩阵版是同一件事的理由。

\subhdr{2.4 操作四步}

\begin{enumerate}
  \item 求齐次基本解组 $\vect{Y}_1,\dots,\vect{Y}_n$（用 L14--15 的三情形）。
  \item 写出 ($\star$)，解出 $c_1',\dots,c_n'$。
  \item 积分得 $c_k(x)$（积分常数取 $0$，因为只要一个特解）。
  \item 代回 $\vect{Y}_p=\sum c_k\vect{Y}_k$，再加齐次通解。
\end{enumerate}

%=====================================================================
\hdr{§3}{参数变易法 II：矩阵版（p11--p13）}

\subhdr{3.1 设形与方程}

设 $\Phi(x)$ 是齐次方程的基本解矩阵，齐次通解是 $\Phi(x)\vect{C}$。把常向量换成向量函数：

\[
\vect{Y}_p(x)=\Phi(x)\vect{C}(x).
\]

求导：

\[
\vect{Y}_p'=\Phi'(x)\vect{C}(x)+\Phi(x)\vect{C}'(x)=A\Phi\vect{C}+\Phi\vect{C}'
=A\vect{Y}_p+\Phi\vect{C}' .
\]

要它满足 (N)，必须

\[
\Phi(x)\vect{C}'(x)=\vect{F}(x)
\ \Longleftrightarrow\
\vect{C}'(x)=\Phi^{-1}(x)\vect{F}(x).\tag{3.1}
\]

\subhdr{3.2 积分与通解（p12）}

积分：

\[
\vect{C}(x)=\int_{x_0}^{x}\Phi^{-1}(s)\vect{F}(s)\,\mathrm{d}s+\vect{C}_1 .
\]

代回：

\[
\vect{Y}_p(x)=\Phi(x)\int_{x_0}^{x}\Phi^{-1}(s)\vect{F}(s)\,\mathrm{d}s,
\qquad
\vect{Y}(x)=\underbrace{\Phi(x)\vect{C}}_{\text{齐次通解}}+\underbrace{\Phi(x)\int_{x_0}^{x}\Phi^{-1}(s)\vect{F}(s)\,\mathrm{d}s}_{\text{一个特解}} .
\]

\textbf{为什么积分下限可以任取？}因为 $\Phi\vect{C}_1$ 已经被齐次通解吸收；取下限 $x_0$ 是为了让特解在 $x=x_0$ 处为零，方便接初值。

\subhdr{3.3 初值问题（p12--p13）}

若给 $\vect{Y}(x_0)=\vect{Y}_0$，则

\[
\vect{Y}(x)=\Phi(x)\Bigl(\Phi^{-1}(x_0)\vect{Y}_0+\int_{x_0}^{x}\Phi^{-1}(s)\vect{F}(s)\,\mathrm{d}s\Bigr).
\]

常系数时取 $\Phi(x)=\mathrm{e}^{Ax}$，上式化为

\[
\boxed{\ \vect{Y}(x)=\mathrm{e}^{A(x-x_0)}\vect{Y}_0+\int_{x_0}^{x}\mathrm{e}^{A(x-s)}\vect{F}(s)\,\mathrm{d}s\ }.
\]

第二项叫\textbf{卷积项}：它把每个时刻 $s$ 的外力 $\vect{F}(s)$ 用 $\mathrm{e}^{A(x-s)}$ 传播到 $x$，再累加。这个公式是这一讲最有用的结论——只要知道 $\mathrm{e}^{Ax}$，非齐次初值问题一步写出。

\subhdr{3.4 为什么 $\Phi^{-1}$ 一定存在}

$\det\Phi(x)=W(x)$ 是基本解组的 Wronskian。由 L13 定理 2.3，$W$ 在区间上处处非零（一旦某点非零就处处非零），所以 $\Phi^{-1}$ 处处存在。课件 p12 写 "$\det\Phi(x)\neq0$ for all admissible $x$"，理由就在这里。

%=====================================================================
\hdr{§4}{待定系数法（p18--p28）}

\subhdr{4.1 三条适用条件（p18）}

\begin{enumerate}
  \item \textbf{$A$ 是常矩阵}。变系数时"猜形"里的指数底数会随 $x$ 变，比较系数无从谈起。
  \item \textbf{$\vect{F}$ 的元素是多项式、指数、正弦余弦，或它们的有限和与积}。其它函数（如 $\ln x$、$1/x$）猜不出来。
  \item \textbf{系数向量待定}：与标量情形不同，这里的"系数"是 $n$ 维\textbf{向量}。
\end{enumerate}

\subhdr{4.2 设形规则（p27--p28）}

\textbf{Case 1}（多项式 × 指数）：

\[
\vect{F}(x)=\mathrm{e}^{rx}\bigl(\alpha_0x^m+\alpha_1x^{m-1}+\cdots+\alpha_m\bigr),\qquad \alpha_i\in\mathbb{R}^n .
\]

设

\[
\vect{Y}_p(x)=\mathrm{e}^{rx}\bigl(\gamma_0x^k+\gamma_1x^{k-1}+\cdots+\gamma_k\bigr),
\qquad
k=\begin{cases}
m, & r\text{ 不是 }A\text{ 的特征值},\\
m+\text{（}r\text{ 的重数）}, & r\text{ 是 }A\text{ 的特征值}.
\end{cases}
\]

\textbf{Case 2}（三角）：

\[
\vect{F}(x)=\mathrm{e}^{rx}\bigl(P_m(x)\cos sx+Q_m(x)\sin sx\bigr),
\]

设 $\vect{Y}_p=\mathrm{e}^{rx}\bigl(\Gamma_k(x)\cos sx+\mathrm{N}_k(x)\sin sx\bigr)$，其中 $\Gamma_k,\mathrm{N}_k$ 是 $k$ 次向量多项式，$k$ 的取法同上——判据是 $r+\mathrm{i}s$ 是不是 $A$ 的特征值。

\subhdr{4.3 "共振"到底在加什么}

规则里 $k=m+\text{重数}$ 的那一项，与标量方程里"共振时乘 $x$"是同一件事。\textbf{为什么必须加}：若 $r$ 是特征值，把 $\mathrm{e}^{rx}\gamma_0x^m$ 代进去，最高次项会给出 $(A-rE)\gamma_0=0$——这个方程非零解要求 $\gamma_0$ 是特征向量，但 $\gamma_0$ 同时还要匹配 $\vect{F}$ 的方向，一般做不到。加一次幂次，让最高次项先被 $(A-rE)$ 吃掉，才能解出系数。

\subhdr{4.4 比较系数得到的是一组矩阵方程}

以 Case 1 的 $m=1$、$r$ 不是特征值为例：设 $\vect{Y}_p=\alpha x+\beta$，代入 $\vect{Y}_p'=A\vect{Y}_p+\vect{F}$ 得

\[
\alpha=A(\alpha x+\beta)+\gamma x+\delta
\ \Longrightarrow\
\begin{cases}
A\alpha+\gamma=0, & (x\text{ 的系数})\\
\alpha=A\beta+\delta, & (\text{常数项})
\end{cases}
\]

先解第一式得 $\alpha=-A^{-1}\gamma$（$r=0$ 不是特征值时 $A$ 可逆），再解第二式得 $\beta=A^{-1}(\alpha-\delta)$。

\textbf{共振时 $A$ 不可逆}，第一式变成 $A\alpha=-\gamma$，它有解的条件是 $-\gamma\in\operatorname{Range}(A)$——这就是课件 p25、p39 里那个"相容性比式"的来源（见 §4.6）。

\subhdr{4.5 与标量待定系数法逐条对照}

这一节的方法不是新东西，是 L8/L11 标量待定系数法的"向量化"。下表把两边逐条对上：

\begin{minipage}{\linewidth}
\tabtop
\begin{tabular}{@{}p{0.18\linewidth}p{0.30\linewidth}p{0.40\linewidth}@{}}
\thead{} & \thead{标量（L8/L11）} & \thead{方程组（本讲）}\\
\end{tabular}
\tabmid
\begin{tabular}{@{}p{0.18\linewidth}p{0.30\linewidth}p{0.40\linewidth}@{}}
方程 & $y''+py'+qy=f$ & $\vect{Y}'=A\vect{Y}+\vect{F}$\\
齐次解 & 两个基 $y_1,y_2$ & $\Phi$ 的列\\
设形 & $y_p=\mathrm{e}^{rx}P_m(x)$ & $\vect{Y}_p=\mathrm{e}^{rx}\bigl(\gamma_0x^k+\cdots+\gamma_k\bigr)$\\
“系数” & 数 $a_i$ & 向量 $\gamma_i\in\mathbb{R}^n$\\
共振判据 & $r$ 是特征根 & $r$ 是 $A$ 的特征值\\
共振处理 & 乘 $x^{(\text{重数})}$ & $k=m+(\text{重数})$\\
定系数 & 比较同次幂 & 比较同次幂，得到\textbf{矩阵方程}\\
解方程 & 一元一次 & $n$ 元一次（$A$ 奇异时要验相容性）\\
\end{tabular}
\tabbot
\end{minipage}

这张表是对照台账，不承载推导；它让你把已经会的标量流程直接搬过来，只改两处：\textbf{系数变向量}、\textbf{方程变矩阵方程}。

\subhdr{4.6 相容性条件：为什么会出现一个比式}

当 $A$ 奇异时，$A\alpha=-\gamma$ 不一定有解。判断办法是看增广矩阵的秩：

\[
\rank(A)=\rank(A\mid-\gamma).
\]

若 $A$ 的秩是 $n-1$，这个条件化成一个标量方程（例如 $2\beta_1+\beta_2=9$）。课件 p25 写 "$\alpha_1+5=-2\alpha_1+8$"、p39 写 "$-5:2=(2\alpha_1-3):(-\alpha_1)$"，都是这个秩条件的具体化——\textbf{它们不是新公式，是"方程要有解"的约束}。

%=====================================================================
\hdr{§5}{两条方法怎么选}

\begin{minipage}{\linewidth}
\tabtop
\begin{tabular}{@{}p{0.20\linewidth}p{0.36\linewidth}p{0.34\linewidth}@{}}
\thead{看什么} & \thead{参数变易} & \thead{待定系数}\\
\end{tabular}
\tabmid
\begin{tabular}{@{}p{0.20\linewidth}p{0.36\linewidth}p{0.34\linewidth}@{}}
$A$ 是否常矩阵 & 不限 & 必须常矩阵\\
$\vect{F}$ 的形状 & 不限 & 只能多项式×指数×三角\\
要不要积分 & 要（可能很难） & 不要，只解代数方程\\
要不要猜 & 不用，有公式 & 要猜，猜错白做\\
$\mathrm{e}^{Ax}$ 已有 & 直接用（最省） & 不必用\\
\end{tabular}
\tabbot
\end{minipage}

这张表是判据台账；两条方法的结论一定一致（最多差一个齐次解），所以\textbf{可以互相验算}。

%=====================================================================
\hdr{§6}{一道题的三种做法（p32--p48）}

课件用同一道题演示了三条路，值得单独说，因为它把"什么时候用哪条"讲活了。题目是

\[
x''=-2x'-5y+3,\qquad y'=x'+2y,\qquad x(0)=0,\ x'(0)=0,\ y(0)=1 .
\]

\textbf{做法一：化成一阶组 + 待定系数}（p32--p40）。令 $x_1=x,\ x_2=x',\ x_3=y$，得 $3\times3$ 常系数系统，外力是常向量 $(0,3,0)^{\mathsf T}$。因为 $0$ 是特征值（课件 p34 算出 $\lambda=0$），按 §4.2 的规则要设 $\vect{X}_p=\alpha t+\beta$（$m=0$，重数 $1$，故 $k=1$）。代进去得到 $A\alpha=0$、$A\beta+\gamma=\alpha$，再由相容性定出 $\alpha=(-6,0,0)$、$\beta=(0,-6,3)$。

\textbf{做法二：化成一阶组 + 参数变易}（p41--p45）。先算 $\Phi$ 与 $\Phi^{-1}$，再算 $\vect{C}'=\Phi^{-1}\vect{F}$，积分后乘回去。课件把 $\Phi^{-1}\vect{F}$ 的三个分量都写出来了，积分也不难——因为外力是常数向量。

\textbf{做法三：化回二阶方程}（p46--p48）。从 $y'=x'+2y$ 解出 $x'=y'-2y$，再求导得 $x''=y''-2y'$，代入第一式得 $y''+y=3$，用标量方法解出 $y$，再回代求 $x$。

\textbf{三条路给出的答案相同}（p45、p48）：

\[
x(t)=-2-6t+2\cos t+6\sin t,\qquad y(t)=3-2\cos t+2\sin t .
\]

\textbf{什么时候选哪条}：题目已经是一阶组、且外力形状简单——走做法一；外力复杂或 $A$ 变系数——走做法二；原题本来就是二阶方程、且能消元——走做法三（少一次求逆）。课件把三条都走一遍，是为了让你看到"同一件事的不同代价"，不是让你每题都走三遍。

%=====================================================================
\hdr{§7}{作业（p49）}

\qbreak
\qstem{（课件 p49）用参数变易法与待定系数法各解一次：
$y_1'=-2y_1-2y_2+2y_3+2x+3$，$y_2'=-3y_1+y_2-4y_3+3x-1$，$y_3'=-3y_1+2y_2-5y_3+3x-2$。}
\par\textbf{考点}：常系数 $3\times3$、多项式外力；两种方法互验。
\par\textbf{关键一步}：特征值 $-1,-2,-3$ 互异；外力是一次多项式且 $0$ 不是特征值，故设 $\vect{Y}_p=\alpha x+\beta$，比较系数得 $A\alpha=-(2,3,3)^{\mathsf T}$、$\alpha=A\beta+(3,-1,-2)^{\mathsf T}$。
\par\textbf{详解}：附录 A.37，第 \pageref{c4:app:hw} 页。

%=====================================================================
\hdr{§8}{易错点}

\begin{enumerate}
  \item \textbf{特解不唯一}。两条方法的结果可以差一个齐次解，别以为算错了（§1.3）。
  \item \textbf{参数变易的积分常数取零}。取非零只是把齐次解重复算进去（§2.4）。
  \item \textbf{待定系数的"系数"是向量}，不是数；每个 $\gamma_k$ 有 $n$ 个分量（§4.1）。
  \item \textbf{共振要乘 $x$（或提高次数）}，判据是"$r$ 是不是特征值"，不是"$\vect{F}$ 里有没有 $x$"（§4.2）。
  \item \textbf{$A$ 奇异时先验相容性}，否则解不出 $\alpha$ 还以为是自己算错（§4.5）。
  \item \textbf{常系数公式里的指数是 $\mathrm{e}^{A(x-s)}$}，不是 $\mathrm{e}^{Ax}\mathrm{e}^{-As}$ 随便写——两者相等，但积分时前者更好用（§3.3）。
\end{enumerate}

%=====================================================================
\hdr{§9}{缺口清单}

\gapitem{代数跳步}{p6 从 $\vect{Y}_p'=A\vect{Y}_p+\vect{F}$ 直接跳到 $\sum c_k'\vect{Y}_k=\vect{F}$。本讲义 §2.2 补出"$c_k$ 那半正好抵消"的一步。}
\gapitem{适用条件缺失}{p6 未说明 ($\star$) 为何有唯一解。本讲义 §2.3 补：系数矩阵是 $\Phi$，行列式非零。}
\gapitem{证明缺口}{p12 “$\det\Phi(x)\neq0$ for all admissible $x$” 未给理由。本讲义 §3.4 补 Wronskian 处处非零。}
\gapitem{边界处理}{p12 积分写成不定积分再加 $\vect{C}_1$，未说明 $\vect{C}_1$ 与齐次通解的关系。本讲义 §3.2 说明取下限 $x_0$ 的理由。}
\gapitem{机制缺口}{p18 只给“只适用于常系数、不能用于任意 $\vect{F}$”，未解释原因。本讲义 §4.1 补三条理由。}
\gapitem{机制缺口}{p23 “for a system, the corresponding form $\alpha x\mathrm{e}^x$ does not work” 未说明为何。本讲义 §4.3 补：$\alpha$ 要同时匹配 $\vect{F}$ 方向与特征向量，一般做不到。}
\gapitem{代数跳步}{p25 “admits solutions if and only if $\alpha_1+5=-2\alpha_1+8$” 直接给条件。本讲义 §4.6 补秩判据 $\rank(A)=\rank(A\mid-\gamma)$。}
\gapitem{记号含糊}{p27--p28 的 $k$ 规则写成 “algebraic multiplicity $+m$”，未说清是“$m$ 加重数”。本讲义 §4.2 改写为 $m+$（重数）。}
\gapitem{代数跳步}{p39 的相容性比式 “$-5:2=(2\alpha_1-3):(-\alpha_1)$” 来源未说明。本讲义 §4.6 统一解释为秩条件。}
\gapitem{数据不一致}{p29 例 5 的矩阵 $A$ 的迹是 $1$，与所标“特征值 $0,\pm5\mathrm{i}$”（其和为零）矛盾；按打印的 $A,\mu,\nu$ 解出的 $\alpha,\beta$ 也与 p31 的答案不符（该矩阵的实际特征值是 $1,\pm5\mathrm{i}$）。本讲义 A.5 改用一个可核对的补例演示三角外力的设形。}
\gapitem{拼接缺口}{p46 化二阶方程时，“$x'=y'-2y$ 再求导得 $x''=y''-2y'$” 两步合并。本讲义 §6 说明这是“解出 $x'$ 再求导”的标准动作。}

%=====================================================================
\hdr{§10}{三句话}

\begin{enumerate}
  \item 非齐次的结构只有一句：通解 $=$ 齐次通解 $+$ 一个特解；齐次部分用前三讲的方法，力气全在特解。
  \item 参数变易法万能：$\vect{C}'=\Phi^{-1}\vect{F}$，积分即得；常系数时化为卷积公式
\[
\vect{Y}(x)=\mathrm{e}^{A(x-x_0)}\vect{Y}_0+\int_{x_0}^{x}\mathrm{e}^{A(x-s)}\vect{F}(s)\,\mathrm{d}s .
\]
  \item 待定系数法快但挑食：常系数 $+$ 特定 $\vect{F}$；设形看"$r$（或 $r+\mathrm{i}s$）是不是特征值"，是就提高次数，比较系数时先解 $A\alpha=\cdots$、再验相容性。
\end{enumerate}

%=====================================================================
\hdr{§11}{与课件的页码对照}

\begin{minipage}{\linewidth}
\tabtop
\begin{tabular}{@{}p{0.12\linewidth}p{0.46\linewidth}p{0.32\linewidth}@{}}
\thead{本讲义} & \thead{内容} & \thead{课件页}\\
\end{tabular}
\tabmid
\begin{tabular}{@{}p{0.12\linewidth}p{0.46\linewidth}p{0.32\linewidth}@{}}
§1 & 解的结构 $\vect{Y}_c+\vect{Y}_p$ & p4\\
§2 & 参数变易 I：逐分量 & p6--p10\\
§3 & 参数变易 II：矩阵与卷积公式 & p11--p16\\
§4 & 待定系数：条件、设形、共振、相容性 & p18--p28\\
§6 & 同一题的三种做法 & p32--p48\\
§7 & 作业 & p49\\
\end{tabular}
\tabbot
\end{minipage}

%=====================================================================
\newpage
\bigpart{III}{幂级数解法 · L18--L20}{换掉“常系数”这一假设：常点用幂级数，正则奇点用 Frobenius 级数。课件 79 页。}
\chapterbanner{5}{幂级数解法}{常点、正则奇点与 Frobenius 方法}{课件 \texttt{18\_power-series.pdf}（37 页）与 \texttt{19-20\_power-series(II).pdf}（42 页），2026-10-20／10-22}
\callout{\textbf{先把这一章的骨架记成五行，后面所有内容都是在解释它们。}
\begin{enumerate}[label=\textcircled{\small\arabic*},leftmargin=2.3em]
  \item \textbf{常点}（$P(x_0)\neq0$）：解是普通幂级数 $y=\sum a_n(x-x_0)^n$，系数由递推逐项定。
  \item \textbf{奇点}分两类：$(x-x_0)p$ 与 $(x-x_0)^2q$ 都在 $x_0$ 解析的叫\textbf{正则奇点}，否则叫非正则。
  \item 正则奇点处设 $y=x^r\sum a_nx^n$，把最低次幂的系数令为零，得到\textbf{指标方程}
        \[ r(r-1)+p_0r+q_0=0,\qquad p_0=\lim_{x\to0}xp(x),\quad q_0=\lim_{x\to0}x^2q(x). \]
  \item 两根之差不是整数：两个解都是 Frobenius 级数；两根相等或差为正整数：第二个解可能带 $\ln x$（也可能不带，要看递推在那个整数步是否自动满足）。
  \item 级数的收敛半径 $\rho$ 等于「展开点到最近奇点（实或复）的距离」；作业题的 $\rho=\infty$，所以级数在整个 $x>0$ 上收敛。
\end{enumerate}}

%=============================================================
\hdr{§0}{课件坐标、记号与本章地图}

\subhdr{0.1 这一章在课程里的位置}

课件 Lecture 1--12 处理的是「能写成初等函数」的方程：可分离、一阶线性、恰当、降阶、常系数、Cauchy--Euler。Lecture 13--17 换成线性方程组的语言，把高阶方程统一成 $Y'=A(x)Y$。Lecture 18--20 面对的是另一类方程：\textbf{系数是函数，而且没有初等解}。

最典型的是课件 p18 给出的两个方程：
\[
\text{Bessel 方程：}\quad x^2y''+xy'+(x^2-\nu^2)y=0,
\qquad
\text{Legendre 方程：}\quad (1-x^2)y''-2xy'+\alpha(\alpha+1)y=0 .
\]
它们的解不是 $\sin$、$\ee^x$、$x^r$ 这类初等函数，而是\textbf{幂级数}本身——Bessel 函数 $J_\nu$、Legendre 多项式 $P_n$ 就是这样定义的。所以这一章不是又学一种「积分技巧」，而是换一种「表示解」的方式：\textbf{把解写成一个收敛的幂级数，并用递推把系数一个个算出来}。

这条方法链只有三步，本讲义按这三步组织：
\begin{enumerate}
  \item \textbf{判断展开点}：$P(x_0)\neq0$ 是常点，$P(x_0)=0$ 是奇点；奇点再分正则、非正则（§2、§3）。
  \item \textbf{决定解的形状}：常点用 $\sum a_n(x-x_0)^n$；正则奇点用 $x^r\sum a_nx^n$，$r$ 由指标方程定（§2、§4）。
  \item \textbf{确定收敛范围}：$\rho$ 由「到最近奇点的距离」给出（§4.6、§5.7）。
\end{enumerate}

\subhdr{0.2 课件逐页对照：哪一页讲什么，本讲义在哪里补}

\begin{center}
\small
\begin{tabular}{@{}p{1.25cm}p{1.55cm}p{7.0cm}p{6.0cm}@{}}
\hline
\thead{课件} & \thead{页} & \thead{课件给了什么} & \thead{本讲义对应节} \\
\hline
L18 & p3--p16 & 幂级数定义、比值法、四则运算、逐项求导、泰勒展开、解析函数、指标平移 & §1（补：比值法的倒数写法、端点检查、为什么能逐项求导）\\
L18 & p17--p28 & 常点/奇点定义、定理 2.1、五步法、通用递推公式 & §2.1--2.3（补：卷积推导、定理证明思路）\\
L18 & p29--p35 & 例 5（含 $\ee^x$ 的初值问题）、例 6（$(x^2+1)y''-4xy'+6y=0$） & §2.5--2.6（补：每步代数、为什么级数截断）\\
L18 & p36 & 作业 1：$y''+(1-x)y'+(x^2-1)y=0$ & 不在本次作业范围；方法同 §2 \\
L19-20 & p3--p7 & 正则/非正则奇点定义、分类例 & §3.1--3.2（补：为什么乘 $(x-x_0)$ 与 $(x-x_0)^2$）\\
L19-20 & p8--p19 & Euler 方程、三种根、$x<0$ 的处理、平移版 & §3.3--3.6（补：重根用降阶推出 $\ln x$）\\
L19-20 & p20--p31 & Frobenius 方法、指标方程、递推、定理 1.1、三情形 & §4（补：分母结构 $(k+r-r_1)(k+r-r_2)$）\\
L19-20 & p32--p34 & 例：$xy''+4y'-xy=0$，只算 $r=0$；$r=-3$ 留作作业 & §4.5（补：$r=-3$ 的完整第二解）\\
L19-20 & p35--p41 & Bessel 方程、$J_\nu$、$Y_\nu$ & §7 \\
L19-20 & p42 & \textbf{作业 1（本次题目）} & \textbf{§5 全解} \\
\hline
\end{tabular}
\end{center}

\aside{本讲义没有照搬「§1 机制—§2/§3 两条主线—§4 等价性」的通用骨架：本章的内容是三个并列的方法（常点、Euler、Frobenius），不是两条路线的对照，硬套骨架反而会把 §3 与 §4 的先后关系讲乱。所以改成「§1 工具 → §2 常点 → §3 奇点与 Euler → §4 Frobenius → §5 作业 → §6 方程组补齐 → §7 Bessel」的线性结构；每一节与课件页的对应关系就是上表。}

\subhdr{0.3 记号}

\begin{center}
\small
\begin{tabular}{@{}p{2.5cm}p{12.8cm}@{}}
\hline
\thead{符号} & \thead{含义} \\
\hline
$x$ & 自变量。教材（Constanda）用 $t$，本讲义统一用 $x$，两者只是字母不同。 \\
$y_1,y_2$ & 方程的两个线性无关解；通解写成 $C_1y_1+C_2y_2$。 \\
$a_n$ & Frobenius 级数 $\sum a_nx^{n+r}$ 的系数；约定 $a_0\neq0$。 \\
$r,r_1,r_2$ & 指标根，规定 $r_1\geq r_2$（实根情形）。 \\
$p_0,q_0$ & $p_0=\lim_{x\to0}xp(x)$，$q_0=\lim_{x\to0}x^2q(x)$，即 $xp$、$x^2q$ 展开式的常数项。 \\
$\rho$ & 定理保证的收敛半径下界：展开点到最近奇点（实或复）的距离。 \\
$W$ & Wronskian：$W=y_1y_2'-y_1'y_2$；系统情形是 $\det\Phi$。 \\
$\Phi,A$ & 基本矩阵与系统矩阵，见 §6。 \\
\hline
\end{tabular}
\end{center}

\aside{教材里 $p_0,q_0$ 的写法是 $\alpha_0,\beta_0$（9.4 节：$tQ/P=\alpha_0+\alpha_1t+\cdots$，$t^2R/P=\beta_0+\beta_1t+\cdots$）。本讲义用 $p_0,q_0$，与课件 p25 一致；对照教材时把 $\alpha_0$ 读成 $p_0$、$\beta_0$ 读成 $q_0$ 即可。}

\subhdr{0.4 本章地图}

\begin{center}
\begin{tikzpicture}[>=Latex,font=\small,
  box/.style={draw=petal,rounded corners=2pt,fill=bandviolet,inner sep=4pt,align=center},
  dec/.style={draw=violet,rounded corners=2pt,fill=white,inner sep=3pt,align=center}]
\node[dec] (q) at (0,0) {$P(x_0)\neq0$ ?};
\node[box,text width=4.0cm] (ord) at (-4.6,-2.0) {常点\\[1pt]$y=\sum a_n(x-x_0)^n$\\[1pt]半径 $\geq$ 到最近奇点};
\node[box,text width=5.6cm] (sing) at (3.4,-2.0) {奇点：再判 $(x-x_0)p$、$(x-x_0)^2q$\\[2pt]正则 $\Rightarrow$ $y=x^r\sum a_nx^n$，指标方程定 $r$\\[2pt]非正则 $\Rightarrow$ 本课不处理};
\draw[->] (q) -- node[above left,inner sep=1pt] {是} (ord);
\draw[->] (q) -- node[above right,inner sep=1pt] {否} (sing);
\end{tikzpicture}\\[2pt]
{\small\color{inkgray}图 1：本章的判别路线。左边一支是 §2 的常点法，右边一支是 §3--§4 的 Frobenius 法。}
\end{center}

这张图是本讲义的全部内容压缩成一张纸。你做题时只要沿着箭头走：先看 $P(x_0)$，再看两个极限是否存在，最后按根的情形选公式。

\subhdr{0.5 读完本章，你应该能回答的三个问题}

\begin{enumerate}
  \item 拿到一个变系数方程，\textbf{展开点属于哪一类}？常点、正则奇点、还是非正则奇点？判断依据是什么？
  \item 在正则奇点处，\textbf{指数 $r$ 从哪来}？为什么 $r_1$ 一定行、$r_2$ 可能出问题？第二个解什么时候要加 $\ln x$？
  \item 算出来的级数\textbf{能用到多远}？$\rho$ 怎么由方程本身读出来，而不是算完级数再猜？
\end{enumerate}

%=============================================================
\hdr{§1}{幂级数工具箱}

这一节对应课件 L18 p3--p16。课件把定义、定理、例子排得很密，但有几处默认你会：比值法为什么写成 $\lim|a_k|/|a_{k+1}|$、端点为什么要单独查、为什么幂级数能逐项求导。本节把这些补上，并把后面反复用到的「指标平移」练熟。

\subhdr{1.1 幂级数、收敛半径与比值法}

\callout{
\textbf{定义（课件 p4--p5）。}形如
\[
\sum_{k=0}^{\infty}a_k(x-x_0)^k=a_0+a_1(x-x_0)+a_2(x-x_0)^2+\cdots
\]
的表达式叫\textbf{幂级数}，$x_0$ 叫展开点。若部分和 $S_n(x)=\sum_{k=0}^n a_k(x-x_0)^k$ 的极限存在，就说级数在 $x$ 处收敛，记其和为 $S(x)$。

若 $\sum_{k=0}^\infty |a_k||x-x_0|^k$ 收敛，就说\textbf{绝对收敛}。使级数在 $|x-x_0|<r$ 内绝对收敛、在 $|x-x_0|>r$ 外发散的 $r$（$0\leq r\leq\infty$）叫\textbf{收敛半径}。
}

三条容易漏的边界事实（课件 p5 的 Remark）：
\begin{itemize}
  \item $r=\infty$：级数对一切实数 $x$ 收敛（例如 $\ee^x$ 的展开）；
  \item $r=0$：级数只在 $x=x_0$ 收敛（例如 $\sum k!x^k$）；
  \item $|x-x_0|=r$ 的两个端点：\textbf{可能收敛、可能发散，必须单独判断}。
\end{itemize}

\callout{
\textbf{定理 1.1（比值法，课件 p6）。}若 $k$ 充分大时 $a_k\neq0$，且
\[
\lim_{k\to\infty}\frac{|a_k|}{|a_{k+1}|}=L,\qquad 0\leq L\leq\infty,
\]
则幂级数的收敛半径为 $r=L$。
}

\aside{注意这个写法是「前项除以后项」，所以 $L$ 直接就是半径，不用再取倒数。如果你习惯写 $\lim|a_{k+1}/a_k|=q$，那么半径是 $r=1/q$（$q=0$ 时 $r=\infty$）。两种写法只差一次倒数，考场上选一种，别混着用。}

\subhdr{1.2 端点：为什么必须单独查}

课件 p6--p7 的例子值得完整走一遍，因为它同时演示了「半径」和「区间」是两件事。

\callout{
\textbf{例 1。}求 $\displaystyle\sum_{n=1}^{\infty}\frac{(x+1)^n}{n\,2^n}$ 的收敛半径与收敛区间。

\textbf{解。}这里 $a_n=\dfrac{1}{n2^n}$，于是
\[
\lim_{n\to\infty}\frac{|a_n|}{|a_{n+1}|}
=\lim_{n\to\infty}\frac{(n+1)2^{n+1}}{n2^n}=2 .
\]
所以 $r=2$：在 $|x+1|<2$ 即 $-3<x<1$ 内绝对收敛，在 $|x+1|>2$ 外发散。端点单独看：
\begin{itemize}
  \item $x=1$：级数变成 $\sum 1/n$，调和级数，\textbf{发散}；
  \item $x=-3$：级数变成 $\sum(-1)^n/n$，交错调和级数，\textbf{收敛}（条件收敛，不绝对收敛）。
\end{itemize}
所以收敛区间是 $[-3,1)$，收敛半径仍是 $2$。
}

\warnbox{
\textbf{易错：}把「收敛半径 $r=2$」直接写成「收敛区间 $(-3,1)$」。半径只告诉你在端点以内绝对收敛、以外发散；端点归谁，必须代进去看。本题两个端点一个收敛一个发散，正好说明端点不能靠半径猜。
}

\subhdr{1.3 四则运算与逐项求导}

课件 p8--p10 给了运算规则。把它们当成「幂级数可以像多项式一样操作」的许可证：

\callout{
\textbf{定理（逐项运算）。}设 $\sum a_k(x-x_0)^k$ 与 $\sum b_k(x-x_0)^k$ 在 $|x-x_0|<r$ 内分别收敛到 $f(x)$、$g(x)$，则在该区间内
\begin{align*}
f(x)\pm g(x)&=\sum_{k=0}^{\infty}(a_k\pm b_k)(x-x_0)^k,\\
f(x)g(x)&=\sum_{n=0}^{\infty}\Bigl(\sum_{k=0}^{n}a_kb_{n-k}\Bigr)(x-x_0)^n,\\
f'(x)&=\sum_{k=1}^{\infty}k\,a_k(x-x_0)^{k-1},\qquad
f''(x)=\sum_{k=2}^{\infty}k(k-1)a_k(x-x_0)^{k-2}.
\end{align*}
若还有 $b_0=g(x_0)\neq0$，则 $f/g$ 也能展成幂级数，但收敛半径可能比 $r$ 小。
}

乘积公式里那个 $\sum_{k=0}^n a_kb_{n-k}$ 叫\textbf{卷积}（Cauchy 乘积）。§2.3 和 §4.2 里，把一个函数级数乘到 $y'$ 或 $y$ 上时，用的就是它：\textbf{两个级数相乘，同幂次的系数是所有下标之和等于该幂次的乘积之和}。

\aside{为什么可以逐项求导？因为幂级数在收敛圆内闭的一致收敛（在任意 $|x-x_0|\leq r'<r$ 上一致收敛），一致收敛的级数可以逐项积分、逐项求导且半径不变。课件 p10 直接给了结论；这条理由在数学分析里证，本讲义只需知道「半径内随便求导」。}

\subhdr{1.4 指标平移：两个动作}

后面每个例子都要做指数平移。课件 p15--p16 的两个例子把动作拆得很清楚，这里把它总结成一句话：\textbf{先把外面的 $x$ 的幂乘进去，再整体升降指标；下限和系数下标要同时改。}

\callout{
\textbf{例 2（升指标）。}把 $\displaystyle\sum_{n=2}^{\infty}a_nx^n$ 写成从 $n=0$ 开始。令 $m=n-2$，则 $n=m+2$，$n=2$ 对应 $m=0$：
\[
\sum_{n=2}^{\infty}a_nx^n=\sum_{m=0}^{\infty}a_{m+2}x^{m+2}=\sum_{n=0}^{\infty}a_{n+2}x^{n+2}.
\]
\textbf{例 3（先乘进去，再降指标）。}把 $\displaystyle x^2\sum_{n=0}^{\infty}(r+n)a_nx^{r+n-1}$ 写成通项含 $x^{r+n}$ 的级数：
\[
x^2\sum_{n=0}^{\infty}(r+n)a_nx^{r+n-1}=\sum_{n=0}^{\infty}(r+n)a_nx^{r+n+1}
=\sum_{n=1}^{\infty}(r+n-1)a_{n-1}x^{r+n}.
\]
}

\warnbox{
平移时最常见的两个错：（1）只改 $\sum$ 的下限，忘了改 $a$ 的下标（$a_n$ 应变成 $a_{n-1}$ 或 $a_{n+1}$）；（2）只改 $a$ 的下标，忘了把 $x$ 的指数也改过来。检查办法：\textbf{代 $n$ 等于新的下限进去，看这一项是否等于原级数的第一项}。
}

\subhdr{1.5 常用展开表}

这张表是「查表用」的台账，不承载推导；每一行都能由泰勒公式在 $x_0=0$ 处算出，收敛范围由比值法得到。

\begin{center}
\small
\begin{tabular}{@{}p{3.6cm}p{8.6cm}p{3.6cm}@{}}
\hline
\thead{函数} & \thead{麦克劳林级数} & \thead{收敛范围} \\
\hline
$\ee^x$ & $\displaystyle\sum_{k=0}^{\infty}\frac{x^k}{k!}$ & $|x|<\infty$ \\
$\sin x$ & $\displaystyle\sum_{k=0}^{\infty}(-1)^k\frac{x^{2k+1}}{(2k+1)!}$ & $|x|<\infty$ \\
$\cos x$ & $\displaystyle\sum_{k=0}^{\infty}(-1)^k\frac{x^{2k}}{(2k)!}$ & $|x|<\infty$ \\
$\sinh x,\ \cosh x$ & $\displaystyle\sum_{k=0}^{\infty}\frac{x^{2k+1}}{(2k+1)!},\ \sum_{k=0}^{\infty}\frac{x^{2k}}{(2k)!}$ & $|x|<\infty$ \\
$\dfrac{1}{1-x}$ & $\displaystyle\sum_{k=0}^{\infty}x^k$ & $|x|<1$ \\
$\ln(1+x)$ & $\displaystyle\sum_{k=1}^{\infty}(-1)^{k-1}\frac{x^k}{k}$ & $|x|<1$ \\
$(1+x)^\alpha$ & $\displaystyle\sum_{k=0}^{\infty}\frac{\alpha(\alpha-1)\cdots(\alpha-k+1)}{k!}x^k$ & $|x|<1$ \\
$\arctan x$ & $\displaystyle\sum_{k=0}^{\infty}(-1)^k\frac{x^{2k+1}}{2k+1}$ & $|x|<1$ \\
\hline
\end{tabular}
\end{center}

读法：左边一列是函数，中间一列是它在 $0$ 点的泰勒级数，右边一列是「这个等式成立的 $x$ 范围」。最后两行的端点不在表内——$\ln(1+x)$ 在 $x=1$ 收敛、$x=-1$ 发散；$\arctan x$ 在两个端点都收敛——但本章只关心开区间内的展开，端点不参与运算。

\subhdr{1.6 解析函数：本章的「合格证」}

\callout{
\textbf{定义（课件 p14）。}若 $f$ 在 $|x-x_0|<r$ 内能展开成收敛的幂级数
\[
f(x)=\sum_{k=0}^{\infty}\frac{f^{(k)}(x_0)}{k!}(x-x_0)^k ,
\]
则称 $f$ 在 $x_0$ \textbf{解析}（analytic）。

解析函数对加、减、乘、除（分母在 $x_0$ 不为零）封闭。$\sin x$、$\ee^x$ 处处解析；$1/x$ 除 $x=0$ 外处处解析；$\tan x$ 除 $\pi/2$ 的奇数倍外处处解析。
}

「解析」这个词在后面只干一件事：\textbf{判断 $p=Q/P$、$q=R/P$ 在展开点附近能不能展成幂级数}。能，就是常点（定理 2.1）；不能但乘上 $(x-x_0)$、$(x-x_0)^2$ 后能，就是正则奇点（§3）。

\warnbox{
课件 p18 写「where $P$, $Q$ and $R$ are analytic \textbf{equations}」，这是笔误，应为 analytic \textbf{functions}。方程没有「解析」一说，解析的是系数函数。
}

%=============================================================
\hdr{§2}{常点邻域：把方程变成递推}

这一节对应课件 L18 p17--p35。核心只有一句话：\textbf{在常点处，解也是解析的；把级数代进去，同幂次的系数令为零，就得到递推}。难的不是思想，是代数——所以本节把两个例子（例 4、例 5）的每一步都写出来。

\subhdr{2.1 常点与奇点}

\callout{
\textbf{定义（课件 p19）。}对方程
\[
P(x)y''+Q(x)y'+R(x)y=0,
\]
若 $P(x_0)\neq0$，称 $x_0$ 为\textbf{常点}（ordinary point）；若 $P(x_0)=0$，称 $x_0$ 为\textbf{奇点}（singular point）。

在常点处可除以 $P$，化成标准形
\[
y''+p(x)y'+q(x)y=0,\qquad p(x)=\frac{Q(x)}{P(x)},\quad q(x)=\frac{R(x)}{P(x)} .
\]
}

判断常点只看 $P$，不看 $Q,R$。例如 $(x^2+1)y''-4xy'+6y=0$ 中 $P=x^2+1$，在实轴上永远不为零，所以每个实数点都是常点；但 $P$ 在 $x=\pm\ii$ 处为零，这两个\textbf{复奇点}会限制实轴上幂级数的收敛半径（§2.6）。

\subhdr{2.2 定理 2.1：解为什么也是解析的}

\callout{
\textbf{定理 2.1（课件 p20）。}设
\[
p(x)=\sum_{n=0}^{\infty}p_n(x-x_0)^n,\qquad q(x)=\sum_{n=0}^{\infty}q_n(x-x_0)^n
\]
在 $|x-x_0|<r$ 内解析。则方程
\[
y''+p(x)y'+q(x)y=0
\]
的\textbf{每一个}解都在 $|x-x_0|<r$ 内解析，即
\[
y(x)=\sum_{k=0}^{\infty}a_k(x-x_0)^k,\qquad |x-x_0|<r .
\]
}

课件没有给证明。这里给一个三行的思路，足够你理解「为什么敢设级数」：

\aside{证明思路（补充）：把 $p,q,y$ 的级数代入，比较 $x^n$ 的系数，得到 $a_{n+2}$ 由 $a_0,\dots,a_{n+1}$ 决定的递推。由 $p,q$ 的收敛性，存在 $M,R$ 使 $|p_n|,|q_n|\leq MR^{-n}$；对递推作归纳可得 $|a_n|\leq C(R')^{-n}$（$R'<R$），于是 $\sum a_n(x-x_0)^n$ 的收敛半径至少 $R'$。令 $R'\to r$ 即得结论。细节不必背，记住「递推 + 系数有界 $\Rightarrow$ 收敛半径不小于系数的解析半径」。}

定理的两个直接推论，是后面反复用的：
\begin{itemize}
  \item 通解有\textbf{两个}任意常数：$a_0=y(x_0)$、$a_1=y'(x_0)$ 可以任取，其余系数被递推唯一决定；
  \item 初值问题 $y(x_0)=y_0,\ y'(x_0)=y_0'$ 在 $|x-x_0|<r$ 内有唯一解，且 $a_0=y_0,\ a_1=y_0'$。
\end{itemize}

\subhdr{2.3 五步法与通用递推公式}

课件 p27--p28 把常点法拆成五步，这是考场上要照着做的流程：

\callout{
\textbf{五步法。}
\begin{enumerate}[label=\textbf{Step \arabic*.},leftmargin=3.4em]
  \item 设 $y=\sum_{k=0}^{\infty}a_kx^k$（若展开点不是 $0$，把 $x$ 换成 $x-x_0$），并写出 $y'$、$y''$ 的级数。
  \item 把 $y,y',y''$ 代回方程，把 $p,q$（或 $Q/P,R/P$）也展成幂级数，\textbf{把所有项整理成同一次幂}。
  \item 令 $x^n$ 的系数为零；这个方程里通常含 $a_n,a_{n+1},a_{n+2}$ 中的几个，解出\textbf{下标最大}的那个，得到\textbf{递推公式}。
  \item 用递推从 $a_0,a_1$ 依次算出 $a_2,a_3,\dots$。
  \item 把系数代回，写出两个线性无关解与通解。
\end{enumerate}
}

当 $p,q$ 本身就是一般幂级数时，第 2 步的「整理」可以用一个通用公式一次完成。课件 p21--p22 给了它，但跳过了卷积的推导，这里补上。

设 $x_0=0$，$p(x)=\sum_{m=0}^\infty p_mx^m$，$q(x)=\sum_{m=0}^\infty q_mx^m$，$y=\sum_{k=0}^\infty a_kx^k$。则
\[
y'=\sum_{k=0}^{\infty}(k+1)a_{k+1}x^k,\qquad
y''=\sum_{k=0}^{\infty}(k+2)(k+1)a_{k+2}x^k .
\]
把 $p$ 与 $y'$ 相乘，用卷积（§1.3）：$x^n$ 的系数是 $\sum_{k=0}^{n}p_{n-k}(k+1)a_{k+1}$。同理 $qy$ 的 $x^n$ 系数是 $\sum_{k=0}^{n}q_{n-k}a_k$。代入方程并令 $x^n$ 的系数为零：
\[
(n+2)(n+1)a_{n+2}+\sum_{k=0}^{n}(k+1)p_{n-k}a_{k+1}+\sum_{k=0}^{n}q_{n-k}a_k=0,
\qquad n=0,1,2,\dots \tag{2.1}
\]
这就是课件 p21 的式子。前三个方程是
\begin{align*}
n=0:\quad & 2a_2+p_0a_1+q_0a_0=0,\\
n=1:\quad & 6a_3+2p_0a_2+(q_0+p_1)a_1+q_1a_0=0,\\
n=2:\quad & 12a_4+3p_0a_3+(q_0+2p_1)a_2+(q_1+p_2)a_1+q_2a_0=0 .
\end{align*}
每一个都解出下标最大的系数：$a_2$ 由 $a_0,a_1$ 定，$a_3$ 由 $a_0,a_1,a_2$ 定，依此类推。

\aside{式 (2.1) 看起来吓人，实际做题时很少直接套：多数题目里 $p,q$ 只有少数几项非零（例如 $p=1$、$q=x$），逐项乘开比套公式快。公式的用处是\textbf{检查}——你算出的 $a_2,a_3$ 应满足 $n=0,1$ 的两个方程。}

\subhdr{2.4 例 4：$y''+y=0$ 的完整演算}

\callout{
\textbf{例 4（课件 p23--p26）。}求 $y''+y=0$ 在 $x=0$ 附近的幂级数解，$-\infty<x<\infty$。
}

设 $y=\sum_{k=0}^\infty a_kx^k$，则 $y''=\sum_{k=2}^\infty k(k-1)a_kx^{k-2}$。把第二个级数的指标升 2（令 $k\to k+2$）：
\[
y''=\sum_{k=0}^{\infty}(k+2)(k+1)a_{k+2}x^k .
\]
代入 $y''+y=0$：
\[
\sum_{k=0}^{\infty}\bigl[(k+2)(k+1)a_{k+2}+a_k\bigr]x^k=0 .
\]
令每个 $x^k$ 的系数为零：
\[
(k+2)(k+1)a_{k+2}+a_k=0
\quad\Longrightarrow\quad
a_{k+2}=-\frac{a_k}{(k+2)(k+1)},\qquad k=0,1,2,\dots \tag{2.2}
\]
这就是递推公式。

\textbf{为什么偶数和奇数解耦。}式 (2.2) 只把 $a_k$ 和 $a_{k+2}$ 连在一起，$k$ 与 $k+2$ 同奇偶，所以\textbf{偶下标只由 $a_0$ 决定，奇下标只由 $a_1$ 决定}，两条链互不相干。这不是巧合，是「方程只含 $y''$ 与 $y$、不含 $y'$」的结果：$y''$ 与 $y$ 都不改变奇偶性。

偶数链：
\[
a_2=-\frac{a_0}{2!},\quad a_4=-\frac{a_2}{3\cdot4}=\frac{a_0}{4!},\quad
a_6=-\frac{a_4}{5\cdot6}=-\frac{a_0}{6!},\quad\dots\quad
a_{2k}=\frac{(-1)^k}{(2k)!}a_0 .
\]
奇数链：
\[
a_3=-\frac{a_1}{3!},\quad a_5=\frac{a_1}{5!},\quad\dots\quad
a_{2k+1}=\frac{(-1)^k}{(2k+1)!}a_1 .
\]
于是
\[
y=a_0\sum_{k=0}^{\infty}\frac{(-1)^k}{(2k)!}x^{2k}
+a_1\sum_{k=0}^{\infty}\frac{(-1)^k}{(2k+1)!}x^{2k+1}
=a_0\cos x+a_1\sin x .
\]
两个级数用比值法都得到 $r=\infty$，所以对一切 $x$ 收敛。$y_1=\cos x$、$y_2=\sin x$ 的 Wronskian
\[
W=\cos x\cos x-(-\sin x)\sin x=1\neq0,
\]
确实线性无关。

\warnbox{
\textbf{注意：}这里 $a_0=y(0)$、$a_1=y'(0)$，不是「$a_0$ 一定是第一个非零项」。若初值给 $y(0)=0$，则 $a_0=0$，偶数链整条为零，解只剩 $\sin$ 一支——这是正常的，不是算错。
}

\subhdr{2.5 例 5：系数含 $\ee^x$ 的初值问题}

\callout{
\textbf{例 5（课件 p29--p32）。}求
\[
y''+\ee^xy'+(1+x^2)y=0,\qquad y(0)=1,\ y'(0)=0
\]
的幂级数解。
}

这里 $p(x)=\ee^x$、$q(x)=1+x^2$ 都在全平面解析，所以由定理 2.1，解在全平面解析，$\rho=\infty$。

设 $y=\sum a_nx^n$。关键是 $\ee^xy'$ 这一项：把 $\ee^x=\sum x^m/m!$ 与 $y'=\sum(k+1)a_{k+1}x^k$ 相乘，用卷积：
\[
\ee^xy'=\sum_{n=0}^{\infty}\Bigl(\sum_{k=0}^{n}\frac{(k+1)a_{k+1}}{(n-k)!}\Bigr)x^n .
\]
另外 $x^2y=\sum_{n=2}^{\infty}a_{n-2}x^n$。把三项合起来，$x^n$ 的系数为零给出
\[
(n+2)(n+1)a_{n+2}+\sum_{k=0}^{n}\frac{(k+1)a_{k+1}}{(n-k)!}+a_n+a_{n-2}=0,
\qquad n\geq0, \tag{2.3}
\]
其中约定 $a_{-1}=a_{-2}=0$（这样 $n=0,1$ 时最后一项自动消失）。逐项写开：
\begin{align*}
n=0:\quad & 2a_2+a_1+a_0=0,\\
n=1:\quad & 6a_3+2a_2+2a_1=0,\\
n=2:\quad & 12a_4+3a_3+3a_2+\tfrac12a_1+a_0=0,\\
n=3:\quad & 20a_5+4a_4+4a_3+a_2+\tfrac76a_1=0,\\
n=4:\quad & 30a_6+5a_5+5a_4+\tfrac32a_3+\tfrac43a_2+\tfrac1{24}a_1=0 .
\end{align*}
初值给 $a_0=1$、$a_1=0$。代入：
\[
a_2=-\frac12,\qquad a_3=\frac16,\qquad a_4=0,\qquad a_5=-\frac1{120},\qquad a_6=\frac{11}{720}.
\]
所以
\[
y(x)=1-\frac12x^2+\frac16x^3-\frac1{120}x^5+\frac{11}{720}x^6+\cdots
\]
注意 $x^4$ 项系数为零，所以写「前五项」时不能把 $x^4$ 算进去——\textbf{「非零项」三个字是有意义的}，本次作业也要求「前五个非零项」。

\aside{验算办法：把上面五项代回 $y''+\ee^xy'+(1+x^2)y$，展开到 $x^3$，残差应为 $O(x^4)$。这类「代回看残差」的检查不依赖任何技巧，是本章最可靠的自检。}

\subhdr{2.6 例 6：级数为什么会截断}

\callout{
\textbf{例 6（课件 p33--p35）。}求
\[
(x^2+1)y''-4xy'+6y=0
\]
在 $x=0$ 附近的级数解（前四个非零项）。
}

$P=x^2+1$，$P(0)=1\neq0$，所以 $x=0$ 是常点。$p=-4x/(x^2+1)$、$q=6/(x^2+1)$ 在 $|x|<1$ 内解析（因为 $x^2+1$ 的零点是 $\pm\ii$，到 $0$ 的距离为 $1$）。设 $y=\sum a_kx^k$，代入：
\[
(x^2+1)\sum_{k=2}^{\infty}k(k-1)a_kx^{k-2}
-4x\sum_{k=1}^{\infty}ka_kx^{k-1}
+6\sum_{k=0}^{\infty}a_kx^k=0 .
\]
第一项拆成两半：$x^2y''$ 升到 $x^k$，$y''$ 本身就是 $x^k$。整理后 $x^k$ 的系数为
\[
\bigl[k(k-1)-4k+6\bigr]a_k+(k+2)(k+1)a_{k+2}=0 .
\]
而 $k(k-1)-4k+6=k^2-5k+6=(k-2)(k-3)$，于是
\[
(k-2)(k-3)a_k+(k+2)(k+1)a_{k+2}=0
\quad\Longrightarrow\quad
a_{k+2}=-\frac{(k-2)(k-3)}{(k+2)(k+1)}a_k . \tag{2.4}
\]
代入 $k=0,1,2,3$：
\[
a_2=-3a_0,\qquad a_3=-\frac13a_1,\qquad a_4=0,\qquad a_5=0 .
\]
从 $k=4$ 起，$(k-2)(k-3)\neq0$，但 $a_4=a_5=0$，递推给出后面全是零。所以级数\textbf{截断}成多项式：
\[
y=a_0(1-3x^2)+a_1\Bigl(x-\frac13x^3\Bigr),\qquad |x|<1 .
\]
直接代回可以验证它确实是解：
\[
y_1=1-3x^2:\quad y_1''=-6,\quad
(x^2+1)(-6)-4x(-6x)+6(1-3x^2)=-6x^2-6+24x^2+6-18x^2=0 .
\]
同理可验 $y_2=x-x^3/3$。

\textbf{为什么截断？}看式 (2.4) 的分子 $(k-2)(k-3)$：它在 $k=2,3$ 时为零，把「$a_4$ 由 $a_2$ 决定、$a_5$ 由 $a_3$ 决定」这两条链\textbf{同时切断}。$a_4=a_5=0$ 之后，后面的递推只能给出零。这种「分子在某个整数处为零」的结构，是幂级数解变成多项式的唯一原因。

\warnbox{
这个例子同时说明：\textbf{实轴上没有奇点，不代表收敛半径是无穷}。$P=x^2+1$ 在实轴上没有零点，但它在 $\pm\ii$ 处为零；从 $0$ 到 $\pm\ii$ 的距离是 $1$，所以级数只在 $|x|<1$ 内保证收敛。复奇点也算数——这是本章最容易忘的一条。
}

\subhdr{2.7 常点法的边界}

课件 p28 最后一句是「The power series method is only applicable when $x=0$ is an ordinary point.」为什么？因为若 $P(x_0)=0$，则 $p=Q/P$、$q=R/P$ 在 $x_0$ 处不解析，代入 $y=\sum a_k(x-x_0)^k$ 后会出现负幂次项，无法与正幂次项合并，系数比较就无从谈起。

但「不能直接用」不等于「没有级数解」：下一节先把奇点分类，再对「正则奇点」把解的形状从整数幂改成 $x^r$ 乘幂级数——这就是 Euler 方程与 Frobenius 方法。

%=============================================================
\hdr{§3}{正则奇点与 Euler 方程}

这一节对应课件 L19--20 p3--p19。它要回答两个问题：奇点能不能「分级」，以及分级之后解长什么样。Euler 方程是分级之后最简单的样本——它的解是 $x^r$、$x^r\ln x$ 或 $x^\lambda\cos(\mu\ln x)$，正好预演了 Frobenius 方法的三种结局。

\subhdr{3.1 正则奇点：奇性「不太强」的精确含义}

\callout{
\textbf{定义（课件 p4）。}对方程 $P(x)y''+Q(x)y'+R(x)y=0$，若 $P(x_0)=0$，称 $x_0$ 为\textbf{奇点}。若还有
\[
\lim_{x\to x_0}(x-x_0)\frac{Q(x)}{P(x)}\quad\text{与}\quad\lim_{x\to x_0}(x-x_0)^2\frac{R(x)}{P(x)}
\]
都有限，则称 $x_0$ 为\textbf{正则奇点}（regular singular point）；否则称为\textbf{非正则奇点}（irregular singular point）。
}

课件 p5 给了一个等价的说法：记 $p=Q/P$、$q=R/P$，则 $x_0$ 正则 $\iff$ $(x-x_0)p(x)$ 与 $(x-x_0)^2q(x)$ 都在 $x_0$ 解析。这个说法比极限好用，因为它直接告诉你「乘完之后是幂级数」。

\aside{为什么偏偏乘 $(x-x_0)$ 和 $(x-x_0)^2$？把方程除以 $P$ 再乘 $(x-x_0)^2$：
\[
(x-x_0)^2y''+(x-x_0)\bigl[(x-x_0)p(x)\bigr]y'+\bigl[(x-x_0)^2q(x)\bigr]y=0 .
\]
如果两个方括号都解析，这个方程就是「Euler 型的三项 + 解析系数」。Euler 方程的解是 $x^r$，系数解析只会把这个解「微调」成 $x^r$ 乘幂级数——这就是 §4 设 $y=x^r\sum a_nx^n$ 的全部理由。乘一次不够、乘三次多余，$(x-x_0)$ 与 $(x-x_0)^2$ 正好把 $p$ 的一阶极点、$q$ 的二阶极点消掉。}

\warnbox{
\textbf{多项式系数的快速判据。}若 $P,Q,R$ 是多项式，$x_0$ 是 $P$ 的 $m$ 阶零点：当 $Q$ 在 $x_0$ 的零点阶数 $\geq m-1$、$R$ 的零点阶数 $\geq m-2$ 时，$x_0$ 是正则奇点；只要有一个不满足，就是非正则。理由：$Q/P$ 的极点阶数至多 $1$，$R/P$ 的极点阶数至多 $2$，正好被乘的因子抵消。
}

\subhdr{3.2 分类的例子（课件 p6--p7）}

\callout{
\textbf{例。}把
\[
(x^2-1)^2y''+(x+1)y'-y=0
\]
的奇点分类。
}

$P=(x^2-1)^2=(x-1)^2(x+1)^2$，奇点为 $x=\pm1$。除以 $P$：
\[
p(x)=\frac{x+1}{(x-1)^2(x+1)^2}=\frac{1}{(x-1)^2(x+1)},\qquad
q(x)=\frac{-1}{(x-1)^2(x+1)^2}.
\]
\textbf{在 $x=1$：}
\[
(x-1)p(x)=\frac{1}{(x-1)(x+1)}\to\infty,
\]
极限不是有限数，所以 $x=1$ 是\textbf{非正则奇点}。

\textbf{在 $x=-1$：}
\[
(x+1)p(x)=\frac{1}{(x-1)^2},\qquad
(x+1)^2q(x)=\frac{-1}{(x-1)^2},
\]
两个都在 $x=-1$ 解析（分母在 $-1$ 处为 $(-2)^2=4\neq0$），所以 $x=-1$ 是\textbf{正则奇点}。

同一个方程的两个奇点，一个正则一个非正则——这说明「正则」不是方程整体的性质，而是\textbf{每个奇点各自}的性质。

\subhdr{3.3 为什么不能再用普通幂级数}

课件 p8 之前先用一个反例说明「非解析解」确实会出现：
\[
x^2y''-6y=0
\]
是 Euler 方程，两个解是 $y_1=x^3$、$y_2=x^{-2}$。$y_2$ 在 $x=0$ 处根本没有有限值，当然不能写成 $\sum a_nx^n$。所以奇点处的解必须允许一个「非整数幂」的因子——这就是 $x^r$。

\subhdr{3.4 Euler 方程：三种根，三种解}

\callout{
\textbf{Euler 方程（课件 p8）。}
\[
ax^2y''+bxy'+cy=0,\qquad x>0 .
\]
设 $y=x^r$，则 $y'=rx^{r-1}$、$y''=r(r-1)x^{r-2}$，代入得
\[
x^r\bigl[ar(r-1)+br+c\bigr]=0 .
\]
因为 $x>0$ 时 $x^r\neq0$，所以必须有
\[
\boxed{\,ar(r-1)+br+c=0\,}
\]
这个二次方程叫\textbf{指标方程}（indicial equation），它的两个根叫\textbf{指标}（exponents）。
}

三种根对应三种解，逐条推导如下。

\callout{
\textbf{情形 1：两个相异实根 $r_1\neq r_2$。}
$y_1=x^{r_1}$、$y_2=x^{r_2}$ 都是解，且 $x^{r_1}/x^{r_2}=x^{r_1-r_2}$ 不是常数，所以线性无关：
\[
y=C_1x^{r_1}+C_2x^{r_2}.
\]

\textbf{情形 2：一对共轭复根 $r=\lambda\pm\ii\mu$（$\mu\neq0$）。}
先把复解写成实函数。用 $x^r=\ee^{r\ln x}$：
\[
x^{\lambda+\ii\mu}=\ee^{(\lambda+\ii\mu)\ln x}
=\ee^{\lambda\ln x}\bigl[\cos(\mu\ln x)+\ii\sin(\mu\ln x)\bigr]
=x^{\lambda}\bigl[\cos(\mu\ln x)+\ii\sin(\mu\ln x)\bigr].
\]
实部与虚部分别是解（方程是实系数的），于是
\[
y=x^{\lambda}\bigl[C_1\cos(\mu\ln x)+C_2\sin(\mu\ln x)\bigr].
\]

\textbf{情形 3：重根 $r_1=r_2=r$。}
只有一个解 $y_1=x^r$。第二个解用降阶（§6 会从方程组角度再推一次；这里用一阶线性方程的积分因子）：
设 $y_2=vy_1$，代入 $y''+\dfrac{b}{ax}y'+\dfrac{c}{ax^2}y=0$，$v$ 的项会消掉，得到
\[
v''+\Bigl(\frac{2r}{x}+\frac{b}{ax}\Bigr)v'=0
\quad\Longrightarrow\quad
\frac{v''}{v'}=-\frac{2r+b/a}{x}.
\]
重根满足 $2r=1-b/a$，所以 $2r+b/a=1$，从而 $v''/v'=-1/x$，积分得 $v'=C/x$，$v=C\ln x$。于是
\[
y=x^r(C_1+C_2\ln x).
\]
}

\begin{center}
\small
\begin{tabular}{@{}p{3.2cm}p{4.6cm}p{7.2cm}@{}}
\hline
\thead{指标根} & \thead{基本解} & \thead{通解} \\
\hline
$r_1\neq r_2$ 实 & $x^{r_1},\ x^{r_2}$ & $C_1x^{r_1}+C_2x^{r_2}$ \\
$r=\lambda\pm\ii\mu$ & $x^\lambda\cos(\mu\ln x),\ x^\lambda\sin(\mu\ln x)$ & $x^\lambda[C_1\cos(\mu\ln x)+C_2\sin(\mu\ln x)]$ \\
$r_1=r_2=r$ & $x^r,\ x^r\ln x$ & $x^r(C_1+C_2\ln x)$ \\
\hline
\end{tabular}
\end{center}

这张表只是把上面三个推导的结论并排抄了一遍；每一行的第二列都能由第一列代回 $ax^2y''+bxy'+cy=0$ 验证。第三行最容易记错：$\ln x$ 乘的是 $x^r$，不是 $x^{r+1}$。

\subhdr{3.5 三个例子}

\callout{
\textbf{例（初值问题，课件 p10--p11）。}解
\[
2x^2y''+3xy'-15y=0,\qquad y(1)=0,\ y'(1)=1 .
\]
\textbf{解。}指标方程
\[
2r(r-1)+3r-15=2r^2+r-15=(2r-5)(r+3)=0,
\]
根为 $r_1=5/2$、$r_2=-3$。通解
\[
y=C_1x^{5/2}+C_2x^{-3}.
\]
代入初值（$x=1$ 时 $x^{5/2}=x^{-3}=1$）：
\[
0=C_1+C_2,\qquad 1=\frac52C_1-3C_2 .
\]
由第一式 $C_2=-C_1$，代入第二式得 $\frac52C_1+3C_1=\frac{11}{2}C_1=1$，所以
\[
C_1=\frac{2}{11},\qquad C_2=-\frac{2}{11},\qquad
y=\frac{2}{11}x^{5/2}-\frac{2}{11}x^{-3}.
\]
}

\warnbox{
课件 p11 把第二式写成「$1=y'(1)=\frac52c_1-3c_3$」，$c_3$ 是笔误，应为 $c_2$。另外这里 $x>0$ 是必须的：$x^{-3}$ 在 $0$ 处无定义，$x^{5/2}$ 对 $x<0$ 不是实数。
}

\callout{
\textbf{例（复根，课件 p13）。}解 $x^2y''+3xy'+4y=0$。

\textbf{解。}$r(r-1)+3r+4=r^2+2r+4=0$，根为 $r=-1\pm\ii\sqrt3$。所以
\[
y=x^{-1}\bigl[C_1\cos(\sqrt3\ln x)+C_2\sin(\sqrt3\ln x)\bigr].
\]
}

\callout{
\textbf{例（重根，课件 p15）。}解 $x^2y''-7xy'+16y=0$。

\textbf{解。}$r(r-1)-7r+16=r^2-8r+16=(r-4)^2=0$，$r=4$。所以
\[
y=x^4(C_1+C_2\ln x).
\]
}

\subhdr{3.6 $x<0$ 怎么办：$\lvert x\rvert$ 的来历}

课件 p16--p17 的处理值得完整看懂，因为作业题只要求 $x>0$，但考试可能问「为什么」。设 $x<0$，令 $\eta=-x>0$，并定义 $u(\eta)=y(-\eta)$。由链式法则
\[
u'(\eta)=-y'(x),\qquad u''(\eta)=y''(x).
\]
代入 $ax^2y''+bxy'+cy=0$，注意 $x=-\eta$：
\[
a\eta^2u''+b\eta u'+cu=0 .
\]
这是关于 $\eta>0$ 的同型 Euler 方程，解是 $\eta^{r}$ 等。再把 $\eta=-x$ 换回来，就得到 $|x|^r$。三种情形合起来：
\[
y=C_1|x|^{r_1}+C_2|x|^{r_2},\qquad
y=|x|^\lambda\bigl[C_1\cos(\mu\ln|x|)+C_2\sin(\mu\ln|x|)\bigr],\qquad
y=|x|^r(C_1+C_2\ln|x|).
\]
这些解在\textbf{任何不包含 $0$ 的区间}上都成立；$x>0$ 与 $x<0$ 是两个独立的区间，解不能从一边「接」到另一边（除非 $r$ 是整数，那时 $|x|^r$ 其实是 $x^r$ 的延拓）。

\subhdr{3.7 平移版 Euler 方程}

若方程是
\[
a(x-x_0)^2y''+b(x-x_0)y'+cy=0,
\]
令 $t=x-x_0$ 就化回原点情形，指标方程不变，解把 $x$ 换成 $|x-x_0|$：
\[
y=C_1|x-x_0|^{r_1}+C_2|x-x_0|^{r_2},\quad\text{等等}.
\]

\callout{
\textbf{例（课件 p19）。}解
\[
3(x+6)^2y''+25(x+6)y'-16y=0
\]
（在不含 $x=-6$ 的区间上）。

\textbf{解。}指标方程 $3r(r-1)+25r-16=3r^2+22r-16=0$，即 $(3r-2)(r+8)=0$，根为 $r_1=2/3$、$r_2=-8$。所以
\[
y=C_1|x+6|^{2/3}+C_2|x+6|^{-8}.
\]
}

%=============================================================
\hdr{§4}{Frobenius 方法：正则奇点处的级数解}

这是本章的核心，对应课件 L19--20 p20--p34。它把 §3 的「$x^r$」与 §2 的「幂级数」拼在一起：\textbf{在正则奇点处，解是 $x^r$ 乘一个幂级数}。本节把课件跳过的三处补上：指标方程怎么来、递推的分母为什么是 $(n+r-r_1)(n+r-r_2)$、以及整数差时第二个解到底长什么样。

\subhdr{4.1 为什么设 $y=x^r\sum a_nx^n$}

设 $x=0$ 是正则奇点。由定义，$xp(x)$ 与 $x^2q(x)$ 在 $0$ 附近解析，设
\[
xp(x)=\sum_{k=0}^{\infty}p_kx^k,\qquad x^2q(x)=\sum_{k=0}^{\infty}q_kx^k,\qquad |x|<\rho .
\]
把方程 $y''+py'+qy=0$ 乘以 $x^2$：
\[
x^2y''+x\bigl(xp(x)\bigr)y'+\bigl(x^2q(x)\bigr)y=0 . \tag{4.1}
\]
如果 $p_k=q_k=0$（$k\geq1$），(4.1) 就是 Euler 方程 $x^2y''+p_0xy'+q_0y=0$，解是 $x^r$。一般的 $p_k,q_k$ 只是给 Euler 方程加了「解析的扰动」，所以自然设想解是「Euler 的解乘一个幂级数」：
\[
y=x^r\sum_{n=0}^{\infty}a_nx^n=\sum_{n=0}^{\infty}a_nx^{n+r},\qquad a_0\neq0 .
\]
这就是 \textbf{Frobenius 级数}。注意 $r$ 不一定是整数，甚至不一定是实数——它由方程在奇点处的「首项」决定。

\subhdr{4.2 代入：指标方程与递推公式}

逐项求导：
\[
y'=\sum_{n=0}^{\infty}(n+r)a_nx^{n+r-1},\qquad
y''=\sum_{n=0}^{\infty}(n+r)(n+r-1)a_nx^{n+r-2}.
\]
三项分别算：
\begin{align*}
x^2y''&=\sum_{n=0}^{\infty}(n+r)(n+r-1)a_nx^{n+r},\\
x\bigl(xp\bigr)y'&=\Bigl(\sum_{k=0}^{\infty}p_kx^k\Bigr)\Bigl(\sum_{m=0}^{\infty}(m+r)a_mx^{m+r}\Bigr)
=\sum_{n=0}^{\infty}\Bigl(\sum_{m=0}^{n}p_{n-m}(m+r)a_m\Bigr)x^{n+r},\\
\bigl(x^2q\bigr)y&=\Bigl(\sum_{k=0}^{\infty}q_kx^k\Bigr)\Bigl(\sum_{m=0}^{\infty}a_mx^{m+r}\Bigr)
=\sum_{n=0}^{\infty}\Bigl(\sum_{m=0}^{n}q_{n-m}a_m\Bigr)x^{n+r}.
\end{align*}
第二个等号用的就是卷积（§1.3）：$x^{n+r}$ 的系数是「两个级数下标相加等于 $n$」的所有乘积之和。

把三式相加，令 $x^{n+r}$ 的系数为零：
\[
(n+r)(n+r-1)a_n+\sum_{m=0}^{n}\bigl[p_{n-m}(m+r)+q_{n-m}\bigr]a_m=0,
\qquad n=0,1,2,\dots \tag{4.2}
\]

\callout{
\textbf{$n=0$：指标方程。}式 (4.2) 在 $n=0$ 时只剩 $m=0$ 一项：
\[
\bigl[r(r-1)+p_0r+q_0\bigr]a_0=0 .
\]
因为 $a_0\neq0$，所以
\[
\boxed{\,r(r-1)+p_0r+q_0=0\,} \tag{4.3}
\]
这就是\textbf{指标方程}。它的两个根 $r_1,r_2$ 叫指标；规定 $r_1\geq r_2$（实根时）。
}

\callout{
\textbf{$n\geq1$：递推公式。}把 (4.2) 中 $m=n$ 的那一项（含 $a_n$）单独拿出来：
\[
\bigl[(n+r)(n+r-1)+(n+r)p_0+q_0\bigr]a_n
=-\sum_{m=0}^{n-1}\bigl[p_{n-m}(m+r)+q_{n-m}\bigr]a_m . \tag{4.4}
\]
只要方括号不为零，就可以解出
\[
a_n=-\frac{1}{(n+r)(n+r-1)+(n+r)p_0+q_0}
\sum_{m=0}^{n-1}\bigl[p_{n-m}(m+r)+q_{n-m}\bigr]a_m .
\]
}

\warnbox{
课件 p24 写「Case 1. If $a_0=0$, then $a_1=a_2=\cdots=0$, resulting in the zero solution」。这句话的本意是「$a_0=0$ 得不到新的解」，但字面上不严谨：若 $a_0=0$，级数从更高的幂次开始，例如 $a_1x^{1+r}+\cdots=x^{r+1}(a_1+\cdots)$，这等价于把指数换成 $r+1$。所以正确的说法是：\textbf{取定指数 $r$ 时约定 $a_0\neq0$ 只是选取「以 $x^r$ 为最低次幂」的那一支解；$a_0=0$ 的级数不是零解，而是另一个指数的解。}做题时不必纠结，按 $a_0=1$ 算即可。
}

\subhdr{4.3 分母的秘密：$(n+r-r_1)(n+r-r_2)$}

式 (4.4) 的分母看起来是个普通二次式，其实它就是指标多项式在 $n+r$ 处的值。记
\[
F(s)=s(s-1)+p_0s+q_0 .
\]
由指标方程，$F(s)=(s-r_1)(s-r_2)$。把 $s=n+r$ 代进去：
\[
(n+r)(n+r-1)+(n+r)p_0+q_0=F(n+r)=(n+r-r_1)(n+r-r_2). \tag{4.5}
\]
这个恒等式是整章的「机关」，它一次性解释了三种情形：

\begin{itemize}
  \item \textbf{取 $r=r_1$}：分母 $(n+r_1-r_1)(n+r_1-r_2)=n\,(n+r_1-r_2)$。$n\geq1$ 时 $n\neq0$；又 $r_1-r_2\geq0$，所以 $n+r_1-r_2>0$。分母永不为零，$r_1$ 这一支\textbf{总能算到底}。
  \item \textbf{取 $r=r_2$}：分母 $(n+r_2-r_1)(n+r_2-r_2)=(n-(r_1-r_2))\,n$。当 $n=r_1-r_2$ 时（若这个差是正整数），分母为零，递推在这一步\textbf{卡住}。
  \item \textbf{若 $r_1-r_2$ 不是整数}：两支都不会碰到零分母，所以两个 Frobenius 级数都存在。
\end{itemize}

卡住的那一步有两种结局。设 $N_0=r_1-r_2$ 是正整数，在 $n=N_0$ 处式 (4.4) 变成
\[
0\cdot a_{N_0}=-\sum_{m=0}^{N_0-1}\bigl[p_{N_0-m}(m+r_2)+q_{N_0-m}\bigr]a_m .
\]
右边若不为零，就矛盾——说明「纯 Frobenius 级数」这一形式不够用，第二个解必须带 $\ln x$；右边若恰好为零，则 $a_{N_0}$ 可以任取，第二个解仍是纯 Frobenius 级数（§4.5 的例子就是这一种）。这就是课件 p27--p28 三种情形的来源。

\subhdr{4.4 Frobenius 定理与三种情形}

\callout{
\textbf{定理 1.1（Frobenius 定理，课件 p31）。}若 $x_0$ 是方程
\[
a_2(x)y''+a_1(x)y'+a_0(x)y=0
\]
的正则奇点，则至少存在一个形如
\[
y=(x-x_0)^r\sum_{n=0}^{\infty}a_n(x-x_0)^n,\qquad a_0\neq0
\]
的级数解，其中 $r=r_1$ 是指标方程的较大根。并且这个级数在
\[
0<x-x_0<\rho
\]
内收敛，$\rho$ 是 $x_0$ 到方程\textbf{最近的其他奇点（实的或复的）}的距离。
}

定理只保证「一个」解。第二个解的存在性与形式由下面三条给出（课件 p27--p28 陈述，证明略——它本质上是 §4.3 分母结构的推论）：

\callout{
\textbf{情形 A：$r_1-r_2$ 不是整数。}
\[
y_2=x^{r_2}\sum_{n=0}^{\infty}b_nx^n,\qquad b_0\neq0 .
\]
两个解都是纯 Frobenius 级数，收敛半径同样是 $\rho$。

\textbf{情形 B：$r_1=r_2$。}
\[
y_2=y_1\ln x+x^{r_2}\sum_{n=0}^{\infty}b_nx^n .
\]

\textbf{情形 C：$r_1-r_2$ 是正整数。}
\[
y_2=a\,y_1\ln x+x^{r_2}\sum_{n=0}^{\infty}b_nx^n,
\]
其中常数 $a$ 可能为零。$a=0$ 时 $y_2$ 仍是纯 Frobenius 级数（没有 $\ln x$）；$a\neq0$ 时必须带对数项。
}

\warnbox{
情形 C 最容易背错：不是「差为正整数就一定带 $\ln x$」，而是「\textbf{可能}带，要看递推在 $n=r_1-r_2$ 那一步是否自动满足」。教材 9.33 的 Remark 明确写了「$c$ may turn out to be zero」，§4.5 就是一个 $c=0$ 的实例。
}

\subhdr{4.5 例：$xy''+4y'-xy=0$ 的两个解}

\callout{
\textbf{例（课件 p32--p34）。}求
\[
xy''+4y'-xy=0,\qquad x>0
\]
在正则奇点 $x=0$ 处的级数解。
}

先化标准形：$y''+\dfrac4xy'-y=0$，所以 $p=4/x$、$q=-1$。于是
\[
xp(x)=4,\qquad x^2q(x)=-x^2,
\]
即 $p_0=4$、$q_0=0$，其余 $p_k=q_k=0$（$k\geq1$）。$x=0$ 是正则奇点，且 $xp$、$x^2q$ 是多项式，所以 $\rho=\infty$。

指标方程：
\[
r(r-1)+4r=r^2+3r=r(r+3)=0\quad\Longrightarrow\quad r_1=0,\ r_2=-3 .
\]

\aside{课件 p33 把这一步写成「$r(r-1)+4r=r+3r=r(r+3)=0$」，中间的「$r+3r$」漏了平方，应为 $r^2+3r$。结论 $r(r+3)=0$ 是对的。}

设 $y=\sum a_nx^{n+r}$，逐项代入：
\begin{align*}
xy''&=\sum_{n=0}^{\infty}(n+r)(n+r-1)a_nx^{n+r-1},\\
4y'&=\sum_{n=0}^{\infty}4(n+r)a_nx^{n+r-1},\\
-xy&=-\sum_{n=0}^{\infty}a_nx^{n+r+1}=-\sum_{n=2}^{\infty}a_{n-2}x^{n+r-1}.
\end{align*}
合并 $x^{n+r-1}$ 的系数：
\[
\bigl[(n+r)(n+r-1)+4(n+r)\bigr]a_n-a_{n-2}=0
\quad\Longrightarrow\quad
(n+r)(n+r+3)a_n=a_{n-2}, \tag{4.6}
\]
其中约定 $a_{-1}=a_{-2}=0$。这正是式 (4.5) 的特例：$(n+r)(n+r+3)=(n+r-r_1)(n+r-r_2)$，因为 $r_1=0,r_2=-3$。

\textbf{第一支 $r=0$。}式 (4.6) 变成 $n(n+3)a_n=a_{n-2}$。$n=0$ 给 $0=0$（指标方程）；$n=1$ 给 $1\cdot4\,a_1=0$，所以 $a_1=0$。于是奇数项全为零，偶数项满足
\[
a_{2k}=\frac{a_{2k-2}}{2k(2k+3)}
=\frac{a_0}{2^k\,k!\,5\cdot7\cdots(2k+3)} .
\]
取 $a_0=1$：
\[
y_1=1+\frac{x^2}{10}+\frac{x^4}{280}+\frac{x^6}{15120}+\cdots
\]
（$2\cdot5=10$，$2\cdot4\cdot5\cdot7=280$，$2\cdot4\cdot6\cdot5\cdot7\cdot9=15120$）。

\textbf{第二支 $r=-3$。}式 (4.6) 变成 $(n-3)n\,a_n=a_{n-2}$。逐项看：
\begin{align*}
n=1:&\quad (-2)(1)a_1=0 \ \Longrightarrow\ a_1=0,\\
n=2:&\quad (-1)(2)a_2=a_0 \ \Longrightarrow\ a_2=-\frac{a_0}{2},\\
n=3:&\quad 0\cdot a_3=a_1=0 \ \Longrightarrow\ a_3\ \text{任意},\\
n=4:&\quad (1)(4)a_4=a_2 \ \Longrightarrow\ a_4=-\frac{a_0}{8},\\
n=5:&\quad (2)(5)a_5=a_3 \ \Longrightarrow\ a_5=\frac{a_3}{10},\\
n=6:&\quad (3)(6)a_6=a_4 \ \Longrightarrow\ a_6=-\frac{a_0}{144}.
\end{align*}
取 $a_0=1$、$a_3=0$，得
\[
y_2=x^{-3}\Bigl(1-\frac{x^2}{2}-\frac{x^4}{8}-\frac{x^6}{144}-\cdots\Bigr).
\]

\textbf{为什么可以取 $a_3=0$，取了之后还少了什么？}若 $a_3\neq0$，那么
\[
x^{-3}\cdot a_3\Bigl(x^3+\frac{x^5}{10}+\cdots\Bigr)=a_3\Bigl(1+\frac{x^2}{10}+\cdots\Bigr)=a_3y_1,
\]
只是把 $y_1$ 混进来，不产生新解。所以「$a_3$ 任意」这一自由度不是第二个独立解，而是一个重复；取 $a_3=0$ 就得到与 $y_1$ 线性无关的那一支。

\textbf{这个例子属于情形 C，但 $a=0$。}两根之差 $r_1-r_2=3$ 是正整数，按定理本应允许 $\ln x$；可是递推在 $n=3$ 那一步右边是 $a_1=0$，方程 $0\cdot a_3=0$ 自动成立，没有矛盾，所以第二个解不需要 $\ln x$。这正是教材 9.33 说的「$c$ may turn out to be zero」。

\warnbox{
\textbf{自检：}把 $y_1=1+x^2/10+x^4/280$ 代回 $xy''+4y'-xy$，展开到 $x^4$，残差应为 $O(x^6)$。这个检查只用求导和合并同类项，不依赖任何定理，建议每次算完都做一遍。
}

\subhdr{4.6 收敛半径 $\rho$：到最近奇点的距离}

定理 1.1 里的 $\rho$ 不是算完级数再量的，而是\textbf{从方程本身读出来的}：它是展开点到方程最近奇点（实或复）的距离。下面三个例子把它说清楚。

\begin{center}
\small
\begin{tabular}{@{}p{6.4cm}p{3.4cm}p{5.8cm}@{}}
\hline
\thead{方程} & \thead{奇点} & \thead{$\rho$（展开点 $0$）} \\
\hline
$xy''+4y'-xy=0$ & 只有 $x=0$ & $\infty$（无其他奇点）\\
$(x^2+1)y''-4xy'+6y=0$ & $x=\pm\ii$ & $1$（到 $\pm\ii$ 的距离）\\
$2x^2y''+(2x^2+x)y'+(3x-1)y=0$ & 只有 $x=0$ & $\infty$（本次作业）\\
\hline
\end{tabular}
\end{center}

表前一句：这张表把三个方程的奇点列出来，$\rho$ 就是「展开点到表里最近那个奇点的距离」。逐行读法：第一行只有 $0$ 一个奇点，没有「其他奇点」可量，所以 $\rho=\infty$；第二行 $P=x^2+1$ 的零点是 $\pm\ii$，$|0-\ii|=1$，所以 $\rho=1$；第三行 $P=2x^2$ 只有零点 $0$，所以 $\rho=\infty$。边界行：$\rho$ 是「至少」保证的半径，实际半径可能更大（例如系数有抵消时），但考试按定理给的下界写。

\begin{center}
\begin{tikzpicture}[>=Latex,font=\small,scale=1.0]
\draw[->] (-2.2,0) -- (2.4,0) node[right] {$\mathrm{Re}$};
\draw[->] (0,-2.0) -- (0,2.2) node[above] {$\mathrm{Im}$};
\draw[rulepurple,dashed] (0,0) circle (1.6);
\fill (0,0) circle (1.6pt) node[below right] {$0$};
\fill (0,1.6) circle (1.6pt) node[right] {$\ii$};
\fill (0,-1.6) circle (1.6pt) node[right] {$-\ii$};
\node at (0.95,0.95) {$\rho=1$};
\end{tikzpicture}\\[2pt]
{\small\color{inkgray}图 2：$(x^2+1)y''-4xy'+6y=0$ 的复奇点 $\pm\ii$ 与收敛圆。实轴上没有奇点，但级数仍只在 $|x|<1$ 内保证收敛。}
\end{center}

\warnbox{
「最小收敛区间」（minimal interval of convergence）这个说法，指的就是定理保证的那个区间：$0<x-x_0<\rho$。它不是「实际能收敛的最大区间」，而是「不用逐项检验就能保证的区间」。做题时先算 $\rho$，再写区间；若 $\rho=\infty$，就写「对整个 $x>x_0$ 成立」。
}

\subhdr{4.7 第二解为什么可能带 $\ln x$：用降阶公式看穿}

§4.4 的三种情形是「结论」，这一节用你学过的降阶公式把它推出来，顺便解释 §4.5 里「$r=-3$ 却没有 $\ln x$」不是巧合。

设 $y_1$ 是 $y''+py'+qy=0$ 的一个解，$W=y_1y_2'-y_1'y_2$。由 Abel 公式（§6.4）
\[
W=C\ee^{-\int p\dd x},
\]
而
\[
\Bigl(\frac{y_2}{y_1}\Bigr)'=\frac{y_1y_2'-y_1'y_2}{y_1^2}=\frac{W}{y_1^2}
=\frac{C\ee^{-\int p}}{y_1^2}.
\]
所以
\[
y_2=y_1\int^x\frac{C\ee^{-\int p(t)\dd t}}{y_1(t)^2}\dd t . \tag{4.7}
\]
这就是降阶公式（L7 学过）。现在把正则奇点处的信息代进去。设
\[
p(x)=\frac{p_0}{x}+\text{(解析部分)},\qquad y_1=x^{r_1}\varphi(x),\quad \varphi(0)\neq0 .
\]
则 $\ee^{-\int p}=x^{-p_0}\psi(x)$（$\psi$ 解析、$\psi(0)\neq0$），$y_1^2=x^{2r_1}\varphi^2$。被积函数是
\[
\frac{C\ee^{-\int p}}{y_1^2}=C\,x^{-p_0-2r_1}\cdot\frac{\psi}{\varphi^2}.
\]
由指标方程 $r_1+r_2=1-p_0$，指数可以改写：
\[
-p_0-2r_1=r_2-r_1-1 .
\]
于是
\[
\frac{C\ee^{-\int p}}{y_1^2}=C\,x^{r_2-r_1-1}\,\Psi(x),\qquad \Psi\ \text{解析},\ \Psi(0)\neq0 . \tag{4.8}
\]

三种情形现在一目了然：
\begin{itemize}
  \item \textbf{$r_1-r_2\notin\Z$}：指数 $r_2-r_1-1$ 不是负整数，逐项积分得 $x^{r_2-r_1}$ 乘解析函数，再乘 $y_1=x^{r_1}\varphi$，得到 $y_2=x^{r_2}\times$（解析函数）。没有 $\ln x$。
  \item \textbf{$r_1=r_2$}：指数是 $-1$，积分给出 $\ln x$。所以重根必带对数。
  \item \textbf{$r_1-r_2=N\in\N$}：指数是 $-N-1$。逐项积分时，$x^{-N-1}$ 项给出 $x^{-N}$，乘 $y_1$ 后是 $x^{r_2}$，仍是纯级数；但如果 $\Psi$ 的展开里\textbf{恰有} $x^N$ 项，那么 $x^{-N-1}\cdot x^N=x^{-1}$，积分就产生 $\ln x$。所以「差为正整数」时，$\ln x$ 出不出现，取决于 $\Psi$ 的第 $N$ 个系数是否为零——这正是 §4.3 里「整数步的右边是否为零」的另一种说法。
\end{itemize}

\textbf{回到 §4.5 的例子。}$xy''+4y'-xy=0$，$r_1=0$、$r_2=-3$、$N=3$，$y_1=1+\frac{x^2}{10}+\cdots$。由 $p=4/x$ 得 $\ee^{-\int p}=x^{-4}$，于是
\[
\frac{\ee^{-\int p}}{y_1^2}=\frac{x^{-4}}{(1+\frac{x^2}{10}+\cdots)^2}
=x^{-4}\Bigl(1-\frac{x^2}{5}+\cdots\Bigr)
=x^{-4}-\frac15x^{-2}+\cdots
\]
只含偶次幂，没有 $x^{-1}$ 项（因为 $N=3$ 是奇数，而 $\Psi$ 只有偶次项）。所以第二个解不带 $\ln x$，与 §4.5 算出的 $y_2=x^{-3}(1-\frac{x^2}{2}-\cdots)$ 一致。

\warnbox{
\textbf{结论：}$\ln x$ 不是「整数差」的必然结果，而是被积函数 (4.8) 里 $x^{-1}$ 项的产物。重根时它必然出现；差为正整数时它可能出现也可能不出现；差非整数时它不出现。把这句话记住，比背三条公式可靠。
}

%=============================================================
\qbreak
\hdr{§5}{作业题全解：$2x^2y''+(2x^2+x)y'+(3x-1)y=0$}
\callout{本节把本章唯一的一道作业题\textbf{完整做完}，用途是示范"一手怎么写"。想要无提示自测，请用配套《习题册》里的同一题（编号 L19--08）：那里只给题干与出处，不给任何步骤。\textbf{先自己做一遍，再回来看这一节}，效果最好。}

\qstem{（课件 p42 第 1 题；教材 Constanda p242 习题 1，书内页码 p234）\textbf{Construct two linearly independent solutions for $x>0$. Compute the first five nonzero terms in each of these solutions and determine the minimal interval of convergence for the power series occurring in them.}
\[
2x^2y''+(2x^2+x)y'+(3x-1)y=0 .
\]}

\textbf{题目要求三件事：}（1）构造两个线性无关解；（2）每个解写出前五个\textbf{非零}项；（3）确定幂级数的\textbf{最小收敛区间}。下面按「判断点 $\to$ 指标方程 $\to$ 递推 $\to$ 两个解 $\to$ 收敛区间 $\to$ 验算」的顺序走。

\subhdr{5.1 第 0 步：化标准形，确认是正则奇点}

方程的首项系数 $P(x)=2x^2$，它在 $x=0$ 有二阶零点，所以 $x=0$ 是奇点；$Q=2x^2+x$、$R=3x-1$ 在 $0$ 处解析。除以 $P$：
\[
y''+\frac{2x^2+x}{2x^2}y'+\frac{3x-1}{2x^2}y=0
\quad\Longrightarrow\quad
p(x)=1+\frac{1}{2x},\qquad q(x)=\frac{3}{2x}-\frac{1}{2x^2}.
\]
再算两个「消奇性」的量：
\[
xp(x)=x+\frac12,\qquad x^2q(x)=\frac{3x-1}{2}=\frac32x-\frac12 .
\]
两个都是多项式，在 $x=0$ 解析，所以 $x=0$ 是\textbf{正则奇点}；并且
\[
p_0=\lim_{x\to0}xp(x)=\frac12,\qquad q_0=\lim_{x\to0}x^2q(x)=-\frac12 .
\]

\warnbox{
\textbf{第一个坑：}指标方程里的 $p_0,q_0$ 是 $xp$、$x^2q$ 的常数项，\textbf{不是} $Q/P$、$R/P$ 的系数。本题若直接拿 $p=1+\frac{1}{2x}$ 的「常数项」$1$ 当 $p_0$，指标方程就错了。记住顺序：先除 $P$ 得 $p,q$，再乘 $x$、$x^2$，最后取 $x\to0$。
}

\subhdr{5.2 第 1 步：指标方程与两根}

把 $p_0=\frac12$、$q_0=-\frac12$ 代入指标方程：
\[
r(r-1)+\frac12r-\frac12=0
\quad\Longrightarrow\quad
2r^2-r-1=0
\quad\Longrightarrow\quad
(2r+1)(r-1)=0 .
\]
两根为
\[
r_1=1,\qquad r_2=-\frac12 .
\]
差 $r_1-r_2=\frac32$ \textbf{不是整数}，属于 §4.4 的\textbf{情形 A}：两个解都是纯 Frobenius 级数，都不需要 $\ln x$。这一点在看递推之前就能确定，是本题的「路线判断」。

\aside{若两根之差是整数，就要先算 $r_1$ 那一支，再看整数步 $n=r_1-r_2$ 处递推是否自动满足（§4.3）。本题差为 $3/2$，没有这个麻烦，直接对两个根各算一支即可。}

\subhdr{5.3 第 2 步：代入，得到递推}

设
\[
y=\sum_{n=0}^{\infty}a_nx^{n+r},\qquad
y'=\sum_{n=0}^{\infty}(n+r)a_nx^{n+r-1},\qquad
y''=\sum_{n=0}^{\infty}(n+r)(n+r-1)a_nx^{n+r-2}.
\]
逐项乘进方程。注意 $2x^2y'$ 与 $3xy$ 会把幂次抬高一次，所以要单独处理：
\begin{align*}
2x^2y''&=\sum_{n=0}^{\infty}2(n+r)(n+r-1)a_nx^{n+r},\\
xy'&=\sum_{n=0}^{\infty}(n+r)a_nx^{n+r},\\
2x^2y'&=\sum_{n=0}^{\infty}2(n+r)a_nx^{n+r+1}
=\sum_{n=1}^{\infty}2(n+r-1)a_{n-1}x^{n+r},\\
3xy&=\sum_{n=0}^{\infty}3a_nx^{n+r+1}
=\sum_{n=1}^{\infty}3a_{n-1}x^{n+r},\\
-y&=-\sum_{n=0}^{\infty}a_nx^{n+r}.
\end{align*}
把 $x^{n+r}$ 的系数合并。$n=0$ 时只有前两项与最后一项有贡献：
\[
\bigl[2r(r-1)+r-1\bigr]a_0=(2r+1)(r-1)a_0=0,
\]
这正是指标方程（与 §5.2 一致）。$n\geq1$ 时：
\[
\bigl[2(n+r)(n+r-1)+(n+r)-1\bigr]a_n
+\bigl[2(n+r-1)+3\bigr]a_{n-1}=0 . \tag{5.1}
\]
记 $s=n+r$。第一个方括号是
\[
2s(s-1)+s-1=2s^2-s-1=(2s+1)(s-1),
\]
第二个方括号是 $2(s-1)+3=2s+1$。于是 (5.1) 变成
\[
(2s+1)\bigl[(s-1)a_n+a_{n-1}\bigr]=0 . \tag{5.2}
\]
因为本题的两个根都使 $2s+1=2(n+r)+1\neq0$（$r=1$ 时是 $2n+3$；$r=-\frac12$ 时是 $2n$），可以约掉：
\[
\boxed{\,a_n=-\frac{a_{n-1}}{n+r-1}\,},\qquad n=1,2,3,\dots \tag{5.3}
\]

\warnbox{
这里的因子 $(2s+1)$ 是题目「特意设计」的：它让递推从二次式退化成一次式，两个根都能一路算到底。如果分母是 $(2s+1)(s-1)$ 而 $(2s+1)$ 在某个 $n$ 处为零，那就要另作讨论。本题不需要——约掉之前先确认它不为零，这一步不能省。
}

\subhdr{5.4 第 3 步：$r_1=1$ 的一支}

把 $r=1$ 代入 (5.3)：
\[
a_n=-\frac{a_{n-1}}{n}.
\]
于是
\[
a_1=-\frac{a_0}{1},\quad
a_2=-\frac{a_1}{2}=\frac{a_0}{2},\quad
a_3=-\frac{a_2}{3}=-\frac{a_0}{6},\quad
a_4=\frac{a_0}{24},\quad
a_5=-\frac{a_0}{120},
\]
通项
\[
a_n=\frac{(-1)^n}{n!}a_0,\qquad n=0,1,2,\dots
\]
取 $a_0=1$：
\[
y_1=x\Bigl(1-x+\frac{x^2}{2}-\frac{x^3}{6}+\frac{x^4}{24}-\cdots\Bigr)
=x-x^2+\frac{x^3}{2}-\frac{x^4}{6}+\frac{x^5}{24}-\cdots \tag{5.4}
\]
括号里正是 $\ee^{-x}$ 的麦克劳林级数，所以这一支有\textbf{闭式}
\[
\boxed{\,y_1=x\ee^{-x}\,}
\]
这不是猜的，是级数求和的结果。闭式给了一个免费的验算：把 $y=x\ee^{-x}$ 直接代回原方程，应该恒等于零。算一遍：
\[
y'=(1-x)\ee^{-x},\qquad y''=(x-2)\ee^{-x},
\]
\begin{align*}
2x^2y''+(2x^2+x)y'+(3x-1)y
&=\ee^{-x}\bigl[2x^2(x-2)+x(2x+1)(1-x)+x(3x-1)\bigr]\\
&=\ee^{-x}\bigl[2x^3-4x^2+x+x^2-2x^3+3x^2-x\bigr]=0 .
\end{align*}
所以 $y_1=x\ee^{-x}$ 确实是解，前五个非零项就是 (5.4) 里的五项。

\subhdr{5.5 第 4 步：$r_2=-\frac12$ 的一支}

把 $r=-\frac12$ 代入 (5.3)：
\[
a_n=-\frac{a_{n-1}}{n-\frac32}=-\frac{2a_{n-1}}{2n-3}.
\]
逐项算（取 $a_0=1$）：
\[
a_1=-\frac{2}{-1}=2,\qquad
a_2=-\frac{2\cdot2}{1}=-4,\qquad
a_3=-\frac{2\cdot(-4)}{3}=\frac83,
\]
\[
a_4=-\frac{2\cdot\frac83}{5}=-\frac{16}{15},\qquad
a_5=-\frac{2\cdot(-\frac{16}{15})}{7}=\frac{32}{105}.
\]
于是
\[
y_2=x^{-1/2}\Bigl(1+2x-4x^2+\frac83x^3-\frac{16}{15}x^4+\frac{32}{105}x^5-\cdots\Bigr). \tag{5.5}
\]
括号内的前五项是
\[
1,\quad 2x,\quad -4x^2,\quad \frac83x^3,\quad -\frac{16}{15}x^4 ;
\]
若把外面的 $x^{-1/2}$ 也算作一项，则前五项是
\[
x^{-1/2},\quad 2x^{1/2},\quad -4x^{3/2},\quad \frac83x^{5/2},\quad -\frac{16}{15}x^{7/2}.
\]
两种数法在作业里都常见，写答案时把括号内、外说清楚即可。

通项可以写成
\[
a_n=\frac{(-1)^{n+1}2^{2n-1}(n-1)!}{(2n-2)!},\qquad n\geq1,
\]
验算：$n=1$ 给 $\frac{1\cdot2\cdot0!}{0!}=2$；$n=2$ 给 $\frac{-8\cdot1!}{2!}=-4$；$n=3$ 给 $\frac{32\cdot2!}{4!}=\frac83$，与上面一致。

\warnbox{
\textbf{第二个坑：}$r_2=-\frac12$ 时 $n+r-1=n-\frac32$，在 $n=1$ 处是 $-\frac12$，符号为负，所以 $a_1=+2$ 而不是 $-2$。很多人在这里把 $-\frac{2}{-1}$ 写成 $-2$，后面整支符号全反。建议先把 $a_1,a_2$ 两项算出来，代回原方程看低阶项是否对得上，再往下算。
}

\subhdr{5.6 第 5 步：最小收敛区间}

按 §4.6，$\rho$ 是展开点到最近奇点的距离。本题
\[
xp(x)=x+\frac12,\qquad x^2q(x)=\frac32x-\frac12
\]
都是多项式，收敛半径都是 $\infty$；方程也只有 $x=0$ 一个奇点，没有「其他奇点」可量。所以
\[
\rho=\infty .
\]
定理保证的区间是 $0<x-0<\infty$，即
\[
\boxed{\,0<x<\infty\,}
\]
在这个区间上，两个级数都收敛。

还可以用比值法独立核对一遍。$y_1$ 的系数是 $1/n!$，$\lim|a_n/a_{n+1}|=\lim(n+1)=\infty$；$y_2$ 的系数满足
\[
\left|\frac{a_n}{a_{n-1}}\right|=\frac{2}{2n-3}\to0
\quad\Longrightarrow\quad
\lim\left|\frac{a_{n-1}}{a_n}\right|=\infty .
\]
两种算法都给出半径 $\infty$，与定理一致。

\warnbox{
\textbf{第三个坑：}「最小收敛区间」不是「收敛半径」。半径是 $\rho=\infty$，但解里有 $x^{-1/2}$，所以解只在 $x>0$ 上有实值意义；答案要写 $0<x<\infty$，不能写 $-\infty<x<\infty$。若题目问的是「幂级数部分 $\sum a_nx^n$ 的收敛区间」，那才是全体实数——把这两个对象分开说。
}

\subhdr{5.7 第 6 步：线性无关性的验算}

两个解都算出来了，还要确认它们线性无关。用 Wronskian：
\[
W=y_1y_2'-y_1'y_2 .
\]
不必硬算（硬算要处理 $x^{-1/2}$ 的级数），用 §6.4 的 Abel 公式更快。对标准形 $y''+py'+qy=0$，
\[
W'=-pW=-\Bigl(1+\frac{1}{2x}\Bigr)W
\quad\Longrightarrow\quad
W=C\ee^{-x}x^{-1/2}.
\]
常数 $C$ 用最低阶项定：$y_1\sim x$，$y_2\sim x^{-1/2}$，于是
\[
W\sim x\cdot\Bigl(-\frac12\Bigr)x^{-3/2}-1\cdot x^{-1/2}
=-\frac32x^{-1/2},
\]
所以 $C=-\frac32$，
\[
W=-\frac32\ee^{-x}x^{-1/2}\neq0\qquad(x>0).
\]
$W$ 在 $x>0$ 上处处不为零，两个解线性无关。

\subhdr{5.8 答案与自检}

\callout{
\textbf{答案。}方程 $2x^2y''+(2x^2+x)y'+(3x-1)y=0$ 在 $x>0$ 上的两个线性无关解为
\[
y_1=x\ee^{-x}=x-x^2+\frac{x^3}{2}-\frac{x^4}{6}+\frac{x^5}{24}-\cdots
\]
\[
y_2=x^{-1/2}\Bigl(1+2x-4x^2+\frac83x^3-\frac{16}{15}x^4+\cdots\Bigr),
\]
通解
\[
y=C_1y_1+C_2y_2,\qquad 0<x<\infty .
\]
最小收敛区间：$0<x<\infty$（$\rho=\infty$，因为 $xp=x+\frac12$ 与 $x^2q=\frac32x-\frac12$ 都是多项式，方程没有其他奇点）。
}

\textbf{自检三问}（答案在 §5.9 末，先自己做）：
\begin{enumerate}
  \item 如果把指标方程误写成 $r(r-1)+p_0r+q_0=0$ 但代入 $p_0=1$、$q_0=0$（即直接用 $Q/P$、$R/P$ 的「常数项」），会得到什么根？错在哪一步？
  \item $y_2$ 的系数 $a_1=2$。如果算成 $a_1=-2$，递推后面会变成什么？用 $n=2$ 的方程检查哪一个对。
  \item 为什么 $y_1=x\ee^{-x}$ 的收敛半径是 $\infty$，而 $y_2$ 的级数部分也是 $\infty$，解却只写在 $x>0$？
\end{enumerate}

\subhdr{5.9 常见错误与自检答案}

\warnbox{
\textbf{本题四个高频错误。}
（1）\textbf{指标方程用错 $p_0,q_0$}：忘记先除 $P$，或忘记乘 $x$、$x^2$。
（2）\textbf{$r_2$ 支的符号}：$n-\frac32$ 在 $n=1$ 为负，$a_1=+2$。
（3）\textbf{把「前五项」数错}：$y_1$ 的 $x^4$ 项系数非零（$-\frac16$），不能跳过；$y_2$ 要把括号内、外说清楚。
（4）\textbf{收敛区间写成 $|x|<\infty$}：$x^{-1/2}$ 只在 $x>0$ 上有实值意义。
}

\textbf{自检答案。}
（1）若取 $p_0=1,q_0=0$，指标方程为 $r(r-1)+r=r^2=0$，得 $r=0$（重根），与正确答案 $1,-\frac12$ 完全不同；错在第一步没有把 $Q/P$、$R/P$ 乘 $x$、$x^2$ 再取极限。
（2）若 $a_1=-2$，由 $a_2=-\frac{2a_1}{1}=4$，与正确值 $-4$ 相反；把 $y_2=x^{-1/2}(1-2x+4x^2-\cdots)$ 代回原方程，$x^{1/2}$ 的系数不为零，所以 $a_1=+2$ 才是对的。
（3）$y_1$ 是整函数 $x\ee^{-x}$，$y_2$ 的幂级数部分也是整函数（比值法给 $\rho=\infty$）；但 $x^{-1/2}$ 在 $x<0$ 不是实数、在 $x=0$ 无定义，所以「解」的实值定义域只能是 $x>0$。收敛半径讲的是级数，定义域讲的是解，两者不是一回事。

%=============================================================
\subhdr{5.10 作业书写模板（可直接照这个顺序写）}

作业要交的是「过程 + 答案」，不是只有最后一行。按下面六块写，老师能一眼看到关键步骤，你自己也不容易漏：

\callout{
\textbf{（1） 定性与标准化。}写出 $P,Q,R$，说明 $x=0$ 是奇点；写出 $p=Q/P$、$q=R/P$，再写 $xp=x+\frac12$、$x^2q=\frac32x-\frac12$，说明二者解析，所以是正则奇点；读出 $p_0=\frac12$、$q_0=-\frac12$。

\textbf{（2） 指标方程。}写 $r(r-1)+\frac12r-\frac12=0$，即 $(2r+1)(r-1)=0$，得 $r_1=1$、$r_2=-\frac12$；写一句「差 $3/2$ 非整数，故两个解均为 Frobenius 级数」。

\textbf{（3） 递推。}设 $y=\sum a_nx^{n+r}$，代入后写 $(2s+1)[(s-1)a_n+a_{n-1}]=0$（$s=n+r$），约去 $2s+1$，得 $a_n=-\frac{a_{n-1}}{n+r-1}$。

\textbf{（4） 第一支。}$r=1$：$a_n=(-1)^n a_0/n!$，写前五项，得 $y_1=x\ee^{-x}=x-x^2+\frac{x^3}{2}-\frac{x^4}{6}+\frac{x^5}{24}-\cdots$；可加一句「代回验证为解」。

\textbf{（5） 第二支。}$r=-\frac12$：$a_n=-\frac{2a_{n-1}}{2n-3}$，写 $a_1=2,a_2=-4,a_3=\frac83,a_4=-\frac{16}{15},a_5=\frac{32}{105}$，得
\[
y_2=x^{-1/2}\Bigl(1+2x-4x^2+\frac83x^3-\frac{16}{15}x^4+\cdots\Bigr).
\]

\textbf{（6） 收敛区间与结论。}写「$xp$、$x^2q$ 为多项式，$\rho=\infty$；又 $x^{-1/2}$ 要求 $x>0$，故最小收敛区间为 $0<x<\infty$」，最后写通解 $y=C_1y_1+C_2y_2$。
}

\warnbox{
\textbf{交卷前 30 秒检查：}（1）指标方程是否与 $n=0$ 的系数方程一致；（2）$a_1$ 的符号；（3）「前五项」是否都是非零项；（4）区间是否写了 $x>0$；（5）有没有一句「两个解线性无关」——可以写 $W\neq0$，也可以写「$y_1/y_2$ 不是常数」。
}

\hdr{§6}{补齐专章：线性微分方程组（回顾）}
\callout{这一节从方程组角度补工具箱：最小线性代数、Wronskian 与 Abel 公式、常系数齐次组的特征值法、矩阵指数、非齐次组的参数变易。\textbf{已读完第 1--4 章、只需幂级数方法的读者可整节跳过}；本章正文有两处直接用到它（§4.7 的降阶公式、§5.7 的 Wronskian 验算），届时回来查 §6.4 与 §6.9 即可。}



\callout{
\textbf{这一节为什么在这里。}本次作业的题干要求「construct two linearly independent solutions」，答案写成「通解 $=C_1y_1+C_2y_2$」，验算用 Wronskian——这三样东西的语言全部来自方程组章（课件 L13--L17）。课件 L18--20 没有再讲一遍，默认你会。你既然说方程组还没学好，这一节就从零把用得到的部分补齐；标「\textbf{本章直接用到}」的小节必须掌握，其余是背景，读懂即可。
}

\subhdr{6.1 从二阶方程到一阶系统（本章直接用到）}

设
\[
y''+p(x)y'+q(x)y=0 .
\]
令 $y_1=y$、$y_2=y'$，则
\[
\begin{cases}
y_1'=y_2,\\
y_2'=-q(x)y_1-p(x)y_2 .
\end{cases}
\]
写成矩阵形式：
\[
\frac{\dd}{\dd x}\begin{pmatrix}y_1\\y_2\end{pmatrix}
=\begin{pmatrix}0&1\\-q(x)&-p(x)\end{pmatrix}
\begin{pmatrix}y_1\\y_2\end{pmatrix},
\qquad\text{即}\quad Y'=A(x)Y . \tag{6.1}
\]
反过来，若一个 $2\times2$ 系统的系数矩阵第一行是 $(0,1)$，那么 $y=y_1$ 满足 $y''+py'+qy=0$：由第一式 $y_2=y_1'$，代入第二式即得。所以\textbf{二阶线性方程与「伴随形式」的一阶系统是同一件事}，只是把「一个二阶方程」换成了「两个一阶方程」。

\aside{不是每个 $2\times2$ 系统都能消成一个二阶方程；能消的恰好是这种「第一行 $(0,1)$」的伴随形式（companion form）。但常系数系统总可以化到若尔当形，逐个分量求解，所以「系统」比「单个二阶方程」更一般。本讲义只需要伴随形式这一支。}

对本题，$p=1+\dfrac{1}{2x}$、$q=\dfrac{3}{2x}-\dfrac{1}{2x^2}$，所以
\[
A(x)=\begin{pmatrix}0&1\\[2pt] \dfrac{1}{2x^2}-\dfrac{3}{2x}&-1-\dfrac{1}{2x}\end{pmatrix}.
\]
注意 $x\to0$ 时 $A(x)$ 的元素无界：\textbf{$x=0$ 是系统的奇点}。这件事在 §6.8 会再用一次。

\subhdr{6.2 最小线性代数工具箱}

这一小节只列后面用到的公式，不证明。

\callout{
\textbf{矩阵与向量。}$A=\begin{pmatrix}a&b\\c&d\end{pmatrix}$，$Y=\begin{pmatrix}u\\v\end{pmatrix}$，
\[
AY=\begin{pmatrix}au+bv\\cu+dv\end{pmatrix},\qquad
\det A=ad-bc,\qquad
A^{-1}=\frac{1}{ad-bc}\begin{pmatrix}d&-b\\-c&a\end{pmatrix}\quad(\det A\neq0).
\]

\textbf{特征值与特征向量。}若存在数 $\lambda$ 和非零向量 $v$ 使
\[
Av=\lambda v,
\]
则 $\lambda$ 叫\textbf{特征值}，$v$ 叫对应的\textbf{特征向量}。把 $Av=\lambda v$ 写成 $(A-\lambda I)v=0$：它有非零解当且仅当
\[
\det(A-\lambda I)=0 .
\]
这是 $\lambda$ 的二次方程（$2\times2$ 时），叫\textbf{特征方程}；它的两个根就是两个特征值（可能重复、可能共轭复）。

\textbf{矩阵指数。}对任意方阵 $A$ 和实数 $x$，定义
\[
\ee^{Ax}=I+Ax+\frac{A^2x^2}{2!}+\frac{A^3x^3}{3!}+\cdots
\]
这个级数对每个元素都绝对收敛，所以 $\ee^{Ax}$ 总是有定义的。
}

\subhdr{6.3 一阶线性系统的基本定理}

\callout{
\textbf{定理 6.1（存在唯一性）。}设 $A(x)$ 与 $F(x)$ 在区间 $[a,b]$ 上连续。则对任意 $x_0\in[a,b]$ 和任意常向量 $Y_0$，初值问题
\[
Y'=A(x)Y+F(x),\qquad Y(x_0)=Y_0
\]
在 $[a,b]$ 上有且只有一个解。
}

这个定理是整章的「地基」：它保证了解的存在、唯一，也保证了解空间是 $n$ 维的。下面只对 $n=2$ 说，$n$ 维完全一样。

\callout{
\textbf{齐次系统的解空间是 $2$ 维的。}设 $Y'=A(x)Y$。
\begin{enumerate}
  \item 若 $Y_1,Y_2$ 是解，则任意线性组合 $c_1Y_1+c_2Y_2$ 也是解（齐次、线性）。
  \item 取定 $x_0$，把每个解映到它的初值 $Y(x_0)\in\R^2$。由存在唯一性，这个映射是一一对应：不同的解给出不同的初值，每个初值也恰好来自一个解。
\end{enumerate}
所以齐次系统的解集与 $\R^2$ 一样是 $2$ 维的：\textbf{只要能找到两个线性无关的解，所有解就都是它们的线性组合}。
}

\callout{
\textbf{基本矩阵。}若 $Y_1,Y_2$ 是两个线性无关的解，把它们按列排成
\[
\Phi(x)=\bigl(Y_1(x)\ \ Y_2(x)\bigr)=\begin{pmatrix}y_{11}&y_{12}\\y_{21}&y_{22}\end{pmatrix},
\]
则 $\Phi$ 叫\textbf{基本矩阵}（fundamental matrix）。它满足
\[
\Phi'=A\Phi,\qquad \det\Phi(x)\neq0,\qquad Y=\Phi C\ \ (C\text{ 为常向量}).
\]
若还要求 $\Phi(x_0)=I$，则称 $\Phi$ 在 $x_0$ \textbf{归一化}；任何基本矩阵右乘 $\Phi(x_0)^{-1}$ 就归一化了。
}

\subhdr{6.4 Wronskian 与 Abel--Liouville 公式（本章直接用到）}

\callout{
\textbf{定义。}$W(x)=\det\Phi(x)$ 叫这组解的 \textbf{Wronskian}。对两个解 $Y_1=(y_{11},y_{21})^{\mathsf T}$、$Y_2=(y_{12},y_{22})^{\mathsf T}$，
\[
W=\det\begin{pmatrix}y_{11}&y_{12}\\y_{21}&y_{22}\end{pmatrix}
=y_{11}y_{22}-y_{12}y_{21}.
\]
}

\callout{
\textbf{定理 6.2（Abel--Liouville）。}若 $\Phi'=A\Phi$，则
\[
W'=(\tr A)\,W,\qquad\text{从而}\qquad
W(x)=W(x_0)\exp\Bigl(\int_{x_0}^{x}\tr A(t)\dd t\Bigr). \tag{6.2}
\]
}

\textbf{$2\times2$ 的直接证明。}记 $Y_1=(u,v)^{\mathsf T}$、$Y_2=(s,t)^{\mathsf T}$，$A=\begin{pmatrix}a&b\\c&d\end{pmatrix}$，则 $W=ut-vs$。由 $Y_1'=AY_1$、$Y_2'=AY_2$：
\[
u'=au+bv,\quad v'=cu+dv,\quad s'=as+bt,\quad t'=cs+dt .
\]
于是
\[
W'=u't+ut'-v's-vs'=(au+bv)t+u(cs+dt)-v(as+bt)-(cu+dv)s .
\]
把四项展开、按 $a,b,c,d$ 归并：
\[
W'=a(ut-vs)+d(ut-vs)=(a+d)W=(\tr A)W .
\]
证毕。

\warnbox{
\textbf{三个推论，做题时天天用。}
（1）$W$ 要么恒为零，要么在区间上处处不为零——因为 $W(x)=W(x_0)\ee^{\int\tr A}$，指数因子永远不为零。
（2）所以判断 $n$ 个解是否线性无关，只需在\textbf{一个点}算 $W$：$W(x_0)\neq0$ 就无关，$W\equiv0$ 就相关。
（3）不要反过来用：$W(x_0)=0$ 只说明这一点处相关，不能直接说「函数线性相关」——但对\textbf{同一条线性方程的解}，由 (6.2) 可知一点为零就处处为零，所以对解组来说这个反向也成立。
}

\textbf{与标量方程的关系（这就是 §5.7 用的公式）。}对 $y''+py'+qy=0$，伴随矩阵 $A=\begin{pmatrix}0&1\\-q&-p\end{pmatrix}$，所以
\[
\tr A=-p,\qquad W'=-pW,\qquad W=W_0\exp\Bigl(-\int p\dd x\Bigr). \tag{6.3}
\]
这正是课件 L8 的 Abel 公式（也叫 Liouville 公式）。§5.7 里直接写 $W=C\ee^{-x}x^{-1/2}$，用的就是 (6.3)，其中 $p=1+\frac{1}{2x}$，$-\int p\dd x=-x-\frac12\ln x$。

\begin{center}
\small
\begin{tabular}{@{}p{6.6cm}p{8.4cm}@{}}
\hline
\thead{标量二阶方程} & \thead{一阶系统} \\
\hline
$y''+py'+qy=0$ & $Y'=A(x)Y$，$A=\begin{pmatrix}0&1\\-q&-p\end{pmatrix}$ \\
两个线性无关解 $y_1,y_2$ & 基本矩阵 $\Phi=(Y_1\ Y_2)$ \\
$W=y_1y_2'-y_1'y_2$ & $W=\det\Phi$ \\
$W'=-pW$ & $W'=(\tr A)W$ \\
$y=C_1y_1+C_2y_2$ & $Y=\Phi C$ \\
\hline
\end{tabular}
\end{center}

这张表把同一件事的两种写法并排：左边一列是你在 L7--L9 学的，右边一列是系统语言。逐行读法：第一行给出矩阵；第二行把两个解按列排；第三行是行列式；第四行把 $W$ 的变化率写成系数矩阵的迹（伴随矩阵的迹正好是 $-p$）；第五行是通解。整张表只有一句话：\textbf{标量方程的 Wronskian 就是伴随系统基本矩阵的行列式}。

\subhdr{6.5 常系数齐次系统：特征值法}

设 $A$ 是常矩阵，解 $Y'=AY$。仿照标量方程试 $y=\ee^{\lambda x}$，这里试
\[
Y=v\ee^{\lambda x}\quad(v\neq0).
\]
代入得 $\lambda v\ee^{\lambda x}=Av\ee^{\lambda x}$，约去 $\ee^{\lambda x}$：
\[
Av=\lambda v .
\]
所以\textbf{$Y=v\ee^{\lambda x}$ 是解 $\iff$ $\lambda$ 是 $A$ 的特征值、$v$ 是对应特征向量}。三种情形：

\callout{
\textbf{情形 1：两个相异实特征值 $\lambda_1\neq\lambda_2$。}对应特征向量 $v_1,v_2$，两个解
\[
Y_1=v_1\ee^{\lambda_1x},\qquad Y_2=v_2\ee^{\lambda_2x}
\]
线性无关，通解 $Y=C_1Y_1+C_2Y_2$。

\textbf{情形 2：一对共轭复特征值 $\lambda=\alpha\pm\ii\beta$（$\beta\neq0$）。}取 $v=u+\ii w$ 对应 $\alpha+\ii\beta$，则
\[
Y=(u+\ii w)\ee^{(\alpha+\ii\beta)x}
=\ee^{\alpha x}\bigl[(u\cos\beta x-w\sin\beta x)+\ii(u\sin\beta x+w\cos\beta x)\bigr].
\]
实部与虚部分别是实解：
\[
Y_1=\ee^{\alpha x}(u\cos\beta x-w\sin\beta x),\qquad
Y_2=\ee^{\alpha x}(u\sin\beta x+w\cos\beta x).
\]

\textbf{情形 3：重特征值 $\lambda$。}若有两个线性无关的特征向量 $v_1,v_2$，照情形 1 写。若只有一个 $v$，需要\textbf{广义特征向量} $w$，满足
\[
(A-\lambda I)w=v,
\]
第二个解取 $Y_2=(vx+w)\ee^{\lambda x}$。
}

\callout{
\textbf{最小例子（情形 3）。}$A=\begin{pmatrix}1&1\\0&1\end{pmatrix}$，特征值 $\lambda=1$（二重），$A-I=\begin{pmatrix}0&1\\0&0\end{pmatrix}$，只有一个特征向量 $v=(1,0)^{\mathsf T}$。解 $(A-I)w=v$：
\[
\begin{pmatrix}0&1\\0&0\end{pmatrix}\begin{pmatrix}w_1\\w_2\end{pmatrix}=\begin{pmatrix}1\\0\end{pmatrix}
\quad\Longrightarrow\quad w_2=1,\ w_1\ \text{任意，取}\ w=(0,1)^{\mathsf T}.
\]
于是
\[
Y=C_1\begin{pmatrix}1\\0\end{pmatrix}\ee^x+C_2\left[\begin{pmatrix}1\\0\end{pmatrix}x+\begin{pmatrix}0\\1\end{pmatrix}\right]\ee^x
=\ee^x\begin{pmatrix}C_1+C_2x\\C_2\end{pmatrix}.
\]
第二分量是常数乘 $\ee^x$，第一分量是 $(C_1+C_2x)\ee^x$——这正是标量重根时的「$x\ee^{\lambda x}$」。
}

\subhdr{6.6 矩阵指数：把三种情形统一}

把 (6.1) 的解写成 $\Phi(x)=\ee^{Ax}$，这不是形式记号，而是真的级数：
\[
\ee^{Ax}=I+Ax+\frac{A^2x^2}{2!}+\cdots
\]
逐项求导得 $\dfrac{\dd}{\dd x}\ee^{Ax}=A\ee^{Ax}=\ee^{Ax}A$，且 $\ee^{A\cdot0}=I$。所以 $\ee^{Ax}$ 就是归一化基本矩阵，任何解都可写成 $Y=\ee^{Ax}Y_0$。

计算 $\ee^{Ax}$ 的三条路：
\begin{enumerate}
  \item \textbf{可对角化}：$A=P\,\mathrm{diag}(\lambda_1,\lambda_2)P^{-1}$，则
        \[
        \ee^{Ax}=P\,\mathrm{diag}(\ee^{\lambda_1x},\ee^{\lambda_2x})P^{-1}.
        \]
  \item \textbf{若尔当块}：$J=\begin{pmatrix}\lambda&1\\0&\lambda\end{pmatrix}=\lambda I+N$，$N^2=0$，所以
        \[
        \ee^{Jx}=\ee^{\lambda x}(I+Nx)=\ee^{\lambda x}\begin{pmatrix}1&x\\0&1\end{pmatrix}.
        \]
  \item \textbf{复特征值}：若 $A$ 的特征值为 $\alpha\pm\ii\beta$，则 $(A-\alpha I)^2=-\beta^2I$，于是
        \[
        \ee^{Ax}=\ee^{\alpha x}\left[\cos(\beta x)\,I+\frac{\sin(\beta x)}{\beta}(A-\alpha I)\right]. \tag{6.4}
        \]
\end{enumerate}

\warnbox{
\textbf{行列式与迹。}由 (6.2) 取 $x_0=0$、$\Phi(0)=I$，得
\[
\det\ee^{Ax}=\ee^{(\tr A)x}.
\]
这条公式有两个用处：一是快速检查 $\ee^{Ax}$ 算对没有（算完取行列式，应等于 $\ee^{(\tr A)x}$）；二是它说明 $\ee^{Ax}$ 永远可逆，所以「基本矩阵」永远不会退化。
}

\callout{
\textbf{例（旋转）。}$A=\begin{pmatrix}0&1\\-1&0\end{pmatrix}$，特征值 $\pm\ii$（$\alpha=0,\beta=1$）。由 (6.4)：
\[
\ee^{Ax}=\cos x\,I+\sin x\,A
=\begin{pmatrix}\cos x&\sin x\\-\sin x&\cos x\end{pmatrix}.
\]
验证：$\det=\cos^2x+\sin^2x=1=\ee^{0\cdot x}$；求导得 $\begin{pmatrix}-\sin x&\cos x\\-\cos x&-\sin x\end{pmatrix}=A\ee^{Ax}$。这个矩阵正是「把向量旋转 $-x$ 角」的矩阵。
}

\subhdr{6.7 非齐次系统：参数变易}

设
\[
Y'=A(x)Y+F(x),
\]
$\Phi$ 是齐次系统 $Y'=AY$ 的基本矩阵。令 $Y=\Phi U$，代入：
\[
\Phi'U+\Phi U'=A\Phi U+F
\quad\Longrightarrow\quad
\Phi U'=F
\quad\Longrightarrow\quad
U'=\Phi^{-1}F .
\]
积分得
\[
Y=\Phi(x)C+\Phi(x)\int_{x_0}^{x}\Phi(t)^{-1}F(t)\dd t . \tag{6.5}
\]
第一项是齐次通解，第二项是一个特解。对标量二阶方程，取 $\Phi=\begin{pmatrix}y_1&y_2\\y_1'&y_2'\end{pmatrix}$，(6.5) 就化为你学过的参数变易公式（L8）：两个积分里的分母正是 $W$。

\subhdr{6.8 把作业题写成系统，并核对 W（本章直接用到）}

本题的 $p,q$ 已在 §6.1 给出，伴随矩阵为
\[
A(x)=\begin{pmatrix}0&1\\[2pt] \dfrac{1}{2x^2}-\dfrac{3}{2x}&-1-\dfrac{1}{2x}\end{pmatrix},
\qquad \tr A=-1-\frac{1}{2x}.
\]
由 (6.2)：
\[
W=W_0\exp\Bigl(-\int\Bigl(1+\frac{1}{2x}\Bigr)\dd x\Bigr)
=W_0\ee^{-x}x^{-1/2}.
\]
这与 §5.7 用 Abel 公式得到的形式一致。再用两个解的领头项定 $W_0$：$y_1\sim x$、$y_2\sim x^{-1/2}$，
\[
W\sim x\cdot\Bigl(-\frac12\Bigr)x^{-3/2}-1\cdot x^{-1/2}=-\frac32x^{-1/2},
\]
所以 $W_0=-\frac32$。$W\neq0$ 说明两个解线性无关，也说明 $\Phi(x_0)$ 可逆、初值问题对任意 $(y_0,y_0')$ 都能定出 $C_1,C_2$。

\textbf{系统视角还解释了「为什么会出现 $x^{-1/2}$」。}$A(x)$ 的元素在 $x=0$ 无界，存在唯一性定理要求 $A$ 在包含 $x_0$ 的区间上连续，所以这个定理\textbf{不能}跨过 $0$ 使用。换句话说，解在 $0$ 附近本来就可能不解析——这正是 §4 要设 $y=x^r\sum a_nx^n$、允许非整数 $r$ 的原因。若 $A$ 在 $0$ 连续（常点），定理 2.1 就会保证解是普通幂级数。

\subhdr{6.9 为什么「两个线性无关解」就够了（本章直接用到）}

给定初值 $(y(x_0),y'(x_0))=(y_0,y_0')$，要解出 $C_1,C_2$：
\[
C_1Y_1(x_0)+C_2Y_2(x_0)=Y_0
\quad\Longleftrightarrow\quad
\Phi(x_0)\begin{pmatrix}C_1\\C_2\end{pmatrix}=Y_0 .
\]
因为 $\det\Phi(x_0)=W(x_0)\neq0$，矩阵 $\Phi(x_0)$ 可逆，所以
\[
\begin{pmatrix}C_1\\C_2\end{pmatrix}=\Phi(x_0)^{-1}Y_0
\]
是唯一的。这就是「通解 $=C_1y_1+C_2y_2$」的严格含义：\textbf{两个线性无关解张成整个 $2$ 维解空间，任何初值都能被它们表示，且表示唯一}。

\subhdr{6.10 三个最小例子与自检}

\callout{
\textbf{例 S1。}$y''-3y'+2y=0$。伴随矩阵 $A=\begin{pmatrix}0&1\\-2&3\end{pmatrix}$，$\tr A=3$，$\det(A-\lambda I)=\lambda^2-3\lambda+2=(\lambda-1)(\lambda-2)$，特征值 $1,2$，特征向量 $(1,1)^{\mathsf T}$、$(1,2)^{\mathsf T}$。基本矩阵
\[
\Phi=\begin{pmatrix}\ee^x&\ee^{2x}\\\ee^x&2\ee^{2x}\end{pmatrix},\qquad
\det\Phi=\ee^{3x}=\ee^{(\tr A)x}.
\]
第一行给出 $y=C_1\ee^x+C_2\ee^{2x}$，与标量解一致。

\textbf{例 S2。}$y''+y=0$。$A=\begin{pmatrix}0&1\\-1&0\end{pmatrix}$，$\ee^{Ax}=\begin{pmatrix}\cos x&\sin x\\-\sin x&\cos x\end{pmatrix}$，第一行给 $y=C_1\cos x+C_2\sin x$。

\textbf{例 S3。}$y''-2y'+y=0$。$A=\begin{pmatrix}0&1\\-1&2\end{pmatrix}$，$\lambda=1$（二重），$v=(1,1)^{\mathsf T}$，$w=(0,1)^{\mathsf T}$。基本矩阵
\[
\Phi=\ee^x\begin{pmatrix}1&x\\1&x+1\end{pmatrix},\qquad \det\Phi=\ee^{2x}.
\]
第一行给 $y=(C_1+C_2x)\ee^x$，正是标量重根的解。
}

\textbf{自检三问}（答案在本节末）：
\begin{enumerate}
  \item 对 $y''+4y=0$，写出伴随矩阵 $A$，用 (6.4) 求 $\ee^{Ax}$，并验证 $\det\ee^{Ax}=1$。
  \item 对 $A=\begin{pmatrix}2&1\\0&2\end{pmatrix}$，求 $\ee^{Ax}$，并说明为什么解里会出现 $x\ee^{2x}$。
  \item 用 (6.2) 验证：若 $A$ 是常矩阵且 $\tr A=0$，则任何基本矩阵的行列式都不随 $x$ 变。
\end{enumerate}

\textbf{自检答案。}
（1）$A=\begin{pmatrix}0&1\\-4&0\end{pmatrix}$，$\lambda=\pm2\ii$，$\alpha=0,\beta=2$，由 (6.4)
\[
\ee^{Ax}=\cos(2x)I+\frac{\sin(2x)}{2}A
=\begin{pmatrix}\cos2x&\frac12\sin2x\\-2\sin2x&\cos2x\end{pmatrix},
\]
$\det=\cos^22x+\sin^22x=1$。
（2）$A=2I+N$，$N=\begin{pmatrix}0&1\\0&0\end{pmatrix}$，$N^2=0$，所以 $\ee^{Ax}=\ee^{2x}(I+Nx)=\ee^{2x}\begin{pmatrix}1&x\\0&1\end{pmatrix}$；右上角的 $x$ 就是「广义特征向量」带来的 $x\ee^{2x}$。
（3）$W'=(\tr A)W=0$，所以 $W$ 为常数；取 $x_0$ 得 $W(x)=W(x_0)$。

\subhdr{6.11 两个补充例子}

\textbf{例 A（相异实特征值）。}解
\[
Y'=\begin{pmatrix}1&1\\4&1\end{pmatrix}Y .
\]
特征方程：
\[
\det(A-\lambda I)=\begin{vmatrix}1-\lambda&1\\4&1-\lambda\end{vmatrix}
=(1-\lambda)^2-4=\lambda^2-2\lambda-3=(\lambda-3)(\lambda+1)=0,
\]
所以 $\lambda_1=3$、$\lambda_2=-1$。

$\lambda_1=3$：$A-3I=\begin{pmatrix}-2&1\\4&-2\end{pmatrix}$，方程 $-2v_1+v_2=0$，取 $v_1=(1,2)^{\mathsf T}$。

$\lambda_2=-1$：$A+I=\begin{pmatrix}2&1\\4&2\end{pmatrix}$，方程 $2v_1+v_2=0$，取 $v_2=(1,-2)^{\mathsf T}$。

于是
\[
Y=C_1\begin{pmatrix}1\\2\end{pmatrix}\ee^{3x}+C_2\begin{pmatrix}1\\-2\end{pmatrix}\ee^{-x},
\]
即 $y_1=C_1\ee^{3x}+C_2\ee^{-x}$、$y_2=2C_1\ee^{3x}-2C_2\ee^{-x}$。检验 $\tr A=2$：基本矩阵的行列式应为 $W_0\ee^{2x}$，而
\[
W=\begin{vmatrix}\ee^{3x}&\ee^{-x}\\2\ee^{3x}&-2\ee^{-x}\end{vmatrix}
=-2\ee^{2x}-2\ee^{2x}=-4\ee^{2x},
\]
与 $\tr A=2$ 一致。

\textbf{例 B（三阶方程化成 $3\times3$ 系统）。}把
\[
y'''-y=0
\]
化成系统。令 $y_1=y$、$y_2=y'$、$y_3=y''$，则
\[
y_1'=y_2,\qquad y_2'=y_3,\qquad y_3'=y'''=y=y_1,
\]
即
\[
Y'=\begin{pmatrix}0&1&0\\0&0&1\\1&0&0\end{pmatrix}Y .
\]
其特征多项式
\[
\det(A-\lambda I)=\begin{vmatrix}-\lambda&1&0\\0&-\lambda&1\\1&0&-\lambda\end{vmatrix}
=-\lambda^3+1=-(\lambda^3-1),
\]
与标量方程的特征方程 $\lambda^3-1=0$ 完全相同（差一个负号，不影响根）。三个特征值是 $1$、$\ee^{2\pi\ii/3}$、$\ee^{-2\pi\ii/3}$，对应
\[
y=C_1\ee^x+\ee^{-x/2}\Bigl[C_2\cos\frac{\sqrt3x}{2}+C_3\sin\frac{\sqrt3x}{2}\Bigr].
\]
这个例子说明：\textbf{$n$ 阶常系数线性方程与 $n\times n$ 伴随系统有同一个特征多项式}，所以「系统章的特征值法」不是新方法，只是把高阶方程的语言换成了矩阵。

\warnbox{
\textbf{本节没讲的三块，先划清边界。}（1）$3\times3$ 以上的\textbf{广义特征向量链}：$g<m$ 时可能要算 $(A-\lambda I)^j w$ 一整条链，本节只讲了 $2\times2$ 的一步链；（2）\textbf{相图与稳定性}（平衡点、鞍点、焦点）：那是系统章后半的内容，与本章的级数解无关；（3）\textbf{用 Laplace 变换解系统}：另一条路线，本章不用。这三块若考试要考，再单独补；就「读懂本讲义 §5、§6.8 的 Wronskian 与通解」而言，上面的内容已经够用。
}

%=============================================================
\hdr{§7}{Bessel 方程：本章方法的最后一次综合}

课件 p35--p41 用 Bessel 方程收尾，不是因为它新，而是因为它把本章所有零件都用了一遍：正则奇点、指标方程、递推、整数差、$\ln x$、收敛半径。

\subhdr{7.1 方程与正则性}

\callout{
\textbf{Bessel 方程（课件 p36）。}
\[
x^2y''+xy'+(x^2-\nu^2)y=0,\qquad \nu\geq0 .
\]
除以 $x^2$：
\[
p(x)=\frac1x,\qquad q(x)=\frac{x^2-\nu^2}{x^2},
\]
于是
\[
xp(x)=1,\qquad x^2q(x)=x^2-\nu^2,
\]
两者都在 $0$ 解析，所以 $x=0$ 是\textbf{正则奇点}；$p_0=1$、$q_0=-\nu^2$。方程没有其他奇点（复平面上也没有），所以 $\rho=\infty$。
}

\subhdr{7.2 指标方程、递推与 $J_\nu$}

指标方程：
\[
r(r-1)+r-\nu^2=r^2-\nu^2=0
\quad\Longrightarrow\quad
r_1=\nu,\quad r_2=-\nu .
\]
设 $y=\sum a_nx^{n+r}$。代入时注意 $x^2y''+xy'$ 合并后是 $\sum(n+r)^2a_nx^{n+r}$，而 $(x^2-\nu^2)y$ 给出 $-\nu^2\sum a_nx^{n+r}$ 与 $\sum a_nx^{n+r+2}$：
\[
\sum_{n=0}^{\infty}\bigl[(n+r)^2-\nu^2\bigr]a_nx^{n+r}+\sum_{n=2}^{\infty}a_{n-2}x^{n+r}=0 .
\]
令 $x^{n+r}$ 的系数为零（约定 $a_{-1}=a_{-2}=0$）：
\[
\bigl[(n+r)^2-\nu^2\bigr]a_n+a_{n-2}=0,\qquad n=0,1,2,\dots \tag{7.1}
\]
$n=0$ 给指标方程；$n=1$ 给 $\bigl[(r+1)^2-\nu^2\bigr]a_1=0$。

取 $r=\nu$。$n=1$ 时 $(2\nu+1)a_1=0$，只要 $\nu\neq-\frac12$ 就有 $a_1=0$；于是奇数项全为零。对偶数项，由 (7.1)：
\[
a_{2k}=-\frac{a_{2k-2}}{(2k+\nu)^2-\nu^2}
=-\frac{a_{2k-2}}{4k(\nu+k)} .
\]
反复代入：
\[
a_{2k}=\frac{(-1)^ka_0}{4^k\,k!\,(\nu+1)(\nu+2)\cdots(\nu+k)}
=\frac{(-1)^k\Gamma(\nu+1)}{4^k\,k!\,\Gamma(\nu+k+1)}a_0 .
\]
取 $a_0=\dfrac{1}{2^\nu\Gamma(\nu+1)}$，得到\textbf{第一类 Bessel 函数}
\[
\boxed{\,J_\nu(x)=\sum_{k=0}^{\infty}\frac{(-1)^k}{k!\,\Gamma(\nu+k+1)}\Bigl(\frac{x}{2}\Bigr)^{2k+\nu}\,} \tag{7.2}
\]
例如 $\nu=0$ 时
\[
J_0(x)=1-\frac{x^2}{4}+\frac{x^4}{64}-\frac{x^6}{2304}+\cdots
\]
（系数由 $a_{2k}=-\dfrac{a_{2k-2}}{4k^2}$ 逐个算出）。

\subhdr{7.3 收敛性与第二解}

式 (7.2) 的相邻系数比：
\[
\left|\frac{a_{2k}}{a_{2k-2}}\right|=\frac{1}{4k(\nu+k)}\to0,
\]
所以级数对一切 $x$ 收敛，$\rho=\infty$。注意 $x^\nu$ 在 $\nu$ 不是整数时只对 $x>0$ 有实值意义，所以解写在 $x>0$ 上。

第二解分两种情形（课件 p41）：
\begin{itemize}
  \item $\nu\notin\N$：$J_\nu$ 与 $J_{-\nu}$ 线性无关，基本解组 $\{J_\nu,J_{-\nu}\}$。
  \item $\nu=n\in\N$：$J_{-n}=(-1)^nJ_n$，两者线性相关，必须另找第二解。课件给出的第二类 Bessel 函数是
  \[
  Y_n(x)=\lim_{\nu\to n}\frac{J_\nu(x)\cos\nu\pi-J_{-\nu}(x)}{\sin\nu\pi},
  \]
  它是一个含 $\ln x$ 的级数（课件 p41 给了展开式）。从本章的角度看，这正是「两根之差 $2n$ 是正整数」的情形 C：第二解带对数项。
\end{itemize}

\aside{本章不需要背 $Y_n$ 的展开式，只需知道：$\nu$ 为整数时第二解带 $\ln x$，来源是 §4.3 的整数差机制。考试若考 Bessel，通常只要求写出 $J_\nu$ 的级数与指标方程。}

\subhdr{7.4 一个能求和的特例：$\nu=\frac12$}

Bessel 级数一般不写成初等函数，但 $\nu=\frac12$ 是例外，它给了一个和 $y_1=x\ee^{-x}$ 同级的「免费验算」。

把 $\nu=\frac12$ 代入 (7.2)。先定 $a_0$：由 $\Gamma(\frac32)=\frac{\sqrt\pi}{2}$，
\[
a_0=\frac{1}{2^{1/2}\Gamma(\frac32)}=\frac{1}{\sqrt2\cdot\frac{\sqrt\pi}{2}}=\sqrt{\frac{2}{\pi}} .
\]
再用递推 $a_{2k}=-\dfrac{a_{2k-2}}{4k(\nu+k)}$，此时 $\nu=\frac12$，$4k(\nu+k)=2k(2k+1)$，所以
\[
a_{2k}=-\frac{a_{2k-2}}{2k(2k+1)} .
\]
逐项：$a_2=-\frac{a_0}{2\cdot3}=-\frac{a_0}{3!}$，$a_4=-\frac{a_2}{4\cdot5}=\frac{a_0}{5!}$，归纳得
\[
a_{2k}=\frac{(-1)^k}{(2k+1)!}a_0 .
\]
于是
\[
J_{1/2}(x)=x^{1/2}\sum_{k=0}^{\infty}\frac{(-1)^k}{(2k+1)!}a_0x^{2k}
=\sqrt{\frac{2}{\pi}}\,x^{1/2}\sum_{k=0}^{\infty}\frac{(-1)^kx^{2k}}{(2k+1)!}
=\sqrt{\frac{2}{\pi x}}\sum_{k=0}^{\infty}\frac{(-1)^kx^{2k+1}}{(2k+1)!}
=\sqrt{\frac{2}{\pi x}}\sin x .
\]
最后一步用的是 $\sin x$ 的麦克劳林级数。这个例子说明：Frobenius 级数不是「只能数值算」的东西，特殊参数下它能求和；同时也提醒你，$x^{1/2}$ 这样的非整数幂在 $x>0$ 之外没有实值意义——这与 §5.6 的讨论是同一件事。

%=============================================================
\hdr{§8}{公式卡与五步流程}

\subhdr{8.1 一页公式}

\callout{
\textbf{（1）判断点。}$P(x_0)\neq0$：常点。$P(x_0)=0$：奇点；再算
\[
\lim_{x\to x_0}(x-x_0)\frac{Q}{P},\qquad \lim_{x\to x_0}(x-x_0)^2\frac{R}{P},
\]
都有限 $\Rightarrow$ 正则奇点；否则非正则。

\textbf{（2）常点。}
\[
y=\sum_{n=0}^{\infty}a_n(x-x_0)^n,
\]
代入后令同幂次系数为零，解出下标最大的系数的递推。

\textbf{（3）正则奇点。}
\[
y=x^r\sum_{n=0}^{\infty}a_nx^n,\qquad a_0\neq0,
\]
\[
r(r-1)+p_0r+q_0=0,\qquad p_0=\lim_{x\to0}xp,\quad q_0=\lim_{x\to0}x^2q .
\]

\textbf{（4）递推。}
\[
a_n=-\frac{1}{(n+r-r_1)(n+r-r_2)}
\sum_{m=0}^{n-1}\bigl[p_{n-m}(m+r)+q_{n-m}\bigr]a_m .
\]

\textbf{（5）第二解。}
$r_1-r_2\notin\Z$：$y_2=x^{r_2}\sum b_nx^n$；
$r_1=r_2$：$y_2=y_1\ln x+x^{r_2}\sum b_nx^n$；
$r_1-r_2\in\N$：$y_2=ay_1\ln x+x^{r_2}\sum b_nx^n$（$a$ 可能为 $0$）。

\textbf{（6）收敛半径。}$\rho=$ 展开点到最近奇点（实或复）的距离；等价地，$\rho=\min\{xp,\ x^2q\ \text{两个级数的收敛半径}\}$。
}

\subhdr{8.2 做题五步}

\begin{enumerate}
  \item \textbf{化标准形}：除以 $P$，读出 $p=Q/P$、$q=R/P$。
  \item \textbf{判点}：算 $xp$、$x^2q$ 在 $x_0$ 的极限；常点走 Taylor，正则奇点走 Frobenius。
  \item \textbf{定指数}：解指标方程，记 $r_1\geq r_2$，看差值是不是整数。
  \item \textbf{算系数}：先算 $r_1$ 支；若差非整数，再算 $r_2$ 支；若差为整数，检查整数步是否自动满足，再决定要不要 $\ln x$。
  \item \textbf{收尾}：写通解、写区间、用 Wronskian 或代回验算。
\end{enumerate}

\subhdr{8.3 三个检查}

\begin{itemize}
  \item \textbf{低阶检查}：把算出的前两三项代回原方程，最低次幂的系数应为零。
  \item \textbf{指数检查}：$n=0$ 的方程必须正好是指标方程；若不等，说明前面某项乘错了。
  \item \textbf{半径检查}：$\rho$ 应等于到最近奇点的距离；若算出的级数比值法给出更小的半径，说明系数算错了。
\end{itemize}

%=============================================================
\hdr{§9}{易错点（每条都能指回一个具体例子）}

\begin{enumerate}[leftmargin=2.2em]
  \item \textbf{把 $Q/P$、$R/P$ 的系数直接当 $p_0,q_0$。}本题 $p=1+\frac{1}{2x}$，常数项是 $1$，但 $p_0=\lim xp=\frac12$。指回 §5.1。
  \item \textbf{指标方程漏掉 $r(r-1)$ 里的 $-r$。}写成 $r^2+p_0r+q_0=0$ 会得到完全不同的根。指回 §4.2 与 §5.2 的对比。
  \item \textbf{把 $a_0=0$ 当成「零解」。}正确的说法是「$a_0=0$ 等价于换一个指数」。指回 §4.2 的警告框。
  \item \textbf{两根差为整数就直接写两个 Frobenius 级数。}要先检查整数步；本题差 $\frac32$ 不用查，但 §4.5 的 $r=-3$ 就是「整数差却没有 $\ln x$」的实例。
  \item \textbf{收敛区间写成 $|x-x_0|<\rho$。}$x^r$ 在 $x<0$ 可能不是实数，答案要按题目要求的 $x>0$ 写。指回 §5.6。
  \item \textbf{「前五项」把零系数算进去。}例 5 的 $a_4=0$，例 6 只有两项非零；题目说「nonzero terms」时要数非零项。指回 §2.5、§2.6。
  \item \textbf{忘记复奇点。}$(x^2+1)y''-4xy'+6y=0$ 在实轴上没有奇点，但 $\pm\ii$ 把半径限制为 $1$。指回 §2.6 与 §4.6。
  \item \textbf{用一点 $W=0$ 判定函数线性相关。}对解组可以（Abel 公式保证要么全零要么全不零），对一般函数不行。指回 §6.4。
  \item \textbf{把 $\ee^{A+B}=\ee^A\ee^B$ 当恒等式。}只有 $A,B$ 可交换时才成立；矩阵指数一般不能拆。指回 §6.6。
\end{enumerate}

%=============================================================
\hdr{§10}{缺口清单：课件哪里跳了步，本讲义补了什么}

\gapitem{代数跳步}{L18 p21 直接给出 $\sum[(n+2)(n+1)a_{n+2}+\sum(k+1)p_{n-k}a_{k+1}+\sum q_{n-k}a_k]x^n=0$，没有说明两个级数相乘后同幂次系数为什么是卷积。补法：§2.3 用 Cauchy 乘积逐项推出，并写出 $n=0,1,2$ 三个方程。}

\gapitem{压缩缺口}{L18 p25 只写「We observe that even and odd coefficients decouple」，没有证明为什么解耦。补法：§2.4 指出递推只连 $a_k$ 与 $a_{k+2}$，二者同奇偶，所以两条链独立。}

\gapitem{术语错误}{L18 p18 写「where $P$, $Q$ and $R$ are analytic \textbf{equations}」，应为 analytic \textbf{functions}。补法：§1.6 末尾订正，并说明「解析」是对系数函数说的。}

\gapitem{适用边界}{L18 p28 只写「The power series method is only applicable when $x=0$ is an ordinary point」，没有解释为什么。补法：§2.7 说明奇点处 $p,q$ 不解析、代入后出现负幂次、无法合并。}

\gapitem{机制缺口}{L19-20 p4--p5 给出正则奇点的极限定义，但没有说明为什么偏偏乘 $(x-x_0)$ 与 $(x-x_0)^2$。补法：§3.1 的旁注把乘完后的方程写成「Euler 型 + 解析系数」，这正是 Frobenius 的出发点。}

\gapitem{平凡解表述}{L19-20 p24 写「If $a_0=0$, then $a_1=a_2=\cdots=0$, resulting in the zero solution」，字面不严谨。补法：§4.2 的警告框说明 $a_0=0$ 等价于把指数换成 $r+N$，不是零解。}

\gapitem{结构缺口}{L19-20 p24 的系数公式分母是 $(k+r)(k+r-1)+(k+r)p_0+q_0$，但没有指出它等于 $(k+r-r_1)(k+r-r_2)$，于是「整数差为什么卡住」无法解释。补法：§4.3 给出恒等式 (4.5) 并用它统一解释三种情形。}

\gapitem{笔误（一）}{L19-20 p33 写「$r(r-1)+4r=r+3r=r(r+3)=0$」，中间漏了平方，应为 $r^2+3r$。补法：§4.5 在推导中订正，并说明结论 $r(r+3)=0$ 仍对。}

\gapitem{作业留白}{L19-20 p34 写「If $r=r_2=-3$, we would actually have obtained two linearly independent solution. (Homework)」，既没给第二解，也没解释为什么没有 $\ln x$。补法：§4.5 完整算出 $y_2=x^{-3}(1-\frac{x^2}{2}-\frac{x^4}{8}-\cdots)$，并用「$n=3$ 时 $0\cdot a_3=0$ 自动成立」解释 $a=0$。}

\gapitem{笔误（二）}{L19-20 p11 的初值方程写「$1=y'(1)=\frac52c_1-3c_3$」，$c_3$ 应为 $c_2$；p15 写「$r(r-1)-7r+16-0$」，应为「$=0$」。补法：§3.5 两处均已订正。}

\gapitem{收敛区间}{L19-20 p31 只写「$\rho$ is the distance from $x_0$ to the nearest other singular point (real or complex)」，没有给算法，也没说复奇点怎么算。补法：§4.6 用三个方程列表说明，并画出 $\pm\ii$ 的收敛圆（图 2）。}

\gapitem{压缩缺口（二）}{L19-20 p22 从两个级数相乘直接跳到 $\sum[\cdots]x^{k+r-2}=0$，中间的下标合并省略。补法：§4.2 把 $x^2y''$、$x(xp)y'$、$(x^2q)y$ 三项分别写成同幂次，再合并。}

%=============================================================
\hdr{§11}{三句话总结}

\begin{enumerate}
  \item \textbf{先判点，再定形}：$P(x_0)\neq0$ 用普通幂级数；$P(x_0)=0$ 且 $xp$、$x^2q$ 解析，用 $x^r\sum a_nx^n$；否则本课不处理。
  \item \textbf{指数看指标方程，第二解看差值}：$r(r-1)+p_0r+q_0=0$；两根差非整数时两支都是纯级数，相等或差正整数时第二解可能带 $\ln x$（也可能不带，看整数步是否自动满足）。
  \item \textbf{半径看奇点}：$\rho$ 是展开点到最近奇点（含复奇点）的距离；本题 $xp=x+\frac12$、$x^2q=\frac32x-\frac12$ 都是多项式，所以 $\rho=\infty$，解在 $0<x<\infty$ 上成立。
\end{enumerate}

%=============================================================
\hdr{§12}{页码对照表}

\begin{center}
\small
\begin{tabular}{@{}p{5.4cm}p{4.2cm}p{5.2cm}@{}}
\hline
\thead{本讲义} & \thead{课件} & \thead{教材 Constanda} \\
\hline
§1 幂级数工具箱 & L18 p3--p16 & 9.1--9.2（p214--228）\\
§2 常点邻域 & L18 p17--p35 & 9.2（p220--228）\\
§3 正则奇点与 Euler & L19-20 p3--p19 & 9.3（p228--230）、4.6 \\
§4 Frobenius & L19-20 p20--p34 & 9.4.1--9.4.3（p231--243）\\
§5 作业题全解 & L19-20 p42 & 9.4.1 习题 1（p234；PDF p242）\\
§6 线性方程组 & L13--L17（本讲义补齐） & 第 7 章（线性系统）\\
§7 Bessel & L19-20 p35--p41 & 9.4 例 9.30--9.37（p237--244）\\
\hline
\end{tabular}
\end{center}

表前一句：这张表把本讲义的每一节对到课件的页码与教材的小节，方便你回查原话。逐行读法：左边是本讲义的节号，中间是课件 PDF 的页码（按 PDF 页数，不是页脚编号），右边是教材的小节与书内页码（括号里是 PDF 页码，两者相差约 8 页）。边界行：教材 9.1--9.2 的页码取自目录，未逐页核对，只作检索起点。
\bigpart{IV}{非线性与稳定性 · L21--L23}{换掉“求解”这个目标：不解方程，画相图定类型；在双曲临界点附近借用线性。课件 61 页。}
\chapterbanner{6}{临界点与稳定性}{相图（phase portrait）、稳定性定义与第一近似（first approximation）}{课件 \texttt{21\_22\_Phase.pdf}（43 页）与 \texttt{22-23\_Linearization.pdf}（18 页），2026-10-29／11-03}
\callout{\textbf{先读这一段。}这份讲义是为四道作业题写的，但它不是"题解集"。它按一条线走：先把方程看成平面上的速度场（§1），再把你学过的线性系统解法读成几何图形（§2），由此得到五种临界点的名字与判据（§3）；然后把"稳定性"从一个感觉变成一个可以验证的定义（§4），得到线性系统的稳定性定理（§5）；最后处理非线性：线性化定理告诉我们什么时候可以只看特征值，什么时候不能（§6），不能的时候用什么办法（§7）。作业的题面与考点在 §8，完整详解在附录 A，另有题面歧义的处理在附录 B。}

%=====================================================================
\hdr{§0}{路线图、记号与读法}

\subhdr{0.1 这批作业在问什么}

本批作业只有两道大题：第 1 题给三个 $2\times2$ 线性系统，要求判断临界点的类型并研究稳定性；第 2 题给一个含二次项的非线性系统，要求研究平凡解的稳定性。四个小题用到的工具并不相同，下表先给出分布。\textbf{正文只给题面与考点，答案统一在附录 A}，你可以先独立做，再翻到对应页对答案。

\begin{minipage}{\linewidth}
\tabtop
\begin{tabular}{@{}llll@{}}
\thead{序号} & \thead{系统} & \thead{要判断的} & \thead{判据与详解位置}\\
\end{tabular}
\tabmid
\begin{tabular}{@{}l l l p{0.40\linewidth}@{}}
1(1) & $x'=3x+2y,\ y'=x+4y$ & 类型、稳定性 & 相异实特征值：§3.2、§3.4；附录 A.38，第 \pageref{c6:app:a1} 页\\
1(2) & $x'=-x,\ y'=-y$ & 类型、稳定性 & 重根与几何重数：§3.5；附录 A.39，第 \pageref{c6:app:a2} 页\\
1(3) & $x'=-x+2y,\ y'=-2x-y$ & 类型、稳定性 & 共轭复特征值：§3.2；附录 A.40，第 \pageref{c6:app:a3} 页\\
2 & 二次非线性系统 & 平凡解的稳定性 & 雅可比矩阵与线性化定理：§6.4；附录 A.41，第 \pageref{c6:app:a4} 页\\
\end{tabular}
\tabbot
\end{minipage}

读这张表只需要注意一件事：前三小题的答案是"矩阵的性质"，第四小题的答案是"矩阵的性质加上一条定理"。前者你已经会算（$L14$--$L18$ 讲过），后者是这一讲真正新增的东西。

\subhdr{0.2 你手上已有什么，这一讲补什么}

你已经学过两段内容：常系数线性系统与矩阵指数（$L13$--$L18$，含特征值三种情形与 $e^{Ax}$），以及幂级数解法（$L18$--$L20$）。本讲\textbf{不需要新的积分技巧}，需要的是三样新语言：

\begin{itemize}
\item \textbf{几何语言}：一条解曲线在平面上是什么样子，什么叫"临界点"，五种类型各自长什么样（§1--§3）。没有这套语言，"类型"三个字就是死记的名词。
\item \textbf{分析语言}：稳定（\EN{stable}）、渐近稳定（\EN{asymptotically stable}）、不稳定（\EN{unstable}），三个词各自的严格含义（§4）。日常说"稳定"其实混着两个概念，必须拆开。
\item \textbf{定理语言}：把非线性系统在临界点作一阶展开，用展开后的线性部分下结论——什么条件下合法、什么条件下失效（§6--§7）。
\end{itemize}

三者的顺序不能颠倒：先有图形，才谈得上"类型"；先有定义，才谈得上"判据"；先有判据，才谈得上"对非线性系统用判据"。

\subhdr{0.3 先认三个身份：时间、状态、规则}

在查表之前，先把三样东西分开。它们不是三个新符号，而是三个身份不同的对象：

\begin{itemize}
\item \textbf{时间 $t$}：唯一的自变量。本讲一律把它读成钟表时间。
\item \textbf{状态 $(x,y)$}：某一时刻系统"在哪里"。$x,y$ 是两个可以独立取值的实数。把它们竖着排成 $\X=(x,y)^{\mathrm T}$，只是排成一列，不做任何计算；上标 $\mathsf T$ 就叫转置。
\item \textbf{规则 $\F=(P,Q)^{\mathsf T}$}：一个只认位置、不认时刻的函数。给它一个位置 $(x,y)$，它返回两个数 $P(x,y)$、$Q(x,y)$，意思是：\textbf{站在这个位置，$x$ 以速度 $P$ 在变，$y$ 以速度 $Q$ 在变。}
\end{itemize}

于是 $\X'=\F(\X)$ 的全部内容是
\[
\frac{\dd\X}{\dd t}=\F(\X)
\qquad\Longleftrightarrow\qquad
\frac{\dd x}{\dd t}=P(x,y),\quad
\frac{\dd y}{\dd t}=Q(x,y).
\]
左边是"状态怎么变"，右边是"在当前位置该往哪变、变多快"。它不是一条要解的代数方程，而是一张地图：\textbf{平面上每个位置都贴了一支箭头}，箭头的方向与长度由 $\F$ 给出。这张图叫\textbf{速度场}（\EN{velocity field}），§1.1 专门讲它。

\subhdr{0.4 记号表（查表用，不要求先背；名词英文见 §13）}

\begin{minipage}{\linewidth}
\tabtop
\begin{tabular}{@{}lll@{}}
\thead{记号} & \thead{含义} & \thead{定义处}\\
\end{tabular}
\tabmid
\begin{tabular}{@{}lll@{}}
$t$ & 自变量；本讲一律当作"时间"来读 & §0.3\\
$x,\ y$（或 $x_1,\ x_2$） & 两个未知函数，也称状态分量 & §0.3\\
$\X=(x,y)^{\mathrm T}$ & 状态向量；$^{\mathrm T}$ 表示转置 & §0.3\\
$\X'=\F(\X)$ & 自治系统（\EN{autonomous system}），$\F=(P,Q)^{\mathrm T}$ & §1.2\\
$A$ & $2\times2$ 常数矩阵，线性系统的系数 & §2.1\\
$\tau=\tr A$,\ $\Delta=\det A$ & 迹（\EN{trace}）与行列式（\EN{determinant}） & §2.2\\
$D=\tau^{2}-4\Delta$ & 特征方程的判别式（\EN{discriminant}） & §2.2\\
$\lambda_1,\ \lambda_2$ & 特征值（\EN{eigenvalue}，特征方程的根） & §2.2\\
$V$ & 实特征向量（\EN{eigenvector}）；复情形记 $w=u+\mathrm{i}v$ & §2.3\\
$\X^{*}$ & 临界点（\EN{critical point}，即平衡点） & §1.4\\
轨线、相图 & 解的像集；全部轨线的集合（\EN{trajectory, phase portrait}） & §1.3\\
平凡解（\EN{trivial solution}） & 零解 $\X(t)\equiv\mathbf{0}$ & §4.5\\
\end{tabular}
\tabbot
\end{minipage}

\subhdr{0.5 一条主线}

整份讲义只走一条线，每一步都是一个定理，不是一条经验：

\begin{enumerate}
\item 把 $\X'=\F(\X)$ 看成平面上的速度场，$\F(\X^{*})=0$ 的点叫临界点（§1）。
\item 在临界点把 $\F$ 作一阶泰勒展开，得到线性系统 $\Y'=A\Y$，$A$ 是雅可比矩阵（§2、§6.2）。
\item 线性系统的解由 $A$ 的特征值完全决定，轨线形状有五种（§2、§3）。
\item 特征值的实部决定稳定性，符号决定"进"还是"出"（§4、§5）。
\item 若所有特征值都不在虚轴上（双曲），则非线性系统在临界点附近与线性系统\textbf{同类}，结论可以搬回去（§6.4）。
\item 若有特征值落在虚轴上，线性化失效，必须看高阶项或另找李雅普诺夫函数（§6.6、§7）。
\end{enumerate}

第 2 题恰好落在第 5 步（特征值是 $1$ 与 $0.5$，都不在虚轴上），所以它\textbf{不需要}李雅普诺夫函数；§7 仍然写出来，是为了让你知道第 6 步长什么样，也知道自己为什么可以跳过它。

%=====================================================================
\hdr{§1}{把方程看成平面上的速度场}

\subhdr{1.1 速度场：每个位置贴一支箭头}

先不看定义，只看一个具体系统：
\[
x'=x,\qquad y'=0 .
\]
（撇号就是对时间求导：$x'=\mathrm{d}x/\mathrm{d}t$，$y'=\mathrm{d}y/\mathrm{d}t$。）它的规则是 $\F(x,y)=(x,0)$，这句话可以逐点读出来：

\begin{itemize}
\item 站在 $(1,1)$：箭头是 $\F(1,1)=(1,0)$，向右，长度 $1$；
\item 站在 $(-2,3)$：箭头是 $\F(-2,3)=(-2,0)$，向左，长度 $2$；
\item 站在 $(0,5)$：箭头是 $\F(0,5)=(0,0)$，\textbf{没有箭头}——这个位置是临界点，§1.4 再讲。
\end{itemize}

把平面上每个位置的箭头都画出来，就是\textbf{速度场}。"场"这个字只说一件事：每个位置都有一个值；"速度"说的是这个值的内容——状态在这里变化得多快、朝哪个方向。箭头从哪来？从规则 $\F$ 来，不需要解方程。

\subhdr{1.2 自治：右端不显含 $t$ 排除了什么}

\textbf{定义 1.1（自治系统）}\ 设 $D$ 是 $\mathbb{R}^{2}$ 中的一个区域，$\F=(P,Q)^{\mathrm T}:D\to\mathbb{R}^{2}$ 连续。若右端不显含 $t$，即系统写成
\[
\X'=\F(\X),\qquad\text{也就是}\quad
\begin{cases}
x'=P(x,y),\\[2pt]
y'=Q(x,y),
\end{cases}
\tag{1.1}
\]
则称它为\textbf{自治系统}（\EN{autonomous system}）。

"右端不显含 $t$"要精确读成一句话：\textbf{同一个位置，箭头永远是同一支}。§1.1 的系统站在 $(1,1)$ 时箭头是 $(1,0)$——今天去看是它，明天去看还是它，这就是"自治"。

但它\textbf{不}意味着箭头不变。同一个系统里，站在 $(1,1)$ 的箭头长 $1$，站在 $(2,1)$ 的箭头长 $2$。系统会移动，位置变了，箭头当然跟着变。

\callout{\textbf{最容易被误读的一处。}"速度场不随时间改变"是错的读法——它听起来像"速度不变"。准确的说法只有一句：\textbf{同一个状态 $(x,y)$，不论什么时刻到达，规则给出的变化率都相同。}至于解 $x(t)$ 自己随时间变，那是显然的，它本来就在动。自治管的是\textbf{规则}，不是\textbf{运动}。}

\subhdr{1.3 轨线、相图：把时间抹掉}

\textbf{定义 1.2（轨线、相图）}\ 设 $\X(t)$，$t\in I$，是 (1.1) 的一个解。集合
\[
\Gamma=\{\X(t):\ t\in I\}\subset\mathbb{R}^{2}
\]
称为一条\textbf{轨线}（\EN{trajectory}，也称\textbf{相轨线}、\EN{orbit}）；全部轨线（连同方向）的集合叫\textbf{相图}（\EN{phase portrait}），它画在\textbf{相平面}（\EN{phase plane}）上。

"解"与"轨线"差在哪，用最小的例子当场兑现。取
\[
x'=-y,\qquad y'=x .
\tag{1.2}
\]
直接验证 $\X(t)=(\cos t,\sin t)$ 是解：
\[
x'=-\sin t=-y,\qquad y'=\cos t=x .
\]
所以这个解告诉你的是\textbf{每个时刻的具体位置}：$t=0$ 在 $(1,0)$，$t=\pi/2$ 在 $(0,1)$，$t=\pi$ 在 $(-1,0)$。

把时间抹掉，只看它走过哪些点：
\[
x^{2}+y^{2}=\cos^{2}t+\sin^{2}t=1,
\]
单位圆。这就是\textbf{轨线}。同一个圆还可以由 $(\cos(t-t_0),\sin(t-t_0))$ 描出，$t_0$ 任意——这正是"同一条轨线有无穷多个解，彼此只差一个时间平移"。

\begin{minipage}{\linewidth}
\tabtop
\begin{tabular}{@{}lll@{}}
\thead{对象} & \thead{在这个例子里是什么} & \thead{回答什么问题}\\
\end{tabular}
\tabmid
\begin{tabular}{@{}l p{0.30\linewidth} p{0.42\linewidth}@{}}
解 & $(\cos t,\sin t)$ & 每个时刻具体在哪里\\
轨线 & 单位圆 $x^{2}+y^{2}=1$ & 一共经过哪些位置\\
相图 & 同心圆族，另有逆时针方向与原点是临界点 & 所有可能的运动放在一起长什么样\\
\end{tabular}
\tabbot
\end{minipage}

"时间参数被抹掉"于是有了具体含义：\textbf{你还看得出走的是哪个圆、朝哪个方向走，但不再标注几点走到哪里。}这也是相图能画成一张静态地图的原因。

\subhdr{1.4 临界点}

\textbf{定义 1.3（临界点）}\ 若 $\X^{*}\in D$ 使 $\F(\X^{*})=\mathbf{0}$，则称 $\X^{*}$ 为系统 (1.1) 的\textbf{临界点}（\EN{critical point}），也称\textbf{平衡点}（\EN{equilibrium point}）、\textbf{驻点}（\EN{stationary point}）；有些教材叫\textbf{奇点}（\EN{singular point}）。

用速度场的语言说，临界点就是\textbf{箭头长度为零}的位置。常值函数 $\X(t)\equiv\X^{*}$ 因此是解，它的轨线退化成一个点。其余轨线在 $\X^{*}$ 处没有确定的切线方向（斜率 $Q/P$ 在分子分母同时为零处失效），所以只能"靠近"或"离开"，不能"穿过"。

\aside{中文教材多用"平衡点"或"奇点"，本课程作业用 critical point，三者同义。作业第 1 题写的是 critical point(s)，复数只是通用写法；当 $\Delta\neq0$ 时临界点只有一个，即原点（见 §1.5）。}

\subhdr{1.5 线性系统的临界点在哪里}

对线性系统 $\X'=A\X$，临界点就是齐次线性方程组 $A\X=\mathbf{0}$ 的解：

\begin{itemize}
\item 若 $\Delta=\det A\neq0$，只有零解，\textbf{临界点唯一，就是原点}；
\item 若 $\Delta=0$，则有一整条直线上的点都是临界点，此时"类型"这个词要分开讨论，不在本讲范围。
\end{itemize}

作业第 1 题三个矩阵的行列式分别是 $10$、$1$、$5$，都不为零，所以三小题都只有一个临界点 $(0,0)$。做题时把这句写出来，是完整性的一部分。

\subhdr{1.6 拿到一张相图，先看这三处}

后面 §3 会把五种临界点各摆一张图，形状完全不同。读任何一张，顺序固定为三步，不要一上来就盯着弯弯曲曲的线：

\begin{enumerate}
\item \textbf{找零箭头}：解 $\F(\X)=\mathbf{0}$。这些点是所有轨线的"汇合处"，也是分类的锚——一张图里可以只有一个，也可以有一串。
\item \textbf{看方向}：挑一个非临界点，把它的 $(x,y)$ 代进 $\F$，看箭头朝哪。这一步只要算术，不用解方程；方向定了，才知道轨线是"进去"还是"出来"。
\item \textbf{看形状}：箭头是直的还是打转？绕不绕圈？最后趋向哪里？三问的答案拼起来，就是 §3 的类型名。
\end{enumerate}

用 (1.2) 走一遍：第一步，零箭头只有原点；第二步，在 $(1,0)$ 处 $\F=(0,1)$，箭头向上，所以是逆时针；第三步，箭头处处与半径垂直，绕圈而不螺旋，轨线是闭合圆。三步合起来，这张相图叫\textbf{中心}（\EN{center}）。

再换一个系统，同样的三步给出另一个名字：
\[
x'=x,\qquad y'=2y .
\]
第一步，零箭头仍是原点；第二步，在 $(1,1)$ 处 $\F=(1,2)$，箭头指向右上方、离开原点；第三步，两条坐标轴上的轨线是直线（$y=0$ 时 $y'\equiv0$，$x=0$ 时 $x'\equiv0$），其余轨线把两式相除、分离变量：
\[
\frac{\dd y}{\dd x}=\frac{2y}{x}
\quad\Longrightarrow\quad
\frac{\dd y}{y}=2\frac{\dd x}{x}
\quad\Longrightarrow\quad
y=Cx^{2},
\]
是一族从原点射出的抛物线。三步合起来，这张相图叫\textbf{结点}（\EN{node}）。

这两个例子说明同一件事：\textbf{类型名是轨线几何的名字，不是矩阵的名字}。§3 的全部工作，就是"不解方程，直接从矩阵读出它是哪一种几何"。

\subhdr{1.7 为什么两条轨线不能相交}

相图是一张"互不相交、带方向"的静态地图——这句话常被一笔带过，其实它需要两个前提：

\begin{enumerate}
\item \textbf{自治}：把运动的起始时刻挪一下，仍然遵守同一套方程。因为右端不含 $t$，若 $\X(t)$ 是解，则 $\X(t-t_0)$ 也是解。
\item \textbf{解的唯一性}（\EN{uniqueness of solution}）：从同一个状态出发，不能有两套不同的后续演化。
\end{enumerate}

两条合起来，论证只有三行。设两条轨线在 $\X_0$ 相交，分别由解 $\X_1(t)$、$\X_2(t)$ 描出，且 $\X_1(t_1)=\X_2(t_2)=\X_0$。由第 1 条，$\tilde{\X}(t)=\X_1(t+t_1-t_2)$ 仍是解，并且 $\tilde{\X}(t_2)=\X_0$。于是 $\tilde{\X}$ 与 $\X_2$ 在同一时刻从同一状态出发，由第 2 条只能相等：
\[
\X_2(t)=\X_1(t+t_1-t_2).
\]
也就是说，两条轨线其实是同一条，只差一个时间平移——所以它们不可能"交叉着穿过"。

\warnbox{\textbf{只有自治，推不出"不相交"。}若右端含 $t$，时间平移就不再保持解，第 1 条的论证直接断掉；若右端只是连续、不满足李普希茨（\EN{Lipschitz}）条件（例如 $x'=\sqrt{\abs{x}}$），从 $x=0$ 出发可以停在原地，也可以走开（$x=t^{2}/4$），两条解在同一处"分岔"。本讲只讨论 $\F\in C^{1}$ 的情形，作业里的右端都是多项式，两个前提自动满足。}

%=====================================================================
\hdr{§2}{线性系统：解的三种形态}

\subhdr{2.1 把两个方程写成一个}

把第 1 题的写法改成矩阵写法（这也是 $L14$--$L15$ 的写法）：
\[
\X'=A\X,\qquad
A=\begin{pmatrix}a_{11}&a_{12}\\ a_{21}&a_{22}\end{pmatrix},\qquad
\X=\begin{pmatrix}x\\ y\end{pmatrix}.
\tag{2.1}
\]
两个写法完全等价：$x'=a_{11}x+a_{12}y$，$y'=a_{21}x+a_{22}y$。写成一个式子的好处是：下面所有结论都只依赖 $A$ 的\textbf{特征值}（\EN{eigenvalue}），与四个元素怎么排布无关。

\subhdr{2.2 特征方程（characteristic equation）}

设 $\X=V\mathrm{e}^{\lambda t}$，$V\neq\mathbf{0}$ 是常向量，代入 (2.1)：
\[
\lambda V\mathrm{e}^{\lambda t}=AV\mathrm{e}^{\lambda t}
\quad\Longrightarrow\quad
AV=\lambda V .
\]
于是 $V$ 是 $A$ 的特征向量（\EN{eigenvector}），$\lambda$ 是特征值。要有非零解，必须
\[
\det(A-\lambda I)=\lambda^{2}-\tau\lambda+\Delta=0,
\qquad
\tau=\tr A=a_{11}+a_{22},\quad \Delta=\det A=a_{11}a_{22}-a_{12}a_{21}.
\tag{2.2}
\]
两根之和为 $\tau$，之积为 $\Delta$（韦达定理），判别式（\EN{discriminant}）记
\[
D=\tau^{2}-4\Delta .
\tag{2.3}
\]
这三个数 $\tau,\Delta,D$ 就是第 1 题的全部计算量：算出它们，类型与稳定性就定了。

\subhdr{2.3 情形一：相异实根}

若 $\lambda_1\neq\lambda_2$ 为实数，对应实特征向量 $V_1,V_2$，则通解（\EN{general solution}）为
\[
\X(t)=c_1 V_1\mathrm{e}^{\lambda_1 t}+c_2 V_2\mathrm{e}^{\lambda_2 t},
\qquad c_1,c_2\in\mathbb{R}.
\tag{2.4}
\]
代入即可验证：$\X'=c_1\lambda_1V_1\mathrm{e}^{\lambda_1t}+c_2\lambda_2V_2\mathrm{e}^{\lambda_2t}=A\X$。$V_1,V_2$ 线性无关（对应不同特征值），所以 (2.4) 是通解。

\subhdr{2.4 情形二：共轭复根（这一步课本常跳过）}

若 $\lambda=\alpha+\mathrm{i}\beta$（$\beta\neq0$）是特征值，因 $A$ 实，$\bar\lambda=\alpha-\mathrm{i}\beta$ 也是。对应的特征向量可以取 $w=u+\mathrm{i}v$（$u,v$ 实向量）。把复解拆成实部（\EN{real part}）与虚部（\EN{imaginary part}）：
\[
w\mathrm{e}^{\lambda t}=(u+\mathrm{i}v)\mathrm{e}^{\alpha t}(\cos\beta t+\mathrm{i}\sin\beta t)
=\mathrm{e}^{\alpha t}\Big[\underbrace{(u\cos\beta t-v\sin\beta t)}_{\text{实部}}
+\mathrm{i}\underbrace{(u\sin\beta t+v\cos\beta t)}_{\text{虚部}}\Big].
\]
实部与虚部各自都是实解（因为方程是实系数的），于是通解为
\[
\X(t)=\mathrm{e}^{\alpha t}\Big[c_1\big(u\cos\beta t-v\sin\beta t\big)
+c_2\big(u\sin\beta t+v\cos\beta t\big)\Big].
\tag{2.5}
\]
这就是"$\mathrm{e}^{\alpha t}$ 控制远近、$\beta t$ 控制旋转"的来历：它不是背下来的，是 $\mathrm{e}^{(\alpha+\mathrm{i}\beta)t}=\mathrm{e}^{\alpha t}(\cos\beta t+\mathrm{i}\sin\beta t)$ 的实部虚部。

\subhdr{2.5 情形三：重根}

若 $\lambda_1=\lambda_2=\lambda=\tau/2$，要分两种，区别在于\textbf{几何重数}（\EN{geometric multiplicity}）$g$——它指线性无关（\EN{linearly independent}）特征向量的个数：

\begin{itemize}
\item $g=2$：有两个线性无关的特征向量 $V_1,V_2$（等价于 $A=\lambda I$），通解
\[
\X=(c_1V_1+c_2V_2)\mathrm{e}^{\lambda t}.
\tag{2.6}
\]
\item $g=1$：只有一个特征向量 $V$。此时第二个解要带一个 $t$：取 $W$ 满足 $(A-\lambda I)W=V$，则
\[
\X=\big(c_1V+c_2(tV+W)\big)\mathrm{e}^{\lambda t}.
\tag{2.7}
\]
\end{itemize}

(2.7) 的验证只要一行，写出来是因为很多人记不住 $W$ 的方程该配 $V$ 还是配 $0$：由 $(A-\lambda I)W=V$ 得 $AW=\lambda W+V$，于是
\[
\X'=\big(c_2V+\lambda(c_1V+c_2(tV+W))\big)\mathrm{e}^{\lambda t},\qquad
A\X=\big(c_1\lambda V+c_2(t\lambda V+\lambda W+V)\big)\mathrm{e}^{\lambda t},
\]
两式相同。作业第 1 题 (2) 正是 $g=2$ 的情形（$A=-I$，任何非零向量都是特征向量），所以用 (2.6)，不会出现 $t$ 项。

\subhdr{2.6 三种形态怎样决定轨线的形状}

把 (2.4)、(2.5)、(2.6)、(2.7) 并排看，可以读出三条规则，后面所有类型名都由它们生成：

\begin{enumerate}
\item \textbf{径向}（离原点多远）由实部控制：$\abs{\mathrm{e}^{\lambda t}}=\mathrm{e}^{(\operatorname{Re}\lambda)t}$。实部为负则整体收缩到原点，为正则整体膨胀。
\item \textbf{角向}（是否旋转）由虚部控制：只有 $\beta\neq0$ 才有 $\cos\beta t,\sin\beta t$，轨线才会绕原点转；$\beta=0$ 时轨线不转，只会"直着走"。
\item \textbf{直线轨线}由特征向量给出：若 $\lambda$ 实、$V$ 实，则 $\X=cV\mathrm{e}^{\lambda t}$ 是一条过原点的直线，方向沿 $V$。两条这样的直线（对应 $V_1,V_2$）是相图的"骨架"，其余轨线都夹在它们之间。
\end{enumerate}

\subhdr{2.7 为什么只看特征值就够了}

读者可能会问：$A$ 的四个元素千变万化，凭什么说类型只有五种？理由是\textbf{可逆线性变换}（\EN{invertible linear transformation}）\textbf{保持相图结构}。

设 $A=PBP^{-1}$（$P$ 可逆），令 $\X=P\Y$，则
\[
\X'=A\X\iff P\Y'=PBP^{-1}P\Y\iff \Y'=B\Y .
\]
也就是说：$A$ 的系统经一个可逆线性变换（拉伸、剪切，不折叠）后，就是 $B$ 的系统。可逆线性变换把直线映成直线、把闭曲线映成闭曲线、把"相切"映成"相切"、把"趋于原点"映成"趋于原点"。所以 $A$ 与 $B$ 的相图\textbf{拓扑等价}（\EN{topologically equivalent}），类型名当然一样。而任意矩阵都可以化到\textbf{若尔当标准形}（\EN{Jordan canonical form}），其对角元正是特征值。于是：
\[
\text{类型与稳定性只取决于特征值（以及重根时的几何重数）}.
\]
这条结论是所有"判别表"的根据；§6 的非线性定理，本质上是它的"局部版本"。

%=====================================================================
%=====================================================================
\hdr{§3}{五种临界点：名字、判据、图形}

\subhdr{3.1 判别式：把两个特征值换成三个数}

由 (2.2)，特征值可以写成
\[
\lambda_{1,2}=\frac{\tau\pm\sqrt{D}}{2},\qquad D=\tau^{2}-4\Delta .
\tag{3.1}
\]
于是判断类型不必真的解方程，只看三个数的符号：$D$ 决定"实还是复"，$\Delta=\lambda_1\lambda_2$ 决定"同号还是异号"，$\tau=\lambda_1+\lambda_2$ 决定"正还是负"。\textbf{下表只是把 (3.1) 与韦达定理读了一遍}，不承载新的推导。

\begin{minipage}{\linewidth}
\tabtop
\begin{tabular}{@{}llll@{}}
\thead{条件} & \thead{特征值} & \thead{类型} & \thead{稳定性}\\
\end{tabular}
\tabmid
\begin{tabular}{@{}p{0.19\linewidth} p{0.20\linewidth} p{0.24\linewidth} p{0.27\linewidth}@{}}
$\Delta<0$ & 实，异号 & 鞍点 & 恒不稳定\\[2pt]
$\Delta>0,\ D>0$ & 实，同号且不等 & 结点 & $\tau<0$ 稳定；$\tau>0$ 不稳定\\[2pt]
$\Delta>0,\ D=0$ & 重根 $\lambda=\tau/2$ & 临界结点（$g=2$）或退化结点（$g=1$） & $\tau<0$ 稳定；$\tau>0$ 不稳定\\[2pt]
$\Delta>0,\ D<0$，\newline $\tau\neq0$ & 复，实部 $\tau/2\neq0$ & 焦点 & $\tau<0$ 稳定；$\tau>0$ 不稳定\\[2pt]
$\Delta>0,\ D<0$，\newline $\tau=0$ & 纯虚 $\pm\mathrm{i}\beta$ & 中心 & 稳定，但非渐近\\[2pt]
$\Delta=0$ & 有零根 & 退化（整条直线都是临界点） & 本题不出现\\
\end{tabular}
\tabbot
\end{minipage}

\textbf{逐行读法}（每一行都对应 $\lambda_{1,2}=(\tau\pm\sqrt D)/2$ 的一次代入）：

\begin{itemize}
\item 第一行：$D=\tau^{2}-4\Delta>0$ 且 $\Delta<0$，故 $\sqrt{D}>\abs{\tau}$，一根为正、一根为负——这是鞍点。
\item 第二行：$D>0$ 给实根，$\Delta>0$ 给同号，$\tau$ 的符号就是两根共同的符号。$\tau>0$ 时两根都为正（都往外跑）。
\item 第三行：$D=0$ 时 $\sqrt{D}=0$，两根重合为 $\tau/2$；此时"结点"的具体形状还要看几何重数 $g$，见 §3.5。
\item 第四行：$D<0$ 给复根 $\tau/2\pm\mathrm{i}\sqrt{\abs{D}}/2$，实部就是 $\tau/2$；实部非零才会既有旋转又有进出。
\item 第五行：$\tau=0$ 使实部为零，只剩纯虚根 $\pm\mathrm{i}\sqrt{\Delta}$，轨线是闭合椭圆——中心。它是这张表里唯一"稳定但非渐近"的行。
\item 第六行：$\Delta=0$ 表示 $A$ 奇异，$A\X=\mathbf{0}$ 有非零解，临界点不再唯一，本章的判据不适用。
\end{itemize}

\textbf{边界行（$D=0$）要单独记住}：$D=0$ 的点恰好落在图 \ref{c6:fig:taud} 的抛物线上，那里既不是"典型的结点"也不是"焦点"，而是重根。作业第 1 题 (2) 就落在这条抛物线上（$\tau=-2$，$\Delta=1$，$D=0$），所以它必须走 §3.5 的分析，不能套"$D>0$ 为结点"那一行。

\subhdr{3.2 五种类型：逐个看图形}

\figblock[fig:kinds]{fig1-kinds.pdf}{0.99}{六种基本相图（本讲插图，数值积分绘制）。浅色为方向场（\EN{direction field}），深色为轨线，虚线为特征方向（\EN{eigendirection}）。其中 (a)(c)(e) 正是作业第 1 题的三个系统：$A=\begin{pmatrix}3&2\\1&4\end{pmatrix}$、$A=\begin{pmatrix}-1&2\\-2&-1\end{pmatrix}$、$A=-I$；(b)(d)(f) 是对照用的标准例子。箭头的方向就是时间增加的方向。}

\subhdr{（1）结点（node）}

条件：$\Delta>0$，$D>0$，两根同号。稳不稳定看 $\tau$：两根都负则所有轨线都进去（稳定），都正则都出来（不稳定）。

形状：两条\textbf{直线轨线}沿 $V_1,V_2$ 的方向；其余轨线是曲线，进入（或离开）原点时\textbf{与"慢方向"相切}。所谓慢方向，就是 $\abs{\lambda}$ 较小的那个特征向量方向。原因是 (2.4) 可以提公因子：
\[
\X=\mathrm{e}^{\lambda_2 t}\Big[c_1V_1\mathrm{e}^{(\lambda_1-\lambda_2)t}+c_2V_2\Big],
\qquad \lambda_1<\lambda_2 .
\]
当 $t\to-\infty$（沿时间倒着看，轨线逼近原点），$\mathrm{e}^{(\lambda_1-\lambda_2)t}\to0$，方括号趋于 $c_2V_2$，所以切线方向就是 $V_2$——较小的特征值。$\tau>0$ 时把 $t\to+\infty$ 代入同一句话，切向仍是 $V_2$（$\lambda_2>\lambda_1$，$\mathrm{e}^{\lambda_2t}$ 长得快）。

图 \ref{c6:fig:kinds}(a) 是 $\lambda=2,5$ 的情形：两条直线轨线分别是 $(-2,1)$ 与 $(1,1)$ 方向，其余轨线离开原点时先贴着 $(-2,1)$（慢方向），走远后才被 $(1,1)$ 拉走。

\subhdr{（2）鞍点（saddle point）}

条件：$\Delta<0$，两根异号。\textbf{恒不稳定}，不必再看 $\tau$。

形状：两条直线轨线，一条进来（沿负特征值方向），一条出去（沿正特征值方向），其余轨线是双曲线型，进来一阵又转出去。图 \ref{c6:fig:kinds}(b) 用的 $A=\operatorname{diag}(1,-1)$，两条直线就是坐标轴。

鞍点在作业里不会出现（第 1 题三个矩阵的 $\Delta$ 都大于零），但它是"$\Delta<0$ 就是不稳定"这条判据的全部理由，考试里常考，值得记住图形。

\subhdr{（3）焦点（focus；课件作 spiral point）}

条件：$\Delta>0$，$D<0$，且 $\tau\neq0$（实部不为零）。稳定性由 $\tau$ 定：$\tau<0$ 稳定，$\tau>0$ 不稳定。

形状：螺旋线，绕原点转无穷多圈，同时半径按 $\mathrm{e}^{\tau t/2}$ 缩放。由 (2.5)，极坐标下半径正比于 $\mathrm{e}^{\tau t/2}$，角度是 $\beta t$ 加常数，两者合起来就是对数螺线（\EN{logarithmic spiral}）——每转一圈，半径乘同一个因子。

图 \ref{c6:fig:kinds}(c) 是作业第 1 题 (3) 的系统：$\lambda=-1\pm2\mathrm{i}$，所以每单位时间转 $2$ 弧度、半径乘 $\mathrm{e}^{-1}$，是"顺时针向内卷"的稳定焦点。

\subhdr{（4）中心（center）}

条件：$\tau=0$，$\Delta>0$（于是 $D=-4\Delta<0$，纯虚根 $\pm\mathrm{i}\sqrt{\Delta}$）。

形状：闭合椭圆，轨线绕原点转圈但半径不变（(2.5) 中 $\mathrm{e}^{\alpha t}=1$）。这是 §1.3 的圆族的一般情形；作一个可逆线性变换就能把圆拉成椭圆，所以一般不写"圆"而写"椭圆"。

中心\textbf{稳定但不渐近稳定}：初值落在哪条椭圆上，就永远待在那条椭圆上，回不到原点。这一点是 §4 的核心例子。

\subhdr{（5）退化情形：重根}

条件：$D=0$，$\lambda_1=\lambda_2=\tau/2$。此时"类型"由\textbf{几何重数} $g$ 决定，两种情形在图上能一眼分辨：

\begin{itemize}
\item $g=2$（$A=\lambda I$）：每个方向都是特征方向，轨线是\textbf{过原点的直线束}，这种叫\textbf{临界结点}（也译星形结点，\EN{proper node / star node}）。(2.6) 里没有 $t$，所以 $y/x$ 是常数。
\item $g=1$：只有一个特征方向，其余轨线都\textbf{切着这一个方向}进出，叫\textbf{退化结点}（\EN{improper node / degenerate node}）。(2.7) 里有 $t$ 因子，轨线会拐弯。
\end{itemize}

作业第 1 题 (2) 是 $A=-I$，属于 $g=2$：图 \ref{c6:fig:kinds}(e) 中所有轨线都是笔直的射线，没有一条弯。

\subhdr{3.3 中英文名称对照}

同一批图形在不同教材里叫法不同，做题时认准"条件"就不怕名字乱。\textbf{下表是一张台账，不承载推导}，只把名字与条件对齐。

\begin{minipage}{\linewidth}
\tabtop
\begin{tabular}{@{}llll@{}}
\thead{中文} & \thead{英文常见写法} & \thead{特征值条件} & \thead{一眼认出的特征}\\
\end{tabular}
\tabmid
\begin{tabular}{@{}l l l p{0.36\linewidth}@{}}
结点 & node & $\Delta>0,\ D>0$ & 两条直线轨线，其余切于慢方向\\
鞍点 & saddle (point) & $\Delta<0$ & 一进一出两条直线，其余双曲\\
焦点 & focus（spiral point） & $\Delta>0,\ D<0,\ \tau\neq0$ & 螺旋，无限多圈\\
中心 & center & $\tau=0,\ \Delta>0$ & 闭合椭圆\\
临界结点 & star node, proper node & $D=0,\ g=2$ & 直线束，各方向皆特征向\\
退化结点 & degenerate node & $D=0,\ g=1$ & 只有一个特征方向\\
\end{tabular}
\tabbot
\end{minipage}

注意英文的 \EN{proper node} 与 \EN{star node} 是同一个东西（对应中文“临界结点”，$g=2$），而 \EN{improper node} 与 \EN{degenerate node} 是另一个东西（对应中文“退化结点”，$g=1$）；普通结点（$\lambda_1\neq\lambda_2$，$\Delta>0,D>0$）就叫 \EN{node}。名字不同不要紧，条件与 $g$ 才是判据。

\subhdr{3.4 一张图看完所有情形：$\tau$–$\Delta$ 平面}

把表 3.1 的条件画在 $(\tau,\Delta)$ 平面上，就是图 \ref{c6:fig:taud}：抛物线 $\Delta=\tau^{2}/4$（即 $D=0$）把上半平面分成"结点区"（抛物线内侧）与"焦点区"（外侧），$\Delta<0$ 的下半平面全是鞍点，纵轴 $\tau=0$ 的正半轴是中心。\textbf{图上的四个点正是这批作业的四个线性部分}，做题时可以先用眼睛定位，再用代数核实。

\figblock[fig:taud]{fig2-taud.pdf}{0.86}{$\tau$–$\Delta$ 判别图（本讲绘制）。纵轴是 $\Delta=\det A$，横轴是 $\tau=\tr A$；粗曲线是 $D=\tau^2-4\Delta=0$。\cnum{1}–\cnum{4} 依次是作业第 1 题 (1)(2)(3) 与第 2 题在原点的雅可比矩阵：\cnum{1}$(7,10)$ 落结点区右上，\cnum{2}$(-2,1)$ 恰在抛物线上，\cnum{3}$(-2,5)$ 落焦点区左上，\cnum{4}$(1.5,0.5)$ 落结点区右下方。}

\begin{minipage}{\linewidth}
\tabtop
\begin{tabular}{@{}lllll@{}}
\thead{题号} & \thead{$\tau$} & \thead{$\Delta$} & \thead{$D$} & \thead{落在哪个区、是什么类型}\\
\end{tabular}
\tabmid
\begin{tabular}{@{}l l l l p{0.46\linewidth}@{}}
1(1) & $7$ & $10$ & $9>0$ & 结点区（$\tau>0$）：不稳定结点\\
1(2) & $-2$ & $1$ & $0$ & 抛物线上的重根：临界结点（$g=2$）\\
1(3) & $-2$ & $5$ & $-16<0$ & 焦点区（$\tau<0$）：稳定焦点\\
2 & $3/2$ & $1/2$ & $1/4>0$ & 结点区（$\tau>0$）：不稳定结点（线性化部分）\\
\end{tabular}
\tabbot
\end{minipage}

这张表是把图 \ref{c6:fig:taud} 的四个点逐个读一遍：横坐标看左右（$\tau$ 的符号），纵坐标看上下（$\Delta$ 的符号），再与抛物线比较看内外。四行都能在图上找到对应的小圆圈，这就是"用眼睛复核代数"的意思。

\subhdr{3.5 重根为什么有两种图形}

§2.5 已经给出两种通解，这里只把差别说透。$g=2$ 时 (2.6) 是
\[
\X=(c_1V_1+c_2V_2)\mathrm{e}^{\lambda t},
\]
括号里是常向量，所以 $\X$ 与某个固定方向共线——所有轨线都是直线。$g=1$ 时 (2.7) 是
\[
\X=\big(c_1V+c_2(tV+W)\big)\mathrm{e}^{\lambda t},
\]
括号里随 $t$ 线性增长，方向慢慢转向 $W$，所以轨线不再是直线，而是"贴着 $V$ 方向"弯进去的曲线。

判 $g$ 的办法只有一条：解 $(A-\lambda I)V=\mathbf{0}$，看有几个线性无关解。对 $2\times2$ 矩阵：
\begin{itemize}
\item 若 $A-\lambda I=0$（即 $A=\lambda I$），$g=2$；
\item 否则 $g=1$，此时要按 (2.7) 求 $W$。
\end{itemize}

\subhdr{3.6 只看 $\Delta$ 会错在哪里}

三个常见误判，都能用图 \ref{c6:fig:taud} 当场驳倒：
\begin{enumerate}
\item \textbf{"$\Delta>0$ 就稳定"}——错。$\Delta>0$ 只说明两根同号，同为正则不稳定（\cnum{1}\cnum{4} 就是）。
\item \textbf{"$\Delta>0$ 就是结点"}——错。要看 $D$：$D<0$ 是焦点（\cnum{3}），$\tau=0$ 是中心。
\item \textbf{"$\tau<0$ 就稳定"}——在 $\Delta>0$ 的前提下对；但 $\Delta<0$ 时 $\tau<0$ 仍是鞍点，恒不稳定。
\end{enumerate}
一句话：\textbf{先看 $\Delta$ 的符号定"鞍点 / 非鞍点"，再看 $D$ 定"实 / 复"，最后看 $\tau$ 定"进 / 出"。}三步顺序不能换。

%=====================================================================
\hdr{§4}{稳定性：把感觉变成定义}

\subhdr{4.1 从图上问三个问题}

看任何一张相图，都可以问三个层层递进的问题：
\begin{enumerate}
\item 初值离临界点很近时，轨线会不会\textbf{跑远}？（跑不远 = 稳定）
\item 不但跑不远，是不是最终\textbf{回到}临界点？（回来 = 渐近稳定）
\item 如果无论把初值取得多近，总有一条轨线\textbf{跑出去}？（= 不稳定）
\end{enumerate}
注意第一个问题问的是"会不会跑远"，不是"跑到哪"。中心就是"跑不远、也不回来"的典型：初值落在半径 $r$ 的椭圆上，轨线永远在那条椭圆上转。

\subhdr{4.2 三个定义}

设 $\X^{*}=\mathbf{0}$ 是系统 (1.1) 的临界点（一般的 $\X^{*}$ 可平移到原点，见 §4.5），记 $\norm{\X}=\sqrt{x^{2}+y^{2}}$。

\textbf{定义 4.1（稳定，Lyapunov stable）}\ 若对任意 $\eps>0$，存在 $\delta>0$，使得只要 $\norm{\X(0)}<\delta$，就有
\[
\norm{\X(t)}<\eps\qquad\text{对一切 }t\ge0,
\]
则称零解（从而临界点）\textbf{稳定}。

\textbf{定义 4.2（渐近稳定，asymptotically stable）}\ 若零解稳定，且存在 $\delta_0>0$，使得只要 $\norm{\X(0)}<\delta_0$，就有
\[
\lim_{t\to+\infty}\X(t)=\mathbf{0},
\]
则称零解\textbf{渐近稳定}。

\textbf{定义 4.3（不稳定，unstable）}\ 若存在 $\eps>0$，使得对任意 $\delta>0$，总存在初值 $\X(0)$ 与时刻 $t_1>0$ 满足
\[
\norm{\X(0)}<\delta,\qquad \norm{\X(t_1)}\ge\eps,
\]
则称零解\textbf{不稳定}。

三个定义的读法：

\begin{itemize}
\item $\eps$ 是"允许的活动范围"，$\delta$ 是"初值能取的范围"。定义说的是：\textbf{你给多大的活动范围，我都能找到一个相应的初值范围}，使轨线不越界。
\item $\delta$ 可以依赖 $\eps$，但\textbf{不能依赖初值}——这正是"一致"的意思。
\item 不稳定是"稳定的否定"：只要存在一个 $\eps$，使得任何 $\delta$ 都关不住某条轨线，就叫不稳定。它不要求"所有轨线都跑到无穷"，只要求"有一条跑出去"。
\end{itemize}

\figblock[fig:epsdelta]{fig3-epsdelta.pdf}{0.93}{$\eps$–$\delta$ 两个圆（本讲绘制）。(a) 稳定且渐近：初值取在 $\delta$ 圆内（半径 $0.5$），轨线全部留在 $\eps$ 圆（半径 $1$）内，并且 $t\to\infty$ 时趋于原点。(b) 不稳定：初值取得极小（半径 $0.06$，远小于 $\delta$），仍有一条轨线越出 $\eps$ 圆——所以无论 $\delta$ 取多小都关不住。}

\subhdr{4.3 三个定义的关系}

\begin{itemize}
\item 渐近稳定 $\Longrightarrow$ 稳定（定义 4.2 的第一句就是"稳定"）。
\item 稳定 $\not\Longrightarrow$ 渐近稳定。在线性系统里（$\Delta\neq0$），唯一的反例来源是\textbf{中心}：$\tau=0$、$\Delta>0$，轨线是闭合椭圆，永远不回到原点，但也永远不跑远。
\item 不稳定 $=$ 非稳定，它不是"渐近稳定"的反义词。逻辑关系是：渐近稳定 $\Longrightarrow$ 稳定，所以不稳定 $\Longrightarrow$ 不渐近稳定。反过来，"不渐近稳定"包含了两种情况——不稳定（如不稳定结点），和"稳定但不渐近"（如中心）。
\end{itemize}

\subhdr{4.4 一维对照：$x'=ax$}

一维情形把三个概念摊得最开：
\[
x(t)=x_0\mathrm{e}^{at}.
\]
\begin{itemize}
\item $a<0$：$x_0$ 再大也趋于 $0$，渐近稳定。
\item $a=0$：$x(t)\equiv x_0$，轨线停在原地，稳定但不渐近。
\item $a>0$：$x_0\neq0$ 时 $\abs{x(t)}\to\infty$，不稳定。
\end{itemize}
二维的判据 (5.1) 就是把这三句话搬到"实部"上：$\operatorname{Re}\lambda<0$ 相当于 $a<0$，$\operatorname{Re}\lambda=0$ 相当于 $a=0$，$\operatorname{Re}\lambda>0$ 相当于 $a>0$。

\subhdr{4.5 什么叫"平凡解"，什么叫"研究临界点的稳定性"}

线性齐次系统 $\X'=A\X$ 永远有零解 $\X(t)\equiv\mathbf{0}$，这个解叫\textbf{平凡解}（\EN{trivial solution}）。同时，$\mathbf{0}$ 又是它的临界点（因为 $A\mathbf{0}=\mathbf{0}$）。所以对这类系统，"平凡解的稳定性"、"原点的稳定性"、"临界点的稳定性"是同一句话。作业第 2 题写的是 "the stability of the trivial solution"，指的就是 $\X=\mathbf{0}$ 这个解。

对非线性系统，$\X=\mathbf{0}$ 是不是临界点要单独验证：把 $\X=\mathbf{0}$ 代入右端，看是否为零向量。第 2 题的两个右端都不含常数项，所以代入得零，原点确为临界点——这一步在详解里必须写出来，它是"平凡解"三个字的前提。

\subhdr{4.6 三个容易混的说法}

\begin{enumerate}
\item \textbf{"稳定"不是"有界"}。稳定说的是"初值小则解小"，是一致性的陈述；只说"解有界"不够，因为界可能依赖初值任意地大。
\item \textbf{"不稳定"不是"所有解都发散"}。以鞍点 $x'=x,\ y'=-y$ 为例，沿 $y$ 轴出发的解（$x\equiv0$）是趋于原点的，有界得很；但沿 $x$ 轴出发的解跑到无穷，所以整个系统不稳定。判不稳定只需\textbf{一条}反例轨线。
\item \textbf{"渐近稳定"比"稳定"多一个"回来"}，而"回来"必须是 $t\to+\infty$ 的极限，不能是"绕了几圈看起来近了"。
\end{enumerate}

%=====================================================================
\hdr{§5}{线性系统的稳定性定理}

\subhdr{5.1 定理（完整陈述）}

\textbf{定理 5.1}\ 设 $A$ 是 $2\times2$ 实常数矩阵，考虑 $\X'=A\X$，$\lambda_1,\lambda_2$ 为其特征值。

\begin{enumerate}
\item[(i)] 若 $\operatorname{Re}\lambda_1<0$ 且 $\operatorname{Re}\lambda_2<0$，则零解\textbf{渐近稳定}（事实上是指数衰减，\EN{exponential decay}：存在 $C,\gamma>0$ 使 $\norm{\X(t)}\le C\mathrm{e}^{-\gamma t}\norm{\X(0)}$）。
\item[(ii)] 若 $\operatorname{Re}\lambda_1>0$ 或 $\operatorname{Re}\lambda_2>0$，则零解\textbf{不稳定}。
\item[(iii)] 若 $\operatorname{Re}\lambda_1\le0$、$\operatorname{Re}\lambda_2\le0$，且实部为零的特征值都满足"代数重数（\EN{algebraic multiplicity}）$=$ 几何重数"（即没有非对角的若尔当块），则零解\textbf{稳定但不渐近稳定}。
\item[(iv)] 其余情形（存在实部为零的重特征值，且对应若尔当块非对角）则零解\textbf{不稳定}。
\end{enumerate}

\subhdr{5.2 证明}

证明不需要新工具，只要把 §2 的四种通解形式逐一看"能不能被 $\eps$ 关住"。

\textbf{情形 $D>0$（相异实根）}：由 (2.4)，
\[
\norm{\X(t)}\le\abs{c_1}\norm{V_1}\mathrm{e}^{\lambda_1t}+\abs{c_2}\norm{V_2}\mathrm{e}^{\lambda_2t}.
\]
若 $\lambda_1,\lambda_2<0$，取 $\gamma=-\max(\lambda_1,\lambda_2)>0$，则右端 $\le C\norm{\X(0)}\mathrm{e}^{-\gamma t}\to0$，得 (i)。若某个 $\lambda_j>0$，沿 $V_j$ 方向取初值（即 $c_j\neq0$，另一系数为零），则 $\norm{\X(t)}=\abs{c_j}\norm{V_j}\mathrm{e}^{\lambda_jt}\to\infty$，对任意 $\eps$ 终将越界，得 (ii)。

\textbf{情形 $D<0$（共轭复根）}：由 (2.5)，$\norm{\X(t)}\le M\mathrm{e}^{\alpha t}\norm{\X(0)}$，$\alpha=\tau/2=\operatorname{Re}\lambda$。$\alpha<0$ 时指数衰减，得 (i)；$\alpha>0$ 时沿任意非零初值都指数增长，得 (ii)；$\alpha=0$ 时 $\norm{\X(t)}$ 只由三角函数决定，$\norm{\X(t)}\le M\norm{\X(0)}$ 且轨线闭合、不趋于零——稳定但不渐近，这正是 (iii) 中"纯虚单根"的情形。

\textbf{情形 $D=0$（重根）}：
\begin{itemize}
\item $g=2$：由 (2.6)，$\norm{\X(t)}\le M\mathrm{e}^{\lambda t}\norm{\X(0)}$，$\lambda=\tau/2$。$\lambda<0$ 得 (i)，$\lambda>0$ 得 (ii)，$\lambda=0$ 得 (iii)（此时 $\X$ 是常向量，稳定不渐近）。
\item $g=1$：由 (2.7)，$\norm{\X(t)}\le M(1+t)\mathrm{e}^{\lambda t}\norm{\X(0)}$。$\lambda<0$ 时 $(1+t)\mathrm{e}^{\lambda t}\to0$，仍得 (i)（只是衰减慢一点，不是指数）；$\lambda>0$ 得 (ii)；$\lambda=0$ 时 $\norm{\X(t)}$ 随 $t$ 线性增长，对任意 $\eps$ 终将越界，得 (iv)。
\end{itemize}

定理证完。注意 (i) 中的"指数衰减"对 $g=1$ 的重根不成立（有 $t$ 因子），但"渐近稳定"仍然成立——所以定理只说渐近稳定，把指数衰减放在括号里作为加强陈述。

\subhdr{5.3 二维的实用判据}

把定理 5.1 与表 3.1 合并，得到只涉及 $\tau,\Delta$ 的判据（二维的 Routh–Hurwitz 判据）：

\textbf{推论 5.2}\ 对 $\X'=A\X$（$\Delta\neq0$）：
\begin{itemize}
\item $\Delta>0$ 且 $\tau<0$ $\iff$ 零解渐近稳定（两根实部都负）；
\item $\Delta>0$ 且 $\tau=0$ $\iff$ 零解稳定但非渐近（中心）；
\item $\Delta>0$ 且 $\tau>0$ $\Longrightarrow$ 零解不稳定；
\item $\Delta<0$ $\Longrightarrow$ 零解不稳定（鞍点）。
\end{itemize}
注意第一条是"当且仅当"：在 $\Delta\neq0$ 的前提下，渐近稳定与 $\tau<0,\Delta>0$ 完全等价，不需要再算 $D$。

\subhdr{5.4 两个边界例子}

\begin{itemize}
\item $A=\begin{pmatrix}0&1\\-1&0\end{pmatrix}$：$\tau=0$，$\Delta=1$，$\lambda=\pm\mathrm{i}$，中心。稳定，但 \textbf{不}渐近稳定。这说明"$\tau=0$"这一行不能并入"$\tau<0$"。
\item $A=\begin{pmatrix}0&1\\0&0\end{pmatrix}$：$\tau=0$，$\Delta=0$，$\lambda=0$ 为二重根且 $g=1$。解为 $x=x_0+y_0t$，$y=y_0$，随 $t$ 线性增长，不稳定。这说明"$\tau=0$"也不能一概算稳定——要看有没有 $t$ 因子（即若尔当块，\EN{Jordan block}）。
\end{itemize}

这两例合起来就是定理 5.1 第 (iii)(iv) 两条的存在理由，也是 §3.5 强调"重根要看 $g$"的原因。

%=====================================================================
\hdr{§6}{非线性系统：临界点、线性化与一条定理}

\subhdr{6.1 非线性的临界点可能不止一个}

对线性系统，临界点就是 $A\X=\mathbf{0}$ 的解，要么唯一（原点），要么一整条直线。对非线性系统，临界点是\textbf{非线性方程组} $\F(\X)=\mathbf{0}$ 的解，可能没有、可能一个、可能好几个。

作业第 2 题的系统是
\[
\begin{cases}
x_1'=x_1-x_1^{2}-x_1x_2,\\[2pt]
x_2'=0.5x_2-0.75x_1x_2-0.25x_2^{2}.
\end{cases}
\tag{6.1}
\]
把 $\X=\mathbf{0}$ 代入右端，两个式子都是 $0$，所以原点确为临界点，$\X(t)\equiv\mathbf{0}$ 确为它的平凡解——题目问的"the trivial solution"就是这个解。至于这个系统还有没有别的临界点，本题不必回答，但 §6.8 会把它算完，因为那是检验你是否真的会用线性化的最好练习。

\subhdr{6.2 线性化（linearization）就是一阶泰勒展开}

设 $\X^{*}$ 是临界点，$\F$ 在 $\X^{*}$ 附近 $C^{1}$。把 $\F$ 在 $\X^{*}$ 处展开到一阶：
\[
\F(\X)=\underbrace{\F(\X^{*})}_{=\,\mathbf{0}}+\Jac(\X^{*})(\X-\X^{*})+o\big(\norm{\X-\X^{*}}\big),
\qquad
\Jac(\X^{*})=\begin{pmatrix}
\dfrac{\partial P}{\partial x}&\dfrac{\partial P}{\partial y}\\[8pt]
\dfrac{\partial Q}{\partial x}&\dfrac{\partial Q}{\partial y}
\end{pmatrix}_{\X=\X^{*}},
\tag{6.2}
\]
其中 $\Jac(\X^{*})$ 叫\textbf{雅可比矩阵}（\EN{Jacobian matrix}，有的书写 Jacobi matrix）。令 $\U=\X-\X^{*}$（把临界点平移到原点），则
\[
\U'=\Jac(\X^{*})\U+o(\norm{\U}),
\tag{6.3}
\]
扔掉高阶项得到的线性系统
\[
\U'=\Jac(\X^{*})\U
\tag{6.4}
\]
称为原系统在 $\X^{*}$ 处的\textbf{线性化系统}（\EN{linearized system}），$\Jac(\X^{*})$ 称为\textbf{线性部分}（\EN{linear part}）。

\callout{\textbf{与上一讲的连接。}你在 $L18$--$L20$ 学的是"把一个函数写成幂级数"：$f(x)=f(0)+f'(0)x+\frac{f''(0)}{2!}x^{2}+\cdots$。式 (6.2) 是同一件事的二维版本：把向量函数写成"常数项 $+$ 一次项 $+$ 二次以上"。区别只有两个：一是这里的常数项恰好为零（因为 $\X^{*}$ 是临界点），二是我们\textbf{只保留一次项}——因为在足够小的邻域里，二次以上的项比一次项小得多，而稳定性与类型只由"一次项占不占主导"决定。这件事什么时候成立，就是 §6.4 的定理要回答的。}

\subhdr{6.3 双曲临界点}

\textbf{定义 6.1（双曲）}\ 若 $\Jac(\X^{*})$ 的特征值全都满足 $\operatorname{Re}\lambda\neq0$，则称 $\X^{*}$ 为\textbf{双曲临界点}（\EN{hyperbolic critical point}）。

"双曲"这个词与双曲函数、双曲几何没有关系，它只是"特征值不落在\textbf{虚轴}（\EN{imaginary axis}）上"的简称。它的用处是：\textbf{双曲时线性化可靠}。

\subhdr{6.4 线性化定理}

\textbf{定理 6.1（李雅普诺夫第一方法，Lyapunov's first method；稳定性部分）}\ 设 $\F\in C^{1}$，$\F(\mathbf{0})=\mathbf{0}$，$A=\Jac(\mathbf{0})$，$\lambda_1,\lambda_2$ 为 $A$ 的特征值。则
\begin{enumerate}
\item[(a)] 若 $\operatorname{Re}\lambda_1<0$ 且 $\operatorname{Re}\lambda_2<0$，则零解\textbf{渐近稳定}；
\item[(b)] 若 $\operatorname{Re}\lambda_1>0$ 或 $\operatorname{Re}\lambda_2>0$，则零解\textbf{不稳定}。
\end{enumerate}

\textbf{定理 6.2（Hartman–Grobman，类型部分）}\ 若 $\mathbf{0}$ 是双曲临界点，则在 $\mathbf{0}$ 的某个邻域（\EN{neighborhood}）内，非线性系统 (6.3) 的相图与线性系统 (6.4) 的相图\textbf{拓扑等价}：存在把两者互相映过去的连续可逆变换，把轨线映成轨线、把方向映成方向。

两个定理分工不同：6.1 管"稳不稳"，6.2 管"叫什么名字"。做作业时两句都要写——只写"渐近稳定"而不说"稳定焦点"，或者只说"结点"而不说"稳定"，都是半截答案。\textbf{第 2 题的答案只需要用到 6.1(b)}：线性部分有正特征值，于是不稳定；再借 6.2 补一句"邻域内呈结点型发散"。

课件把用一次项判稳定性这件事叫\textbf{第一近似}（\EN{first approximation}）：\textbf{第一近似系统}（\EN{system in the first approximation}）就是 (6.4)，而第一近似判不出结论的情形就是 §6.6。答“由第一近似可知……”与答“由定理 6.1 可知……”是同一句话的两种说法，考试里两种都算对。

\subhdr{6.5 定理 6.1 为什么成立（证明思路）}

把 (6.3) 写成
\[
\X'=A\X+\mathbf{R}(\X),\qquad \norm{\mathbf{R}(\X)}\le\eps\norm{\X}\ \ (\norm{\X}\le\rho),
\tag{6.5}
\]
这里 $\eps$ 可以取得任意小（$\mathbf{R}$ 是 $o(\norm{\X})$ 的直接翻译）。由常数变易公式，
\[
\X(t)=\mathrm{e}^{At}\X(0)+\int_0^{t}\mathrm{e}^{A(t-s)}\mathbf{R}\big(\X(s)\big)\dd s .
\tag{6.6}
\]
关键是一个矩阵估计：若所有特征值实部 $<0$，则存在 $K\ge1$、$\gamma>0$ 使
\[
\norm{\mathrm{e}^{At}}\le K\mathrm{e}^{-\gamma t}\qquad(t\ge0).
\tag{6.7}
\]
（理由：把 $A$ 化到若尔当标准形，每个块 $\mathrm{e}^{\lambda t}$ 前面至多带一个 $t^{k}$ 因子，而 $t^{k}\mathrm{e}^{\lambda t}\le C\mathrm{e}^{(\operatorname{Re}\lambda+\gamma)t}$，块只有有限个，取最大值即可。）把 (6.7) 代入 (6.6)：
\[
\norm{\X(t)}\le K\mathrm{e}^{-\gamma t}\norm{\X(0)}+K\eps\int_0^{t}\mathrm{e}^{-\gamma(t-s)}\norm{\X(s)}\dd s .
\]
两边乘 $\mathrm{e}^{\gamma t}$，记 $\varphi(t)=\mathrm{e}^{\gamma t}\norm{\X(t)}$，得
\[
\varphi(t)\le K\norm{\X(0)}+K\eps\int_0^{t}\varphi(s)\dd s ,
\]
由格朗沃尔（\EN{Gronwall}）不等式，$\varphi(t)\le K\norm{\X(0)}\mathrm{e}^{K\eps t}$，即
\[
\norm{\X(t)}\le K\norm{\X(0)}\mathrm{e}^{-({\gamma-K\eps})t}.
\]
取 $\eps<\gamma/K$，指数仍为负，于是 $\X(t)\to\mathbf{0}$，且衰减是指数的——(a) 得证。(b) 的证明思路类似：若某 $\operatorname{Re}\lambda>0$，沿对应的（广义）特征方向取初值，线性解已经指数增长，非线性项只贡献 $o(\norm{\X})$，不足以把它拉回来（严格的论证用"不稳定流形"或切塔耶夫方法，本讲不展开）。

\subhdr{6.6 为什么必须"双曲"：三个线性部分完全相同的例子}

三个系统，线性部分都是 $A=\begin{pmatrix}0&-1\\1&0\end{pmatrix}$，即 $\lambda=\pm\mathrm{i}$、$\tau=0$、$\Delta=1$，线性化是"中心"，按定理 6.1 属于\textbf{无结论}的情形。可是它们的真实行为完全不同。取 $V=x^{2}+y^{2}=r^{2}$，沿轨线求导：
\[
\dot V=\frac{\dd}{\dd t}(x^{2}+y^{2})=2(xx'+yy').
\]

\textbf{(i)} $x'=-y-x(x^{2}+y^{2})$，$y'=x-y(x^{2}+y^{2})$。代入：
\[
\dot V=2\big[x(-y-xr^{2})+y(x-yr^{2})\big]
=2\big[-xy-x^{2}r^{2}+xy-y^{2}r^{2}\big]=-2r^{4}<0 .
\]
$r>0$ 时严格为负，所以 $r$ 单调递减、趋于 $0$：\textbf{渐近稳定}，相图是向内卷的螺旋（稳定焦点）。

\textbf{(ii)} $x'=-y+x(x^{2}+y^{2})$，$y'=x+y(x^{2}+y^{2})$。同样代入，两项的符号翻转：
\[
\dot V=2\big[-xy+x^{2}r^{2}+xy+y^{2}r^{2}\big]=+2r^{4}>0 ,
\]
$r$ 单调递增：\textbf{不稳定}，相图是向外张的螺旋（不稳定焦点）。

\textbf{(iii)} $x'=-y$，$y'=x$。此时
\[
\dot V=2(-xy+yx)=0 ,
\]
$r$ 恒定，轨线是圆：\textbf{中心}，稳定但不渐近。

三个系统的线性部分一模一样，结论却有三种。原因就写在 (6.3) 的 $o(\norm{\U})$ 里：当特征值实部为零时，一次项既不收缩也不放大，二次项的量级不再"小到可以忽略"，它\textbf{决定}了命运。这就是定理 6.1 必须要求 $\operatorname{Re}\lambda\neq0$ 的理由。

\aside{注意："线性化无结论"不等于"结论与线性化相反"，也不等于"系统一定不稳定"。单摆（§6.7）在最低点附近就是"无结论但实际是中心"的例子。}

\subhdr{6.7 例 1：单摆的两个临界点}

\textbf{例 1}\ 无阻尼单摆（\EN{undamped pendulum}，取 $g=\ell=1$）满足 $x''=-\sin x$，令 $y=x'$ 化为一阶系统
\[
x'=y,\qquad y'=-\sin x .
\]
\textbf{求临界点并判断类型。}

\textbf{解}\ 临界点满足 $y=0$、$\sin x=0$，即 $(k\pi,0)$，$k\in\mathbb{Z}$。雅可比矩阵
\[
\Jac(x,y)=\begin{pmatrix}0&1\\-\cos x&0\end{pmatrix}
\quad\Longrightarrow\quad
\Jac(k\pi,0)=\begin{pmatrix}0&1\\(-1)^{k+1}&0\end{pmatrix},
\]
特征方程 $\lambda^{2}-(-1)^{k+1}=0$，即 $\lambda^{2}=(-1)^{k+1}$。

$k$ 为奇数（如 $(\pi,0)$，摆球在最高点）：$\lambda^{2}=1$，$\lambda=\pm1$ 异号，$\Delta=-1<0$，是\textbf{鞍点}，不稳定——物理上，摆球停在最高点，任何微小扰动都会让它落下。

$k$ 为偶数（如 $(0,0)$，摆球在最低点）：$\lambda=\pm\mathrm{i}$，线性化是中心，\textbf{定理 6.1 不适用}。真实行为确实是中心（因为系统有能量守恒），但要证明它需要另找守恒量（\EN{conserved quantity}）$E=\frac{1}{2}y^{2}+1-\cos x$，不能靠线性化。这个例子说明：无结论时，要么另找方法，要么承认判断不了——不能硬套。

\subhdr{6.8 例 2：把这个二次系统算到底}

本题只问原点，但"多临界点"是这一章最容易漏的地方，所以把同一个系统算完。由 (6.1) 令右端为零：
\[
x_1(1-x_1-x_2)=0,\qquad x_2(0.5-0.75x_1-0.25x_2)=0 .
\tag{6.8}
\]
两个因子各取一个零，得四组解：
\[
(0,0),\qquad (1,0),\qquad (0,2),\qquad (0.5,\,0.5).
\]
（第四组来自 $1-x_1-x_2=0$ 与 $0.5-0.75x_1-0.25x_2=0$ 联立：把 $x_2=1-x_1$ 代入第二式得 $0.5-0.75x_1-0.25+0.25x_1=0$，即 $0.25=0.5x_1$，故 $x_1=0.5$，$x_2=0.5$。）

一般点的雅可比矩阵由 (6.2) 逐项求偏导：
\[
\Jac(x_1,x_2)=
\begin{pmatrix}
1-2x_1-x_2 & -x_1\\[4pt]
-0.75x_2 & 0.5-0.75x_1-0.5x_2
\end{pmatrix}.
\tag{6.9}
\]
逐点代入，得到下表。\textbf{表前一句}：这张表只是把 (6.9) 在四个点上各算一次特征值，再按表 3.1 归类。

\begin{minipage}{\linewidth}
\tabtop
\begin{tabular}{@{}lllll@{}}
\thead{临界点} & \thead{$\Jac$} & \thead{特征值} & \thead{类型} & \thead{稳定性}\\
\end{tabular}
\tabmid
\begin{tabular}{@{}l l l l l@{}}
$(0,0)$ & $\operatorname{diag}(1,\,0.5)$ & $1,\ 0.5$ & 结点 & 不稳定\\
$(1,0)$ & $\begin{pmatrix}-1&-1\\0&-0.25\end{pmatrix}$ & $-1,\ -0.25$ & 结点 & 渐近稳定\\
$(0,2)$ & $\begin{pmatrix}-1&0\\-1.5&-0.5\end{pmatrix}$ & $-1,\ -0.5$ & 结点 & 渐近稳定\\
$(0.5,0.5)$ & $\begin{pmatrix}-0.5&-0.5\\-0.375&-0.125\end{pmatrix}$ & $\approx0.159,\ \approx-0.784$ & 鞍点 & 不稳定\\
\end{tabular}
\tabbot
\end{minipage}

\textbf{逐行读法}：第一行 $\tau=1.5>0$、$\Delta=0.5>0$、$D=0.25>0$，落在图 \ref{c6:fig:taud} 的结点区右上方；第二、三行 $\tau<0$、$\Delta>0$、$D>0$，落在结点区左上方；第四行 $\Delta=-1/8<0$，$\tau=-5/8$，落在鞍点区（$\Delta<0$ 与 $\tau$ 的符号无关）。第四行的特征值是 $(-5\pm\sqrt{57})/16$，一正一负，正是 $\Delta<0$ 的必然结果，不必真算出来也能判鞍点。

\figblock[fig:hw2]{fig4-hw2.pdf}{0.99}{第 2 题系统的相图（本讲数值绘制）。(a) 全局：四个临界点各带类型标签，$\times$ 为鞍点；鞍点的两条分界线（\EN{separatrix}）把平面分成几个"流域"。(b) 原点放大：初值取得很小，轨线仍全部离开原点，且先沿 $x_2$ 轴（$\lambda=0.5$，增长较慢）离开，再被 $x_1$ 方向（$\lambda=1$，增长较快）拉走。}

注意 (a) 与 (b) 的对照：从\textbf{局部}（\EN{local}）看，原点是"不稳定结点"；从\textbf{全局}（\EN{global}）看，从原点附近出发的轨线并不会跑到无穷，而是被另外两个稳定结点 $(1,0)$、$(0,2)$"接住"，或者绕过鞍点 $(0.5,0.5)$ 转向。这提醒一件事：\textbf{稳定性是局部性质}，定理 6.1 只说"原点附近"，不说"整个平面"。

\subhdr{6.9 定理的适用边界（做题时要自觉写出来）}

\begin{enumerate}
\item \textbf{局部}：结论只对 $\X^{*}$ 的某个邻域成立。远处的行为可能完全不同（例 2 就是）。
\item \textbf{光滑性}：要求 $\F\in C^{1}$。多项式自动满足；含 $\abs{x}$、$\sqrt{x}$ 之类的要小心。
\item \textbf{双曲}：$\operatorname{Re}\lambda\neq0$。只要有一个特征值在虚轴上，6.1 就"不说话"（不是"说错"）。
\item \textbf{结论的内容}：给的是"稳定性"与"类型名"，不是解的具体表达式。非线性系统一般解不出来，也不需要解出来。
\end{enumerate}

%=====================================================================
\hdr{§7}{线性化失效时：李雅普诺夫第二方法（选读）}

\subhdr{7.1 想法}

§6.6 的三个例子用的都是同一个技巧：找一个函数 $V(x,y)$，它沿轨线\textbf{单调下降}，于是轨线只能往"低处"走。这个 $V$ 像能量：真实系统的能量会耗散，所以状态只能朝能量更低的方向走。李雅普诺夫把这件事抽象成一条判据，好处是\textbf{不需要解方程}。

\subhdr{7.2 定理}

\textbf{定义 7.1（正定）}\ 设 $V$ 在原点附近连续可微，$V(\mathbf{0})=0$。若对 $\X\neq\mathbf{0}$ 有 $V(\X)>0$，称 $V$ \textbf{正定}（\EN{positive definite}）；若 $V(\X)\ge0$，称\textbf{半正定}（\EN{positive semi-definite}）。

\textbf{定理 7.1（李雅普诺夫第二方法，Lyapunov's second method）}\ 设 $\F\in C^{1}$，$\F(\mathbf{0})=\mathbf{0}$，$V$ 在原点某邻域 $U$ 上连续可微、正定，记
\[
\dot V(\X)=\nabla V(\X)\cdot\F(\X)=\frac{\partial V}{\partial x}P+\frac{\partial V}{\partial y}Q .
\]
\begin{enumerate}
\item[(a)] 若 $U$ 上 $\dot V\le0$（半负定，\EN{negative semi-definite}），则零解\textbf{稳定}；
\item[(b)] 若 $U\setminus\{\mathbf{0}\}$ 上 $\dot V<0$（负定，\EN{negative definite}），则零解\textbf{渐近稳定}；
\item[(c)] 若存在区域使 $V>0$ 且 $\dot V>0$（切塔耶夫条件，\EN{Chetaev condition}），则零解\textbf{不稳定}。
\end{enumerate}

读法：(a) 说"V 不增，所以出不去"；(b) 说"V 严格减，只能往原点去"；(c) 是不稳定的判据，方向反过来用。

\subhdr{7.3 二次型（quadratic form）怎么取}

最常用的是 $V=ax^{2}+bxy+cy^{2}$。它正定的充要条件（西尔维斯特判据，\EN{Sylvester's criterion}）是
\[
a>0,\qquad 4ac-b^{2}>0 .
\tag{7.1}
\]
最省事的选择是 $V=x^{2}+y^{2}=r^{2}$（即 $a=c=1,b=0$），因为
\[
\dot V=2(xP+yQ)
\]
只要算两个乘法。若 $V=r^{2}$ 算出来 $\dot V$ 不定号，再考虑带交叉项 $bxy$ 的二次型；本讲只用到 $r^{2}$。

\subhdr{7.4 例：线性化完全失效的系统}

\textbf{例 2}\ 判断 $x'=-x^{3}$，$y'=-y^{3}$ 的零解稳定性。

\textbf{解}\ 右端在原点的一阶偏导全为零，$\Jac(\mathbf{0})=\mathbf{0}$，特征值 $\lambda_{1}=\lambda_{2}=0$，线性化给出的是"原点附近什么都不动"，完全失效。取 $V=x^{2}+y^{2}$，
\[
\dot V=2\big[x(-x^{3})+y(-y^{3})\big]=-2(x^{4}+y^{4})<0\qquad(\X\neq\mathbf{0}).
\]
由定理 7.1(b)，零解\textbf{渐近稳定}。

这道题还能分离变量验算：$x'=-x^{3}$ 得 $\frac{\dd}{\dd t}\big(\frac{1}{x^{2}}\big)=2$，故 $x^{2}(t)=\dfrac{x_0^{2}}{1+2x_0^{2}t}\to0$，衰减像 $t^{-1/2}$，不是指数——与定理 5.1 里"指数衰减"的说法对照，可以看出"渐近稳定"确实比"指数稳定"弱。

\subhdr{7.5 本题为什么不需要它}

第 2 题的 $\Jac(\mathbf{0})=\operatorname{diag}(1,\,0.5)$，特征值 $1$ 与 $0.5$ 都不为零，是\textbf{双曲}临界点，定理 6.1(b) 直接给出"不稳定"。此时再去找李雅普诺夫函数是多余的：定理 6.1 已经给出了结论，而且比第二方法更强（第二方法要猜 $V$，猜不出来不代表稳定）。

什么时候才必须用第二方法：线性化后出现 $\operatorname{Re}\lambda=0$（中心、零特征值），即 §6.6 的三种情形；或者系统不是多项式、难以线性化。这就是 §7 存在的理由。

%=====================================================================
\needspace{0.32\textheight}
\hdr{§8}{作业题面与考点}

本批作业共 \textbf{2 道大题}：第 1 题含 3 个小问（三个 $2\times2$ 线性系统），第 2 题 1 问（一个二次非线性系统）。全部对应 Lecture 22--23，截止\textbf{第十一周周二 13:00}（按第 1 讲 9 月 1 日、周一为周首推算，约为 11 月 10 日；以课程页面为准）。下面按"一题一块"给出题面与考点，\textbf{详解统一在附录 A}，题面下面标了页码。建议先把四题各写一遍再翻附录，否则你会觉得"都懂"，但手不一定写得出来。

\qbreak
\subhdr{第 1 题 1)：相异实根}
\qstem{（课件 L22--23 作业 第 1 题 1)）\ $x'=3x+2y$，\ $y'=x+4y$。\textbf{Determine the type of the critical point and study the stability.}}
\textbf{考点}：$\tau,\Delta,D$ 三个数的计算（§2.2）；"相异实根、同号"对应结点（§3.2）；稳定性的符号判据（§5.3）。
\par\textbf{详解}：附录 A.38，第 \pageref{c6:app:a1} 页。

\qbreak
\subhdr{第 1 题 2)：重根}
\qstem{（课件 L22--23 作业 第 1 题 2)）\ $x'=-x$，\ $y'=-y$。\textbf{Determine the type of the critical point and study the stability.}}
\textbf{考点}：重根与几何重数（§2.5、§3.5）；"临界结点"与"退化结点"的区别；$\tau<0$ 且 $g=2$ 时的稳定性。
\par\textbf{详解}：附录 A.39，第 \pageref{c6:app:a2} 页。

\qbreak
\subhdr{第 1 题 3)：共轭复根}
\qstem{（课件 L22--23 作业 第 1 题 3)）\ $x'=-x+2y$，\ $y'=-2x-y$。\textbf{Determine the type of the critical point and study the stability.}}
\textbf{考点}：共轭复特征值与实解的写法（§2.4）；$D<0$ 且 $\tau\neq0$ 对应焦点（§3.2）；"稳定焦点"与"中心"的分界（§4.3）。
\par\textbf{详解}：附录 A.40，第 \pageref{c6:app:a3} 页。

\qbreak
\subhdr{第 2 题：二次非线性系统}
\qstem{（课件 L22--23 作业 第 2 题）\ $x_1'=x_1-x_1^{2}-x_1x_2$，\ $x_2'=0.5x_2-0.75x_1x_2-0.25x_2^{2}$。\textbf{Investigate the stability of the trivial solution.}}
\textbf{考点}：平凡解与临界点（§4.5）；雅可比矩阵与线性化（§6.2）；线性化定理的适用条件（双曲）与结论（§6.4）；为什么本题\textbf{不需要}李雅普诺夫函数（§7.5）。
\par\textbf{详解}：附录 A.41，第 \pageref{c6:app:a4} 页。

\callout{\textbf{一个提醒。}第 1 题问的是"types of critical point(s) \textbf{and} the stability"，是两件事：先定类型（结点/鞍点/焦点/中心/临界结点），再定稳定性（稳定/渐近稳定/不稳定），并用一句话给出依据（哪条定理、哪个特征值）。只写"不稳定"或只写"结点"都算答了一半。}

%=====================================================================
\hdr{§9}{易错点：每条都能指回某一道题}

\begin{enumerate}
\item \textbf{只算 $\Delta$ 就下结论。}$\Delta>0$ 只排除鞍点；还要看 $D$ 定"实/复"、看 $\tau$ 定"进/出"（§3.6）。第 1 题 (1) 与 (3) 的 $\tau$ 都是负的，但一个是结点、一个是焦点，差别全在 $D$ 的符号上。
\item \textbf{把重根一律叫"结点"。}重根要分 $g=2$（临界结点，直线束）与 $g=1$（退化结点，切着唯一方向拐弯）。第 1 题 (2) 是 $g=2$，答案该写"临界结点"（§3.5）。
\item \textbf{把"稳定"当成"渐近稳定"。}中心是稳定的，但初值在哪条椭圆上就永远待在那条椭圆上，回不来（§4.3）。
\item \textbf{把"不稳定"说成"趋于无穷"。}不稳定只要求\textbf{存在}一条越界轨线；第 2 题原点附近的轨线其实被别的临界点接住了，但它们确实离开了原点，故仍是不稳定（§4.6、图 \ref{c6:fig:hw2}）。
\item \textbf{非线性系统直接套特征值，不验证双曲。}第 2 题可以套，是因为 $1$ 与 $0.5$ 都不在虚轴上；若线性化给出 $\pm\mathrm{i}$，必须停下来说明"无结论"（§6.6）。
\item \textbf{把 $x_1x_2$、$x_1^{2}$ 当线性项。}线性化只保留\textbf{一次}项：$x_1^{2}$ 与 $x_1x_2$ 在原点处对偏导没有贡献（$x_1=0$ 时偏导为零），所以 $\Jac(\mathbf{0})$ 是对角的（§6.2 与 A.4）。
\item \textbf{忘记说明"唯一临界点"。}线性系统要先说 $\Delta\neq0$ 故临界点只有原点，再谈类型；否则"the critical point"这个说法没有依据（§1.5）。
\item \textbf{复特征值只写三角函数，丢了 $\mathrm{e}^{\alpha t}$。}第 1 题 (3) 若丢掉 $\mathrm{e}^{-t}$，就会把稳定焦点误判成中心（§2.4）。
\end{enumerate}

%=====================================================================
\hdr{§10}{本讲义补齐了哪些跳步}

本讲课件尚未入库，因此这里不做"课件页码 + 原文短语"式的缺口表（没有课件就不能凭空引用页码）。下面列的是按标准教材讲法\textbf{最容易跳过}的 10 处，以及本讲义补在哪一节——每一条都对应上文的一个具体位置，可以逐条对照阅读。

\gapitem{形式改写（复根通解）}{多数教材直接给出 (2.5) 式，不解释它来自 $\mathrm{e}^{(\alpha+\mathrm{i}\beta)t}$ 的实部与虚部。补在 §2.4。}
\gapitem{代数跳步（重根的 $W$）}{教材给出 $(A-\lambda I)W=V$ 却不验算，学生常把右端写成 $0$。补在 §2.5，含一行验证。}
\gapitem{机制缺口（类型为什么只有五种）}{教材列完表格不说明"为什么相似变换不改变类型"。补在 §2.7。}
\gapitem{机制缺口（结点切向）}{"轨线与慢方向相切"通常只给结论。补在 §3.2，用提公因子的极限说明。}
\gapitem{压缩缺口（判别表）}{表 3.1 常作为结论直接背。本讲义在表后逐行给出 $\lambda=(\tau\pm\sqrt D)/2$ 的代入过程（§3.1）。}
\gapitem{定义缺口（稳定性的三个词）}{"稳定"与"渐近稳定"常被混用。补在 §4.2--§4.3，含 $\eps$–$\delta$ 定义与"中心稳定但不渐近"这个反例。}
\gapitem{证明缺口（稳定性定理）}{定理 5.1 的 (iii)(iv) 两条几乎总被略去。补在 §5.2，并用两个 $2\times2$ 反例收尾（§5.4）。}
\gapitem{机制缺口（线性化为什么合法）}{线性化定理常被当作"显然"。补在 §6.5，给出 (6.6)--(6.7) 的估计链与格朗沃尔不等式。}
\gapitem{机制缺口（双曲条件不可省）}{三个"线性部分相同、结论不同"的例子补在 §6.6，每一个都算出 $\dot V$。}
\gapitem{适用区间缺口（局部 vs 全局）}{"稳定性是局部性质"常被忽略。补在 §6.8 与 §6.9，用第 2 题系统的全局相图（图 \ref{c6:fig:hw2}）说明。}

%=====================================================================
\hdr{§11}{三句话总结}

\begin{enumerate}
\item \textbf{类型}由线性部分 $A$（或 $\Jac(\X^{*})$）的特征值决定，读法是三步：$\Delta$ 的符号定"鞍点/非鞍点"，$D=\tau^{2}-4\Delta$ 定"实根/复根"，$\tau$ 的符号定"进/出"。
\item \textbf{稳定性}只看实部：全部 $<0$ 渐近稳定；有一个 $>0$ 不稳定；全部 $\le0$ 且零实部无重根则稳定但不渐近（中心）。
\item \textbf{非线性}：在双曲临界点（实部全不为零）处，线性化的结论可以原样搬回；只要有一个特征值落在虚轴上，线性化就失效，必须另找方法（守恒量、李雅普诺夫函数，或如实说明判断不了）。
\end{enumerate}

%=====================================================================
\hdr{§12}{四步用法与速查}

\subhdr{12.1 拿到一道题的四个动作}

\begin{enumerate}
\item \textbf{找临界点}：解 $\F(\X)=\mathbf{0}$。线性系统先看 $\Delta$：$\Delta\neq0$ 则只有原点；非线性系统把解全部列出（可能不止一个，见 §6.8）。
\item \textbf{写线性部分}：线性系统直接取 $A$；非线性系统算 $\Jac(\X^{*})$（(6.2) 或 (6.9) 的形式）。
\item \textbf{算三个数}：$\tau=\tr A$，$\Delta=\det A$，$D=\tau^{2}-4\Delta$；再解特征值（若需要写通解）。
\item \textbf{写结论}：类型（表 3.1）+ 稳定性（$\tau$ 的符号）+ 依据（定理号与条件，如"$\Jac$ 特征值 $1,0.5$ 双曲，由定理 6.1(b)"）。
\end{enumerate}

\subhdr{12.2 一句话模板}

\begin{itemize}
\item 线性题：\ "$\Delta=\cdots\neq0$，故唯一临界点为原点；$\tau=\cdots$，$\Delta=\cdots$，$D=\cdots$，故为\underline{\hspace{2em}}（类型）；又 $\tau\ \square\ 0$，故\underline{\hspace{2em}}（稳定性）。"
\item 非线性题：\ "$\X=\mathbf{0}$ 是临界点（代入右端为零）；$\Jac(\mathbf{0})=\cdots$，特征值 $\cdots$，均非零（双曲）；由定理 6.1\underline{\hspace{0.6em}}，平凡解\underline{\hspace{2em}}。"
\end{itemize}

\subhdr{12.3 自检清单（交卷前逐条打勾）}

\begin{itemize}
\item 临界点是否\underline{全部}找到？是否说明了"唯一"或"共有几个"？
\item $\tau,\Delta,D$ 是否重算过一遍？特征值是否为特征方程的根（代入验算）？
\item 类型名与 $(\tau,\Delta,D)$ 是否对得上表 3.1？重根是否查了 $g$？
\item 稳定性是否与类型一致（鞍点必不稳定、中心必非渐近）？
\item 非线性题是否写了"双曲"或"特征值非零"这句前提？是否引了定理号？
\item 答案是否同时包含"类型"与"稳定性"两项（题目问了两件事）？
\end{itemize}

\subhdr{12.4 页码与章节对照}

\begin{minipage}{\linewidth}
\tabtop
\begin{tabular}{@{}lll@{}}
\thead{本讲义} & \thead{内容} & \thead{对应}\\
\end{tabular}
\tabmid
\begin{tabular}{@{}l l p{0.42\linewidth}@{}}
§1--§2 & 相平面、临界点、线性系统三种解 & 第 1 题全部；Constanda §7.3--§7.4\\
§3 & 五种类型与判别表、$\tau$–$\Delta$ 图 & 第 1 题 (1)(2)(3)\\
§4--§5 & 稳定性定义与线性判据 & 第 1 题全部、第 2 题的结论表述\\
§6 & 线性化、双曲、定理 6.1/6.2 & 第 2 题\\
§7 & 李雅普诺夫第二方法 & 第 2 题的"备用工具"（本题不用）\\
§13 & 名词中英对照（考前查词） & 全章；题干原句的译法\\
附录 A & 四道题的完整详解 & 第 1 题 (1)(2)(3)、第 2 题\\
附录 B & 题面歧义与替代读法 & 第 2 题\\
\end{tabular}
\tabbot
\end{minipage}

%=====================================================================
\hdr{§13}{名词中英对照（考前查用）}

\textbf{这一节解决什么问题。}课件与考卷都是英文，讲义是中文；考场上真正卡住的常常不是方法，而是题干里那个词。下面按用途分成五组，只做对照，不重复推导。正文里每个词在首次出现处也已用灰色斜体注过一次（例如“临界点（\EN{critical point}）”），两处可以互相印证。全书级的简表在 §0.5，那里每章只收一两个词；本节把第 6 章的名词收全。

\subhdr{13.1 相图与点型}

\begin{minipage}{\linewidth}
{\small
\tabtop
\begin{tabular}{@{}>{\raggedright\arraybackslash}p{0.19\linewidth}>{\raggedright\arraybackslash}p{0.27\linewidth}>{\raggedright\arraybackslash}p{0.19\linewidth}>{\raggedright\arraybackslash}p{0.27\linewidth}@{}}
\thead{中文} & \thead{英文} & \thead{中文} & \thead{英文}\\
\end{tabular}
\tabmid
\begin{tabular}{@{}>{\raggedright\arraybackslash}p{0.19\linewidth}>{\raggedright\arraybackslash}p{0.27\linewidth}>{\raggedright\arraybackslash}p{0.19\linewidth}>{\raggedright\arraybackslash}p{0.27\linewidth}@{}}
相图 & \EN{phase portrait} & 轨线 & \EN{trajectory (orbit)}\\
相平面 & \EN{phase plane} & 速度场 & \EN{velocity field}\\
方向场 & \EN{direction field} & 自治系统 & \EN{autonomous system}\\
临界点 & \EN{critical point} & 平衡点 & \EN{equilibrium point}\\
奇点 & \EN{singular point} & 驻点 & \EN{stationary point}\\
特征方向 & \EN{eigendirection} & 结点 & \EN{node}\\
鞍点 & \EN{saddle point} & 焦点 & \EN{focus (spiral point)}\\
中心 & \EN{center} & 临界结点 & \EN{proper node (star node)}\\
退化结点 & \EN{improper node (degenerate node)} & 分界线 & \EN{separatrix}\\
源 & \EN{source} & 汇 & \EN{sink}\\
\end{tabular}
\tabbot
}
\end{minipage}

\subhdr{13.2 稳定性、线性化与方法}

\begin{minipage}{\linewidth}
{\small
\tabtop
\begin{tabular}{@{}>{\raggedright\arraybackslash}p{0.19\linewidth}>{\raggedright\arraybackslash}p{0.27\linewidth}>{\raggedright\arraybackslash}p{0.19\linewidth}>{\raggedright\arraybackslash}p{0.27\linewidth}@{}}
\thead{中文} & \thead{英文} & \thead{中文} & \thead{英文}\\
\end{tabular}
\tabmid
\begin{tabular}{@{}>{\raggedright\arraybackslash}p{0.19\linewidth}>{\raggedright\arraybackslash}p{0.27\linewidth}>{\raggedright\arraybackslash}p{0.19\linewidth}>{\raggedright\arraybackslash}p{0.27\linewidth}@{}}
稳定 & \EN{stable} & 渐近稳定 & \EN{asymptotically stable}\\
不稳定 & \EN{unstable} & 平凡解 & \EN{trivial solution}\\
指数衰减 & \EN{exponential decay} & 有界 & \EN{bounded}\\
雅可比矩阵 & \EN{Jacobian matrix} & 线性化 & \EN{linearization}\\
线性化系统 & \EN{linearized system} & 线性部分 & \EN{linear part}\\
第一近似 & \EN{first approximation} & 第一近似系统 & \EN{system in the first approximation}\\
双曲临界点 & \EN{hyperbolic critical point} & 虚轴 & \EN{imaginary axis}\\
拓扑等价 & \EN{topologically equivalent} & 邻域 & \EN{neighborhood}\\
李雅普诺夫第一方法 & \EN{Lyapunov's first method} & 李雅普诺夫第二方法 & \EN{Lyapunov's second method}\\
正定 & \EN{positive definite} & 负定 & \EN{negative definite}\\
半正定 & \EN{positive semi-definite} & 半负定 & \EN{negative semi-definite}\\
二次型 & \EN{quadratic form} & 西尔维斯特判据 & \EN{Sylvester's criterion}\\
守恒量 & \EN{conserved quantity} & 格朗沃尔不等式 & \EN{Gronwall's inequality}\\
局部 & \EN{local} & 全局 & \EN{global}\\
单摆 & \EN{pendulum} & 逻辑斯蒂方程 & \EN{logistic equation}\\
未扰解 & \EN{unperturbed solution} & 受扰解 & \EN{perturbed solution}\\
\end{tabular}
\tabbot
}
\end{minipage}

{\small\color{inkgray}最后一行的\EN{unperturbed solution}与\EN{perturbed solution}是课件定义稳定性时的用语（要研究的那个解叫未扰解，从附近初值出发的解叫受扰解）；本书 §4.2 直接用 $\X^{*}$ 与 $\X(t)$ 写，没有引入这两个词，认词即可。\par}

\subhdr{13.3 矩阵与判据}

\begin{minipage}{\linewidth}
{\small
\tabtop
\begin{tabular}{@{}>{\raggedright\arraybackslash}p{0.19\linewidth}>{\raggedright\arraybackslash}p{0.27\linewidth}>{\raggedright\arraybackslash}p{0.19\linewidth}>{\raggedright\arraybackslash}p{0.27\linewidth}@{}}
\thead{中文} & \thead{英文} & \thead{中文} & \thead{英文}\\
\end{tabular}
\tabmid
\begin{tabular}{@{}>{\raggedright\arraybackslash}p{0.19\linewidth}>{\raggedright\arraybackslash}p{0.27\linewidth}>{\raggedright\arraybackslash}p{0.19\linewidth}>{\raggedright\arraybackslash}p{0.27\linewidth}@{}}
特征值 & \EN{eigenvalue} & 特征向量 & \EN{eigenvector}\\
特征方程 & \EN{characteristic equation} & 迹 & \EN{trace}\\
行列式 & \EN{determinant} & 判别式 & \EN{discriminant}\\
几何重数 & \EN{geometric multiplicity} & 代数重数 & \EN{algebraic multiplicity}\\
若尔当标准形 & \EN{Jordan canonical form} & 若尔当块 & \EN{Jordan block}\\
通解 & \EN{general solution} & 线性无关 & \EN{linearly independent}\\
可逆线性变换 & \EN{invertible linear transformation} & 二重根 & \EN{double root}\\
初值 & \EN{initial value} & 唯一性 & \EN{uniqueness of solution}\\
\end{tabular}
\tabbot
}
\end{minipage}

\needspace{4\baselineskip}
\subhdr{13.4 六组容易混的词}

\begin{itemize}
\item \EN{stable}（稳定）与 \EN{asymptotically stable}（渐近稳定）：前者只要求不跑远，后者还要求回来（§4.3）。
\item \EN{proper node}（临界结点）与 \EN{improper node}（退化结点）：前者各方向都是特征方向（$g=2$），后者只有一个特征方向（$g=1$）。作业第 1 题 (2) 是前者。
\item \EN{node}（结点）与 \EN{saddle point}（鞍点）：两者靠 $\Delta$ 的符号分家，与 $\tau$ 的符号无关；$\tau$ 只管进出（§3.6）。
\item \EN{source}（源）与 \EN{sink}（汇）：课件讲结点时按 $\tau$ 的符号取名——$\tau>0$ 叫 source，$\tau<0$ 叫 sink；本书正文写作“不稳定结点”“稳定结点”，同义。
\item \EN{hyperbolic}（双曲）与 \EN{hyperbola}（双曲线）：前者是特征值实部不为零的条件，与双曲线没有关系（§6.3）。
\item \EN{local}（局部）与 \EN{global}（全局）：定理 6.1 只保证临界点附近的事，远处的行为可以完全不同（§6.8、§6.9）。
\end{itemize}

\subhdr{13.5 题干原句}

\qstem{（课件 L22--23 作业 第 1 题）\textbf{Determine the type of the critical point and study the stability.}}
译：判定临界点的类型，并研究其稳定性。\textbf{两件事都要写}：先给类型（结点/鞍点/焦点/中心/临界结点/退化结点），再给稳定性（稳定/渐近稳定/不稳定），最后补一句依据（哪条定理、哪个特征值）。第 1 题的三个小问都是这一句；课件把它写成 \EN{critical point(s)}，括号里的复数只是通用写法，$\Delta\neq0$ 时线性系统只有一个临界点（§1.5）。

\qstem{（课件 L22--23 作业 第 2 题）\textbf{Investigate the stability of the trivial solution.}}
译：研究平凡解的稳定性。关键词是 \EN{trivial solution}（平凡解，即 $\X\equiv\mathbf{0}$），不是一般解；第 2 题问的正是原点。

本讲义与考卷的英文标题：本章为 \EN{Critical Points and Stability}，两讲课件分别是 \EN{Phase Portrait}（相图）与 \EN{Linearization}（线性化）；课件里“第一近似下的稳定性”写作 \EN{stability in the first approximation}，与书中"线性化定理"是同一件事（§6.4）。

\flowerrule

%=====================================================================
\needspace{0.30\textheight}
\clearpage
\settags{附录 A\quad 全部题目详解}{微分方程 · 期中之后}
\hdr{附录 A}{全部题目详解（A.1--41）}

\callout{本附录是六章正文所引“详解”的汇总：每一小节先复述题干与出处，再给完整推导。小节号按章顺序连续编号；正文里“附录 A.$n$，第 $m$ 页”的页码由 \LaTeX{} 自动解析，重排版后仍然准确。}
\subhdr{第 1 章（L13）· 各题详解}

每块自成一页（块前用 \verb|\qbreak| 保证整块不跨页），按“题干 → 考点 → 关键一步 → 详解 → 答案 → 自检”排。题干里的条件是完整复述，不回头也能做。

\qpagebreak
\subhdr{A.1\quad 例 1：验证通解与常数的等价写法}
\label{c1:app:a1}

\qbreak
\qstem{（课件 p10 例 1）$\dfrac{\mathrm{d}y_1}{\mathrm{d}x}=y_2$，$\dfrac{\mathrm{d}y_2}{\mathrm{d}x}=-y_1$。验证 $y_1=c_1\cos(x-c_2)$、$y_2=-c_1\sin(x-c_2)$ 是解，并说明常数可以换写法。}

\textbf{考点}\ 解的定义（代回恒成立）；通解里常数的等价性。

\textbf{关键一步}\ 分别求导，代回右端；再用和角公式把 $c_1\cos(x-c_2)$ 展开成 $A\cos x+B\sin x$。

\textbf{详解}\ 先验证第一式：$y_1'=-c_1\sin(x-c_2)$，而右端 $y_2=-c_1\sin(x-c_2)$，两者相同。再验证第二式：$y_2'=-c_1\cos(x-c_2)$，右端 $-y_1=-c_1\cos(x-c_2)$，也相同。所以这组函数是解。

再看常数：由和角公式

\[
c_1\cos(x-c_2)=c_1\cos c_2\cos x+c_1\sin c_2\sin x=A\cos x+B\sin x,
\]

其中 $A=c_1\cos c_2$，$B=c_1\sin c_2$。反过来，给定 $A,B$（不同时为零），取 $c_1=\sqrt{A^2+B^2}$、$c_2=\arctan(B/A)$ 即可还原。两种写法覆盖的解集合相同，都是两个任意常数。

\textbf{但两者的“常数”不是同一组数}：$A,B$ 就是 $x=0$ 处的初值 $y_1(0),y_2(0)$；而 $c_1,c_2$ 是振幅与相位，$c_1=\sqrt{y_1(x_0)^2+y_2(x_0)^2}$，$c_2=x_0-\arctan\bigl(-y_2(x_0)/y_1(x_0)\bigr)$。答卷上写哪一种都行，但不能两种混着用。

\textbf{答案}\ 两式代回均恒成立；$c_1\cos(x-c_2)$ 与 $A\cos x+B\sin x$ 等价，常数个数都是 2。

\textbf{自检}\ 取 $c_1=1,c_2=0$：$y_1=\cos x$，$y_2=-\sin x$。$y_1'=-\sin x=y_2$ \ding{51}；$y_2'=-\cos x=-y_1$ \ding{51}。

\qpagebreak
\subhdr{A.2\quad 例 2：通解、定义域与奇解}
\label{c1:app:a2}

\qbreak
\qstem{（课件 p10 例 2）$\dfrac{\mathrm{d}y_1}{\mathrm{d}x}=2\sqrt{y_1}$，$\dfrac{\mathrm{d}y_2}{\mathrm{d}x}=y_1$。验证通解 $y_1=(x+c_1)^2$、$y_2=\tfrac13(x+c_1)^3+c_2$，以及奇解 $y_1=0,\ y_2=c$。}

\textbf{考点}\ 非线性方程的定义域限制；“奇解”为什么确实不能由通解取常数得到。

\textbf{关键一步}\ 第一式只含 $y_1$，先解它；检查 $\sqrt{(x+c_1)^2}=|x+c_1|$ 与 $x+c_1$ 的符号；最后看 $y_1\equiv0$ 是否可能。

\textbf{详解}\ 第一式分离变量：$\mathrm{d}y_1/(2\sqrt{y_1})=\mathrm{d}x$，积分得 $\sqrt{y_1}=x+c_1$（取正支），平方得 $y_1=(x+c_1)^2$。代回验算：$y_1'=2(x+c_1)$，而 $2\sqrt{y_1}=2|x+c_1|$，两者相等\textbf{仅当} $x+c_1\ge0$。所以这条解只在 $[-c_1,+\infty)$ 上成立——这是非线性方程特有的。

第二式：$y_2'=y_1=(x+c_1)^2$，积分得 $y_2=\tfrac13(x+c_1)^3+c_2$。验算 $y_2'=(x+c_1)^2=y_1$ \ding{51}。

再看奇解：$y_1\equiv0$ 时第一式两边都是 $0$，成立；$y_2\equiv c$ 时第二式也是 $0=0$，成立。它能否由通解取常数得到？要 $y_1=(x+c_1)^2\equiv0$ 对区间内所有 $x$ 成立，必须 $x+c_1=0$ 对所有 $x$ 成立——不可能。所以它是通解之外的新解，即奇解。

\textbf{答案}\ 通解在 $x\ge-c_1$ 上成立；奇解为 $y_1\equiv0,\ y_2\equiv c$。

\textbf{自检}\ 把 $(y_1,y_2)=(0,5)$ 代入：$y_1'=0=2\sqrt0$ \ding{51}；$y_2'=0=y_1$ \ding{51}。

\qpagebreak
\subhdr{A.3\quad 例 3：写成矩阵}
\label{c1:app:a3}

\qbreak
\qstem{（课件 p14 例 3）$x'=y-z$，$y'=x+y$，$z'=x+z$，未知函数为 $x(t),y(t),z(t)$。写成矩阵形式并判断是否齐次。}

\textbf{考点}\ 逐行读系数；齐次判断。

\textbf{关键一步}\ 按 $x,y,z$ 的顺序，把每个方程里三个系数依次抄成一行；没出现的变量系数写 $0$。

\textbf{详解}\ 第一式 $x'=0\cdot x+1\cdot y+(-1)\cdot z$，系数行 $(0,1,-1)$；第二式 $y'=1\cdot x+1\cdot y+0\cdot z$，行 $(1,1,0)$；第三式 $z'=1\cdot x+0\cdot y+1\cdot z$，行 $(1,0,1)$。于是

\[
\mat{x'\\ y'\\ z'}=\mat{0&1&-1\\ 1&1&0\\ 1&0&1}\mat{x\\ y\\ z}.
\]

右端没有只依赖 $t$ 的项，故\textbf{齐次}。

\textbf{答案}\ $A=\mat{0&1&-1\\1&1&0\\1&0&1}$，齐次线性系统。

\textbf{自检}\ 取 $(x,y,z)=(1,0,0)$：右端 $(0,1,1)$，矩阵乘同样得 $(0,1,1)$ \ding{51}。

\qpagebreak
\subhdr{A.4\quad 例 4：验证两解与叠加}
\label{c1:app:a4}

\qbreak
\qstem{（课件 p18 例 4）$\vect{Y}'=\mat{1&1\\4&1}\vect{Y}$，$\vect{Y}_1=\mat{1\\2}\mathrm{e}^{3x}$，$\vect{Y}_2=\mat{1\\-2}\mathrm{e}^{-x}$。验证两者是解，从而 $c_1\vect{Y}_1+c_2\vect{Y}_2$ 是解。}

\textbf{考点}\ 解的定义；叠加原理。

\textbf{关键一步}\ 对 $\vect{Y}_k=\vect{V}_k\mathrm{e}^{\lambda_k x}$，只需验证 $A\vect{V}_k=\lambda_k\vect{V}_k$——指数因子两边可约掉。

\textbf{详解}\ 第一式：

\[
A\vect{Y}_1=\mat{1&1\\4&1}\mat{1\\2}\mathrm{e}^{3x}=\mat{3\\6}\mathrm{e}^{3x}=3\mat{1\\2}\mathrm{e}^{3x}=3\vect{Y}_1 ,
\]

而 $\vect{Y}_1'=3\vect{Y}_1$，故 $A\vect{Y}_1=\vect{Y}_1'$ \ding{51}。第二式：

\[
A\vect{Y}_2=\mat{1&1\\4&1}\mat{1\\-2}\mathrm{e}^{-x}=\mat{-1\\2}\mathrm{e}^{-x}=(-1)\mat{1\\-2}\mathrm{e}^{-x}=\vect{Y}_2' .
\]

两式都是齐次方程的解，由性质 2.1，任意线性组合仍是解。

\textbf{答案}\ 两者都是解；$c_1\vect{Y}_1+c_2\vect{Y}_2$ 是通解（两解线性无关：Wronskian $=\det\mat{\mathrm{e}^{3x}&\mathrm{e}^{-x}\\2\mathrm{e}^{3x}&-2\mathrm{e}^{-x}}=-4\mathrm{e}^{2x}\neq0$）。

\textbf{自检}\ 取 $c_1=c_2=1$：$\vect{Y}=\vect{Y}_1+\vect{Y}_2$。在 $x=0$ 处 $\vect{Y}(0)=(1,2)+(1,-2)=(2,0)$；求导得 $\vect{Y}'(0)=3(1,2)-1(1,-2)=(2,8)$；而右端

\[
A\vect{Y}(0)=\mat{1&1\\4&1}\mat{2\\0}=\mat{2\\8},
\]

两边相同 \ding{51}。这也说明：$A\vect{Y}_k=\lambda_k\vect{Y}_k$ 成立时，组合的每一项各按自己的 $\lambda$ 走，互不干扰。

\qpagebreak
\subhdr{A.5\quad 例 5：Wronskian 与通解}
\label{c1:app:a5}

\qbreak
\qstem{（课件 p23 例 5）$\mathrm{d}y_1/\mathrm{d}x=y_2$，$\mathrm{d}y_2/\mathrm{d}x=-y_1$；$\vect{Y}_1=(\cos x,-\sin x)^{\mathsf T}$，$\vect{Y}_2=(\sin x,\cos x)^{\mathsf T}$。算 $W$ 并写通解。}

\textbf{考点}\ Wronskian 的计算；由基本解组写通解。

\textbf{关键一步}\ 按列排：第一列 $\vect{Y}_1$，第二列 $\vect{Y}_2$，算 $2\times2$ 行列式。

\textbf{详解}\

\[
W(x)=\det\mat{\cos x & \sin x\\ -\sin x & \cos x}
=\cos^2x-(-\sin x)\sin x=\cos^2x+\sin^2x=1\neq0 .
\]

两个解在每一点都不为零，故线性无关，构成基本解组；由定理 2.5，

\[
\vect{Y}=C_1\vect{Y}_1+C_2\vect{Y}_2
\ \Longrightarrow\
y_1=C_1\cos x+C_2\sin x,\quad y_2=-C_1\sin x+C_2\cos x .
\]

\textbf{答案}\ $W\equiv1$；通解如上。

\textbf{自检}\ 取 $C_1=1,C_2=0$：$y_1=\cos x,\ y_2=-\sin x$，与 A.1 的验证一致 \ding{51}。初值 $y_1(0)=C_1,\ y_2(0)=C_2$，说明常数就是初值，任意初值都能取到。

\qpagebreak
\subhdr{A.6\quad 例 6：一般线性方程的五步变换}
\label{c1:app:a6}

\qbreak
\qstem{（课件 p26--p29 例 6）$b_0(t)y^{(n)}+b_1(t)y^{(n-1)}+\cdots+b_{n-1}(t)y'+b_n(t)y=g(t)$，$b_0\neq0$，初值 $y(t_0)=c_0,\dots,y^{(n-1)}(t_0)=c_{n-1}$。化为 $\vect{Y}'=A(t)\vect{Y}+\vect{F}(t)$。}

\textbf{考点}\ 归一化；变量替换；系数矩阵的构造（尤其最后一行的倒序）。

\textbf{关键一步}\ 先除 $b_0$ 把最高阶导数解出来，再设 $y_k=y^{(k-1)}$；最后一行按“$y_1$ 乘 $a_n$”的顺序排。

\textbf{详解}\ 第一步归一化：$y^{(n)}=a_1y^{(n-1)}+\cdots+a_{n-1}y'+a_ny+f$，其中 $a_j=-b_j/b_0$，$f=g/b_0$。

第二步设 $y_1=y,\ y_2=y',\ \dots,\ y_n=y^{(n-1)}$，于是 $y_k'=y_{k+1}$（$k=1,\dots,n-1$）。

第三步写最后一式：$y_n'=y^{(n)}=a_1y_n+a_2y_{n-1}+\cdots+a_ny_1+f$，按 $y_1,\dots,y_n$ 的顺序重排：

\[
y_n'=a_ny_1+a_{n-1}y_2+\cdots+a_1y_n+f .
\]

第四步拼矩阵：前 $n-1$ 行是移位，最后一行是上面的系数：

\[
A=\begin{pmatrix}
0&1&0&\cdots&0\\
0&0&1&\cdots&0\\
\vdots& & &\ddots&\vdots\\
0&0&0&\cdots&1\\
a_n&a_{n-1}&a_{n-2}&\cdots&a_1
\end{pmatrix},\quad
\vect{F}=\mat{0\\ \vdots\\ 0\\ f},\quad
\vect{Y}(t_0)=\mat{c_0\\ \vdots\\ c_{n-1}}.
\]

\textbf{答案}\ 如上；关键是最后一行按 $y_1,\dots,y_n$ 记为 $(a_n,\dots,a_1)$。

\textbf{自检（$n=3$）}\ 取 $y'''=y''+2y'+3y$，即 $a_1=1,a_2=2,a_3=3$，最后一行应为 $(3,2,1)$：

\[
A=\mat{0&1&0\\0&0&1\\3&2&1}.
\]

验证：第一行给 $y_1'=y_2$ \ding{51}；第二行给 $y_2'=y_3$ \ding{51}；第三行给 $y_3'=3y_1+2y_2+y_3=3y+2y'+y''$ \ding{51}。

\qpagebreak
\subhdr{A.7\quad 例 7：含 $t$ 的二阶方程}
\label{c1:app:a7}

\qbreak
\qstem{（课件 p30--p31 例 7）$t^4\dfrac{\mathrm{d}^2x}{\mathrm{d}t^2}+\sin t\dfrac{\mathrm{d}x}{\mathrm{d}t}-4x=\ln t$，$t>0$。化为一阶线性组并写成矩阵形式。}

\textbf{考点}\ 归一化后设变量；变系数矩阵 $A(t)$。

\textbf{关键一步}\ 先除 $t^4$：$x''=\dfrac{4}{t^4}x-\dfrac{\sin t}{t^4}x'+\dfrac{\ln t}{t^4}$；再令 $y=x'$。

\textbf{详解}\ 设 $y=x'$，则 $y'=x''$，于是

\[
\begin{cases}
x'=y,\\[3pt]
y'=\dfrac{4}{t^4}x-\dfrac{\sin t}{t^4}y+\dfrac{\ln t}{t^4}.
\end{cases}
\]

写成矩阵：

\[
\mat{x'\\ y'}=\mat{0&1\\[2pt] \dfrac{4}{t^4}&-\dfrac{\sin t}{t^4}}\mat{x\\ y}+\mat{0\\[2pt]\dfrac{\ln t}{t^4}}.
\]

右端有只依赖 $t$ 的项，故\textbf{非齐次}。

\textbf{答案}\ 如上；$A(t)$ 的两行分别来自 $x'=y$ 与归一化后的 $x''$。

\textbf{自检}\ 取 $t=1$：$x''=4x-(\sin1)x'+\ln1=4x-0.8415x'$，与矩阵在 $t=1$ 的取值一致 \ding{51}。

\qpagebreak
\subhdr{A.8\quad 作业：把 $y'''=y''-x^2(y')^2$ 化为一阶组}
\label{c1:app:a8}

\qbreak
\qstem{（课件 p37 作业 I）把方程 $y'''=y''-x^2(y')^2$ 化为等价的一阶方程组。课件原注：It's not a linear system。}

\textbf{考点}\ 一般变换在非线性方程上的用法；为什么不能写成 $A\vect{Y}$。

\textbf{关键一步}\ 设 $y_1=y,\ y_2=y',\ y_3=y''$；前两式照抄，第三式把 $y'''$ 用原方程替换。

\textbf{详解}\ 由定义：$y_1'=y_2$，$y_2'=y_3$。第三式：$y_3'=y'''$，而原方程给 $y'''=y''-x^2(y')^2$，用新变量代掉 $y''$ 与 $y'$：

\[
y_3'=y_3-x^2y_2^2 .
\]

于是

\[
\begin{cases}
y_1'=y_2,\\[3pt]
y_2'=y_3,\\[3pt]
y_3'=y_3-x^2y_2^2 .
\end{cases}
\]

\textbf{为什么拼不成 $A(x)\vect{Y}$？}右端对 $y_2$ 是二次的（$y_2^2$），不是线性组合，所以只能写成 $\vect{Y}'=\vect{F}(x,\vect{Y})$，其中

\[
\vect{F}(x,\vect{Y})=\bigl(y_2,\ y_3,\ y_3-x^2y_2^2\bigr)^{\mathsf T}.
\]

第一、二行仍是“移位”形式，与 §6.4 的模板一样；非线性的部分只落在第三行。

\textbf{答案}\ 如上三式；是非线性系统，不能写成矩阵乘向量。

\textbf{自检}\ 在 $x=0$ 处非线性项消失，系统退化为 $y'''=y''$，其解为 $y=c_1+c_2x+c_3\mathrm{e}^{x}$。代进系统：$y_1=c_1+c_2x+c_3\mathrm{e}^x$，$y_2=c_2+c_3\mathrm{e}^x$，$y_3=c_3\mathrm{e}^x$。三式验证：$y_1'=c_2+c_3\mathrm{e}^x=y_2$ \ding{51}；$y_2'=c_3\mathrm{e}^x=y_3$ \ding{51}；$y_3'=c_3\mathrm{e}^x=y_3-x^2y_2^2\big|_{x=0}$ \ding{51}。

\flowerrule

{\small\color{inkgray}
依据课件：\texttt{13\_DE\_systems.pdf}（37 页，2026-09-24，W.Y., DLUT）。页码为该 PDF 页序。
缺口 10 条见 §10；示意图两张为本文重绘（弹簧-质量、并联 LRC），元件与端口按课件 Figure 1、Figure 2 的约定标注。}
\subhdr{第 2 章（L14--15）· 例题与作业详解}

每块自成一页，按“题干 → 考点 → 关键一步 → 详解 → 答案 → 自检”排。题号后标课件页码。

\qbreak
\subhdr{A.9\quad 例 1：三个相异实根（p11--p15）}
\label{c2:app:ex1}
\qbreak
\qstem{（课件 p11 例 1）解 $\vect{Y}'=A\vect{Y}$，
$A=\mat{1&-1&1\\1&1&-1\\2&-1&0}$。}
\textbf{考点}\ Case 1：特征方程、特征向量、通解。
\textbf{关键一步}\ 特征多项式是三次的，试根 $\lambda=1$ 命中，再降为二次。
\textbf{详解}\ 特征方程
\[
\det(A-\lambda E)=\begin{vmatrix}1-\lambda&-1&1\\1&1-\lambda&-1\\2&-1&-\lambda\end{vmatrix}
=(2-\lambda)(1-\lambda)(\lambda+1)=0,
\]
故 $\lambda_1=2$，$\lambda_2=1$，$\lambda_3=-1$。
对 $\lambda_1=2$：$A-2E=\mat{-1&-1&1\\1&-1&-1\\2&-1&-2}$，行化简得 $(1,0,1)^{\mathsf T}$ 方向，取 $\vect{V}_1=(1,0,1)^{\mathsf T}$。检验 $A\vect{V}_1=(2,0,2)^{\mathsf T}=2\vect{V}_1$。
对 $\lambda_2=1$：$A-E=\mat{0&-1&1\\1&0&-1\\2&-1&-1}$，解得 $\vect{V}_2=(1,1,1)^{\mathsf T}$；$A\vect{V}_2=(1,1,1)^{\mathsf T}$。
对 $\lambda_3=-1$：$A+E=\mat{2&-1&1\\1&2&-1\\2&-1&1}$，解得 $\vect{V}_3=(-1,3,5)^{\mathsf T}$；$A\vect{V}_3=(1,-3,-5)^{\mathsf T}=-\vect{V}_3$。
\textbf{答案}\ $\vect{Y}=c_1\mat{1\\0\\1}\mathrm{e}^{2x}+c_2\mat{1\\1\\1}\mathrm{e}^{x}+c_3\mat{-1\\3\\5}\mathrm{e}^{-x}$，即三个分量 $y_1=c_1\mathrm{e}^{2x}+c_2\mathrm{e}^{x}-c_3\mathrm{e}^{-x}$、$y_2=c_2\mathrm{e}^{x}+3c_3\mathrm{e}^{-x}$、$y_3=c_1\mathrm{e}^{2x}+c_2\mathrm{e}^{x}+5c_3\mathrm{e}^{-x}$。
\textbf{自检}\ 取 $c_1=1$ 其余为零：$\vect{Y}=(1,0,1)^{\mathsf T}\mathrm{e}^{2x}$，$\vect{Y}'=(2,0,2)^{\mathsf T}\mathrm{e}^{2x}=A\vect{Y}$ \ding{51}。

\qpagebreak
\subhdr{A.10\quad 例 2：一对共轭复根（p17--p19）}
\label{c2:app:ex2}
\qbreak
\qstem{（课件 p17 例 2）求 $\vect{Y}'=A\vect{Y}$ 的实值基本解组，$A=\mat{-1/2&1\\-1&-1/2}$。}
\textbf{考点}\ Case 2：复特征值、实值化、Wronskian。
\textbf{关键一步}\ 特征方程是 $(\lambda+\tfrac12)^2+1=0$；取 $\lambda=-\tfrac12+\mathrm{i}$，特征向量 $(1,\mathrm{i})^{\mathsf T}$。
\textbf{详解}\ $\det(A-\lambda E)=(-\tfrac12-\lambda)^2+1=0$，故 $\lambda=-\tfrac12\pm\mathrm{i}$，即 $a=-\tfrac12$、$b=1$。
对 $\lambda=-\tfrac12+\mathrm{i}$，$(A-\lambda E)\vect{V}=0$ 化为 $-\mathrm{i}v_1+v_2=0$，取 $\vect{V}=(1,\mathrm{i})^{\mathsf T}=\vect{P}+\mathrm{i}\vect{Q}$，其中 $\vect{P}=(1,0)^{\mathsf T}$、$\vect{Q}=(0,1)^{\mathsf T}$。
代入 §4.3 的公式：
\[
\vect{U}=\mathrm{e}^{-x/2}\bigl(\vect{P}\cos x-\vect{Q}\sin x\bigr)=\mathrm{e}^{-x/2}\mat{\cos x\\-\sin x},
\]
\[
\vect{W}=\mathrm{e}^{-x/2}\bigl(\vect{P}\sin x+\vect{Q}\cos x\bigr)=\mathrm{e}^{-x/2}\mat{\sin x\\\cos x}.
\]
\textbf{答案}\ 基本解组为 $\vect{U},\vect{W}$；Wronskian
\[
W=\det\mat{\mathrm{e}^{-x/2}\cos x&\mathrm{e}^{-x/2}\sin x\\-\mathrm{e}^{-x/2}\sin x&\mathrm{e}^{-x/2}\cos x}=\mathrm{e}^{-x}\neq0 .
\]
\textbf{自检}\ 直接代 $\vect{U}$：$U_1'=\mathrm{e}^{-x/2}(-\tfrac12\cos x-\sin x)$，而 $-\tfrac12U_1+U_2=\mathrm{e}^{-x/2}(-\tfrac12\cos x-\sin x)$ \ding{51}。

\qpagebreak
\subhdr{A.11\quad 例 3：实根与复根混合（p22--p25）}
\label{c2:app:ex3}
\qbreak
\qstem{（课件 p22 例 3）解 $\vect{Y}'=A\vect{Y}$，$A=\mat{2&1&0\\1&3&-1\\-1&2&3}$。}
\textbf{考点}\ 同一矩阵里 Case 1 与 Case 2 并用。
\textbf{关键一步}\ 特征多项式 $(2-\lambda)(\lambda^2-6\lambda+10)$；二次因子给一对复根 $3\pm\mathrm{i}$。
\textbf{详解}\ $\lambda_1=2$：$A-2E=\mat{0&1&0\\1&1&-1\\-1&2&1}$，解得 $\vect{V}_1=(1,0,1)^{\mathsf T}$，检验 $A\vect{V}_1=(2,0,2)^{\mathsf T}$ \ding{51}。
$\lambda=3+\mathrm{i}$：$A-(3+\mathrm{i})E$ 化简后取 $\vect{V}=\mat{1\\1\\2}+\mathrm{i}\mat{0\\1\\-1}$，即 $\vect{P}=(1,1,2)^{\mathsf T}$、$\vect{Q}=(0,1,-1)^{\mathsf T}$，$a=3$、$b=1$。
\[
\vect{U}=\mathrm{e}^{3x}\bigl(\vect{P}\cos x-\vect{Q}\sin x\bigr),\qquad
\vect{W}=\mathrm{e}^{3x}\bigl(\vect{P}\sin x+\vect{Q}\cos x\bigr).
\]
\textbf{答案}\ $\vect{Y}=c_1\vect{V}_1\mathrm{e}^{2x}+c_2\vect{U}+c_3\vect{W}$，分量式为
$y_1=c_1\mathrm{e}^{2x}+\mathrm{e}^{3x}(c_2\cos x+c_3\sin x)$、
$y_2=\mathrm{e}^{3x}[c_2(\cos x-\sin x)+c_3(\cos x+\sin x)]$、
$y_3=c_1\mathrm{e}^{2x}+\mathrm{e}^{3x}[c_2(2\cos x+\sin x)+c_3(2\sin x-\cos x)]$。
\textbf{自检}\ 取 $c_2=1$、其余为零，看 $x=0$：$\vect{U}(0)=\vect{P}=(1,1,2)$，而
\[
\vect{U}'(0)=3\vect{P}-\vect{Q}=(3,3,6)-(0,1,-1)=(3,2,7),
\]
另一边 $A\vect{P}=(2+1,\ 1+3-2,\ -1+2+6)=(3,2,7)$，两边相同 \ding{51}。注意这里用的是 $\vect{U}=\mathrm{e}^{3x}(\vect{P}\cos x-\vect{Q}\sin x)$ 的导数在 $0$ 处等于 $3\vect{P}-\vect{Q}$，不是 $3\vect{P}+\vect{Q}$——符号错了这一验就对不上。

\qpagebreak
\subhdr{A.12\quad 例 4：四维弹簧-质量（p27--p29）}
\label{c2:app:ex4}
\qbreak
\qstem{（课件 p27 例 4）解四维系统，$A=\mat{0&0&1&0\\0&0&0&1\\-2&3/2&0&0\\4/3&-3&0&0}$。}
\textbf{考点}\ 纯虚特征值；四维复根的实值化。
\textbf{关键一步}\ 特征方程只含 $\lambda^2$，令 $\mu=\lambda^2$ 解二次方程。
\textbf{详解}\ $\det(A-\lambda E)=\lambda^4+5\lambda^2+4=(\lambda^2+1)(\lambda^2+4)$，故 $\lambda=\pm\mathrm{i},\pm2\mathrm{i}$。
$\lambda=\mathrm{i}$ 的特征向量为 $\vect{V}_1=(3,2,3\mathrm{i},2\mathrm{i})^{\mathsf T}$。逐行验算 $A\vect{V}_1$：第一行取 $y_3=3\mathrm{i}$；第二行取 $y_4=2\mathrm{i}$；第三行 $-2\cdot3+\tfrac32\cdot2=-6+3=-3$；第四行 $\tfrac43\cdot3-3\cdot2=4-6=-2$。于是 $A\vect{V}_1=(3\mathrm{i},2\mathrm{i},-3,-2)^{\mathsf T}=\mathrm{i}\vect{V}_1$ \ding{51}。
取实部虚部得两条实解；$\lambda=2\mathrm{i}$ 的特征向量 $\vect{V}_3=(3,-4,6\mathrm{i},-8\mathrm{i})^{\mathsf T}$ 再给两条。
\textbf{答案}\ 通解由四条实解线性组合而成（课件 p29）：
\[
\vect{Y}=c_1\mat{3\cos t\\2\cos t\\-3\sin t\\-2\sin t}+c_2\mat{3\sin t\\2\sin t\\3\cos t\\2\cos t}
+c_3\mat{3\cos2t\\-4\cos2t\\-6\sin2t\\8\sin2t}+c_4\mat{3\sin2t\\-4\sin2t\\6\cos2t\\-8\cos2t}.
\]
\textbf{自检}\ 取 $c_1=1$ 其余 $0$、$t=0$：$\vect{Y}(0)=(3,2,0,0)$，$\vect{Y}'(0)=(0,0,-3,-2)$；而 $A\vect{Y}(0)$ 第一行 $=y_3=0$、第二行 $=y_4=0$、第三行 $=-6+3=-3$、第四行 $=4-6=-2$ \ding{51}。

\qpagebreak
\subhdr{A.13\quad 例 5：$g=m=2$ 的重根（p33--p36）}
\label{c2:app:ex5}
\qbreak
\qstem{（课件 p33 例 5）解 $\vect{Y}'=A\vect{Y}$，$A=\mat{4&-1&-1\\1&2&-1\\1&-1&2}$。}
\textbf{考点}\ Case 3a：重根但特征向量够用。
\textbf{关键一步}\ 先算 $\rank(A-3E)$：它是 $1$，所以 $g=3-1=2=m$。
\textbf{详解}\ 特征方程 $(3-\lambda)(\lambda-3)(\lambda-2)=0$，即 $(\lambda-3)^2(\lambda-2)=0$，故 $\lambda_1=2$（$m_1=1$）、$\lambda_2=3$（$m_2=2$）。
$\lambda_1=2$：$A-2E=\mat{2&-1&-1\\1&0&-1\\1&-1&1}$，解得 $\vect{V}_1=(1,1,1)^{\mathsf T}$，检验 $A\vect{V}_1=(2,2,2)^{\mathsf T}$ \ding{51}。
$\lambda_2=3$：$A-3E=\mat{1&-1&-1\\1&-1&-1\\1&-1&-1}$，秩 $1$，故 $g=2$；方程只约束 $v_1=v_2+v_3$，取 $(v_2,v_3)=(1,0)$ 与 $(0,1)$ 得 $\vect{V}_2=(1,1,0)^{\mathsf T}$、$\vect{V}_3=(1,0,1)^{\mathsf T}$。检验 $A\vect{V}_2=(3,3,0)^{\mathsf T}$、$A\vect{V}_3=(3,0,3)^{\mathsf T}$ \ding{51}。
\textbf{答案}\ $\vect{Y}=c_1\mat{1\\1\\1}\mathrm{e}^{2x}+c_2\mat{1\\1\\0}\mathrm{e}^{3x}+c_3\mat{1\\0\\1}\mathrm{e}^{3x}$。
\textbf{自检}\ 三个解在 $x=0$ 的向量为 $(1,1,1),(1,1,0),(1,0,1)$，行列式 $=-1\neq0$，故构成基本解组 \ding{51}。

\qpagebreak
\subhdr{A.14\quad 例 6：$g=1<m=2$ 的重根（p39--p42）}
\label{c2:app:ex6}
\qbreak
\qstem{（课件 p39 例 6）解 $\vect{Y}'=A\vect{Y}$，$A=\mat{1&-1\\1&3}$。}
\textbf{考点}\ Case 3b 的最小例子：二重根、只有一个特征向量。
\textbf{关键一步}\ $(A-2E)=\mat{-1&-1\\1&1}$ 秩 $1$，$g=2-1=1<m=2$，要接一个广义特征向量。
\textbf{详解}\ $\det(A-\lambda E)=(\lambda-2)^2$，$\lambda=2$（$m=2$）。
特征向量：$v_1+v_2=0$，取 $\vect{V}_1=(1,-1)^{\mathsf T}$。
广义特征向量：解 $(A-2E)\vect{V}^{(1)}=\vect{V}_1$：
\[
\mat{-1&-1\\1&1}\mat{w_1\\w_2}=\mat{1\\-1}
\Longrightarrow -w_1-w_2=1,
\]
取 $w_1=0$ 得 $w_2=-1$，即 $\vect{V}^{(1)}=(0,-1)^{\mathsf T}$。
\textbf{答案}\ $\vect{Y}=c_1\mathrm{e}^{2x}\mat{1\\-1}+c_2\Bigl[x\mathrm{e}^{2x}\mat{1\\-1}+\mathrm{e}^{2x}\mat{0\\-1}\Bigr]$。
\textbf{自检}\ 取 $c_2=1,c_1=0$：$\vect{Y}=(x,\,-x-1)^{\mathsf T}\mathrm{e}^{2x}$。求导
\[
\vect{Y}'=\mathrm{e}^{2x}\bigl[(1,1)^{\mathsf T}+2(x,-x-1)^{\mathsf T}\bigr]=\mathrm{e}^{2x}(1+2x,\ -3-2x)^{\mathsf T}.
\]
另一边逐行算 $A\vect{Y}$：第一行 $=y_1-y_2=x-(-x-1)=2x+1$；第二行 $=y_1+3y_2=x-3x-3=-2x-3$。两边相同 \ding{51}。

\qpagebreak
\subhdr{A.15\quad 复习例 (1)：三重根、$g=1$（p46--p51）}
\label{c2:app:rev1}
\qbreak
\qstem{（课件 p46 复习例 (1)）解 $\vect{Y}'=A\vect{Y}$，$A=\mat{4&-1&0\\3&1&-1\\1&0&1}$。}
\textbf{考点}\ 长度 2 的链（$m=3$、$g=1$ 时链长 $m-g=2$）。
\textbf{关键一步}\ 先求 $\vect{V}_1$，再解两次 $(A-2E)\vect{V}^{(1)}=\vect{V}_1$、$(A-2E)\vect{V}^{(2)}=\vect{V}^{(1)}$。
\textbf{详解}\ $\det(A-\lambda E)=-(\lambda-2)^3$，$\lambda=2$（$m=3$）。$A-2E=\mat{2&-1&0\\3&-1&-1\\1&0&-1}$，秩 $2$，故 $g=3-2=1$。特征向量 $\vect{V}_1=(1,2,1)^{\mathsf T}$。
链：$(A-2E)\vect{V}^{(1)}=\vect{V}_1$ 取 $\vect{V}^{(1)}=(1,1,0)^{\mathsf T}$（检验：第一行 $2-1=1$ \ding{51}，第二行 $3-1=2$ \ding{51}，第三行 $1=1$ \ding{51}）。
$(A-2E)\vect{V}^{(2)}=\vect{V}^{(1)}$ 取 $\vect{V}^{(2)}=(0,-1,0)^{\mathsf T}$（第一行 $1=1$？：$2\cdot0-(-1)=1$ \ding{51}；第二行 $1=1$ \ding{51}；第三行 $0$ \ding{51}）。
\textbf{答案}\ $\vect{Y}=c_1\vect{V}_1\mathrm{e}^{2x}+c_2(\vect{V}_1x+\vect{V}^{(1)})\mathrm{e}^{2x}+c_3\bigl(\vect{V}_1\tfrac{x^2}{2}+\vect{V}^{(1)}x+\vect{V}^{(2)}\bigr)\mathrm{e}^{2x}$。
\textbf{自检}\ 用 $x=0$：三条解在 $0$ 处为 $(1,2,1),(1,1,0),(0,-1,0)$，行列式非零 \ding{51}。

\qpagebreak
\subhdr{A.16\quad 复习例 (2)：实根 + 重根（p52--p57）}
\label{c2:app:rev2}
\qbreak
\qstem{（课件 p52 复习例 (2)）解 $\vect{Y}'=A\vect{Y}$，$A=\mat{-1&1&-2\\4&1&0\\2&1&-1}$。}
\textbf{考点}\ 一个单根与一个重根并存；重根 $g=1<m=2$。
\textbf{关键一步}\ 特征方程 $(1-\lambda^2)(\lambda+1)=0$，即 $\lambda=1$ 与 $\lambda=-1$（$m=2$）。
\textbf{详解}\ $\lambda_1=1$：$A-E=\mat{-2&1&-2\\4&0&0\\2&1&-2}$，解得 $\vect{V}_1=(0,2,1)^{\mathsf T}$。
$\lambda_2=-1$（$m=2$）：$A+E=\mat{0&1&-2\\4&2&0\\2&1&0}$，秩 $2$，$g=3-2=1$；特征向量 $\vect{V}_2=(-1,2,1)^{\mathsf T}$。
广义：解 $(A+E)\vect{V}^{(1)}=\vect{V}_2$，即 $w_2-2w_3=-1$、$4w_1+2w_2=2$、$2w_1+w_2=1$，取 $w_3=0$ 得 $w_2=-1$、$w_1=1$，或按课件取 $\vect{V}^{(1)}=(0,1,1)^{\mathsf T}$（检验：$1-2=-1$ \ding{51}，$2=2$ \ding{51}，$1=1$ \ding{51}）。
\textbf{答案}\ $\vect{Y}=c_1\mat{0\\2\\1}\mathrm{e}^{x}+c_2\mat{-1\\2\\1}\mathrm{e}^{-x}+c_3\Bigl[\mat{-1\\2\\1}x+\mat{0\\1\\1}\Bigr]\mathrm{e}^{-x}$。
\textbf{自检}\ 三个解在 $x=0$ 为 $(0,2,1),(-1,2,1),(0,1,1)$，行列式非零 \ding{51}。

\qpagebreak
\subhdr{A.17\quad 复习例 (3)：三重根、$g=2$（p58--p63）}
\label{c2:app:rev3}
\qbreak
\qstem{（课件 p58 复习例 (3)）解 $\vect{Y}'=A\vect{Y}$，$A=\mat{0&-4&0\\1&-4&0\\1&-2&-2}$。}
\textbf{考点}\ $g=2$、$m=3$：两个特征向量，但只有一个能接上链。
\textbf{关键一步}\ 先取两个特征向量 $\vect{V}_1=(2,1,0)^{\mathsf T}$、$\vect{V}_2=(0,0,1)^{\mathsf T}$；再检验哪个在 $\operatorname{Range}(A+2E)$ 里（比较秩）。
\textbf{详解}\ $\det(A-\lambda E)=-(\lambda+2)^3$，$\lambda=-2$（$m=3$）。$A+2E=\mat{2&-4&0\\1&-2&0\\1&-2&0}$ 秩 $1$，$g=3-1=2$；特征向量由 $v_1=2v_2$ 给出：$\vect{V}_1=(2,1,0)^{\mathsf T}$、$\vect{V}_2=(0,0,1)^{\mathsf T}$。
\textbf{关键处}：$(A+2E)\vect{W}=\vect{V}$ 要有解，需 $\rank(A+2E\mid\vect{V})=\rank(A+2E)=1$。
取 $\vect{V}=\vect{V}_2=(0,0,1)$：增广矩阵第三行给 $w_1-2w_2=1$，与第一、二行的 $0$ 矛盾，无解。
取 $\vect{V}=\vect{V}_1+\vect{V}_2=(2,1,1)$：满足条件，解 $(A+2E)\vect{W}=(2,1,1)$ 得 $\vect{W}=(1,0,0)^{\mathsf T}$（检验：$(2,1,1)$ \ding{51}）。
\textbf{答案}\ $\vect{Y}=c_1\mat{0\\0\\1}\mathrm{e}^{-2x}+c_2\mat{2\\1\\1}\mathrm{e}^{-2x}+c_3\Bigl[\mat{1\\0\\0}+x\mat{2\\1\\1}\Bigr]\mathrm{e}^{-2x}$。
\textbf{自检}\ 三个向量 $(0,0,1),(2,1,1),(1,0,0)$ 线性无关 \ding{51}；且 $(A+2E)(1,0,0)^{\mathsf T}=(2,1,1)^{\mathsf T}$ \ding{51}。

\qpagebreak
\subhdr{A.18\quad 作业第 1 题（p64）}
\label{c2:app:hw1}
\qbreak
\qstem{（课件 p64 第 1 题）$A=\mat{1&1&-3\\3&2&-3\\1&0&1}$，解 $\vect{Y}'=A\vect{Y}$。}
\textbf{考点}\ 一重根 + 二重根（链长 1）。
\textbf{关键一步}\ 特征多项式 $\lambda^3-4\lambda^2+5\lambda-2=(\lambda-1)^2(\lambda-2)$；对 $\lambda=1$ 算 $g$。
\textbf{详解}\ $\lambda_1=1$（$m=2$）、$\lambda_2=2$（$m=1$）。
$\lambda=1$：$A-E=\mat{0&1&-3\\3&1&-3\\1&0&0}$，秩 $2$，$g=3-2=1<m$；特征向量由 $v_1=0$、$v_2=3v_3$ 得 $\vect{V}_1=(0,3,1)^{\mathsf T}$。
广义：$(A-E)\vect{W}=\vect{V}_1$，即 $w_1=1$、$w_2-3w_3=0$、$3w_1+w_2-3w_3=3$，取 $w_3=0$ 得 $\vect{W}=(1,0,0)^{\mathsf T}$。
$\lambda=2$：$A-2E=\mat{-1&1&-3\\3&0&-3\\1&0&-1}$，解得 $\vect{V}_2=(1,4,1)^{\mathsf T}$。
\textbf{答案}\ $\vect{Y}=c_1\mat{0\\3\\1}\mathrm{e}^{x}+c_2\Bigl[\mat{1\\0\\0}+x\mat{0\\3\\1}\Bigr]\mathrm{e}^{x}+c_3\mat{1\\4\\1}\mathrm{e}^{2x}$。
\textbf{自检}\ $(A-E)(1,0,0)^{\mathsf T}=(0,3,1)^{\mathsf T}$ \ding{51}；$A(1,4,1)^{\mathsf T}=(1+4-3,3+8-3,1+1)=(2,8,2)=2(1,4,1)$ \ding{51}。

\qpagebreak
\subhdr{A.19\quad 作业第 2 题（p64）}
\label{c2:app:hw2}
\qbreak
\qstem{（课件 p64 第 2 题）$A=\mat{-2&1\\-5&2}$，解 $\vect{Y}'=A\vect{Y}$。}
\textbf{考点}\ 纯虚根；实值化。
\textbf{关键一步}\ 迹 $=0$、行列式 $=1$，特征方程 $\lambda^2+1=0$。
\textbf{详解}\ $\lambda=\pm\mathrm{i}$，$a=0$、$b=1$。
对 $\lambda=\mathrm{i}$：$A-\mathrm{i}E=\mat{-2-\mathrm{i}&1\\-5&2-\mathrm{i}}$，方程 $(-2-\mathrm{i})v_1+v_2=0$，取 $\vect{V}=(1,2+\mathrm{i})^{\mathsf T}=\vect{P}+\mathrm{i}\vect{Q}$，$\vect{P}=(1,2)$、$\vect{Q}=(0,1)$。
\[
\vect{U}=\vect{P}\cos x-\vect{Q}\sin x=\mat{\cos x\\2\cos x-\sin x},\qquad
\vect{W}=\vect{P}\sin x+\vect{Q}\cos x=\mat{\sin x\\2\sin x+\cos x}.
\]
\textbf{答案}\ $\vect{Y}=c_1\mat{\cos x\\2\cos x-\sin x}+c_2\mat{\sin x\\2\sin x+\cos x}$。
\textbf{自检}\ 取 $\vect{U}$ 在 $x=0$：$\vect{U}(0)=(1,2)$，$\vect{U}'(0)=(0,-1)$；$A\vect{U}(0)=(-2+2,-5+4)=(0,-1)$ \ding{51}。

\qpagebreak
\subhdr{A.20\quad 作业第 3 题（p64）}
\label{c2:app:hw3}
\qbreak
\qstem{（课件 p64 第 3 题）$A=\mat{-2&-1&1\\-3&-3&2\\-4&-3&2}$，解 $\vect{Y}'=A\vect{Y}$。}
\textbf{考点}\ 三重根、$g=1$：链长 2。
\textbf{关键一步}\ 特征多项式 $(\lambda+1)^3$；$A+E$ 的秩为 $2$，故 $g=1$。
\textbf{详解}\ $\lambda=-1$（$m=3$）。$A+E=\mat{-1&-1&1\\-3&-2&2\\-4&-3&3}$，化简得秩 $2$；特征向量由 $v_2=v_3$、$v_1=0$ 给出，$\vect{V}=(0,1,1)^{\mathsf T}$。
$(A+E)\vect{W}=\vect{V}$：取 $\vect{W}=(-1,-3,-4)^{\mathsf T}$（检验：$1+3-4=0$ \ding{51}；$3+6-8=1$ \ding{51}；$4+9-12=1$ \ding{51}）。
$(A+E)\vect{U}=\vect{W}$：取 $\vect{U}=(1,0,0)^{\mathsf T}$（检验：$(-1,-3,-4)$ \ding{51}）。
\textbf{答案}\ $\vect{Y}=c_1\vect{V}\mathrm{e}^{-x}+c_2(\vect{W}+x\vect{V})\mathrm{e}^{-x}+c_3\bigl(\vect{U}+x\vect{W}+\tfrac{x^2}{2}\vect{V}\bigr)\mathrm{e}^{-x}$，其中 $\vect{V}=(0,1,1)$、$\vect{W}=(-1,-3,-4)$、$\vect{U}=(1,0,0)$。
\textbf{自检}\ 三个向量线性无关（行列式 $=-1$）\ding{51}；链的两步已逐行验算 \ding{51}。

\qpagebreak
\subhdr{A.21\quad 作业第 4 题（p64）}
\label{c2:app:hw4}
\qbreak
\qstem{（课件 p64 第 4 题）$A=\mat{1&1\\9&1}$，解 $\vect{Y}'=A\vect{Y}$。}
\textbf{考点}\ 相异实根，Case 1 的最基本操作。
\textbf{关键一步}\ 迹 $2$、行列式 $-8$，特征方程 $\lambda^2-2\lambda-8=0$。
\textbf{详解}\ $\lambda_1=4$、$\lambda_2=-2$。
$\lambda=4$：$A-4E=\mat{-3&1\\9&-3}$，$-3v_1+v_2=0$，取 $\vect{V}_1=(1,3)^{\mathsf T}$。
$\lambda=-2$：$A+2E=\mat{3&1\\9&3}$，$3v_1+v_2=0$，取 $\vect{V}_2=(1,-3)^{\mathsf T}$。
\textbf{答案}\ $\vect{Y}=c_1\mat{1\\3}\mathrm{e}^{4x}+c_2\mat{1\\-3}\mathrm{e}^{-2x}$。
\textbf{自检}\ $A(1,3)^{\mathsf T}=(1+3,9+3)=(4,12)=4(1,3)$ \ding{51}；$A(1,-3)^{\mathsf T}=(1-3,9-3)=(-2,6)=-2(1,-3)$ \ding{51}。

\flowerrule

{\small\color{inkgray}
依据课件：\texttt{14\_15\_Linear-DE-systems-with-constant-coefficients.pdf}（64 页，2026-09-29 / 10-08，W.Y., DLUT）。
缺口 10 条见 §12；附录 A.18--A.13 的四题解答已用符号计算核对（特征值、广义特征向量与代回残差）。}
\subhdr{第 3 章（L16--17）· 例题与作业详解}

每块自成一页，按“题干 → 考点 → 关键一步 → 详解 → 答案 → 自检”排。

\qbreak
\subhdr{A.22\quad 例 1：对角化算 $\mathrm{e}^{Ax}$（p13--p15）}
\label{c3:app:ex1}
\qbreak
\qstem{（课件 p13 例 1）$A=\mat{1&-1&1\\1&1&-1\\2&-1&0}$，求 $\mathrm{e}^{Ax}$。}
\par\textbf{考点}\ 对角化路线：$A=V\Lambda V^{-1}$。
\par\textbf{关键一步}\ 特征值 $2,1,-1$ 互异，故可对角化；特征向量按列排成 $V$，再求 $V^{-1}$。
\par\textbf{详解}\ 特征向量（L14--15 例 1 已算）：$\vect{V}_1=(1,0,1)^{\mathsf T}$、$\vect{V}_2=(1,1,1)^{\mathsf T}$、$\vect{V}_3=(-1,3,5)^{\mathsf T}$，对应 $\lambda_1=2,\lambda_2=1,\lambda_3=-1$。于是
\[
V=\mat{1&1&-1\\0&1&3\\1&1&5},\qquad
\Lambda=\mat{2&0&0\\0&1&0\\0&0&-1},\qquad
V^{-1}=\mat{1/3&-1&2/3\\1/2&1&-1/2\\-1/6&0&1/6}.
\]
\textbf{注意}：课件 p14 把 $V^{-1}$ 的第 1 行第 2 列印成 $0$，正确值是 $-1$；用 $V\cdot V^{-1}=E$ 一乘就能看出。
\[
\mathrm{e}^{Ax}=V\mathrm{e}^{\Lambda x}V^{-1}
=V\mat{\mathrm{e}^{2x}&0&0\\0&\mathrm{e}^{x}&0\\0&0&\mathrm{e}^{-x}}V^{-1}.
\]
\par\textbf{答案}\ 上式即 $\mathrm{e}^{Ax}$；通解 $\vect{Y}=\mathrm{e}^{Ax}\vect{C}=c_1\vect{V}_1\mathrm{e}^{2x}+c_2\vect{V}_2\mathrm{e}^{x}+c_3\vect{V}_3\mathrm{e}^{-x}$。
\par\textbf{自检}\ 在 $x=0$：$\mathrm{e}^{A\cdot0}=V E V^{-1}=E$ \ding{51}。对 $x$ 求导在 $0$ 处应得 $A$，等价于 $V\Lambda V^{-1}=A$，而 $\Lambda=V^{-1}AV$ 正是构造，\ding{51}。

\qpagebreak
\subhdr{A.23\quad 例 2：多项式降阶（p20）}
\label{c3:app:ex2}
\qbreak
\qstem{（课件 p20 例 2）对 $A=\mat{3&1\\1&2}$，把 $p(A)=A^4+3A^3+2A^2+A+E$ 降到一次。}
\par\textbf{考点}\ Cayley--Hamilton 的降阶用法。
\par\textbf{关键一步}\ 先算 $\Delta(s)=\det(A-sE)=s^2-5s+5$，再做多项式除法取余。
\par\textbf{详解}\ 长除：
\[
\frac{s^4+3s^3+2s^2+s+1}{s^2-5s+5}=s^2+8s+37+\frac{146s-184}{s^2-5s+5}.
\]
即 $p(s)=(s^2+8s+37)\Delta(s)+(146s-184)$。由 C--H，$\Delta(A)=[0]$，故
\[
p(A)=146A-184E=\mat{438&146\\146&292}-\mat{184&0\\0&184}=\mat{254&146\\146&108}.
\]
\par\textbf{答案}\ $p(A)=\mat{254&146\\146&108}$。
\par\textbf{自检}\ 用特征值验：$\lambda_{1,2}=(5\pm\sqrt5)/2$，直接代入 $p(\lambda)=\lambda^4+3\lambda^3+2\lambda^2+\lambda+1$ 与 $146\lambda-184$，两者相等（代数恒等式）。

\qpagebreak
\subhdr{A.24\quad 例 4：$2\times2$ 实根（p24--p25）}
\label{c3:app:ex3}
\qbreak
\qstem{（课件 p24 例 4）$A=\mat{1&1\\9&1}$，求 $\mathrm{e}^{Ax}$ 与通解。}
\par\textbf{考点}\ C--H 的最小例子：$n=2$，两个相异实根。
\par\textbf{关键一步}\ 设 $\mathrm{e}^{Ax}=\alpha_1Ax+\alpha_0E$，用 $\mathrm{e}^{4x}$ 与 $\mathrm{e}^{-2x}$ 两个方程定 $\alpha_1x,\alpha_0$。
\par\textbf{详解}\ $A$ 的特征值：$\lambda^2-2\lambda-8=0$，$\lambda=4,-2$。$Ax$ 的特征值是 $4x,-2x$。
\[
\mathrm{e}^{4x}=4x\alpha_1+\alpha_0,\qquad \mathrm{e}^{-2x}=-2x\alpha_1+\alpha_0 .
\]
两式相减得 $6x\alpha_1=\mathrm{e}^{4x}-\mathrm{e}^{-2x}$，相加得 $2\alpha_0=\mathrm{e}^{4x}+2\mathrm{e}^{-2x}$：
\[
x\alpha_1=\frac{\mathrm{e}^{4x}-\mathrm{e}^{-2x}}{6},\qquad
\alpha_0=\frac{\mathrm{e}^{4x}+2\mathrm{e}^{-2x}}{3}.
\]
\[
\mathrm{e}^{Ax}=\alpha_1Ax+\alpha_0E
=\frac16\mat{3\mathrm{e}^{4x}+3\mathrm{e}^{-2x}&\mathrm{e}^{4x}-\mathrm{e}^{-2x}\\9\mathrm{e}^{4x}-9\mathrm{e}^{-2x}&3\mathrm{e}^{4x}+3\mathrm{e}^{-2x}}.
\]
\par\textbf{答案}\ 上式；通解 $\vect{Y}=\tilde c_1\mathrm{e}^{4x}\binom{1}{3}+\tilde c_2\mathrm{e}^{-2x}\binom{1}{-3}$。
\par\textbf{自检}\ $x=0$ 时矩阵为 $\frac16\mat{6&0\\0&6}=E$ \ding{51}。再验导数：$(1,1)$ 元素 $\frac16(3\mathrm{e}^{4x}+3\mathrm{e}^{-2x})$ 在 $0$ 处导数为 $\frac16(12-6)=1$，与 $A_{11}=1$ 一致 \ding{51}；$(1,2)$ 元素 $\frac16(\mathrm{e}^{4x}-\mathrm{e}^{-2x})$ 的导数为 $\frac16(4+2)=1=A_{12}$ \ding{51}；$(2,1)$ 元素给 $\frac16(36+18)=9=A_{21}$ \ding{51}。

\qpagebreak
\subhdr{A.25\quad 例 5：纯虚根与旋转（p26--p27）}
\label{c3:app:ex4}
\qbreak
\qstem{（课件 p26 例 5）$A=\mat{0&1\\-1&0}$，求 $\mathrm{e}^{Ax}$ 与通解。}
\par\textbf{考点}\ 纯虚特征值；$\mathrm{e}^{Ax}$ 是旋转矩阵。
\par\textbf{关键一步}\ $Ax$ 的特征值是 $\pm\mathrm{i}x$；用欧拉公式把 $\alpha_1x,\alpha_0$ 解成 $\sin x,\cos x$。
\par\textbf{详解}\ 特征方程 $\lambda^2+1=0$，$\lambda=\pm\mathrm{i}$；$Ax$ 的特征值 $\pm\mathrm{i}x$。
\[
\mathrm{e}^{\mathrm{i}x}=\mathrm{i}x\alpha_1+\alpha_0,\qquad \mathrm{e}^{-\mathrm{i}x}=-\mathrm{i}x\alpha_1+\alpha_0 .
\]
两式相加：$\alpha_0=\cos x$；相减除以 $2\mathrm{i}$：$x\alpha_1=\sin x$。于是
\[
\mathrm{e}^{Ax}=\alpha_0E+\alpha_1Ax=\mat{\cos x&\sin x\\-\sin x&\cos x}.
\]
\par\textbf{答案}\ $\mathrm{e}^{Ax}=\mat{\cos x&\sin x\\-\sin x&\cos x}$，即绕原点顺时针旋转 $x$ 的矩阵（行列式 $1$，正交）。
通解 $\vect{Y}=\mathrm{e}^{Ax}\vect{C}=c_1\binom{\cos x}{-\sin x}+c_2\binom{\sin x}{\cos x}$。
\textbf{注意}：课件 p27 末行把第二项写成 $\binom{-\sin x}{\cos x}$，符号错了；直接乘 $\mathrm{e}^{Ax}\binom{c_1}{c_2}$ 得 $(c_1\cos x+c_2\sin x,\ -c_1\sin x+c_2\cos x)$，第二项应是 $\binom{\sin x}{\cos x}$。
\par\textbf{自检}\ $x=0$ 得 $E$ \ding{51}；求导在 $0$ 处得 $\mat{0&1\\-1&0}=A$ \ding{51}。

\qpagebreak
\subhdr{A.26\quad 例 7：三重根（p33--p35）}
\label{c3:app:ex5}
\qbreak
\qstem{（课件 p33 例 7）$A=\mat{0&-4&0\\1&-4&0\\1&-2&-2}$，求 $\mathrm{e}^{Ax}$。}
\par\textbf{考点}\ 重根时补导数方程。
\par\textbf{关键一步}\ $A$ 的特征值 $-2$（三重），故 $Ax$ 的特征值 $-2x$（三重）；要写 $\mathrm{e}^{-2x}=r(-2x)$、$\mathrm{e}^{-2x}=r'(-2x)$、$\mathrm{e}^{-2x}=r''(-2x)$。
\par\textbf{详解}\ 设 $r(\lambda)=\alpha_2\lambda^2+\alpha_1\lambda+\alpha_0$。三个方程：
\[
\mathrm{e}^{-2x}=4x^2\alpha_2-2x\alpha_1+\alpha_0,\quad
\mathrm{e}^{-2x}=-4x\alpha_2+\alpha_1,\quad
\mathrm{e}^{-2x}=2\alpha_2 .
\]
由第三式 $\alpha_2=\tfrac12\mathrm{e}^{-2x}$；代入第二式 $\alpha_1=(1+2x)\mathrm{e}^{-2x}$；再代第一式 $\alpha_0=(1+2x+2x^2)\mathrm{e}^{-2x}$。
代回 $\mathrm{e}^{Ax}=\alpha_2A^2x^2+\alpha_1Ax+\alpha_0E$，整理得
\[
\mathrm{e}^{Ax}=\mathrm{e}^{-2x}\mat{2x+1&-4x&0\\x&1-2x&0\\x&-2x&1}.
\]
\par\textbf{答案}\ 上式。
\par\textbf{自检}\ $x=0$ 得 $E$ \ding{51}；求导在 $0$ 处：$(1,1)$ 元素 $2-2=0$ \ding{51}，$(2,1)$ 元素 $1$ \ding{51}，$(3,2)$ 元素 $-2$ \ding{51}，与 $A$ 同。

\qpagebreak
\subhdr{A.27\quad 例 8：$0$ 与 $\pm\mathrm{i}$ 并存（p36--p38）}
\label{c3:app:ex6}
\qbreak
\qstem{（课件 p36 例 8）$A=\mat{0&1&0\\0&-2&-5\\0&1&2}$，求 $\mathrm{e}^{Ax}$。}
\par\textbf{考点}\ 实根与纯虚根混合；$\alpha_k$ 用三个方程定。
\par\textbf{关键一步}\ $A$ 的特征值 $0,\pm\mathrm{i}$；$Ax$ 的特征值 $0,\pm\mathrm{i}x$。
\par\textbf{详解}\ 设 $r(\lambda)=\alpha_2\lambda^2+\alpha_1\lambda+\alpha_0$。三个方程：
\[
1=\alpha_0,\qquad \mathrm{e}^{\mathrm{i}x}=-x^2\alpha_2+\mathrm{i}x\alpha_1+\alpha_0,\qquad \mathrm{e}^{-\mathrm{i}x}=-x^2\alpha_2-\mathrm{i}x\alpha_1+\alpha_0 .
\]
后两式相加得 $2\cos x=-2x^2\alpha_2+2\alpha_0$，故 $x^2\alpha_2=1-\cos x$；相减得 $2\mathrm{i}\sin x=2\mathrm{i}x\alpha_1$，故 $x\alpha_1=\sin x$。
代回（课件 p36 已算好 $A^2$ 的表达式）：
\[
\mathrm{e}^{Ax}=\mat{1&-2+2\cos x+\sin x&-5+5\cos x\\0&\cos x-2\sin x&-5\sin x\\0&\sin x&\cos x+2\sin x}.
\]
\par\textbf{答案}\ 上式。
\par\textbf{自检}\ $x=0$ 得 $E$ \ding{51}；求导在 $0$ 处：$(1,2)$ 元素 $0+1=1$ \ding{51}，$(2,3)$ 元素 $-5$ \ding{51}，$(3,3)$ 元素 $2$ \ding{51}。

\qpagebreak
\subhdr{A.28\quad 作业第 1 题}
\label{c3:app:hw1}
\qbreak
\qstem{（课件 p44 第 1 题）$A=\mat{-6&-4\\8&6}$，求 $\mathrm{e}^{Ax}$ 与通解。}
\par\textbf{考点}\ C--H，$2\times2$ 相异实根。
\par\textbf{关键一步}\ 特征值 $\pm2$，$Ax$ 的特征值 $\pm2x$；设 $\mathrm{e}^{Ax}=\alpha_1Ax+\alpha_0E$。
\par\textbf{详解}\ $\operatorname{tr}A=0$、$\det A=-36+32=-4$，特征方程 $\lambda^2-4=0$，$\lambda=\pm2$。
\[
\mathrm{e}^{2x}=2x\alpha_1+\alpha_0,\qquad \mathrm{e}^{-2x}=-2x\alpha_1+\alpha_0 .
\]
解得 $x\alpha_1=\tfrac14(\mathrm{e}^{2x}-\mathrm{e}^{-2x})=\tfrac12\sinh2x$，$\alpha_0=\tfrac12(\mathrm{e}^{2x}+\mathrm{e}^{-2x})=\cosh2x$。于是
\[
\mathrm{e}^{Ax}=\alpha_1Ax+\alpha_0E
=\frac14(\mathrm{e}^{2x}-\mathrm{e}^{-2x})A+\cosh2x\,E
=\mat{2\mathrm{e}^{-2x}-\mathrm{e}^{2x}&\mathrm{e}^{-2x}-\mathrm{e}^{2x}\\2\mathrm{e}^{2x}-2\mathrm{e}^{-2x}&2\mathrm{e}^{2x}-\mathrm{e}^{-2x}}.
\]
\par\textbf{答案}\ 上式；通解为列向量的线性组合（两个特征向量 $(1,-2)$ 与 $(1,-1)$ 分别配 $\mathrm{e}^{2x}$、$\mathrm{e}^{-2x}$）。
\par\textbf{自检}\ $x=0$ 得 $E$ \ding{51}；求导在 $0$ 处：第一行 $( -4-2,\ -2-2)=(-6,-4)$，与 $A$ 第一行同 \ding{51}。

\qpagebreak
\subhdr{A.29\quad 作业第 2 题}
\label{c3:app:hw2}
\qbreak
\qstem{（课件 p44 第 2 题）$A=\mat{-4&4\\-5&4}$，求 $\mathrm{e}^{Ax}$ 与通解。}
\par\textbf{考点}\ 纯虚根，$\mathrm{e}^{Ax}$ 含 $\cos2x$ 与 $\sin2x$。
\par\textbf{关键一步}\ 特征值 $\pm2\mathrm{i}$；用 $\mathrm{e}^{\pm2\mathrm{i}x}$ 的两个方程定 $\alpha$。
\par\textbf{详解}\ $\operatorname{tr}A=0$、$\det A=-16+20=4$，特征方程 $\lambda^2+4=0$，$\lambda=\pm2\mathrm{i}$。
\[
\mathrm{e}^{2\mathrm{i}x}=2\mathrm{i}x\alpha_1+\alpha_0,\qquad \mathrm{e}^{-2\mathrm{i}x}=-2\mathrm{i}x\alpha_1+\alpha_0 .
\]
得 $\alpha_0=\cos2x$，$x\alpha_1=\tfrac{1}{2}\sin2x$。于是
\[
\mathrm{e}^{Ax}=\alpha_0E+\alpha_1Ax
=\mat{\cos2x-2\sin2x&2\sin2x\\-\tfrac52\sin2x&\cos2x+2\sin2x}.
\]
\par\textbf{答案}\ 上式；通解 $\vect{Y}=\mathrm{e}^{Ax}\vect{C}$。
\par\textbf{自检}\ $x=0$ 得 $E$ \ding{51}；求导在 $0$ 处第一行 $( -4,\ 4)$、第二行 $(-5,\ 4)$，与 $A$ 同 \ding{51}。

\qpagebreak
\subhdr{A.30\quad 作业第 3 题}
\label{c3:app:hw3}
\qbreak
\qstem{（课件 p44 第 3 题）$A=\mat{-24&0&0\\4&-2&0\\1&0&2}$，求 $\mathrm{e}^{Ax}$ 与通解。}
\par\textbf{考点}\ $3\times3$、三个相异实根；下三角阵的指数仍是下三角。
\par\textbf{关键一步}\ 下三角阵的特征值就是对角元 $-24,-2,2$；代三个方程解 $\alpha_2,\alpha_1,\alpha_0$。
\par\textbf{详解}\ 设 $\mathrm{e}^{Ax}=\alpha_2A^2x^2+\alpha_1Ax+\alpha_0E$，三个方程：
\[
\mathrm{e}^{-24x}=(-24x)^2\alpha_2-24x\alpha_1+\alpha_0,\quad
\mathrm{e}^{-2x}=4x^2\alpha_2-2x\alpha_1+\alpha_0,\quad
\mathrm{e}^{2x}=4x^2\alpha_2+2x\alpha_1+\alpha_0 .
\]
解得
\[
\alpha_2x^2=\frac{2\mathrm{e}^{-24x}+11\mathrm{e}^{2x}-13\mathrm{e}^{-2x}}{1144},\qquad
\alpha_1x=\frac{\mathrm{e}^{2x}-\mathrm{e}^{-2x}}{4},
\]
而 $\alpha_0$ 由第二式回代：$\alpha_0=\mathrm{e}^{-2x}-4x^2\alpha_2+2x\alpha_1$。把这三者代回 $\mathrm{e}^{Ax}=\alpha_2A^2x^2+\alpha_1Ax+\alpha_0E$ 并合并同类项，得
直接代回并整理（用 $A^2$ 的显式）得
\[
\mathrm{e}^{Ax}=\mat{\mathrm{e}^{-24x}&0&0\\[2pt]\dfrac{2\mathrm{e}^{-2x}-2\mathrm{e}^{-24x}}{11}&\mathrm{e}^{-2x}&0\\[4pt]\dfrac{\mathrm{e}^{2x}-\mathrm{e}^{-24x}}{26}&0&\mathrm{e}^{2x}}.
\]
\par\textbf{答案}\ 上式；通解 $\vect{Y}=\mathrm{e}^{Ax}\vect{C}$。
\par\textbf{自检}\ $x=0$ 得 $E$ \ding{51}；求导在 $0$ 处：$(2,1)$ 元素 $\dfrac{-4+48}{11}=4$ \ding{51}，$(3,1)$ 元素 $\dfrac{2+24}{26}=1$ \ding{51}，与 $A$ 同。

\flowerrule

{\small\color{inkgray}
依据课件：\texttt{16\_17\_matrix\_exponential.pdf}（44 页，2026-10-13，W.Y., DLUT）。页码为该 PDF 页序。
缺口 10 条见 §9；附录中的 $\mathrm{e}^{Ax}$ 均用符号计算核对过导数与初值。}
\subhdr{第 4 章（L17--18）· 例题与作业详解}

每块自成一页，按“题干 → 考点 → 关键一步 → 详解 → 答案 → 自检”排。

\qbreak
\subhdr{A.31\quad 例 1：参数变易（逐分量）（p7--p10）}
\label{c4:app:ex1}
\qbreak
\qstem{（课件 p7 例 1）$y_1'=\tfrac73y_1+\tfrac23y_2-5\mathrm{e}^x$，$y_2'=\tfrac43y_1+\tfrac53y_2-8\mathrm{e}^x$，$y_1(0)=1,\ y_2(0)=4$。用参数变易法解。}
\par\textbf{考点}\ 逐分量参数变易的全流程。
\par\textbf{关键一步}\ 先求齐次基本解组，再解 $c_1'\vect{Y}_1+c_2'\vect{Y}_2=\vect{F}$。
\par\textbf{详解}\ $A=\mat{7/3&2/3\\4/3&5/3}$，特征方程 $\lambda^2-4\lambda+3=0$，$\lambda=1,3$；特征向量 $\vect{V}_1=(-1,2)^{\mathsf T}$、$\vect{V}_2=(1,1)^{\mathsf T}$，故
\[
\vect{Y}_1=\mathrm{e}^{x}\mat{-1\\2},\qquad \vect{Y}_2=\mathrm{e}^{3x}\mat{1\\1}.
\]
设 $\vect{Y}_p=c_1\vect{Y}_1+c_2\vect{Y}_2$，由 ($\star$)：
\[
-c_1'\mathrm{e}^x+c_2'\mathrm{e}^{3x}=-5\mathrm{e}^x,\qquad
2c_1'\mathrm{e}^x+c_2'\mathrm{e}^{3x}=-8\mathrm{e}^x .
\]
两式相减：$3c_1'\mathrm{e}^x=-3\mathrm{e}^x$，故 $c_1'=-1$；代回得 $c_2'\mathrm{e}^{3x}=-6\mathrm{e}^x$，故 $c_2'=-6\mathrm{e}^{-2x}$。积分（常数取零）：$c_1=-x$，$c_2=3\mathrm{e}^{-2x}$。于是
\[
\vect{Y}_p=-x\mathrm{e}^x\mat{-1\\2}+3\mathrm{e}^{-2x}\mathrm{e}^{3x}\mat{1\\1}=\mathrm{e}^x\mat{x+3\\-2x+3}.
\]
\textbf{答案}\ 通解 $\vect{Y}=\vect{Y}_p+c_1\vect{Y}_1+c_2\vect{Y}_2$；代初值 $(1,4)$ 得 $c_1=1,c_2=-1$，故
\[
\vect{Y}=\mathrm{e}^x\mat{x+2\\-2x+5}-\mathrm{e}^{3x}\mat{1\\1}.
\]
\par\textbf{自检}\ 代 $x=0$：$(2,5)-(1,1)=(1,4)$ \ding{51}。代回第一式：$y_1=(\mathrm{e}^x(x+2))'-\mathrm{e}^{3x}=\mathrm{e}^x(x+3)-3\mathrm{e}^{3x}$，而右端 $\tfrac73y_1+\tfrac23y_2-5\mathrm{e}^x=\mathrm{e}^x[\tfrac73(x+2)+\tfrac23(-2x+5)-5]+\mathrm{e}^{3x}[-\tfrac73-\tfrac23]=\mathrm{e}^x(x+3)-3\mathrm{e}^{3x}$ \ding{51}。

\qpagebreak
\subhdr{A.32\quad 例 2：参数变易（矩阵 / $\mathrm{e}^{Ax}$）（p14--p16）}
\label{c4:app:ex2}
\qbreak
\qstem{（课件 p14 例 2）同一题，改用 $\mathrm{e}^{Ax}$ 与卷积公式。}
\par\textbf{考点}\ 卷积公式 $\vect{Y}=\mathrm{e}^{Ax}\vect{Y}_0+\int_0^x\mathrm{e}^{A(x-s)}\vect{F}(s)\mathrm{d}s$。
\par\textbf{关键一步}\ 先用 C--H 求 $\mathrm{e}^{Ax}$（$A$ 的特征值 $1,3$），再代公式。
\par\textbf{详解}\ 由 C--H：$\mathrm{e}^{Ax}=\alpha_1Ax+\alpha_0E$，$\mathrm{e}^{x}=x\alpha_1+\alpha_0$、$\mathrm{e}^{3x}=3x\alpha_1+\alpha_0$，解得 $x\alpha_1=\tfrac12(\mathrm{e}^{3x}-\mathrm{e}^x)$、$\alpha_0=\tfrac12(3\mathrm{e}^x-\mathrm{e}^{3x})$。代回整理：
\[
\mathrm{e}^{Ax}=\frac13\mat{2\mathrm{e}^{3x}+\mathrm{e}^x&\mathrm{e}^{3x}-\mathrm{e}^x\\2\mathrm{e}^{3x}-2\mathrm{e}^x&\mathrm{e}^{3x}+2\mathrm{e}^x}.
\]
卷积项 $\int_0^x\mathrm{e}^{A(x-s)}\vect{F}(s)\mathrm{d}s$ 的被积函数是 $\mathrm{e}^{A(x-s)}\binom{-5}{-8}\mathrm{e}^{s}$，逐项积分（指数为 $\mathrm{e}^{2s}$、$\mathrm{e}^{4s}$ 型），结果与例 1 的 $\vect{Y}_p$ 相同。
\par\textbf{答案}\ 同 A.1：$\vect{Y}=\mathrm{e}^x\binom{x+2}{-2x+5}-\mathrm{e}^{3x}\binom{1}{1}$。
\par\textbf{自检}\ 两条路给出同一答案，说明 $\vect{C}'=\mathrm{e}^{-Ax}\vect{F}$ 与逐分量解 $c_k'$ 等价 \ding{51}。

\qpagebreak
\subhdr{A.33\quad 例 3：待定系数（多项式外力）（p19--p21）}
\label{c4:app:ex3}
\qbreak
\qstem{（课件 p19 例 3）$\vect{Y}'=\mat{1&2\\3&2}\vect{Y}+x\binom{2}{-4}$，求通解。}
\par\textbf{考点}\ Case 1 且 $r=0$ 不是特征值：设 $\vect{Y}_p=\alpha x+\beta$。
\par\textbf{关键一步}\ 比较系数得 $A\alpha=-\gamma$、$\alpha=A\beta+\delta$，其中 $\gamma=(2,-4)^{\mathsf T}$、$\delta=\vect{0}$。
\par\textbf{详解}\ $A$ 的特征值：$\lambda^2-3\lambda-4=0$，$\lambda=-1,4$；特征向量 $(1,-1)$、$(2,3)$。
设 $\vect{Y}_p=\alpha x+\beta$，$\vect{Y}_p'=\alpha$，代入：
\[
\alpha=A\alpha x+A\beta+\gamma x
\ \Longrightarrow\
A\alpha=-\gamma,\quad \alpha=A\beta .
\]
第一式：$\mat{1&2\\3&2}\alpha=\binom{-2}{4}$，解得 $\alpha=\binom{3}{-5/2}$。
第二式：$A\beta=\alpha$，即 $\mat{1&2\\3&2}\beta=\binom{3}{-5/2}$，解得 $\beta=\binom{-11/4}{23/8}$。
\par\textbf{答案}\ $\vect{Y}=c_1\mathrm{e}^{-x}\binom{1}{-1}+c_2\mathrm{e}^{4x}\binom{2}{3}+x\binom{3}{-5/2}+\binom{-11/4}{23/8}$。
\par\textbf{自检}\ 把 $\vect{Y}_p$ 代回：$\vect{Y}_p'=\binom{3}{-5/2}$；$A\vect{Y}_p+\vect{F}=\binom{3}{-5/2}x+\binom{0}{0}+x\binom{2}{-4}$？逐项：$A\alpha x+A\beta+\gamma x=(-2,4)^{\mathsf T}x+(3,-5/2)^{\mathsf T}+(2,-4)^{\mathsf T}x=(3,-5/2)^{\mathsf T}$ \ding{51}。

\qpagebreak
\subhdr{A.34\quad 例 4：共振（$r$ 是特征值）（p22--p26）}
\label{c4:app:ex4}
\qbreak
\qstem{（课件 p22 例 4）$A=\mat{7/3&2/3\\4/3&5/3}$，$\vect{F}=\binom{-5}{-8}\mathrm{e}^x$，$\vect{Y}(0)=\binom14$。用待定系数法解。}
\par\textbf{考点}\ 共振：$r=1$ 是特征值，必须设 $\vect{Y}_p=\mathrm{e}^x(\alpha x+\beta)$。
\par\textbf{关键一步}\ 比较系数得 $\alpha=A\alpha$（$\alpha$ 是特征向量）与 $\beta+\alpha=A\beta+\gamma$；后者是 $(A-E)\beta=\alpha-\gamma$，要验相容性。
\par\textbf{详解}\ 齐次解 $c_1\mathrm{e}^x\binom{-1}{2}+c_2\mathrm{e}^{3x}\binom11$。
设 $\vect{Y}_p=\mathrm{e}^x(\alpha x+\beta)$，求导：$\vect{Y}_p'=\mathrm{e}^x(\alpha x+\beta+\alpha)$；右端 $A\vect{Y}_p+\vect{F}=\mathrm{e}^x(A\alpha x+A\beta+\gamma)$。比较：
\[
x\mathrm{e}^x:\ \alpha=A\alpha;\qquad \mathrm{e}^x:\ \beta+\alpha=A\beta+\gamma .
\]
第一式说明 $\alpha$ 是 $\lambda=1$ 的特征向量，故 $\alpha=k\binom{-1}{2}$。第二式即 $(A-E)\beta=\alpha-\gamma$。这里 $\rank(A-E)=1$，相容性要求 $\alpha-\gamma$ 与 $(A-E)$ 的行空间正交，具体化为
\[
\alpha_1+5=-2\alpha_1+8\ \Longrightarrow\ \alpha_1=1,
\]
故 $\alpha=\binom{1}{-2}$。代回解出 $2\beta_1+\beta_2=9$，取 $\beta=\binom44$ 即得 $\beta=\binom{4}{1}$。
\par\textbf{答案}\ $\vect{Y}=c_1\mathrm{e}^x\binom{-1}{2}+c_2\mathrm{e}^{3x}\binom11+x\mathrm{e}^x\binom{1}{-2}+\mathrm{e}^x\binom41$；代初值得 $c_1=2,c_2=-1$，即 $\vect{Y}=\mathrm{e}^x\binom{x+2}{-2x+5}-\mathrm{e}^{3x}\binom11$（与 A.1 相同）。
\par\textbf{自检}\ 把 $c_1=2,c_2=-1$ 代回初值条件：
\[
2\binom{-1}{2}+(-1)\binom11+\binom41=\binom{-2-1+4}{4-1+1}=\binom14 \ \text{\ding{51}}.
\]
即 $c_1\binom{-1}{2}+c_2\binom11+\binom41=\binom14$ 解出 $c_1=2,c_2=-1$，与上面一致。

\qpagebreak
\subhdr{A.35\quad 三角外力：一个可核对的补例}
\label{c4:app:ex5}
\qbreak
\qstem{（本讲义补例；课件 p29 例 5 的数据自相矛盾，说明见 §9）解 $\vect{Y}'=A\vect{Y}+\vect{F}$，$A=\mat{1&0\\0&2}$，$\vect{F}=\binom{\cos2x}{\sin2x}$。}
\par\textbf{考点}\ Case 2：$r+\mathrm{i}s=2\mathrm{i}$ 不是特征值（特征值是 $1,2$），设 $\vect{Y}_p=\alpha\cos2x+\beta\sin2x$。
\par\textbf{关键一步}\ 求导得 $-2\alpha\sin2x+2\beta\cos2x$，与 $A(\alpha\cos2x+\beta\sin2x)+\vect{F}$ 比较 $\cos$、$\sin$ 的系数。
\par\textbf{详解}\ 比较系数：
\[
\cos2x:\ 2\beta=A\alpha+\binom{1}{0};\qquad \sin2x:\ -2\alpha=A\beta+\binom{0}{1}.
\]
写成分量（设 $\alpha=(a_1,a_2)$、$\beta=(b_1,b_2)$）：
\[
2b_1=a_1+1,\quad 2b_2=2a_2;\qquad -2a_1=b_1,\quad -2a_2=2b_2+1 .
\]
第一组代第二组：$b_1=-2a_1$ 代入 $2b_1=a_1+1$ 得 $-4a_1=a_1+1$，即 $a_1=-\tfrac15$，$b_1=\tfrac25$；$b_2=a_2$ 代入 $-2a_2=2a_2+1$ 得 $a_2=-\tfrac14$，$b_2=-\tfrac14$。
\par\textbf{答案}\
\[
\vect{Y}_p=\binom{-\tfrac15\cos2x+\tfrac25\sin2x}{-\tfrac14\cos2x-\tfrac14\sin2x},
\]
通解再加齐次部分 $c_1\mathrm{e}^{x}\binom10+c_2\mathrm{e}^{2x}\binom01$。
\par\textbf{自检}\ 代回第一个分量：$\vect{Y}_{p,1}'=\tfrac25\cos2x+\tfrac45\sin2x$；右端 $y_{p,1}+\cos2x=-\tfrac15\cos2x+\tfrac25\sin2x+\cos2x=\tfrac45\cos2x+\tfrac25\sin2x$ \ding{51}。第二个分量：$\vect{Y}_{p,2}'=\tfrac12\sin2x-\tfrac12\cos2x$；右端 $2y_{p,2}+\sin2x=-\tfrac12\cos2x-\tfrac12\sin2x+\sin2x=-\tfrac12\cos2x+\tfrac12\sin2x$ \ding{51}。

\qpagebreak
\subhdr{A.36\quad 例 6：三种做法（p32--p48）}
\label{c4:app:ex6}
\qbreak
\qstem{（课件 p32 例 6）$x''=-2x'-5y+3$，$y'=x'+2y$，$x(0)=0,x'(0)=0,y(0)=1$。}
\par\textbf{考点}\ 一题三解：系统 + 待定系数、系统 + 参数变易、化二阶方程。
\par\textbf{关键一步}\ 三种做法的公共起点是"把二阶方程化成一阶组"或"把组化成二阶方程"，二选一。
\par\textbf{详解}\ \textbf{做法一}：令 $x_1=x,x_2=x',x_3=y$，得 $3\times3$ 系统，$A$ 的特征值 $0,\pm\mathrm{i}$。外力是常向量，$0$ 是特征值，故设 $\vect{X}_p=\alpha t+\beta$，比较系数得 $A\alpha=0$、$A\beta+\gamma=\alpha$，相容性给出 $\alpha=(-6,0,0)$、$\beta=(0,-6,3)$。
\textbf{做法二}：同一系统，用 $\Phi$ 与 $\vect{C}'=\Phi^{-1}\vect{F}$，积分后乘回，得同一 $\vect{X}_p$。
\textbf{做法三}：由 $y'=x'+2y$ 解出 $x'=y'-2y$，求导得 $x''=y''-2y'$，代入第一式得 $y''+y=3$；解 $y=c_1\cos t+c_2\sin t+3$，回代求 $x$。
\par\textbf{答案}\ $x=-2-6t+2\cos t+6\sin t$，$y=3-2\cos t+2\sin t$。
\par\textbf{自检}\ $x(0)=0$、$x'(0)=-6+6=0$、$y(0)=1$ \ding{51}；且 $y'=2\sin t+2\cos t$，$x'+2y=(-6-2\sin t+6\cos t)+(6-4\cos t+4\sin t)=2\sin t+2\cos t$ \ding{51}。

\qpagebreak
\subhdr{A.37\quad 作业：两种方法各解一次}
\label{c4:app:hw}
\qbreak
\qstem{（课件 p49）$y_1'=-2y_1-2y_2+2y_3+2x+3$，$y_2'=-3y_1+y_2-4y_3+3x-1$，$y_3'=-3y_1+2y_2-5y_3+3x-2$。}
\par\textbf{考点}\ $3\times3$ 常系数、多项式外力，两条方法互验。
\par\textbf{关键一步}\ 特征值 $-1,-2,-3$ 互异；$0$ 不是特征值，设 $\vect{Y}_p=\alpha x+\beta$。
\par\textbf{详解}\ $A=\mat{-2&-2&2\\-3&1&-4\\-3&2&-5}$，特征值 $-1,-2,-3$，特征向量
\[
\lambda=-1:\ \vect{V}_1=\mat{-2\\5\\4},\qquad \lambda=-2:\ \vect{V}_2=\mat{-1\\3\\3},\qquad \lambda=-3:\ \vect{V}_3=\mat{0\\1\\1}.
\]
\textbf{待定系数法}：设 $\vect{Y}_p=\alpha x+\beta$，外力写成 $\gamma x+\delta$，$\gamma=(2,3,3)^{\mathsf T}$、$\delta=(3,-1,-2)^{\mathsf T}$。比较系数：
\[
A\alpha=-\gamma,\qquad \alpha=A\beta+\delta .
\]
第一式解得 $\alpha=(1,0,0)^{\mathsf T}$；第二式解得 $\beta=(0,1,0)^{\mathsf T}$。故 $\vect{Y}_p=(x,1,0)^{\mathsf T}$。
\textbf{参数变易法}：取 $\Phi=V\mathrm{e}^{\Lambda x}$，$V=(\vect{V}_1,\vect{V}_2,\vect{V}_3)$。计算 $\Phi^{-1}\vect{F}$ 得
\[
\Phi^{-1}\vect{F}=\mat{\mathrm{e}^{x}\\(-2x-5)\mathrm{e}^{2x}\\9(x+1)\mathrm{e}^{3x}},
\]
积分：$\vect{C}=\bigl(\mathrm{e}^x,\,-(x+2)\mathrm{e}^{2x},\,(3x+2)\mathrm{e}^{3x}\bigr)^{\mathsf T}$。乘回去：
\[
\Phi\vect{C}=\mat{x\\1\\0},
\]
与待定系数法\textbf{完全相同}——两条方法在这题上给出同一个特解。
\par\textbf{答案}\
\[
\vect{Y}=c_1\mathrm{e}^{-x}\mat{-2\\5\\4}+c_2\mathrm{e}^{-2x}\mat{-1\\3\\3}+c_3\mathrm{e}^{-3x}\mat{0\\1\\1}+\mat{x\\1\\0}.
\]
\par\textbf{自检}\ 代 $\vect{Y}_p=(x,1,0)$：$\vect{Y}_p'=(1,0,0)$，而 $A\vect{Y}_p+\vect{F}=\mat{-2x-2+2x+3\\-3x+1+3x-1\\-3x+2+3x-2}=(1,0,0)$ \ding{51}。

\flowerrule

{\small\color{inkgray}
依据课件：\texttt{17\_18\_Nonhomogeneous-linear-systems.pdf}（49 页，2026-10-15 / 10-20，W.Y., DLUT）。页码为该 PDF 页序。
缺口 10 条见 §9；作业两法的结果已用符号计算核对（特解相同、代回残差为零）。}
\subhdr{第 6 章（L21--L23）· 四道题的完整详解}

每道题自成一页（块前用 \verb|\qbreak| 保证整块不跨页）：块首复述题干与关键条件，然后五段——考点、关键一步、详解、答案、自检。答案里的每个数都能在正文里找到出处，若找不到，说明是本附录新算的，步骤也写全了。

下面这张表只是索引，\textbf{不含答案}：做完一题，翻到对应页，再核对。

\begin{minipage}{\linewidth}
\tabtop
\begin{tabular}{@{}llll@{}}
\thead{题号} & \thead{要判断什么} & \thead{回顾正文的哪一节} & \thead{页码}\\
\end{tabular}
\tabmid
\begin{tabular}{@{}l p{0.30\linewidth} p{0.36\linewidth} l@{}}
第 1 题 1) & 相异实根的类型与稳定性 & §2.3 通解、§3.2 结点、§5.3 判据 & 第 \pageref{c6:app:a1} 页\\
第 1 题 2) & 重根的类型（$g$ 的作用）与稳定性 & §2.5 重根、§3.5 两种重根 & 第 \pageref{c6:app:a2} 页\\
第 1 题 3) & 复根的类型与稳定性 & §2.4 实解、§3.2 焦点、§4.3 与中心的区别 & 第 \pageref{c6:app:a3} 页\\
第 2 题 & 非线性系统平凡解 & §6.2 雅可比、§6.4 定理 6.1 & 第 \pageref{c6:app:a4} 页\\
\end{tabular}
\tabbot
\end{minipage}

%---------------------------------------------------------------------
\qpagebreak
\qbreak
\subhdr{A.38\quad 第 1 题 1)：不稳定结点}
\label{c6:app:a1}
\qstem{（课件 L22--23 作业 第 1 题 1)；条件摘要：$x'=3x+2y$，$y'=x+4y$，系数矩阵按 $x,y$ 的顺序读出）\textbf{求临界点的类型，并研究其稳定性。}}

\textbf{考点}\ 线性系统的临界点（§1.5）；$\tau,\Delta,D$ 与类型（§3.1）；结点与稳定性的符号判据（§3.2、§5.3）。

\textbf{关键一步}\ 先算 $\Delta$，确认"临界点是唯一的"；再用 $\tau,\Delta,D$ 三个数代替解特征方程。这一步走通，全题就是代数。

\textbf{详解}\
第一步，找临界点。解
\[
3x+2y=0,\qquad x+4y=0,\qquad
\det\begin{pmatrix}3&2\\1&4\end{pmatrix}=12-2=10\neq0,
\]
齐次方程组只有零解，所以\textbf{唯一}临界点是 $(0,0)$。

第二步，算三个数：
\[
\tau=3+4=7,\qquad \Delta=3\cdot4-2\cdot1=10,\qquad D=7^{2}-4\cdot10=9>0 .
\]
特征方程 $\lambda^{2}-7\lambda+10=0$，即 $(\lambda-2)(\lambda-5)=0$，$\lambda_1=2$，$\lambda_2=5$。

第三步，求特征向量（写通解时要用，也为作图提供"骨架"）。
对 $\lambda=2$：$A-2I=\begin{pmatrix}1&2\\1&2\end{pmatrix}$，$x+2y=0$，取 $V_1=(-2,1)^{\mathrm T}$。
对 $\lambda=5$：$A-5I=\begin{pmatrix}-2&2\\1&-1\end{pmatrix}$，$x=y$，取 $V_2=(1,1)^{\mathrm T}$。
通解为
\[
\X=c_1\begin{pmatrix}-2\\1\end{pmatrix}\mathrm{e}^{2t}+c_2\begin{pmatrix}1\\1\end{pmatrix}\mathrm{e}^{5t}.
\]

第四步，定类型。$\Delta=10>0$、$D=9>0$、两根同号且不等，按表 3.1 是\textbf{结点}；$\tau=7>0$ 两根都为正，轨线都往外走，故是\textbf{不稳定结点}（图 \ref{c6:fig:kinds}(a)）。

第五步，定稳定性。$\lambda_2=5>0$，由定理 5.1(ii) 零解\textbf{不稳定}。也可直接指认一条越界轨线：取初值 $(1,1)$，解为 $\X=(1,1)^{\mathrm T}\mathrm{e}^{5t}$，范数 $\sqrt2\,\mathrm{e}^{5t}\to\infty$，对任意 $\eps$ 终将越界。

\textbf{答案}\ 唯一临界点 $(0,0)$；类型为\textbf{不稳定结点}；零解\textbf{不稳定}。

\textbf{自检}\ 把 $\X=(1,1)^{\mathrm T}\mathrm{e}^{5t}$ 代回原方程：$x'=y'=5\mathrm{e}^{5t}$；而 $3x+2y=5\mathrm{e}^{5t}$，$x+4y=5\mathrm{e}^{5t}$，两式都对。再看初值反解：$c_1(-2)+c_2=x_0$，$c_1+c_2=y_0$，系数行列式 $-3\neq0$，任意初值都可表示。

%---------------------------------------------------------------------
\qpagebreak
\qbreak
\subhdr{A.39\quad 第 1 题 2)：临界结点}
\label{c6:app:a2}
\qstem{（课件 L22--23 作业 第 1 题 2)；条件摘要：$x'=-x$，$y'=-y$，$A=-I$，$\tau=-2$，$\Delta=1$）\textbf{求临界点的类型，并研究其稳定性。}}

\textbf{考点}\ 重根、几何重数 $g$（§2.5、§3.5）；"临界结点"的判定；$\tau<0$ 且 $g=2$ 时的稳定性（定理 5.1(i)）。

\textbf{关键一步}\ $D=0$ 说明是重根，此时不能止步于"结点"，必须再算 $A-\lambda I$ 的秩，看 $g$ 是 $2$ 还是 $1$——这一步决定了类型名，也决定了要不要出现 $t$ 项。

\textbf{详解}\
第一步，临界点。$A=-I$，$\Delta=\det(-I)=1\neq0$，故唯一临界点是 $(0,0)$。

第二步，特征值与重数。
\[
\tau=-1-1=-2,\qquad \Delta=(-1)(-1)-0=1,\qquad D=4-4=0 .
\]
特征方程 $\lambda^{2}+2\lambda+1=(\lambda+1)^{2}=0$，故 $\lambda=-1$ 是\textbf{二重根}（\EN{double root}）。

第三步，几何重数。$A-\lambda I=A+I=0$（零矩阵），方程 $(A-\lambda I)V=\mathbf{0}$ 对\textbf{任意} $V$ 都成立，所以线性无关的特征向量有两个，$g=2$。这属于 §2.5 的第一种情形，通解取 (2.6)：
\[
x=c_1\mathrm{e}^{-t},\qquad y=c_2\mathrm{e}^{-t}
\]
（这里直接取 $V_1=(1,0)^{\mathrm T}$、$V_2=(0,1)^{\mathrm T}$，因为两个坐标方向都是特征方向）。

第四步，看轨线。$c_1\neq0$ 时
\[
\frac{y(t)}{x(t)}=\frac{c_2}{c_1}=\text{常数},
\]
所以每条轨线都是\textbf{过原点的射线}，这正是图 \ref{c6:fig:kinds}(e) 的"直线束"。这种图形叫\textbf{临界结点}（星形结点，proper node）。

第五步，定稳定性。$\lambda=-1<0$，且 $\norm{\X(t)}=\sqrt{c_1^{2}+c_2^{2}}\,\mathrm{e}^{-t}\to0$，由定理 5.1(i) 零解\textbf{渐近稳定}（还是指数衰减）。

\textbf{答案}\ 唯一临界点 $(0,0)$；类型为\textbf{临界结点}（$g=2$ 的重根）；零解\textbf{渐近稳定}。

\textbf{自检}\ 取 $(x_0,y_0)=(1,2)$，则 $x=\mathrm{e}^{-t}$、$y=2\mathrm{e}^{-t}$，恒有 $y=2x$（轨线是直线 $y=2x$ 的一段），且 $t\to\infty$ 时趋于原点。再验证方程：$x'=-\mathrm{e}^{-t}=-x$，$y'=-2\mathrm{e}^{-t}=-y$，两式都对。

\textbf{一个对照}\ 若本题的 $A$ 是 $\begin{pmatrix}-1&1\\0&-1\end{pmatrix}$（同样 $\lambda=-1$ 二重，但 $g=1$），通解会含 $t\mathrm{e}^{-t}$ 项，类型要写"退化结点"；稳定性仍是渐近稳定。所以"$g$"只改类型名，不改稳定性——但类型名是题目要问的，不能省。

%---------------------------------------------------------------------
\qpagebreak
\qbreak
\subhdr{A.40\quad 第 1 题 3)：稳定焦点}
\label{c6:app:a3}
\qstem{（课件 L22--23 作业 第 1 题 3)；条件摘要：$x'=-x+2y$，$y'=-2x-y$，系数矩阵按 $x,y$ 的顺序读出）\textbf{求临界点的类型，并研究其稳定性。}}

\textbf{考点}\ 共轭复根与实解 (2.5)（§2.4）；$D<0$、$\tau\neq0$ 对应焦点（§3.2）；"稳定焦点"与"中心"的分界（§4.3）。

\textbf{关键一步}\ 算出 $\tau$ 与 $\Delta$ 后先看 $D$：$D<0$ 说明有旋转；再看 $\tau/2$ 的符号，它决定螺旋是向内还是向外。两句话就能定位，不必先解方程。

\textbf{详解}\
第一步，临界点。$\Delta=(-1)(-1)-2\cdot(-2)=1+4=5\neq0$，唯一临界点 $(0,0)$。

第二步，三个数：
\[
\tau=-1-1=-2,\qquad \Delta=5,\qquad D=4-20=-16<0 .
\]
特征方程 $\lambda^{2}+2\lambda+5=0$，$\lambda=-1\pm2\mathrm{i}$。即 $\alpha=-1$，$\beta=2$。

第三步，特征向量。对 $\lambda=-1+2\mathrm{i}$，
\[
A-\lambda I=\begin{pmatrix}-2\mathrm{i}&2\\-2&-2\mathrm{i}\end{pmatrix},
\qquad -2\mathrm{i}x+2y=0\ \Longrightarrow\ y=\mathrm{i}x .
\]
取 $x=1$，得 $w=(1,\mathrm{i})^{\mathrm T}=u+\mathrm{i}v$，其中 $u=(1,0)^{\mathrm T}$，$v=(0,1)^{\mathrm T}$。

第四步，写实解。代入 (2.5)：
\[
x=\mathrm{e}^{-t}\big(c_1\cos2t+c_2\sin2t\big),\qquad
y=\mathrm{e}^{-t}\big(-c_1\sin2t+c_2\cos2t\big).
\]
在 $t=0$ 处 $x=c_1$，$y=c_2$，所以 $c_1,c_2$ 就是初值。

第五步，定类型与稳定性。$\Delta>0$、$D<0$、$\tau/2=-1\neq0$，按表 3.1 是\textbf{焦点}；$\tau=-2<0$ 故\textbf{稳定焦点}。半径
\[
r(t)=\sqrt{x^{2}+y^{2}}=\mathrm{e}^{-t}\sqrt{c_1^{2}+c_2^{2}}\to0,
\]
角度以角速度 $2$ 增加（从 $(1,0)$ 出发时 $y'=-x=-1<0$，向下走，故为顺时针），所以是顺时针向内卷的螺旋，由定理 5.1(i) 零解\textbf{渐近稳定}。

\textbf{答案}\ 唯一临界点 $(0,0)$；类型为\textbf{稳定焦点}；零解\textbf{渐近稳定}（指数衰减，每单位时间半径乘 $\mathrm{e}^{-1}$）。

\textbf{自检}\ 取 $(c_1,c_2)=(1,0)$：$x=\mathrm{e}^{-t}\cos2t$，$y=-\mathrm{e}^{-t}\sin2t$。代第一式：
\[
x'=-\mathrm{e}^{-t}\cos2t-2\mathrm{e}^{-t}\sin2t,\qquad
-x+2y=-\mathrm{e}^{-t}\cos2t-2\mathrm{e}^{-t}\sin2t,
\]
两边相同；第二式同理。再检查 $t=0$ 时 $(x,y)=(1,0)$，与取定的初值一致。

%---------------------------------------------------------------------
\qpagebreak
\qbreak
\subhdr{A.41\quad 第 2 题：平凡解不稳定}
\label{c6:app:a4}
\qstem{（课件 L22--23 作业 第 2 题）\textbf{条件摘要}：$x_1'=x_1-x_1^{2}-x_1x_2$，$x_2'=0.5x_2-0.75x_1x_2-0.25x_2^{2}$；平凡解指 $\X\equiv\mathbf{0}$。\textbf{研究平凡解的稳定性。}}

\textbf{考点}\ 平凡解与临界点（§4.5）；雅可比矩阵 (6.9)（§6.2）；线性化定理 6.1(b) 及其"双曲"前提（§6.3、§6.4）；为什么不需要李雅普诺夫函数（§7.5）。

\textbf{关键一步}\ 先确认 $\mathbf{0}$ 确是临界点（否则"平凡解"无从谈起），再只保留一次项算 $\Jac(\mathbf{0})$；只要特征值全不为零，线性化定理就能直接下结论，不需再碰二次项。

\textbf{详解}\
第一步，确认平凡解。把 $x_1=x_2=0$ 代入右端，两个式子都是 $0$，所以 $\X(t)\equiv\mathbf{0}$ 是解，且原点是临界点。

第二步，线性化。记两个右端分别为 $P=x_1-x_1^{2}-x_1x_2$、$Q=0.5x_2-0.75x_1x_2-0.25x_2^{2}$，逐项求偏导：
\[
P_{x_1}=1-2x_1-x_2,\qquad P_{x_2}=-x_1,\qquad
Q_{x_1}=-0.75x_2,\qquad Q_{x_2}=0.5-0.75x_1-0.5x_2 .
\]
在原点代入 $x_1=x_2=0$，得到
\[
\Jac(\mathbf{0})=\begin{pmatrix}P_{x_1}&P_{x_2}\\Q_{x_1}&Q_{x_2}\end{pmatrix}_{\mathbf{0}}
=\begin{pmatrix}1&0\\0&0.5\end{pmatrix}.
\]
注意二次项在原点处的偏导全为零（$x_1^{2}$ 的偏导是 $2x_1$，$x_1x_2$ 对 $x_1$ 的偏导是 $x_2$，代 $0$ 都得零），所以 $\Jac(\mathbf{0})$ 恰好就是"线性部分的系数矩阵"。

第三步，特征值。对角矩阵，$\lambda_1=1$，$\lambda_2=0.5$。

第四步，判断。两个特征值都不为零，故原点是\textbf{双曲}临界点，定理 6.1 适用；$\lambda_1=1>0$，由定理 6.1(b)，平凡解\textbf{不稳定}。再由定理 6.2（Hartman--Grobman），邻域内的相图与线性系统 $\U'=\operatorname{diag}(1,0.5)\U$ 拓扑等价，即\textbf{不稳定结点}型：轨线沿两个坐标方向离开原点，$x_1$ 方向增长快（因子 $\mathrm{e}^{t}$），$x_2$ 方向增长慢（因子 $\mathrm{e}^{0.5t}$），见图 \ref{c6:fig:hw2}(b)。

第五步，说明不需要第二方法。$\Jac(\mathbf{0})$ 的特征值非零，属于定理 6.1 的适用情形，结论已经确定；再用李雅普诺夫函数既不能加强结论，也不是必需（§7.5）。

\textbf{答案}\ 平凡解 $\X\equiv\mathbf{0}$ \textbf{不稳定}（既非稳定，也非渐近稳定）；线性化部分为不稳定结点（$\lambda=1,\ 0.5$）。

\textbf{自检}\ 取初值 $(0.01,\,0)$：此时 $x_2\equiv0$，方程组退化为 $x_1'=x_1(1-x_1)$（逻辑斯蒂方程，\EN{logistic equation}），其解
\[
x_1(t)=\frac{1}{1+99\,\mathrm{e}^{-t}},
\]
从 $0.01$ 单调增到 $1$——确实离开了原点（距离趋于 $1$），最后停在另一个临界点 $(1,0)$。这同时说明"不稳定"不等于"跑到无穷"：它只是"离开了原点"（§4.6）。另一个初值 $(0,\,0.01)$ 类似，会被 $(0,2)$ 接住。

%=====================================================================
\newpage

\settags{附录 B\quad 题面歧义说明}{微分方程 · 期中之后}
\hdr{附录 B}{末项的题面歧义与替代读法}

这道题是从手机截图转录的，末项 $0.25x_2^{2}$ 上的平方号很小，早期 OCR 曾把它读成线性的 $0.25x_2$。本讲义按\textbf{原图}（放大两倍后可见上标 $2$）取 $0.25x_2^{2}$。下面说明：即使按另一种读法，结论也不变——这一点值得写进答案，因为它说明结论对题面的小歧义是稳健的。

\subhdr{B.1 两种读法与各自的雅可比矩阵}

读法一（本讲义采用）：$x_2'=0.5x_2-0.75x_1x_2-0.25x_2^{2}$。$\Jac(\mathbf{0})=\operatorname{diag}(1,\,0.5)$，特征值 $1$、$0.5$。

读法二：$x_2'=0.5x_2-0.75x_1x_2-0.25x_2$，即末项线性。合并同类项得 $x_2'=0.25x_2-0.75x_1x_2$，于是 $\Jac(\mathbf{0})=\operatorname{diag}(1,\,0.25)$，特征值 $1$、$0.25$。

两种读法的特征值都\textbf{为正、互不相等}，故都落在表 3.1 的"结点、$\tau>0$"一行：\textbf{不稳定结点}，平凡解\textbf{不稳定}。差别只在第二个特征值的大小（$0.5$ 与 $0.25$），它影响轨线"扁"的程度，不影响类型与稳定性。

\subhdr{B.2 连系数读错也不影响结论}

若交叉项系数 $0.75$ 被读成别的正数（例如 $0.25$），$\Jac(\mathbf{0})$ 仍然是 $\operatorname{diag}(1,\,0.5)$——因为 $\partial(x_1x_2)/\partial x_1=x_2$ 在原点为零，交叉项根本进不了一次项。所以这道题的"结论"只依赖两个线性系数（$1$ 与 $0.5$），与二次项的系数无关。这也是一个可写进答案的观察：\textbf{二次项只影响高阶修正，不影响线性化结论}。

\subhdr{B.3 建议的写法}

按原图作答，并在答案里把雅可比矩阵写出来：
\[
\Jac(\mathbf{0})=\begin{pmatrix}1&0\\0&0.5\end{pmatrix},\qquad \lambda_1=1>0,\ \lambda_2=0.5>0,
\]
然后引定理 6.1(b) 得"不稳定"。若担心老师原题是线性末项，可以加一句："若末项为 $0.25x_2$（线性），则第二特征值为 $0.25$，结论（不稳定结点）不变。"这句话既显示你检查过题面，也不会被判为答非所问。

\subhdr{B.4 关于重根那一问的名称}

有的教材不区分"临界结点"与普通结点，一律叫 node；本课程作业只要求"type"，写\ "结点（临界结点/星形结点，$g=2$）"\ 最稳妥：既给出中文名，也给出英文别名与判定依据。第 1 题 (1) 同理，写"结点（不稳定结点）"即可，不必纠结 proper/improper 的译名。

\flowerrule

\end{document}

\flowerrule
{\small\color{inkgray}
合并说明：本讲义由六份分册讲义（2026-10-28 至 2026-11-13）合并而成，正文推导、例题与附录详解均保留原稿；合并只做了章题、节号前缀、附录连续编号与标签去重。
各章正文的“详解”页码由 \LaTeX{} 交叉引用自动生成，重排版后会同步更新。
制表：2026-11-18。}

\end{document}
